\chapter{全球AI智能体安全研究现状}
\DefTblrTemplate{conthead-text}{default}{（续表）}
\DefTblrTemplate{contfoot-text}{default}{续下页}

本章以2025--2026年为时间窗口期，围绕AI智能体安全领域展开全域现状调研，系统梳理国内外学术研究成果、产业端主流安全产品解决方案、全球各国监管政策与合规规制体系，并结合近两年爆发的典型重大安全事件进行复盘分析，全面剖析领域技术、产业、治理及风险等方面的发展现状。

% 本章主要内容包括学术界研究现状（以顶尖高校与国家实验室为主）、工业界技术产品现状（以头部科技企业和主流产品为主）各国智能体安全法律法规以及全球AI硬件和芯片底座安全研究。本章主要聚焦北美、欧洲与我国，横向对比三大区域在AI智能体安全领域的学术研究、产业、政策与底层硬件等方面差异。

\section{全球AI智能体安全学术研究现状}
本节主要围绕北美、欧洲及我国的顶尖高校与国家级实验室进行AI智能体安全研究的调研与分析。
% 学术研究方面，北美高校较早把智能体放入真实 Web、代码库和工具调用环境中评测，红队测试、透明度披露和企业级安全接口也较活跃。
% 欧洲高校和研究机构更重视可信 AI、隐私保护、形式化验证、法律责任和人类监督，AgentDojo 等工作则提供了可复现的提示注入评测样本。
% 国内高校和科研院所主要从大模型安全、网络空间安全、软件工程和行业应用进入这一议题，研究对象正在从内容安全扩展到工具安全、智能体评测和安全运营。

\subsection{北美AI智能体安全学术研究现状}
% 北美作为全球AI智能体安全领域核心研究高地，依托顶尖高校、国家级实验室与跨学科研究平台，建成分层、体系化、实战导向的完整研究布局。
本节梳理北美地区AI智能体安全的研究主体、研究方向和标志性学术成果，覆盖美国麻省理工学院、斯坦福大学、加州大学伯克利分校、卡内基梅隆大学等顶尖高校，以及加拿大多伦多大学、麦吉尔大学、英属哥伦比亚大学等顶尖高校，核心发现如下（详见表~\ref{tab:us_top_institution}）：

\begin{itemize}[tabitem]
\item \textbf{研究资源集中于头部高校，各校研究方向存在明确区分。}
MIT开展底层形式化验证、网络攻防智能体、通用AI风险分类研究\cite{mitRus2026,mitCSSGroup2026,mitAIRiskRepo2026}；斯坦福开展提示注入防御、权限校验、行为溯源、模型透明性评估\cite{stanfordHAI2026,stanfordCRFM2026}；加州伯克利开展函数调用评测、上下文安全理论、线下攻防实验\cite{berkeleyBAIR2026,contextsecurity2026}；卡内基梅隆大学搭建交互式智能体仿真环境，研发代码安全奖励模型\cite{seccodeprm2026}。

现有研究划分基础理论、系统防护、场景应用三层体系：理论产出上下文安全四维属性理论\cite{contextsecurity2026}、DemPO对齐算法\cite{dempo2026}、WARP检索污染攻击\cite{warp2026}；系统工具包含JAI沙箱\cite{stanfordJai2026}、IsolateGPT插件隔离、Progent权限代理\cite{progent2025}、WaVe WebAssembly验证工具；研究场景覆盖代码、Web、机器人、网络攻防四类智能体；评测工具包含Gorilla、BFCL、APIBench、SWE-bench、WebArena、SafeArena\cite{safearena2025}、HELM、MIT风险库\cite{mitAIRiskRepo2026}、FMTI透明度指数，分别用于函数调用、代码修复、网页交互、风险分类与模型透明性量化测试。

\item \textbf{大量公开实验观测到自主智能体可实现端到端链式攻击，多款自动化红队平台已实现漏洞挖掘、验证、修复全流程自动化。}
卡内基梅隆大学复现Equifax完整数据泄露链路，智能体自主完成漏洞利用、恶意程序部署、数据外传\cite{equifaxsim2025}；纽约大学PromptLock可无人工干预执行勒索攻击，EnIGMA网安智能体攻防性能优于传统安全系统\cite{promptlock2025}；多伦多大学实验观测AI蠕虫可动态更新攻击策略、跨主机横向传播\cite{aiworm2026}。
%
普渡大学ASTRA系统针对代码智能体漏洞检出率超过90\%\cite{astra2025}；伯克利CyberGym可自动挖掘零日漏洞并完成CVE申报，覆盖漏洞挖掘、危害验证、修复迭代完整流程\cite{cybergym2025}。

\item \textbf{研究同时覆盖技术、政策、运营多维度，北美高校与国家级机构同步产出治理规范与校内安全处置机制。}
哈佛Kempner中心、普林斯顿CITP开展智能体商业部署边界、平台风控责任相关研究\cite{berkmanklein2026,princetonCITP2026}；加拿大多家国家级机构发布可信AI、生成式AI治理文档，相关内容被官方政策引用\cite{vectorResponsibleAI2025,caisi2026}。
卡内基梅隆大学设立校级AI安全响应团队AISIRT，专职处理生成式AI安全漏洞与突发事件\cite{cmuAISIRT2026}；伊利诺伊大学香槟分校将智能体安全流程并入原有网络安全漏洞管理、应急响应体系\cite{uiucCSL2026}。
\end{itemize}

\begingroup

\begin{center}

% \nopagebreak[4]
\enlargethispage{1cm}
\footnotesize



\begin{longtblr}[
caption = {北美顶尖高校AI智能体安全学术研究汇总},
label   = {tab:us_top_institution},
]
{
colspec = {|m{0.8cm}
|m{1.2cm}
|m{2.0cm}
|m{\dimexpr
  \linewidth-0.8cm-1.2cm-2.0cm-8\tabcolsep
\relax}|},
rowhead = 1,
}

\hline
\textbf{国家} & \textbf{高校} & \textbf{实验室 /\newline 研究团队} & \textbf{AI智能体安全核心研究方向} \\
\hline
\SetCell[r=1]{c} 美国
& \SetCell[r=4]{c} 麻省理工学院\newline (MIT)
& DRL~\cite{mitRus2026}
& \textbf{软体机器人安全控制与人机交互约束：}Daniela Rus团队研发具备数学安全保障的软体机器人控制系统，约束机器人形变、环境适配与人机交互过程中的违规行为~\cite{mitSoftRobotSafety2026}。 \\
\cline{3-4}
& & ALFA Lab~\cite{mitOReilly2026}
& \textbf{智能体网络攻防鲁棒性与威胁情报知识图谱：}Una-May O'Reilly团队基于CybORG平台完成LLM智能体攻防能力评测，构建融合MITRE ATT\&CK、CWE、CVE的威胁知识图谱BRON~\cite{mitALFAResearch2026}。 \\
\cline{3-4}
& & MIT FutureTech~\cite{mitAIRiskRepo2026}
& \textbf{AI风险统一分类与治理框架：}搭建包含1725条条目的AI风险仓库，采用因果、领域双维度分类体系，为智能体风险审计与标准化治理提供分类依据~\cite{mitAIRiskRepo2026}。 \\
\cline{3-4}
& & CSS Group~\cite{mitCSSGroup2026}
& \textbf{形式化验证与底层系统安全：}依托形式化验证技术开发FSCQ安全文件系统、CryptDB加密数据库，从底层硬件存储层保障智能体数据读写安全~\cite{mitCSSGroup2026}。 \\
\cline{2-4}
& \SetCell[r=3]{c} 斯坦福大学\newline (Stan-\newline ford)
& HAI \& CAIS~\cite{stanfordHAI2026}~\cite{stanfordCAIS2026} \allowbreak ~\cite{stanfordKochenderfer2026}
& \textbf{提示注入防护与智能体资产风险分级：}搭建多层越狱防护、插件权限校验体系，提出高风险场景智能体合规溯源与资产分级评估方案；参与《Agents of Chaos》实验，证实自主智能体存在数据泄露、恶意系统操作、身份冒充等原生安全隐患~\cite{stanfordKochenderfer2026,agentsofchaos2026}。 \\
\cline{3-4}
& & SCS Group~\cite{stanfordSCS2026}
& \textbf{智能体文件系统隔离沙箱：}研发jai轻量级Linux沙箱，依托目录权限管控、写时复制机制，抵御智能体误删、篡改本地文件的安全风险~\cite{stanfordJai2026}。 \\
\cline{3-4}
& & CRFM~\cite{stanfordCRFM2026}
& \textbf{大模型风险与能力统一评测：}推出HELM标准化评测框架，实现多任务多维度模型能力与风险量化对比；提出FMTI指数，量化评估模型全生命周期研发、训练、部署透明度水平。 \\
\cline{2-4}
& \SetCell[r=4]{c} 加州大学伯克利分校\newline (UC Berkeley)
& BAIR~\cite{berkeleyBAIR2026}~\cite{pieterAbbeelHomepage2026}
& \textbf{多智能体强化学习决策与函数调用评测体系：}研究机器人自适应协同强化学习策略，支撑具身智能体协同决策；构建Gorilla、BFCL开源基准，标准化量化智能体函数选择、参数生成能力~\cite{pieterAbbeelHomepage2026,bfcl2025}。 \\
\cline{3-4}
& & Sky Lab~\cite{berkeleySky2026}
& \textbf{智能体工具调用运行时隔离：}从操作系统层优化智能体工具调度、任务可观测性，搭建轻量化运行时安全隔离平台~\cite{berkeleySky2026}。 \\
\cline{3-4}
& & CHAI~\cite{berkeleyCHAI2026}
& \textbf{人类偏好对齐与智能体可控性：}基于人机协作理论设计意图对齐算法，提升自主智能体行为可解释性与人工干预可控能力~\cite{berkeleyCHAI2026}。 \\
\cline{3-4}
& & RDI~\cite{dawnSongBerkeley2026}
& \textbf{工具调用权限管控、代码智能体攻防与攻击分类理论：}Progent权限代理系统将工具调用攻击成功率由39.9\%降至1.0\%；CyberGym平台自动挖掘34个软件零日漏洞，验证代码智能体原生攻击风险；提出上下文安全四维属性模型，重构智能体攻击风险完整分类体系~\cite{progent2025,cybergym2025,contextsecurity2026}。 \\
\cline{2-4}
& \SetCell[r=3]{c} 卡内基梅隆大学\newline (CMU)
% & LTI
% & \textbf{Web场景智能体交互评测：}搭建WebArena、VisualWebArena真实网页交互环境，完成文本、多模态智能体交互式任务执行能力标准化评测。 \\
% \cline{3-4}
& CyLab~\cite{cmuCyLab2026}
& \textbf{代码智能体安全对齐与自主链式网络攻击验证：}提出SecCodePRM过程奖励模型，抑制代码生成阶段高危漏洞产出；复现Equifax数据泄露完整漏洞链，证实LLM智能体可自主完成全流程网络攻击~\cite{seccodeprm2026,equifaxsim2025}。 \\
\cline{3-4}
& & Robotics Institute~\cite{cmuRI2026}
& \textbf{多智能体博弈与任务分配：}基于博弈论设计分布式任务调度机制，规避多智能体协同过程中的策略冲突与资源争抢风险~\cite{cmuRI2026}。 \\
\cline{3-4}
& & AISIRT~\cite{cmuAISIRT2026}
& \textbf{生成式AI安全事件响应：}校级专项研究团队，专职开展AI智能体漏洞研判、安全应急处置与风险闭环管控工作~\cite{cmuAISIRT2026}。 \\
\cline{2-4}
& \SetCell[r=4]{c} 普渡\newline 大学\newline (Pur-\newline due)
& ASSET~\cite{purdueASSET2026}
& \textbf{代码智能体推理能力评测：}提出CoRe静态分析基准，量化评测大模型识别、推理代码漏洞的风险判别能力~\cite{purdueASSET2026}。 \\
\cline{3-4}
& & PurCL~\cite{astra2025}
& \textbf{代码智能体攻防实战与隐性漏洞挖掘：}团队斩获Amazon Nova AI Challenge攻击赛道全球第一，验证自动化红队攻击能力；ASTRA系统攻击成功率超90\%，定向挖掘代码助手多步骤链式隐性漏洞~\cite{purcl2025nova,astra2025}。 \\
\cline{3-4}
& & IDEAS~\cite{purdueIDEAS2026}
& \textbf{拒止环境机器人协同导航：}获美军项目资助，研究GPS拒止场景下空地多机器人智能协同定位与安全导航技术~\cite{ideasdod2025}。 \\
\cline{3-4}
& & CERIAS~\cite{purdueCERIAS2026}
& \textbf{组织级智能体安全工程化：}将智能体安全管控嵌入企业漏洞管理、应急演练、合规培训全业务流程，形成标准化安全工程落地方案~\cite{purdueCERIAS2026}。 \\
\cline{2-4}
& \SetCell[r=4]{c} 哈佛\newline 大学\newline (Har-\newline vard)
& Kempner Institute~\cite{kakadeHomepage2026}
& \textbf{自主智能体有害生成机制：}拆解大模型推理规划底层逻辑，统一揭示自主智能体主动生成有害内容的内在机理~\cite{kempnerHarmfulness2026}。 \\
\cline{3-4}
& & PT Project~\cite{privacytoolsharvard2026}
& \textbf{差分隐私数据保护工具：}开源OpenDP轻量化差分隐私框架，为智能体跨主体数据交互提供标准化隐私防护能力~\cite{privacytoolsharvard2026}。 \\
\cline{3-4}
& & James Mickens团队~\cite{mickensharvard2026}
& \textbf{Agentic AI安全治理战略：}系统梳理自主智能体全维度网络安全威胁体系，提出产学研多方协同长效安全治理框架~\cite{mickensharvard2026}。 \\
\cline{3-4}
& & Berkman Klein Center~\cite{berkmanklein2026}
& \textbf{高风险行业智能体制度边界：}从社会技术交叉视角，明确金融、医疗等高风险领域智能体落地的合规红线与监管权责边界~\cite{berkmanklein2026}。 \\
\cline{2-4}
& \SetCell[r=4]{c} 康奈尔大学\newline (Cor-\newline nell)
& Cornell AI Alignment实验室~\cite{cornellLevine2026}
& \textbf{AI对齐基础理论研究：}依托严格数学理论搭建自主智能体价值对齐、行为硬约束底层原理完整理论框架~\cite{cornellLevine2026}。 \\
\cline{3-4}
& & Vitaly Shmatikov团队~\cite{shmatikovHomepage2026}
& \textbf{检索型智能体投毒攻击与嵌入向量信息泄露：}提出WARP网页投毒方法，验证检索增强智能体易遭受外部内容篡改攻击；vec2vec算法证实跨模型嵌入向量可还原原始文本语义，暴露向量检索智能体隐私泄露风险~\cite{warp2026,vec2vec2025}。 \\
\cline{3-4}
& & Martin T. Wells团队~\cite{dempo2026}
& \textbf{多用户价值对齐冲突消解：}设计DemPO优化算法，解决多主体交互场景下智能体训练存在的人群价值偏好偏差冲突问题~\cite{dempo2026}。 \\
\cline{3-4}
& & Thomas Ristenpart团队~\cite{ristenpartcornell2026}
& \textbf{LLM路由系统安全风险：}深度剖析大模型流量调度、任务路由底层架构存在的原生安全漏洞与各类权限绕过攻击路径~\cite{reroutingllm2025}。 \\
\cline{2-4}
& \SetCell[r=4]{c} 佐治亚理工学院\newline (Georgia Tech.)
& SSS Lab~\cite{gatechSCP2026}
& \textbf{软件系统漏洞挖掘实战：}面向智能体底层依赖库、运行时组件开展定向漏洞挖掘与攻防验证实验~\cite{gatechSCP2026}。 \\
\cline{3-4}
& & TR Lab~\cite{gatechSCP2026}
& \textbf{自主机器人环境安全落地：}研究服务、工业机器人智能体在复杂人居场景下的安全决策、动态风险规避机制~\cite{gatechSCP2026}。 \\
\cline{3-4}
& & BEES Lab~\cite{gatechSCP2026}
& \textbf{AI安全风险成因量化分析：}基于大规模全网流量测量，量化解析智能体安全事件现实诱因、传播路径与演化规律~\cite{gatechSCP2026}。 \\
\cline{3-4}
& & Team Atlanta~\cite{gatechAIxCC2026}
& \textbf{LLM自主漏洞攻防系统：}Atlantis框架实现全自动无人工干预漏洞挖掘与修复，斩获DARPA AI Cyber Challenge半决赛冠军~\cite{gatechAIxCC2026}。 \\
\cline{2-4}
& \SetCell[r=4]{c} 伊利诺伊大学厄巴纳-香槟分校\newline (UIUC)
& CSL~\cite{uiucCSL2026}
& \textbf{可信AI智能体综合研究：}融合机器学习、底层系统安全、软件工程技术，搭建端到端全链路可信智能体技术体系~\cite{uiucCSL2026}。 \\
\cline{3-4}
& & ITI~\cite{uiucITI2026}
& \textbf{关键基础设施智能体弹性防护：}面向智能电网等工业关键设施，研发防护型智能体，提升基础设施网络抵御高强度攻击的弹性能力~\cite{uiucITI2026}。 \\
\cline{3-4}
& & SL Lab~\cite{uiucBoLi2026}
& \textbf{智能体对抗鲁棒性增强：}提出具备理论证明的鲁棒学习算法，抵御针对智能体视觉、文本感知层的对抗样本攻击~\cite{uiucBoLi2026}。 \\
\cline{3-4}
& & Heng Ji团队~\cite{uiucHengJi2026}
& \textbf{代码化智能体动作安全表示：}CodeAct框架采用可执行代码替代纯文本动作指令，显著提升智能体工具调用成功率与执行安全边界可控性。 \\
\cline{2-4}
& \SetCell[r=4]{c} 普林斯顿大学\newline (Princ-\newline eton)
& NLP团队~\cite{narasimhanHomepage2026}
& \textbf{代码修复智能体评测与优化：}构建SWE-bench真实软件工程基准，设计SWE-agent专用交互接口，提升智能体自动修复代码漏洞的稳定性与成功率。 \\
\cline{3-4}
& & CITP~\cite{princetonCITP2026}
& \textbf{AI能力夸大风险辨析：}系统拆解智能体商业宣传中的能力夸大、虚假宣传问题，提供标准化风险甄别与核验方法论。 \\
\cline{3-4}
& & Aleksandra Korolova团队~\cite{princetonKorolova2026}
& \textbf{LLM决策系统公平性审计：}基于差分隐私理论范式，量化审计简历筛选、信用评估等落地场景下LLM智能体的算法公平性偏差~\cite{princetonKorolova2026}。 \\
\cline{3-4}
& & S\&P Group~\cite{princetonSecPriv2026}
& \textbf{智能体应用层安全防护：}结合现代密码学、网络攻防技术，构建面向业务场景智能体的上层应用完整安全防护体系~\cite{princetonSecPriv2026}。 \\
\cline{2-4}
& \SetCell[r=4]{c} 纽约\newline 大学\newline (NYU)
& Muhammad Shafique团队~\cite{shafiqueHomepage2026}
& \textbf{网安智能体自动化攻防、自主勒索攻击与多智能体越狱攻击：}EnIGMA网安智能体解题能力达到传统模型三倍，适配全品类CTF攻防场景；PromptLock验证LLM可独立完成完整勒索软件攻击链路；MetaCipher依靠多智能体协同高效绕过主流大模型安全防护护栏~\cite{promptlock2025,metacipher2025}。 \\
\cline{3-4}
& & AITP Lab~\cite{sucholutskyAITP2026}
& \textbf{思维伙伴智能体人机交互：}研究共生式协同思维伙伴系统，优化人机联合决策的安全性、交互流畅度与风险可控性~\cite{sucholutskyAITP2026}。 \\
\cline{3-4}
& & AgentAuditor团队~\cite{agentauditor2025}
& \textbf{通用智能体安全审计方法：}依托安全经验记忆库，使LLM智能体实现接近人类安全专家水平的全维度智能体风险评估能力~\cite{agentauditor2025}。 \\
\cline{3-4}
& & Quanyan Zhu团队~\cite{stackelbergDefense2025}
& \textbf{攻防博弈智能体防御建模：}基于斯塔克尔伯格博弈理论量化攻防双方策略均衡，推导面向智能体的最优主动防御部署方案~\cite{stackelbergDefense2025}。 \\
\hline
\SetCell[r=1]{c} 加拿大
& \SetCell[r=4]{c} 多伦多大学\newline (Toro-\newline nto)
& SRI~\cite{schwartzReisman2026}
& \textbf{智能体全周期治理框架与国家级AI安全指南：}搭建覆盖研发、部署、运维的智能体风险管控与责任追溯体系，规范企业级智能体落地流程；牵头编制的AI治理指南被加拿大政府纳入中小企业AI部署标准工具包~\cite{schwartzReisman2026,vectorResponsibleAI2025}。 \\
\cline{3-4}
& & CleverHans Lab~\cite{cleverhans2026}
& \textbf{AI蠕虫自主传播风险与隐私协作学习框架：}实验证实开源大模型可生成自适应AI蠕虫，依托终端资源实现跨设备自主扩散；提出CaPC Learning多方协作学习框架，在联合训练过程中实现无损隐私保护~\cite{aiworm2026}。 \\
\cline{3-4}
& & TISL~\cite{tisl2026}
& \textbf{具身智能体场景仿真：}优化机器人感知、环境表征算法，搭建高保真物理仿真环境用于具身智能体安全训练与风险验证~\cite{tisl2026}。 \\
\cline{3-4}
& & DORL~\cite{dorl2026}
& \textbf{多智能体强化学习优化：}研究动态开放环境下多智能体协同优化算法，降低复杂任务协作中的冲突与安全失控风险~\cite{dorl2026}。 \\
\cline{2-4}
& \SetCell[r=4]{c} 麦吉尔大学\newline (McGill)
& McGill NLP Group~\cite{sivaReddy2026}
& \textbf{Web智能体安全评测基准：}推出SafeArena评测基准，覆盖正常、恶意两类交互场景，标准化量化Web智能体抗误导、抗各类注入攻击能力~\cite{safearena2025,safearenaGithub2025}。 \\
\cline{3-4}
& & Mila~\cite{mila2026}
& \textbf{全球前沿AI风险治理：}牵头编撰《国际AI安全报告》，主导高能力通用智能体全球协同治理规则体系研究~\cite{internationalAISafetyReport2025}。 \\
\cline{3-4}
& & IVADO R3AI~\cite{sivaReddy2026}
& \textbf{长期AI安全专项资助：}设立专项科研经费，持续支持2024–2030年度智能体安全全方向课题攻关~\cite{sivaReddy2026}。 \\
\cline{3-4}
& & Jackie Cheung团队~\cite{mcgillCheung2026}
& \textbf{文本生成智能体内容安全：}研究大模型文本生成阶段有害内容过滤、生成内容溯源与可追溯取证技术~\cite{mcgillCheung2026}。 \\
\cline{2-4}
& \SetCell[r=4]{c} 英属哥伦比亚大学\newline (UBC)
& S\&P Group~\cite{ubcSPG2026}
& \textbf{部署型智能体对抗检测：}研发深度神经网络异常识别算法，实时在线识别智能体运行阶段遭遇的各类对抗样本攻击~\cite{ubcSPG2026}。 \\
\cline{3-4}
& & Mathias Lécuyer团队~\cite{ubcLecuyer2026}
& \textbf{无重训练隐私审计与多步骤防御认证鲁棒性：}PANORAMIA框架无需模型微调即可快速完成智能体全链路隐私风险检测；自适应平滑算法为多层防御智能体提供可证明的安全鲁棒保障。 \\
\cline{3-4}
& & Victoria Lemieux团队~\cite{ubcLemieux2026}
& \textbf{AI安全公共资源治理：}结合公地治理理论，构建兼顾原住民权益视角的AI数据资产与全域安全治理框架~\cite{ubcLemieux2026}。 \\
\cline{3-4}
& & CAISI~\cite{caisi2026}
& \textbf{国家级AI安全认证体系：}联合加拿大多家顶尖科研机构，制定标准化智能体安全分级评估与合规认证完整规范~\cite{caisi2026}。 \\
\hline
\end{longtblr}
\end{center}
\endgroup

% \begingroup

% \begin{center}

% % % \nopagebreak[4]
% \enlargethispage{1cm}
% \footnotesize

% \begin{longtable}{|>{\rule{0pt}{2pt}}m{0.8cm}
% |>{\rule{0pt}{2pt}}m{1.2cm}
% |>{\rule{0pt}{2pt}}m{2.0cm}
% |>{\rule{0pt}{2pt}}m{\dimexpr
%   \linewidth-0.8cm-1.2cm-2.0cm-8\tabcolsep
% \relax}|}
% % \begin{longtable}{|>{\rule{0pt}{2pt}}m{0.8cm}
% % |>{\rule{0pt}{2pt}\let\multirow\SchoolMultirow}m{1.2cm}
% % |>{\rule{0pt}{2pt}}m{2.0cm}
% % |>{\rule{0pt}{2pt}}m{\dimexpr\linewidth - 0.8cm - 1.2cm - 2.0cm - 8\tabcolsep\relax}|}
% \caption{北美顶尖高校AI智能体安全学术研究汇总}\label{tab:us_top_institution}\\
% \hline
% \textbf{国家} & \textbf{高校} & \textbf{实验室 /\newline 研究团队} & \textbf{AI智能体安全核心研究方向} \\
% \hline
% \endfirsthead
% \multicolumn{4}{l}{\tablename\ \thetable\ -- 续上页}\\
% \hline
% \textbf{国家} & \textbf{高校} & \textbf{实验室 /\newline 研究团队} & \textbf{AI智能体安全核心研究方向} \\
% \hline
% \endhead
% \hline
% \multicolumn{4}{r}{续下页}\\
% \endfoot
% \hline
% \endlastfoot

% \multirow{56}{=}{美国}
% & \multirow{4}{=}{麻省理工学院\\(MIT)}
% & DRL~\cite{mitRus2026}
% & \textbf{软体机器人安全控制与人机交互约束：}Daniela Rus团队研发具备数学安全保障的软体机器人控制系统，约束机器人形变、环境适配与人机交互过程中的违规行为~\cite{mitSoftRobotSafety2026}。 \\
% \cline{3-4}
% & & ALFA Lab~\cite{mitOReilly2026}
% & \textbf{智能体网络攻防鲁棒性与威胁情报知识图谱：}Una-May O'Reilly团队基于CybORG平台完成LLM智能体攻防能力评测，构建融合MITRE ATT\&CK、CWE、CVE的威胁知识图谱BRON~\cite{mitALFAResearch2026}。 \\
% \cline{3-4}
% & & MIT FutureTech~\cite{mitAIRiskRepo2026}
% & \textbf{AI风险统一分类与治理框架：}搭建包含1725条条目的AI风险仓库，采用因果、领域双维度分类体系，为智能体风险审计与标准化治理提供分类依据~\cite{mitAIRiskRepo2026}。 \\
% \cline{3-4}
% & & CSS Group~\cite{mitCSSGroup2026}
% & \textbf{形式化验证与底层系统安全：}依托形式化验证技术开发FSCQ安全文件系统、CryptDB加密数据库，从底层硬件存储层保障智能体数据读写安全~\cite{mitCSSGroup2026}。 \\
% \cline{2-4}
% & \multirow{3}{=}{斯坦福大学\\(Stan-\\ford)}
% & HAI \& CAIS~\cite{stanfordHAI2026}~\cite{stanfordCAIS2026} \allowbreak ~\cite{stanfordKochenderfer2026}
% & \textbf{提示注入防护与智能体资产风险分级：}搭建多层越狱防护、插件权限校验体系，提出高风险场景智能体合规溯源与资产分级评估方案；参与《Agents of Chaos》实验，证实自主智能体存在数据泄露、恶意系统操作、身份冒充等原生安全隐患~\cite{stanfordKochenderfer2026,agentsofchaos2026}。 \\
% \cline{3-4}
% & & SCS Group~\cite{stanfordSCS2026}
% & \textbf{智能体文件系统隔离沙箱：}研发jai轻量级Linux沙箱，依托目录权限管控、写时复制机制，抵御智能体误删、篡改本地文件的安全风险~\cite{stanfordJai2026}。 \\
% \cline{3-4}
% & & CRFM~\cite{stanfordCRFM2026}
% & \textbf{大模型风险与能力统一评测：}推出HELM标准化评测框架，实现多任务多维度模型能力与风险量化对比；提出FMTI指数，量化评估模型全生命周期研发、训练、部署透明度水平。 \\
% \cline{2-4}
% & \multirow{4}{=}{加州大学伯克利分校\\(UC Berkeley)}
% & BAIR~\cite{berkeleyBAIR2026}~\cite{pieterAbbeelHomepage2026}
% & \textbf{多智能体强化学习决策与函数调用评测体系：}研究机器人自适应协同强化学习策略，支撑具身智能体协同决策；构建Gorilla、BFCL开源基准，标准化量化智能体函数选择、参数生成能力~\cite{pieterAbbeelHomepage2026,bfcl2025}。 \\
% \cline{3-4}
% & & Sky Lab~\cite{berkeleySky2026}
% & \textbf{智能体工具调用运行时隔离：}从操作系统层优化智能体工具调度、任务可观测性，搭建轻量化运行时安全隔离平台~\cite{berkeleySky2026}。 \\
% \cline{3-4}
% & & CHAI~\cite{berkeleyCHAI2026}
% & \textbf{人类偏好对齐与智能体可控性：}基于人机协作理论设计意图对齐算法，提升自主智能体行为可解释性与人工干预可控能力~\cite{berkeleyCHAI2026}。 \\
% \cline{3-4}
% & & RDI~\cite{dawnSongBerkeley2026}
% & \textbf{工具调用权限管控、代码智能体攻防与攻击分类理论：}Progent权限代理系统将工具调用攻击成功率由39.9\%降至1.0\%；CyberGym平台自动挖掘34个软件零日漏洞，验证代码智能体原生攻击风险；提出上下文安全四维属性模型，重构智能体攻击风险完整分类体系~\cite{progent2025,cybergym2025,contextsecurity2026}。 \\
% \cline{2-4}
% & \multirow{3}{=}{卡内基梅隆大学\\(CMU)}
% % & LTI
% % & \textbf{Web场景智能体交互评测：}搭建WebArena、VisualWebArena真实网页交互环境，完成文本、多模态智能体交互式任务执行能力标准化评测。 \\
% % \cline{3-4}
% & CyLab~\cite{cmuCyLab2026}
% & \textbf{代码智能体安全对齐与自主链式网络攻击验证：}提出SecCodePRM过程奖励模型，抑制代码生成阶段高危漏洞产出；复现Equifax数据泄露完整漏洞链，证实LLM智能体可自主完成全流程网络攻击~\cite{seccodeprm2026,equifaxsim2025}。 \\
% \cline{3-4}
% & & Robotics Institute~\cite{cmuRI2026}
% & \textbf{多智能体博弈与任务分配：}基于博弈论设计分布式任务调度机制，规避多智能体协同过程中的策略冲突与资源争抢风险~\cite{cmuRI2026}。 \\
% \cline{3-4}
% & & AISIRT~\cite{cmuAISIRT2026}
% & \textbf{生成式AI安全事件响应：}校级专项研究团队，专职开展AI智能体漏洞研判、安全应急处置与风险闭环管控工作~\cite{cmuAISIRT2026}。 \\
% \cline{2-4}

% & \multirow{4}{=}{普渡\\大学\\(Pur-\\due)}
% & ASSET~\cite{purdueASSET2026}
% & \textbf{代码智能体推理能力评测：}提出CoRe静态分析基准，量化评测大模型识别、推理代码漏洞的风险判别能力~\cite{purdueASSET2026}。 \\*

% \nobreakcline{3-4}
% & & PurCL~\cite{astra2025}
% & \textbf{代码智能体攻防实战与隐性漏洞挖掘：}团队斩获Amazon Nova AI Challenge攻击赛道全球第一，验证自动化红队攻击能力；ASTRA系统攻击成功率超90\%，定向挖掘代码助手多步骤链式隐性漏洞~\cite{purcl2025nova,astra2025}。 \\*
% \nobreakcline{3-4}
% & & IDEAS~\cite{purdueIDEAS2026}
% & \textbf{拒止环境机器人协同导航：}获美军项目资助，研究GPS拒止场景下空地多机器人智能协同定位与安全导航技术~\cite{ideasdod2025}。 \\*
% \cline{3-4}
% & & CERIAS~\cite{purdueCERIAS2026}
% & \textbf{组织级智能体安全工程化：}将智能体安全管控嵌入企业漏洞管理、应急演练、合规培训全业务流程，形成标准化安全工程落地方案~\cite{purdueCERIAS2026}。 \\*
% \cline{2-4}
% & \multirow{4}{=}{哈佛\\大学\\(Har-\\vard)}
% & Kempner Institute~\cite{kakadeHomepage2026}
% & \textbf{自主智能体有害生成机制：}拆解大模型推理规划底层逻辑，统一揭示自主智能体主动生成有害内容的内在机理~\cite{kempnerHarmfulness2026}。 \\
% \cline{3-4}
% & & PT Project~\cite{privacytoolsharvard2026}
% & \textbf{差分隐私数据保护工具：}开源OpenDP轻量化差分隐私框架，为智能体跨主体数据交互提供标准化隐私防护能力~\cite{privacytoolsharvard2026}。 \\
% \cline{3-4}
% & & James Mickens团队~\cite{mickensharvard2026}
% & \textbf{Agentic AI安全治理战略：}系统梳理自主智能体全维度网络安全威胁体系，提出产学研多方协同长效安全治理框架~\cite{mickensharvard2026}。 \\
% \cline{3-4}
% & & Berkman Klein Center~\cite{berkmanklein2026}
% & \textbf{高风险行业智能体制度边界：}从社会技术交叉视角，明确金融、医疗等高风险领域智能体落地的合规红线与监管权责边界~\cite{berkmanklein2026}。 \\
% \cline{2-4}
% & \multirow{4}{=}{康奈尔大学\\(Cor-\\nell)}
% & Cornell AI Alignment实验室~\cite{cornellLevine2026}
% & \textbf{AI对齐基础理论研究：}依托严格数学理论搭建自主智能体价值对齐、行为硬约束底层原理完整理论框架~\cite{cornellLevine2026}。 \\
% \cline{3-4}
% & & Vitaly Shmatikov团队~\cite{shmatikovHomepage2026}
% & \textbf{检索型智能体投毒攻击与嵌入向量信息泄露：}提出WARP网页投毒方法，验证检索增强智能体易遭受外部内容篡改攻击；vec2vec算法证实跨模型嵌入向量可还原原始文本语义，暴露向量检索智能体隐私泄露风险~\cite{warp2026,vec2vec2025}。 \\
% \cline{3-4}
% & & Martin T. Wells团队~\cite{dempo2026}
% & \textbf{多用户价值对齐冲突消解：}设计DemPO优化算法，解决多主体交互场景下智能体训练存在的人群价值偏好偏差冲突问题~\cite{dempo2026}。 \\
% \cline{3-4}
% & & Thomas Ristenpart团队~\cite{ristenpartcornell2026}
% & \textbf{LLM路由系统安全风险：}深度剖析大模型流量调度、任务路由底层架构存在的原生安全漏洞与各类权限绕过攻击路径~\cite{reroutingllm2025}。 \\
% \cline{2-4}
% & \multirow{4}{=}{佐治亚理工学院\\(Georgia Tech.)}
% & SSS Lab~\cite{gatechSCP2026}
% & \textbf{软件系统漏洞挖掘实战：}面向智能体底层依赖库、运行时组件开展定向漏洞挖掘与攻防验证实验~\cite{gatechSCP2026}。 \\
% \cline{3-4}
% & & TR Lab~\cite{gatechSCP2026}
% & \textbf{自主机器人环境安全落地：}研究服务、工业机器人智能体在复杂人居场景下的安全决策、动态风险规避机制~\cite{gatechSCP2026}。 \\
% \cline{3-4}
% & & BEES Lab~\cite{gatechSCP2026}
% & \textbf{AI安全风险成因量化分析：}基于大规模全网流量测量，量化解析智能体安全事件现实诱因、传播路径与演化规律~\cite{gatechSCP2026}。 \\
% \cline{3-4}
% & & Team Atlanta~\cite{gatechAIxCC2026}
% & \textbf{LLM自主漏洞攻防系统：}Atlantis框架实现全自动无人工干预漏洞挖掘与修复，斩获DARPA AI Cyber Challenge半决赛冠军~\cite{gatechAIxCC2026}。 \\
% \cline{2-4}
% & \multirow{4}{=}[1.5\baselineskip]{伊利诺伊大学厄巴纳-香槟分校\\(UIUC)}
% & CSL~\cite{uiucCSL2026}
% & \textbf{可信AI智能体综合研究：}融合机器学习、底层系统安全、软件工程技术，搭建端到端全链路可信智能体技术体系~\cite{uiucCSL2026}。 \\
% \cline{3-4}
% & & ITI~\cite{uiucITI2026}
% & \textbf{关键基础设施智能体弹性防护：}面向智能电网等工业关键设施，研发防护型智能体，提升基础设施网络抵御高强度攻击的弹性能力~\cite{uiucITI2026}。 \\
% \cline{3-4}
% & & SL Lab~\cite{uiucBoLi2026}
% & \textbf{智能体对抗鲁棒性增强：}提出具备理论证明的鲁棒学习算法，抵御针对智能体视觉、文本感知层的对抗样本攻击~\cite{uiucBoLi2026}。 \\
% \cline{3-4}
% & & Heng Ji团队~\cite{uiucHengJi2026}
% & \textbf{代码化智能体动作安全表示：}CodeAct框架采用可执行代码替代纯文本动作指令，显著提升智能体工具调用成功率与执行安全边界可控性。 \\
% \cline{2-4}
% & \multirow{4}{=}[-1.0\baselineskip]{普林斯顿大学\\(Princ-\\eton)}
% & NLP团队~\cite{narasimhanHomepage2026}
% & \textbf{代码修复智能体评测与优化：}构建SWE-bench真实软件工程基准，设计SWE-agent专用交互接口，提升智能体自动修复代码漏洞的稳定性与成功率。 \\*
% \nobreakcline{3-4}

% & & CITP~\cite{princetonCITP2026}
% & \textbf{AI能力夸大风险辨析：}系统拆解智能体商业宣传中的能力夸大、虚假宣传问题，提供标准化风险甄别与核验方法论。 \\*
% \nobreakcline{3-4}

% & & Aleksandra Korolova团队~\cite{princetonKorolova2026}
% & \textbf{LLM决策系统公平性审计：}基于差分隐私理论范式，量化审计简历筛选、信用评估等落地场景下LLM智能体的算法公平性偏差~\cite{princetonKorolova2026}。 \\*
% \cline{3-4}

% & & S\&P Group~\cite{princetonSecPriv2026}
% & \textbf{智能体应用层安全防护：}结合现代密码学、网络攻防技术，构建面向业务场景智能体的上层应用完整安全防护体系~\cite{princetonSecPriv2026}。 \\*
% \cline{2-4}

% & \multirow{4}{=}{纽约\\大学\\(NYU)}
% & Muhammad Shafique团队~\cite{shafiqueHomepage2026}
% & \textbf{网安智能体自动化攻防、自主勒索攻击与多智能体越狱攻击：}EnIGMA网安智能体解题能力达到传统模型三倍，适配全品类CTF攻防场景；PromptLock验证LLM可独立完成完整勒索软件攻击链路；MetaCipher依靠多智能体协同高效绕过主流大模型安全防护护栏~\cite{promptlock2025,metacipher2025}。 \\
% \cline{3-4}
% & & AITP Lab~\cite{sucholutskyAITP2026}
% & \textbf{思维伙伴智能体人机交互：}研究共生式协同思维伙伴系统，优化人机联合决策的安全性、交互流畅度与风险可控性~\cite{sucholutskyAITP2026}。 \\
% \cline{3-4}
% & & AgentAuditor团队~\cite{agentauditor2025}
% & \textbf{通用智能体安全审计方法：}依托安全经验记忆库，使LLM智能体实现接近人类安全专家水平的全维度智能体风险评估能力~\cite{agentauditor2025}。 \\
% \cline{3-4}
% & & Quanyan Zhu团队~\cite{stackelbergDefense2025}
% & \textbf{攻防博弈智能体防御建模：}基于斯塔克尔伯格博弈理论量化攻防双方策略均衡，推导面向智能体的最优主动防御部署方案~\cite{stackelbergDefense2025}。 \\
% \hline
% \multirow{12}{=}{加拿大}
% & \multirow{4}{=}{多伦多大学\\(Toro-\\nto)}
% & SRI~\cite{schwartzReisman2026}
% & \textbf{智能体全周期治理框架与国家级AI安全指南：}搭建覆盖研发、部署、运维的智能体风险管控与责任追溯体系，规范企业级智能体落地流程；牵头编制的AI治理指南被加拿大政府纳入中小企业AI部署标准工具包~\cite{schwartzReisman2026,vectorResponsibleAI2025}。 \\
% \cline{3-4}
% & & CleverHans Lab~\cite{cleverhans2026}
% & \textbf{AI蠕虫自主传播风险与隐私协作学习框架：}实验证实开源大模型可生成自适应AI蠕虫，依托终端资源实现跨设备自主扩散；提出CaPC Learning多方协作学习框架，在联合训练过程中实现无损隐私保护~\cite{aiworm2026}。 \\
% \cline{3-4}
% & & TISL~\cite{tisl2026}
% & \textbf{具身智能体场景仿真：}优化机器人感知、环境表征算法，搭建高保真物理仿真环境用于具身智能体安全训练与风险验证~\cite{tisl2026}。 \\
% \cline{3-4}
% & & DORL~\cite{dorl2026}
% & \textbf{多智能体强化学习优化：}研究动态开放环境下多智能体协同优化算法，降低复杂任务协作中的冲突与安全失控风险~\cite{dorl2026}。 \\
% \cline{2-4}
% & \multirow{4}{=}{麦吉尔大学\\(McGill)}
% & McGill NLP Group~\cite{sivaReddy2026}
% & \textbf{Web智能体安全评测基准：}推出SafeArena评测基准，覆盖正常、恶意两类交互场景，标准化量化Web智能体抗误导、抗各类注入攻击能力~\cite{safearena2025,safearenaGithub2025}。 \\
% \cline{3-4}
% & & Mila~\cite{mila2026}
% & \textbf{全球前沿AI风险治理：}牵头编撰《国际AI安全报告》，主导高能力通用智能体全球协同治理规则体系研究~\cite{internationalAISafetyReport2025}。 \\
% \cline{3-4}
% & & IVADO R3AI~\cite{sivaReddy2026}
% & \textbf{长期AI安全专项资助：}设立专项科研经费，持续支持2024–2030年度智能体安全全方向课题攻关~\cite{sivaReddy2026}。 \\
% \cline{3-4}
% & & Jackie Cheung团队~\cite{mcgillCheung2026}
% & \textbf{文本生成智能体内容安全：}研究大模型文本生成阶段有害内容过滤、生成内容溯源与可追溯取证技术~\cite{mcgillCheung2026}。 \\
% \cline{2-4}
% & \multirow{4}{=}{英属哥伦比亚大学\\(UBC)}
% & S\&P Group~\cite{ubcSPG2026}
% & \textbf{部署型智能体对抗检测：}研发深度神经网络异常识别算法，实时在线识别智能体运行阶段遭遇的各类对抗样本攻击~\cite{ubcSPG2026}。 \\
% \cline{3-4}
% & & Mathias Lécuyer团队~\cite{ubcLecuyer2026}
% & \textbf{无重训练隐私审计与多步骤防御认证鲁棒性：}PANORAMIA框架无需模型微调即可快速完成智能体全链路隐私风险检测；自适应平滑算法为多层防御智能体提供可证明的安全鲁棒保障。 \\
% \cline{3-4}
% & & Victoria Lemieux团队~\cite{ubcLemieux2026}
% & \textbf{AI安全公共资源治理：}结合公地治理理论，构建兼顾原住民权益视角的AI数据资产与全域安全治理框架~\cite{ubcLemieux2026}。 \\
% \cline{3-4}
% & & CAISI~\cite{caisi2026}
% & \textbf{国家级AI安全认证体系：}联合加拿大多家顶尖科研机构，制定标准化智能体安全分级评估与合规认证完整规范~\cite{caisi2026}。 
% \end{longtable}
% \end{center}
% \endgroup

% \begin{figure}[htbp]
% \centering
% \includegraphics[width=0.48\textwidth]{figures/north_america_academic_share.png}
% \includegraphics[width=0.48\textwidth]{figures/north_america_direction_share.png}
% \caption{北美高校样本占比（\textcolor{red}{这个是指论文吗？基于第二章论文个数的}）与研究方向占比（\textcolor{red}{这个是指什么呢？}）}
% \end{figure}

% % 从公开资料看，北美研究也存在若干限制。首先，部分工作受产品能力牵引，更多关注“当前 Agent 能完成哪些任务”和“现有 Agent 如何被攻击”，对需求阶段的权限建模、组织流程设计、长期记忆退役和跨组织责任链讨论相对有限。其次，企业平台掌握大量真实运行日志和事故样本，但公开披露范围有限，学术界难以完整复现实验环境。因此，北美成果主要提供技术前沿和评测基准参照；高风险行业部署仍需结合治理框架、行业流程和组织责任设计。


% \begin{table}[!ht]
% \centering
% \renewcommand{\arraystretch}{0.7}
% \caption{北美智能体安全研究方向样本占比\textcolor{red}{这个是指什么呢？怎么算的？}）}
% \begin{tabular}{p{5cm}p{2cm}p{7cm}}
% \hline
% 研究方向 & 占比 & 说明 \\
% \hline
% Web/代码执行智能体评测 & 24\% & WebArena、SWE-bench、真实工作流评测较强 \\
% 工具调用与 MCP/插件安全 & 20\% & 工具描述、API 调用、沙箱和供应链风险突出 \\
% 提示注入与上下文安全 & 18\% & 间接提示注入、恶意网页、邮件和文档攻击 \\
% 模型透明度与治理评测 & 15\% & FMTI、HELM、系统卡和安全披露 \\
% 多智能体与具身智能体安全 & 13\% & 机器人、多智能体通信和协作失效 \\
% 隐私、身份与运行审计 & 10\% & 企业身份、日志、权限和数据泄露 \\
% \hline
% \end{tabular}
% \end{table}

\subsection{欧洲AI智能体安全学术研究现状}
欧洲是全球AI安全治理与基础理论研究核心策源地，依托《欧盟AI法案》顶层规制、国家级科研院所与头部高校，形成法规驱动、重理论、重隐私与合规导向的智能体安全研究体系，与北美偏重实战攻防、工程落地的路线形成明显区分。欧洲AI智能体研究主要围绕合规约束、可证明安全、隐私防护展开，覆盖软件智能体、具身机器人、信息物理系统、去中心化多智能体等应用场景\cite{mpiEUAIAct2026}。
% 欧洲科研力量呈现分层协同布局：英国牛津、剑桥、谷歌 DeepMind 与艾伦·图灵研究所主攻AGI安全与治理\cite{deepmindControlRoadmap2026,turingGal2026}；德国CISPA、马克斯·普朗克研究所深耕底层攻击机理与机器学习安全\cite{cispaFritz2026,mpiAbdelnabi2026}；瑞士ETH Zurich、EPFL专注安全评测与系统防护\cite{ethSPYLab2026,payerHomepage2026}；法国、荷兰、比利时机构集中攻关形式化验证、隐私技术与法规工程转化\cite{inriaWhy3ProofInUse2023,tudelftErkin2026,cosicKuLeuven2026}。
% 整体形成“攻击机理挖掘—形式化安全证明—合规落地—跨域风险管控”完整研究链路，在间接提示注入、时序后门、可证明安全控制、法规适配等方向具备全球领先成果\cite{cispaIndirectPI2023,mpiTimeBombs2026,tumSafeLLMRobots2026}。
本节梳理主要欧洲国家（包括英国、德国、瑞士、法国、荷兰与比利时）顶尖高校、国家级研究院及企业实验室的研究团队、研究方向与标志性成果，核心发现如下（详见表~\ref{tab:europe_top_institution}）：

\begin{itemize}[tabitem]
\item \textbf{跨国分工布局清晰，技术与合规深度落地。} 英国依托牛津大学、剑桥大学主攻AGI对齐、人机权责划分与长期风险监管\cite{wooldridgeHomepage2026,oxfordAIGI2026,cambridgeCFI2026}；德国聚焦底层攻击机理与硬件侧信道安全\cite{cispaFritz2026,tumSigl2026}；瑞士专攻安全评测基准与安全强化学习\cite{as2025actsafe}；荷兰研究隐私重构攻击\cite{tudelftErkin2026}、比利时KU Leuven构建适配GDPR的安全成熟度模型。整体依托《欧盟AI法案》与GDPR，将学术成果转化为可量化合规指标、商用智能体披露标准与跨境算力管控方案，实现科研、立法、产业审计一体化闭环\cite{mpiEUAIAct2026,citipDataProtectionByDesign2025}。

\item \textbf{基础威胁机理领先，提出多项原生攻击范式。} 德国CISPA完整定义并分类间接提示注入攻击，厘清不可信外部数据带来的指令混淆风险\cite{cispaAgentSafety2026}；慕尼黑工大验证软提示嵌入逃逸与评测工具随机性缺陷，揭露现有红队评估体系的可靠性漏洞\cite{tumCoinFlipSafety2026}；帝国理工学院提出Checkpoint-GCG，证明微调防御存在固有结构性漏洞，推翻单一模型即可完成安全防护的传统结论\cite{checkpointGCG2025}。此外，马克斯·普朗克研究所提出时序记忆炸弹后门，挖掘智能体长时序交互下的隐性攻击机理，完善智能体动态威胁体系\cite{mpiTimeBombs2026}。

\item \textbf{安全体系完善，动态评测工具谱系齐全。} 软件层由牛津、剑桥基于CHERI隔离与Why3工具实现智能体模块隔离与演绎证明\cite{cambridgeCHERI2026}；物理层依托慕尼黑工大可达性分析、ETH Zurich ActSafe框架界定LLM机器人安全边界，实现物理执行层可证明安全保障\cite{tumSafeLLMRobots2026,as2025actsafe}；群体层完成多机协同碰撞约束与人机轨迹风险量化，完善多智能体协同安全验证体系\cite{gao2025provablySafe}；评测层面主打动态对抗检验，依托ETH AgentDojo、BaxBench、Foolbox等标杆工具，覆盖通用、代码、视觉、法律与长时序规划细分场景，构建完备的智能体安全评测基准矩阵\cite{baxbench2026,mpiLegalAgentBenchmark2026}。
\end{itemize}


\begingroup

% 允许 longtblr 在跨行单元格内部自动分页
% 对新版 tabularray 生效；旧版不存在该变量时不会报错
% \ExplSyntaxOn
% \cs_if_exist:NT \lTblrCellBreakBool
%   {
%     \bool_set_true:N \lTblrCellBreakBool
%   }
% \ExplSyntaxOff

\begin{center}
% \nopagebreak[4]
%\enlargethispage{1cm}
\footnotesize
\setlength{\tabcolsep}{3pt}
% \renewcommand{\arraystretch}{0.7}
\begin{longtblr}[
caption = {欧洲顶尖高校AI智能体安全学术研究汇总},
label   = {tab:europe_top_institution},
]
{
colspec = {|m{0.8cm}|m{1.2cm}|m{2.5cm}|m{\dimexpr\linewidth - 0.8cm - 1.2cm - 2.5cm - 10\tabcolsep\relax}|},
rowhead = 1,
}
\hline
\textbf{国家} & \textbf{高校} & \textbf{实验室/研究团队} & \textbf{核心智能体安全研究方向} \\
\hline
\SetCell[r=1]{c} 英国
& \SetCell[r=4]{c} 牛津\newline 大学(Oxford)
& MAS~\cite{wooldridgeHomepage2026}
& \textbf{多智能体系统理论与LLM智能体失效模式：}Michael Wooldridge团队将经典多智能体系统理论应用于LLM智能体失效模式、对齐挑战与责任行为研究，2026年主办LaMAS 2026研讨会聚焦该核心议题~\cite{oxfordLaMAS2026}。 \\
\cline{3-4}
& & Institute for Ethics in AI~\cite{oxfordHarcourt2026}
& \textbf{AI伦理治理与多元价值冲突：}Edward Harcourt团队从伦理、责任、治理与社会影响维度研究AI部署风险，通过《Value Kaleidoscope》探究AI系统处理人类多元价值观、权利与义务冲突的机制。 \\
\cline{3-4}
& & AI Governance Initiative~\cite{oxfordAIGI2026}
& \textbf{国际AI治理与算力监管协调：}Robert Trager团队研究跨国AI治理体系设计，聚焦算力治理与跨境监管协调，2024年发布白皮书提出算力提供方跨境报告机制，防范监管套利。 \\
\cline{3-4}
& & Automated Verification~\cite{oxfordKwiatkowska2026}
& \textbf{深度神经网络形式化安全验证：}Marta Kwiatkowska团队研究DNN形式化安全验证与鲁棒性证明，基于博弈论近似技术为智能体决策核心组件提供可证明安全保障。 \\
\cline{2-4}
& \SetCell[r=4]{c} 剑桥\newline 大学\newline (Cam-\newline bridge)
& CFI~\cite{cambridgeCFI2026,cambridgeCave2026}
& \textbf{智能体透明度评估与安全披露审计：}Stephen Cave团队构建AI Agent Index，结构化评估30款主流智能体的能力、生态交互与安全披露情况，指出当前智能体普遍缺少系统安全测试与第三方评估材料~\cite{cambridgeAgentIndex2026}。 \\
\cline{3-4}
& & Security Group~\cite{cambridgeSecurityGroup2026}
& \textbf{间接提示注入防御与权限约束架构：}提出规划与检索模块分离的安全架构，抵御智能体间接提示注入攻击；基于CHERI能力硬件隔离体系，为智能体工具调用、文件访问提供底层权限安全支撑~\cite{cambridgeFoersterAgentSecurity2026,cambridgeCHERI2026}。 \\
\cline{3-4}
& & Prorok Lab~\cite{cambridgeProrokLab2026}
& \textbf{多机器人可证明安全导航与协同：}融合控制屏障函数、在线环境估计与去中心化规划，提出未知环境下多智能体安全导航方法，为机器人群体碰撞规避、任务可达性提供形式化保证~\cite{gao2025provablySafe}。 \\
\cline{3-4}
& & BIRL~\cite{cambridgeBIRL2026}
& \textbf{具身智能体物理执行层安全边界：}从软体机器人、触觉感知与人机交互切入，将智能体安全研究从软件规划层延伸至物理执行层的接触安全、动作边界与环境适配领域~\cite{birlPublications2026}。 \\
\cline{2-4}
& \SetCell[r=4]{c} 帝国理工学院\newline (ICL)
& APSS~\cite{apssLab2026}
& \textbf{神经符号辩论与合规性可解释判断：}依托NeSyDebates项目，基于神经符号辩论机制实现大模型政策法规合规性的可解释、可争议、可执行安全判定，覆盖生成式AI与物联网系统安全场景~\cite{apssNeSyDebates2026}。 \\
\cline{3-4}
& & AISP~\cite{imperialAISPLab2026}
& \textbf{提示注入防御审计与成员推断风险：}提出Checkpoint-GCG审计微调类提示注入防御机制，证实模型侧防御无法根除外部数据引发的间接提示注入风险；基于LLM成员推断技术构建企业知识库智能体数据泄露评估体系~\cite{checkpointGCG2025,llmMembershipInference2025}。 \\
\cline{3-4}
& & SAIL~\cite{imperialSAIL2026}
& \textbf{神经网络形式化验证与安全强化学习：}研发Venus验证工具，检测神经网络可达性与对抗鲁棒性，为自动驾驶、群体智能体、强化学习系统提供形式化安全保障~\cite{venusVerifier2026}。 \\
\cline{3-4}
& & AIRL~\cite{imperialAIRL2026}
& \textbf{机器人损伤恢复与自主适应能力：}提出基于Quality-Diversity的智能试错算法，使机器人在未知损伤、环境扰动下快速生成补偿性行为，提升开放环境下具身智能体的自主韧性。 \\
\cline{2-4}
& \SetCell[r=4]{c} 伦敦大学学院\newline (UCL)
& S2Lab~\cite{uclS2Lab2026}
& \textbf{机器学习安全评测方法论陷阱梳理：}系统总结恶意软件检测、漏洞挖掘、二进制分析等安全任务中机器学习评测的常见误区，构建可信评测规范体系。 \\
\cline{3-4}
& & ISR Group~\cite{uclISec2026}
& \textbf{智能体自主漏洞利用与多智能体渗透测试：}构建大模型智能体合约漏洞生成框架，实现智能合约自主分析与可执行攻击；研发MAPTA多智能体Web渗透测试系统，支撑智能体攻防能力量化评估~\cite{gervaisA1Exploit2025,uclGoogleFunding2026}。 \\
\cline{3-4}
& & CAI~\cite{uclAICentre2026}
& \textbf{负责任生成模型开发与信息泄露防护：}牵头英国生成式AI Hub，构建生成模型负责任开发与评估工具体系；提出Aspective Agentic AI框架，通过动态信息视域约束抑制多智能体系统信息泄露风险~\cite{uclGenAIHub2026,aspectiveAgenticAI2025}。 \\
\cline{3-4}
& & UCLIC~\cite{uclIC2026}
& \textbf{人本可用安全与隐私伤害辅助报告：}研究对话式AI辅助用户上报隐私伤害的机制，其人本安全研究成果落地英国政府与行业安全规范~\cite{uclHumanCentredSecurity2026}。 \\
\hline
& \SetCell[r=4]{c} 爱丁堡大学\newline (Edin-\newline burgh)
& AIAI~\cite{edinburghAIAI2026}
& \textbf{长时序智能体规划与协作任务管理：}基于O-Plan/I-X系列规划技术，整合任务表示、动态重规划与跨主体协作机制，奠定长时序智能体任务调度与协同管理的系统基础~\cite{tateOPlan1996}。 \\
\cline{3-4}
& & CAIAA~\cite{edinburghAssistiveAutonomy2026}
& \textbf{人机协作环境下安全鲁棒自主决策：}融合自动驾驶、机器人学习与人机对话技术，构建辅助自主平台，研究不确定性场景下人机协同智能体的安全决策机制。 \\
\cline{3-4}
& & SPT Group~\cite{edinburghSPT2026}
& \textbf{可验证安全证据与信息流控制：}基于携带证明代码与形式化协议验证技术，为智能体外部代码执行、插件调用、分布式服务访问提供可验证安全证据与信息流管控方案~\cite{aspinallMRG2005}。 \\
\cline{3-4}
& & CTF~\cite{edinburghCTF2026}
& \textbf{负责任AI生态治理分歧梳理：}依托BRAID项目联动学界、产业与监管机构，系统厘清AI安全、伦理、透明度、问责与治理体系的重叠与冲突边界~\cite{braidProgramme2026,gyevnar2026responsibleDivides}。 \\
\cline{2-4}
% & \SetCell[r=4]{c} Google DeepMind
% & AGI Safety\newline and Alignment~\cite{deepmindRohinShah2026}
% & \textbf{AGI技术安全与智能体行为监测护栏：}发布GDM AI Control Roadmap v0.1，借鉴网络安全内部威胁防护思路，构建智能体工具访问控制与实时行为监测安全护栏体系~\cite{deepmindControlRoadmap2026}。 \\
% \cline{3-4}
% & & Frontier Safety\newline and Governance~\cite{deepmindDafoe2026}
% & \textbf{前沿模型危险能力评估与多智能体风险基金：}开展前沿模型危害能力定级，设立千万级专项基金，研究大规模多智能体交互中的诈骗、提示注入与网络攻击风险~\cite{deepmindMultiAgentFund2026}。 \\
% \cline{3-4}
% & & Mechanistic\newline Interpretability~\cite{deepmindNeelNanda2026}
% & \textbf{机械可解释性与隐藏目标识别：}通过神经网络逆向工程识别模型隐藏目标与欺骗倾向，开源Gemma Scope稀疏自编码器工具集支撑安全社区可解释性研究~\cite{deepmindGemmaScope2024}。 \\
% \cline{3-4}
% & & 具身智能体\newline 研究方向~\cite{deepmindRoboCat2023}
% & \textbf{具身智能体跨任务迁移与动作限制：}基于RoboCat与SIMA模型验证语言指令与物理动作的安全鸿沟，提出动作边界约束与人工接管机制防控具身智能体执行风险~\cite{deepmindRoboCat2023}。 \\
% \cline{2-4}
% & \SetCell[r=4]{c} Alan Turing Institute
% & Safe and\newline Ethical AI~\cite{turingGal2026}
% & \textbf{大规模LLM数据投毒后门研究：}完成迄今最大规模LLM数据投毒实验，证实仅250份恶意文档即可在任意规模大模型中植入持久后门，为训练数据安全管控提供依据~\cite{turingPoisoning2025}。 \\
% \cline{3-4}
% & & Defence and National\newline Security / DARe~\cite{turingLASR2026}
% & \textbf{受限环境多机器人韧性与对抗性威胁研究：}依托820万英镑LASR实验室，联合GCHQ等机构研究通信脆弱、资源受限场景下边缘AI与多机器人系统的对抗性风险与韧性增强技术~\cite{turingLASR2026}。 \\
% \cline{3-4}
% & & Public Policy\newline Programme~\cite{turingBright2026}
% & \textbf{在线安全与AI公共政策立法咨询：}为英欧政府提供算法问责、在线有害内容检测领域技术支撑，推动AI安全相关立法落地~\cite{turingBright2026}。 \\
% \cline{3-4}
% & & Human-AI Interfaces\newline and Robotics~\cite{turingVijayakumar2026}
% & \textbf{人机协作责任边界与控制权移交：}研究人机协同机器人的技能共享控制机制，界定自主系统与人类操作员之间的安全责任边界与动态控制权切换规则~\cite{turingVijayakumar2026}。 \\
\hline
\SetCell[r=1]{c} 德国
% & \SetCell[r=4]{c} CISPA Helmholtz Center
% & Trustworthy ML\newline and Cybersecurity~\cite{cispaFritz2026}
% & \textbf{间接提示注入攻击识别与智能体部署安全：}2023年首次系统定义并分类间接提示注入攻击，厘清LLM应用数据与指令边界漏洞；2026年提出开放式智能体部署前安全优先的核心准则~\cite{cispaIndirectPI2023,cispaAgentSafety2026}。 \\
% \cline{3-4}
% & & Systems Security~\cite{cispaHolz2026}
% & \textbf{系统层安全与固件模糊测试：}提出新型固件模糊测试方案，解决中断处理导致的测试覆盖率偏低难题，提升嵌入式设备底层漏洞挖掘效率~\cite{cispaFirmwareFuzzing2025}。 \\
% \cline{3-4}
% & & Web Security~\cite{cispaStock2026}
% & \textbf{第三方服务监控盲区测量：}提出LeakyLinks框架，量化分析URL扫描服务在浏览器生态中形成的隐私监控盲区，为Web端智能体隐私防护提供依据~\cite{cispaLeakyLinks2026}。 \\
% \cline{3-4}
% & & Trustworthy ML\newline (LLM Safety)~\cite{cispaYangZhang2026}
% & \textbf{仅标签成员推理攻击与训练数据泄露：}实现低查询量增强型仅标签成员推理攻击，揭示LLM训练数据泄露的新型攻击路径，适配轻量化智能体安全审计场景~\cite{cispaMembershipInference2025}。 \\
% \cline{2-4}
& \SetCell[r=4]{c} 慕尼黑工业大学(TUM)
& CPS~\cite{tumAlthoff2026}
& \textbf{可达性分析与LLM控制机器人安全保证：}基于可达性分析理论，构建LLM驱动机器人的形式化验证框架，提供数学可证明的物理执行安全保障~\cite{tumSafeLLMRobots2026}。 \\
\cline{3-4}
& & DAML~\cite{tumGuennemann2026}
& \textbf{嵌入空间对齐绕过与红队评估鲁棒性：}证实嵌入空间软提示攻击可绕过LLM安全对齐与遗忘机制；证明LLM-as-a-Judge在红队测试中评估鲁棒性接近随机猜测，质疑现有评估范式可靠性~\cite{tumCoinFlipSafety2026}。 \\
\cline{3-4}
& & SPRLS ~\cite{tumSchoellig2026}
& \textbf{学习型控制系统安全性能基准平台：}开发safe-control-gym开源基准，实现传统控制、强化学习等智能体控制系统安全性能的定量对比与标准化评测。 \\
\cline{3-4}
& & SIT ~\cite{tumSigl2026}
& \textbf{硬件侧信道攻击与模型参数窃取：}验证物理访问场景下，侧信道分析可窃取神经网络加速器的核心参数，为边缘部署智能体的硬件安全防护提供依据~\cite{tumSideChannelNN2026}。 \\
\cline{2-4}
% & \SetCell[r=4]{c} 马克斯·普朗克智能系统研究所(MPI-IS)
% & AI Safety\newline and Alignment~\cite{mpiAndriushchenko2026}
% & \textbf{自主LLM智能体对齐与能力评估：}聚焦自主式大模型智能体价值对齐研究，构建严格的前沿模型能力分级评估体系，推进AI风险公众科普与认知引导~\cite{mpiAndriushchenko2026}。 \\
% \cline{3-4}
% & & COMPASS~\cite{mpiAbdelnabi2026}
% & \textbf{隐性记忆隐藏通道与时序后门：}提出时序性"记忆炸弹"后门机制，仅在特定交互序列下激活恶意行为，揭示跨智能体隐蔽通信与基准数据集污染风险~\cite{mpiTimeBombs2026}。 \\
% \cline{3-4}
% & & Empirical\newline Inference~\cite{mpiScholkopf2026}
% & \textbf{因果学习理论与AI监管科学支撑：}深耕因果学习基础理论，团队多人入选联合国、欧盟AI法案专家小组，为国际AI监管提供科学理论支撑~\cite{mpiEUAIAct2026}。 \\
% \cline{3-4}
% & & Social Foundations\newline of Computation~\cite{mpiHardt2026}
% & \textbf{法律智能体真实任务表现评测：}参与构建Legal Agent Benchmark基准，聚焦真实法律场景量化评测专业领域智能体的任务执行能力与安全合规性~\cite{mpiLegalAgentBenchmark2026}。 \\
% \cline{2-4}
& \SetCell[r=2]{c} 图宾根大学
& Bethge Lab~\cite{tuebingenBethge2026}
& \textbf{对抗鲁棒性评测工具与基准建设：}开发Foolbox开源工具包，标准化对抗样本生成与模型鲁棒性量化流程，成为视觉智能体安全评测的通用基准工具。 \\
\cline{3-4}
& & AVG ~\cite{tuebingenGeiger2026}
& \textbf{自动驾驶3D场景理解基准数据集：}维护KITTI系列权威自动驾驶数据集，提出TransFuser端到端感知模型，成为具身驾驶智能体场景理解的基线标准。 \\
\hline
\SetCell[r=1]{c} 瑞士
& \SetCell[r=4]{c} 苏黎世联邦理工学院(ETH Zurich)
& SPY Lab~\cite{ethSPYLab2026,tramerHomepage2026}
& \textbf{智能体提示注入攻防动态评测环境：}构建AgentDojo评测平台，集成97项真实任务与629条安全用例，统一量化评估LLM智能体工具调用、外部数据读取中的间接提示注入攻防性能~\cite{agentdojoGithub2026}。 \\
\cline{3-4}
& & SRI Lab~\cite{ethSRILab2026,ethVechev2026}
& \textbf{代码智能体安全关键任务生成评测：}提出BaxBench基准，通过392项高危后端编程任务评测代码智能体生成内容的正确性与安全性，证实顶尖模型仍存在大量高危漏洞~\cite{baxbench2026}。 \\
\cline{3-4}
& & LAS Group~\cite{ethLAS2026,ethKrause2026}
& \textbf{安全强化学习中的主动探索约束：}提出ActSafe算法，融合主动学习与硬安全约束，管控强化学习智能体在未知环境探索中的风险暴露范围~\cite{as2025actsafe}。 \\
\cline{3-4}
& & RSL~\cite{ethRSL2026,ethHutter2026}
& \textbf{腿式机器人复杂地形自主导航安全：}基于强化学习与机载感知，实现ANYmal系列机器人复杂地形稳定移动，将智能体安全体系延伸至物理设备运动边界与环境适配场景~\cite{ethAnyMal2026}。 \\
\cline{2-4}
& \SetCell[r=4]{c} 洛桑联邦理工学院(EPFL)
& HexHive~\cite{payerHomepage2026}
& \textbf{固件移植与边缘设备漏洞检测：}提出EmbedFuzz固件移植技术，实现微控制器固件高性能原生模糊测试，支撑边缘智能体插件与控制模块的底层漏洞挖掘。 \\
\cline{3-4}
& & SPRING Lab~\cite{springLab2026}
& \textbf{去中心化隐私保护与敏感数据风险通知：}研发DP3T去中心化追踪系统，以最小数据收集原则实现隐私风险通知，为医疗、政务敏感数据处理智能体提供隐私设计范式。 \\
\cline{3-4}
& & DEDIS Lab~\cite{epflDEDIS2026}
& \textbf{去中心化信任基础设施与多机构协同审计：}构建Cothority/OmniLedger分布式信任架构，分散信任节点权限，支撑跨机构多智能体协同的身份认证与不可篡改审计追踪~\cite{epflOmniLedger2026}。 \\
\cline{3-4}
& & VITA Lab~\cite{epflVITA2026}
& \textbf{人机共存场景轨迹预测与碰撞风险规避：}提出Social-Transmotion视觉轨迹预测方法，提升公共空间移动智能体对行人意图与碰撞风险的预判能力，优化人机共存安全决策。 \\
\hline
% \SetCell[r=4]{c} 法国
% & \SetCell[r=4]{c} Inria
% & Prosecco~\cite{inriaProsecco2026}
% & \textbf{密码协议形式化验证与后量子安全分析：}基于ProVerif工具完成Signal后量子密钥协商协议PQXDH的自动化形式化验证，筑牢智能体加密通信底层安全基础~\cite{inriaPQXDH2024}。 \\
% \cline{3-4}
% & & PESTO~\cite{inriaPesto2026}
% & \textbf{隐私性质符号化验证与协议攻击发现：}实现安全协议隐私属性的符号化自动验证，发现新型哈希漏洞攻击模式，获USENIX Security 2023杰出论文奖~\cite{inriaHashGoneBad2023}。 \\
% \cline{3-4}
% & & Privatics~\cite{inriaCastelluccia2026}
% & \textbf{数据隐私保护与算法认知安全威胁：}研究大规模数据隐私脱敏技术，重点分析AI智能体操纵人类决策的认知层安全威胁与防御路径~\cite{inriaCastelluccia2026}。 \\
% \cline{3-4}
% & & Toccata~\cite{inriaToccata2026}
% & \textbf{演绎验证平台与关键系统安全工具链：}研发Why3演绎验证框架，集成至航空级SPARK/Ada工具链，支撑高可信智能体核心控制软件的形式化证明~\cite{inriaWhy3ProofInUse2023}。 \\
% \hline
\SetCell[r=1]{c} 荷兰
& \SetCell[r=4]{c} 代尔夫特理工大学(TU Delft)
& II Group~\cite{tudelftJonker2026}
& \textbf{智能体协商与人机混合智能系统：}深耕自主智能体自动协商机制与人机协同认知建模，获2024年ACM/SIGAI自主智能体领域终身成就奖~\cite{tudelftJonker2026}。 \\
\cline{3-4}
& & CS Group~\cite{tudelftErkin2026}
& \textbf{隐私求和协议的图结构重构攻击：}发现可证明安全的隐私求和协议在并行执行时存在图结构泄露漏洞，为隐私保护型多智能体集群部署提供风险警示。 \\
\cline{3-4}
& & AMR Lab~\cite{tudelftAlonsoMora2026}
& \textbf{LLM融合导航与有界概率运动规划：}融合大语言模型与模型预测控制实现机器人个性化安全导航；提出带碰撞概率上界的场景化规划方法，保障机器人群体运动安全~\cite{tudelftHeyRobot2026,tudelftScenarioPlanning2025}。 \\
\cline{3-4}
& & GC~\cite{tudelftVanEeten2026}
& \textbf{关键基础设施网络安全治理测量：}基于全网实测数据量化分析ISP与终端设备对Mirai僵尸网络的清除效果，支撑物联网智能体集群治理策略优化。 \\
\cline{2-4}
& \SetCell[r=2]{c} 阿姆斯特丹大学(UvA)
& VIS/QUVA Lab~\cite{uvaSnoek2026}
& \textbf{具身多智能体协作规划优化：}提出CaPo协同规划算法，优化具身场景下多智能体的任务分配与动作协同效率，降低协作冲突风险~\cite{uvaCaPo2025}。 \\
\cline{3-4}
& & IRLab~\cite{uvaDeRijke2026}
& \textbf{检索增强智能体信息可靠性理论：}构建检索增强生成（RAG）智能体的信息可信度评估理论体系，解决外部检索内容带来的事实性与安全风险问题~\cite{uvaDeRijke2026}。 \\
\hline
\SetCell[r=1]{c} 比利时
& \SetCell[r=4]{c} 鲁汶大学(KU Leuven)
& COSIC~\cite{cosicKuLeuven2026}
& \textbf{密码学与敏感数据访问技术控制：}基于哈希函数、隐私增强技术构建智能体敏感数据访问的底层密码管控机制，支撑合规化数据权限治理~\cite{cosicKuLeuven2026}。 \\
\cline{3-4}
& & DistriNet（软件安全）~\cite{distrinetJoosen2026}
& \textbf{LLM聊天机器人冒充与网络民主威胁防御：}系统研究AI聊天机器人冒充人类干预公共讨论的攻击路径，构建针对性检测与防御体系~\cite{distrinetLLMChatbot2026}。 \\
\cline{3-4}
& & DistriNet（AI可信性）~\cite{distrinetRimmer2026}
& \textbf{入侵检测场景下数据驱动AI可信度：}分析无监督数据环境下AI模型在恶意代码检测中的可信度缺陷，提出风险感知型智能体检测方案~\cite{distrinetRimmer2026}。 \\
\cline{3-4}
& & CiTiP~\cite{citipValcke2026}
& \textbf{GDPR合规转化与陪伴型智能体数据保护：}将GDPR合规条款转化为可量化的网络安全成熟度模型；针对性解析陪伴型聊天机器人的设计阶段数据保护合规要点~\cite{citipDataProtectionByDesign2025}。 \\
\hline
\end{longtblr}
\end{center}
\endgroup



% \begin{center}
% % \nopagebreak[4]
% \enlargethispage{1cm}
% \footnotesize
% \setlength{\tabcolsep}{3pt}
% % \renewcommand{\arraystretch}{0.7}
% \begin{longtable}{|>{\rule{0pt}{2pt}}m{0.8cm}
% |>{\rule{0pt}{2pt}}m{1.2cm}
% |>{\rule{0pt}{2pt}}m{3cm}
% |>{\rule{0pt}{2pt}}m{\dimexpr\linewidth - 0.8cm - 1.2cm - 3cm - 8\tabcolsep\relax}|}

% \caption{欧洲顶尖高校AI智能体安全学术研究汇总}\label{tab:europe_top_institution}\\
% \hline
% \textbf{国家} & \textbf{高校} & \textbf{实验室/研究团队} & \textbf{核心智能体安全研究方向} \\
% \hline
% \endfirsthead
% \multicolumn{4}{l}{\tablename\ \thetable\ -- 续上页}\\
% \hline
% \textbf{国家} & \textbf{高校} & \textbf{实验室/研究团队} & \textbf{核心智能体安全研究方向} \\
% \hline
% \endhead
% \hline
% \multicolumn{4}{r}{续下页}\\
% \endfoot
% \hline
% \endlastfoot

% \multirow{28}{=}{英国}
% & \multirow{4}{=}{牛津\\大学(Oxford)}
% & MAS~\cite{wooldridgeHomepage2026}
% & \textbf{多智能体系统理论与LLM智能体失效模式：}Michael Wooldridge团队将经典多智能体系统理论应用于LLM智能体失效模式、对齐挑战与责任行为研究，2026年主办LaMAS 2026研讨会聚焦该核心议题~\cite{oxfordLaMAS2026}。 \\
% \cline{3-4}
% & & Institute for Ethics in AI~\cite{oxfordHarcourt2026}
% & \textbf{AI伦理治理与多元价值冲突：}Edward Harcourt团队从伦理、责任、治理与社会影响维度研究AI部署风险，通过《Value Kaleidoscope》探究AI系统处理人类多元价值观、权利与义务冲突的机制。 \\
% \cline{3-4}
% & & AI Governance Initiative~\cite{oxfordAIGI2026}
% & \textbf{国际AI治理与算力监管协调：}Robert Trager团队研究跨国AI治理体系设计，聚焦算力治理与跨境监管协调，2024年发布白皮书提出算力提供方跨境报告机制，防范监管套利。 \\
% \cline{3-4}
% & & Automated Verification~\cite{oxfordKwiatkowska2026}
% & \textbf{深度神经网络形式化安全验证：}Marta Kwiatkowska团队研究DNN形式化安全验证与鲁棒性证明，基于博弈论近似技术为智能体决策核心组件提供可证明安全保障。 \\
% \cline{2-4}

% & \multirow{4}{=}{剑桥\\大学\\(Cam-\\bridge)}
% & CFI~\cite{cambridgeCFI2026,cambridgeCave2026}
% & \textbf{智能体透明度评估与安全披露审计：}Stephen Cave团队构建AI Agent Index，结构化评估30款主流智能体的能力、生态交互与安全披露情况，指出当前智能体普遍缺少系统安全测试与第三方评估材料~\cite{cambridgeAgentIndex2026}。 \\
% \cline{3-4}
% & & Security Group~\cite{cambridgeSecurityGroup2026}
% & \textbf{间接提示注入防御与权限约束架构：}提出规划与检索模块分离的安全架构，抵御智能体间接提示注入攻击；基于CHERI能力硬件隔离体系，为智能体工具调用、文件访问提供底层权限安全支撑~\cite{cambridgeFoersterAgentSecurity2026,cambridgeCHERI2026}。 \\
% \cline{3-4}
% & & Prorok Lab~\cite{cambridgeProrokLab2026}
% & \textbf{多机器人可证明安全导航与协同：}融合控制屏障函数、在线环境估计与去中心化规划，提出未知环境下多智能体安全导航方法，为机器人群体碰撞规避、任务可达性提供形式化保证~\cite{gao2025provablySafe}。 \\
% \cline{3-4}
% & & BIRL~\cite{cambridgeBIRL2026}
% & \textbf{具身智能体物理执行层安全边界：}从软体机器人、触觉感知与人机交互切入，将智能体安全研究从软件规划层延伸至物理执行层的接触安全、动作边界与环境适配领域~\cite{birlPublications2026}。 \\
% \cline{2-4}

% & \multirow{4}{=}{帝国理工学院\\(ICL)}
% & APSS~\cite{apssLab2026}
% & \textbf{神经符号辩论与合规性可解释判断：}依托NeSyDebates项目，基于神经符号辩论机制实现大模型政策法规合规性的可解释、可争议、可执行安全判定，覆盖生成式AI与物联网系统安全场景~\cite{apssNeSyDebates2026}。 \\
% \cline{3-4}
% & & AISP~\cite{imperialAISPLab2026}
% & \textbf{提示注入防御审计与成员推断风险：}提出Checkpoint-GCG审计微调类提示注入防御机制，证实模型侧防御无法根除外部数据引发的间接提示注入风险；基于LLM成员推断技术构建企业知识库智能体数据泄露评估体系~\cite{checkpointGCG2025,llmMembershipInference2025}。 \\
% \cline{3-4}
% & & SAIL~\cite{imperialSAIL2026}
% & \textbf{神经网络形式化验证与安全强化学习：}研发Venus验证工具，检测神经网络可达性与对抗鲁棒性，为自动驾驶、群体智能体、强化学习系统提供形式化安全保障~\cite{venusVerifier2026}。 \\
% \cline{3-4}
% & & AIRL~\cite{imperialAIRL2026}
% & \textbf{机器人损伤恢复与自主适应能力：}提出基于Quality-Diversity的智能试错算法，使机器人在未知损伤、环境扰动下快速生成补偿性行为，提升开放环境下具身智能体的自主韧性。 \\
% \cline{2-4}

% & \multirow{4}{=}{伦敦大学学院\\(UCL)}
% & S2Lab~\cite{uclS2Lab2026}
% & \textbf{机器学习安全评测方法论陷阱梳理：}系统总结恶意软件检测、漏洞挖掘、二进制分析等安全任务中机器学习评测的常见误区，构建可信评测规范体系。 \\
% \cline{3-4}
% & & ISR Group~\cite{uclISec2026}
% & \textbf{智能体自主漏洞利用与多智能体渗透测试：}构建大模型智能体合约漏洞生成框架，实现智能合约自主分析与可执行攻击；研发MAPTA多智能体Web渗透测试系统，支撑智能体攻防能力量化评估~\cite{gervaisA1Exploit2025,uclGoogleFunding2026}。 \\
% \cline{3-4}
% & & CAI~\cite{uclAICentre2026}
% & \textbf{负责任生成模型开发与信息泄露防护：}牵头英国生成式AI Hub，构建生成模型负责任开发与评估工具体系；提出Aspective Agentic AI框架，通过动态信息视域约束抑制多智能体系统信息泄露风险~\cite{uclGenAIHub2026,aspectiveAgenticAI2025}。 \\
% \cline{3-4}
% & & UCLIC~\cite{uclIC2026}
% & \textbf{人本可用安全与隐私伤害辅助报告：}研究对话式AI辅助用户上报隐私伤害的机制，其人本安全研究成果落地英国政府与行业安全规范~\cite{uclHumanCentredSecurity2026}。 \\
% \cline{2-4}

% & \multirow{4}{=}{爱丁堡大学\\(Edin-\\burgh)}
% & AIAI~\cite{edinburghAIAI2026}
% & \textbf{长时序智能体规划与协作任务管理：}基于O-Plan/I-X系列规划技术，整合任务表示、动态重规划与跨主体协作机制，奠定长时序智能体任务调度与协同管理的系统基础~\cite{tateOPlan1996}。 \\
% \cline{3-4}
% & & CAIAA~\cite{edinburghAssistiveAutonomy2026}
% & \textbf{人机协作环境下安全鲁棒自主决策：}融合自动驾驶、机器人学习与人机对话技术，构建辅助自主平台，研究不确定性场景下人机协同智能体的安全决策机制。 \\
% \cline{3-4}
% & & SPT Group~\cite{edinburghSPT2026}
% & \textbf{可验证安全证据与信息流控制：}基于携带证明代码与形式化协议验证技术，为智能体外部代码执行、插件调用、分布式服务访问提供可验证安全证据与信息流管控方案~\cite{aspinallMRG2005}。 \\
% \cline{3-4}
% & & CTF~\cite{edinburghCTF2026}
% & \textbf{负责任AI生态治理分歧梳理：}依托BRAID项目联动学界、产业与监管机构，系统厘清AI安全、伦理、透明度、问责与治理体系的重叠与冲突边界~\cite{braidProgramme2026,gyevnar2026responsibleDivides}。 \\
% \cline{2-4}

% % & \multirow{4}{=}{Google DeepMind}
% % & AGI Safety\newline and Alignment~\cite{deepmindRohinShah2026}
% % & \textbf{AGI技术安全与智能体行为监测护栏：}发布GDM AI Control Roadmap v0.1，借鉴网络安全内部威胁防护思路，构建智能体工具访问控制与实时行为监测安全护栏体系~\cite{deepmindControlRoadmap2026}。 \\
% % \cline{3-4}
% % & & Frontier Safety\newline and Governance~\cite{deepmindDafoe2026}
% % & \textbf{前沿模型危险能力评估与多智能体风险基金：}开展前沿模型危害能力定级，设立千万级专项基金，研究大规模多智能体交互中的诈骗、提示注入与网络攻击风险~\cite{deepmindMultiAgentFund2026}。 \\
% % \cline{3-4}
% % & & Mechanistic\newline Interpretability~\cite{deepmindNeelNanda2026}
% % & \textbf{机械可解释性与隐藏目标识别：}通过神经网络逆向工程识别模型隐藏目标与欺骗倾向，开源Gemma Scope稀疏自编码器工具集支撑安全社区可解释性研究~\cite{deepmindGemmaScope2024}。 \\
% % \cline{3-4}
% % & & 具身智能体\newline 研究方向~\cite{deepmindRoboCat2023}
% % & \textbf{具身智能体跨任务迁移与动作限制：}基于RoboCat与SIMA模型验证语言指令与物理动作的安全鸿沟，提出动作边界约束与人工接管机制防控具身智能体执行风险~\cite{deepmindRoboCat2023}。 \\
% % \cline{2-4}

% % & \multirow{4}{=}{Alan Turing Institute}
% % & Safe and\newline Ethical AI~\cite{turingGal2026}
% % & \textbf{大规模LLM数据投毒后门研究：}完成迄今最大规模LLM数据投毒实验，证实仅250份恶意文档即可在任意规模大模型中植入持久后门，为训练数据安全管控提供依据~\cite{turingPoisoning2025}。 \\
% % \cline{3-4}
% % & & Defence and National\newline Security / DARe~\cite{turingLASR2026}
% % & \textbf{受限环境多机器人韧性与对抗性威胁研究：}依托820万英镑LASR实验室，联合GCHQ等机构研究通信脆弱、资源受限场景下边缘AI与多机器人系统的对抗性风险与韧性增强技术~\cite{turingLASR2026}。 \\
% % \cline{3-4}
% % & & Public Policy\newline Programme~\cite{turingBright2026}
% % & \textbf{在线安全与AI公共政策立法咨询：}为英欧政府提供算法问责、在线有害内容检测领域技术支撑，推动AI安全相关立法落地~\cite{turingBright2026}。 \\
% % \cline{3-4}
% % & & Human-AI Interfaces\newline and Robotics~\cite{turingVijayakumar2026}
% % & \textbf{人机协作责任边界与控制权移交：}研究人机协同机器人的技能共享控制机制，界定自主系统与人类操作员之间的安全责任边界与动态控制权切换规则~\cite{turingVijayakumar2026}。 \\
% \hline
% \multirow{14}{=}{德国}
% % & \multirow{4}{=}{CISPA Helmholtz Center}
% % & Trustworthy ML\newline and Cybersecurity~\cite{cispaFritz2026}
% % & \textbf{间接提示注入攻击识别与智能体部署安全：}2023年首次系统定义并分类间接提示注入攻击，厘清LLM应用数据与指令边界漏洞；2026年提出开放式智能体部署前安全优先的核心准则~\cite{cispaIndirectPI2023,cispaAgentSafety2026}。 \\
% % \cline{3-4}
% % & & Systems Security~\cite{cispaHolz2026}
% % & \textbf{系统层安全与固件模糊测试：}提出新型固件模糊测试方案，解决中断处理导致的测试覆盖率偏低难题，提升嵌入式设备底层漏洞挖掘效率~\cite{cispaFirmwareFuzzing2025}。 \\
% % \cline{3-4}
% % & & Web Security~\cite{cispaStock2026}
% % & \textbf{第三方服务监控盲区测量：}提出LeakyLinks框架，量化分析URL扫描服务在浏览器生态中形成的隐私监控盲区，为Web端智能体隐私防护提供依据~\cite{cispaLeakyLinks2026}。 \\
% % \cline{3-4}
% % & & Trustworthy ML\newline (LLM Safety)~\cite{cispaYangZhang2026}
% % & \textbf{仅标签成员推理攻击与训练数据泄露：}实现低查询量增强型仅标签成员推理攻击，揭示LLM训练数据泄露的新型攻击路径，适配轻量化智能体安全审计场景~\cite{cispaMembershipInference2025}。 \\
% % \cline{2-4}

% & \multirow{4}{=}{慕尼黑工业大学(TUM)}
% & CPS~\cite{tumAlthoff2026}
% & \textbf{可达性分析与LLM控制机器人安全保证：}基于可达性分析理论，构建LLM驱动机器人的形式化验证框架，提供数学可证明的物理执行安全保障~\cite{tumSafeLLMRobots2026}。 \\
% \cline{3-4}
% & & DAML~\cite{tumGuennemann2026}
% & \textbf{嵌入空间对齐绕过与红队评估鲁棒性：}证实嵌入空间软提示攻击可绕过LLM安全对齐与遗忘机制；证明LLM-as-a-Judge在红队测试中评估鲁棒性接近随机猜测，质疑现有评估范式可靠性~\cite{tumCoinFlipSafety2026}。 \\
% \cline{3-4}
% & & SPRLS ~\cite{tumSchoellig2026}
% & \textbf{学习型控制系统安全性能基准平台：}开发safe-control-gym开源基准，实现传统控制、强化学习等智能体控制系统安全性能的定量对比与标准化评测。 \\
% \cline{3-4}
% & & SIT ~\cite{tumSigl2026}
% & \textbf{硬件侧信道攻击与模型参数窃取：}验证物理访问场景下，侧信道分析可窃取神经网络加速器的核心参数，为边缘部署智能体的硬件安全防护提供依据~\cite{tumSideChannelNN2026}。 \\
% \cline{2-4}

% % & \multirow{4}{=}{马克斯·普朗克智能系统研究所(MPI-IS)}
% % & AI Safety\newline and Alignment~\cite{mpiAndriushchenko2026}
% % & \textbf{自主LLM智能体对齐与能力评估：}聚焦自主式大模型智能体价值对齐研究，构建严格的前沿模型能力分级评估体系，推进AI风险公众科普与认知引导~\cite{mpiAndriushchenko2026}。 \\
% % \cline{3-4}
% % & & COMPASS~\cite{mpiAbdelnabi2026}
% % & \textbf{隐性记忆隐藏通道与时序后门：}提出时序性"记忆炸弹"后门机制，仅在特定交互序列下激活恶意行为，揭示跨智能体隐蔽通信与基准数据集污染风险~\cite{mpiTimeBombs2026}。 \\
% % \cline{3-4}
% % & & Empirical\newline Inference~\cite{mpiScholkopf2026}
% % & \textbf{因果学习理论与AI监管科学支撑：}深耕因果学习基础理论，团队多人入选联合国、欧盟AI法案专家小组，为国际AI监管提供科学理论支撑~\cite{mpiEUAIAct2026}。 \\
% % \cline{3-4}
% % & & Social Foundations\newline of Computation~\cite{mpiHardt2026}
% % & \textbf{法律智能体真实任务表现评测：}参与构建Legal Agent Benchmark基准，聚焦真实法律场景量化评测专业领域智能体的任务执行能力与安全合规性~\cite{mpiLegalAgentBenchmark2026}。 \\
% % \cline{2-4}

% & \multirow{2}{=}{图宾根大学}
% & Bethge Lab~\cite{tuebingenBethge2026}
% & \textbf{对抗鲁棒性评测工具与基准建设：}开发Foolbox开源工具包，标准化对抗样本生成与模型鲁棒性量化流程，成为视觉智能体安全评测的通用基准工具。 \\
% \cline{3-4}
% & & AVG ~\cite{tuebingenGeiger2026}
% & \textbf{自动驾驶3D场景理解基准数据集：}维护KITTI系列权威自动驾驶数据集，提出TransFuser端到端感知模型，成为具身驾驶智能体场景理解的基线标准。 \\
% \hline

% \multirow{8}{=}{瑞士}
% & \multirow{4}{=}{苏黎世联邦理工学院(ETH Zurich)}
% & SPY Lab~\cite{ethSPYLab2026,tramerHomepage2026}
% & \textbf{智能体提示注入攻防动态评测环境：}构建AgentDojo评测平台，集成97项真实任务与629条安全用例，统一量化评估LLM智能体工具调用、外部数据读取中的间接提示注入攻防性能~\cite{agentdojoGithub2026}。 \\
% \cline{3-4}
% & & SRI Lab~\cite{ethSRILab2026,ethVechev2026}
% & \textbf{代码智能体安全关键任务生成评测：}提出BaxBench基准，通过392项高危后端编程任务评测代码智能体生成内容的正确性与安全性，证实顶尖模型仍存在大量高危漏洞~\cite{baxbench2026}。 \\
% \cline{3-4}
% & & LAS Group~\cite{ethLAS2026,ethKrause2026}
% & \textbf{安全强化学习中的主动探索约束：}提出ActSafe算法，融合主动学习与硬安全约束，管控强化学习智能体在未知环境探索中的风险暴露范围~\cite{as2025actsafe}。 \\
% \cline{3-4}
% & & RSL~\cite{ethRSL2026,ethHutter2026}
% & \textbf{腿式机器人复杂地形自主导航安全：}基于强化学习与机载感知，实现ANYmal系列机器人复杂地形稳定移动，将智能体安全体系延伸至物理设备运动边界与环境适配场景~\cite{ethAnyMal2026}。 \\
% \cline{2-4}

% & \multirow{4}{=}{洛桑联邦理工学院(EPFL)}
% & HexHive~\cite{payerHomepage2026}
% & \textbf{固件移植与边缘设备漏洞检测：}提出EmbedFuzz固件移植技术，实现微控制器固件高性能原生模糊测试，支撑边缘智能体插件与控制模块的底层漏洞挖掘。 \\*
% \nobreakcline{3-4}
% & & SPRING Lab~\cite{springLab2026}
% & \textbf{去中心化隐私保护与敏感数据风险通知：}研发DP3T去中心化追踪系统，以最小数据收集原则实现隐私风险通知，为医疗、政务敏感数据处理智能体提供隐私设计范式。 \\*
% \cline{3-4}
% & & DEDIS Lab~\cite{epflDEDIS2026}
% & \textbf{去中心化信任基础设施与多机构协同审计：}构建Cothority/OmniLedger分布式信任架构，分散信任节点权限，支撑跨机构多智能体协同的身份认证与不可篡改审计追踪~\cite{epflOmniLedger2026}。 \\*
% \cline{3-4}
% & & VITA Lab~\cite{epflVITA2026}
% & \textbf{人机共存场景轨迹预测与碰撞风险规避：}提出Social-Transmotion视觉轨迹预测方法，提升公共空间移动智能体对行人意图与碰撞风险的预判能力，优化人机共存安全决策。 \\*
% \hline

% % \multirow{4}{=}{法国}
% % & \multirow{4}{=}{Inria}
% % & Prosecco~\cite{inriaProsecco2026}
% % & \textbf{密码协议形式化验证与后量子安全分析：}基于ProVerif工具完成Signal后量子密钥协商协议PQXDH的自动化形式化验证，筑牢智能体加密通信底层安全基础~\cite{inriaPQXDH2024}。 \\
% % \cline{3-4}
% % & & PESTO~\cite{inriaPesto2026}
% % & \textbf{隐私性质符号化验证与协议攻击发现：}实现安全协议隐私属性的符号化自动验证，发现新型哈希漏洞攻击模式，获USENIX Security 2023杰出论文奖~\cite{inriaHashGoneBad2023}。 \\
% % \cline{3-4}
% % & & Privatics~\cite{inriaCastelluccia2026}
% % & \textbf{数据隐私保护与算法认知安全威胁：}研究大规模数据隐私脱敏技术，重点分析AI智能体操纵人类决策的认知层安全威胁与防御路径~\cite{inriaCastelluccia2026}。 \\
% % \cline{3-4}
% % & & Toccata~\cite{inriaToccata2026}
% % & \textbf{演绎验证平台与关键系统安全工具链：}研发Why3演绎验证框架，集成至航空级SPARK/Ada工具链，支撑高可信智能体核心控制软件的形式化证明~\cite{inriaWhy3ProofInUse2023}。 \\
% % \hline

% \multirow{6}{=}{荷兰}
% & \multirow{4}{=}{代尔夫特理工大学(TU Delft)}
% & II Group~\cite{tudelftJonker2026}
% & \textbf{智能体协商与人机混合智能系统：}深耕自主智能体自动协商机制与人机协同认知建模，获2024年ACM/SIGAI自主智能体领域终身成就奖~\cite{tudelftJonker2026}。 \\
% \cline{3-4}
% & & CS Group~\cite{tudelftErkin2026}
% & \textbf{隐私求和协议的图结构重构攻击：}发现可证明安全的隐私求和协议在并行执行时存在图结构泄露漏洞，为隐私保护型多智能体集群部署提供风险警示。 \\
% \cline{3-4}
% & & AMR Lab~\cite{tudelftAlonsoMora2026}
% & \textbf{LLM融合导航与有界概率运动规划：}融合大语言模型与模型预测控制实现机器人个性化安全导航；提出带碰撞概率上界的场景化规划方法，保障机器人群体运动安全~\cite{tudelftHeyRobot2026,tudelftScenarioPlanning2025}。 \\
% \cline{3-4}
% & & GC~\cite{tudelftVanEeten2026}
% & \textbf{关键基础设施网络安全治理测量：}基于全网实测数据量化分析ISP与终端设备对Mirai僵尸网络的清除效果，支撑物联网智能体集群治理策略优化。 \\
% \cline{2-4}

% & \multirow{2}{=}{阿姆斯特丹大学(UvA)}
% & VIS/QUVA Lab~\cite{uvaSnoek2026}
% & \textbf{具身多智能体协作规划优化：}提出CaPo协同规划算法，优化具身场景下多智能体的任务分配与动作协同效率，降低协作冲突风险~\cite{uvaCaPo2025}。 \\
% \cline{3-4}
% & & IRLab~\cite{uvaDeRijke2026}
% & \textbf{检索增强智能体信息可靠性理论：}构建检索增强生成（RAG）智能体的信息可信度评估理论体系，解决外部检索内容带来的事实性与安全风险问题~\cite{uvaDeRijke2026}。 \\
% \hline

% \multirow{4}{=}{比利时}
% & \multirow{4}{=}{鲁汶大学(KU Leuven)}
% & COSIC~\cite{cosicKuLeuven2026}
% & \textbf{密码学与敏感数据访问技术控制：}基于哈希函数、隐私增强技术构建智能体敏感数据访问的底层密码管控机制，支撑合规化数据权限治理~\cite{cosicKuLeuven2026}。 \\
% \cline{3-4}
% & & DistriNet（软件安全）~\cite{distrinetJoosen2026}
% & \textbf{LLM聊天机器人冒充与网络民主威胁防御：}系统研究AI聊天机器人冒充人类干预公共讨论的攻击路径，构建针对性检测与防御体系~\cite{distrinetLLMChatbot2026}。 \\
% \cline{3-4}
% & & DistriNet（AI可信性）~\cite{distrinetRimmer2026}
% & \textbf{入侵检测场景下数据驱动AI可信度：}分析无监督数据环境下AI模型在恶意代码检测中的可信度缺陷，提出风险感知型智能体检测方案~\cite{distrinetRimmer2026}。 \\
% \cline{3-4}
% & & CiTiP~\cite{citipValcke2026}
% & \textbf{GDPR合规转化与陪伴型智能体数据保护：}将GDPR合规条款转化为可量化的网络安全成熟度模型；针对性解析陪伴型聊天机器人的设计阶段数据保护合规要点~\cite{citipDataProtectionByDesign2025}。 
% \end{longtable}
% \end{center}
% \endgroup

% \begin{table}[htbp]
% \centering
% \caption{欧洲智能体安全研究方向样本占比}
% \begin{tabular}{p{5cm}p{2cm}p{7cm}}
% \hline
% 研究方向 & 占比 & 说明 \\
% \hline
% 可信 AI、解释性与人类监督 & 24\% & 高风险场景、监督边界、可解释轨迹 \\
% 隐私保护与数据治理 & 22\% & GDPR、数据最小化、企业数据边界 \\
% 形式化验证与系统安全 & 18\% & 策略约束、程序分析、可信执行 \\
% 提示注入与工具调用评测 & 16\% & AgentDojo 等动态攻防环境 \\
% 多智能体与机器人安全 & 12\% & 协作、具身智能体、工业控制 \\
% 政策治理与责任链 & 8\% & AI Act、责任归因、影响评估 \\
% \hline
% \end{tabular}
% \end{table}

% \begin{figure}[htbp]
% \centering
% \includegraphics[width=0.48\textwidth]{figures/europe_academic_share.png}
% \includegraphics[width=0.48\textwidth]{figures/europe_direction_share.png}
% \caption{欧洲高校/机构样本占比与研究方向占比}
% \end{figure}


\subsection{我国AI智能体安全学术研究现状}
区别于北美偏向前沿攻防与防务赋能、欧洲侧重法理规制与可证明基础理论的研究路径，我国研究更聚焦大模型原生智能体的工程化安全问题，以评测基准构建、工具调用风险防控、中文场景安全对齐、产业级落地防护为核心主线。
%
% 我国AI智能体安全研究主体呈现“顶尖高校引领、国家级科研院所攻坚、新型研发机构赋能、头部实验室产业化转化”的立体化布局。高校层面以清华、北大、上交、浙大、南大、复旦为核心，在智能体能力评测、安全对齐、工具学习安全、多模态防护等方向形成标杆成果\cite{tsinghuaTHUDM2026,pkuAlignment2026,sjtuAISecLab2026}；科研院所以中国科学院自动化所、计算所、信工所等主体，主攻底层系统安全、多智能体博弈攻防、伪造检测与国家级安全测评平台研发\cite{pandaguardCASIA2026,chenweiICT2026,huSonglinGalexy2024}；北京智源、上海AI实验室、北京通用AI研究院等新型研发机构立足通用AI底座，构建智能体原生安全框架、分级评测体系与产业级安全操作系统\cite{flagevalBAAI2024,safework2025,bigaiTongTest2024}；区域实验室与地方高校则聚焦垂直行业场景，补齐端侧智能体、行业合规安全的研究短板。
%
本节梳理我国高校、科研院所、新型研发机构的主要团队、研究方向和标志性成果，主要发现如下 （详见表~\ref{tab:china_ai_agent_sec}）：

\begin{itemize}[tabitem]
\item \textbf{研究资源集聚分层，头部机构形成差异化分工壁垒。}我国AI智能体安全研究以清华、北大与中国科学院体系为核心研发主体，在基础评测、安全对齐、底层系统安全领域形成引领性优势\cite{pandaguardCASIA2026}；上交、浙大、南大、复旦构成第二梯队，聚焦工程化落地与隐私安全\cite{aicertZJU2026,cosecNJU2026,jadeWhitzard2026}；
网安强校深耕底层漏洞与端侧防护，中国科学院主攻国家级测评平台与硬件底层安全\cite{chenweiICT2026,ictCASOfficial2026}，新型研发机构聚焦通用智能体原生安全与产业转化\cite{flagopenBAAI2026,intershannon2026}。

\item \textbf{技术导向偏重工程实用，构建本土化特色赛道。}我国锚定产业刚需，主要研究包括覆盖通用、交互、场景化的智能体安全双维评测体系\cite{csebenchmark2025}、基于国产开源底座的工具学习与RAG安全方案\cite{openmmlabSHLab2026}、适配中文语境的价值对齐范式，以及规模化落地的多模态与具身智能体防护体系\cite{reinEAD2025}。依托自主算力架构，我国在硬件级端侧可信执行、权限隔离领域形成独有技术特色，\cite{ictCASOfficial2026,flagopenBAAI2026}。

\item \textbf{整体呈并跑-跟跑格局，基础理论存在核心短板。}我国在开源工具链安全、产业级测评平台、安全标准转化领域紧跟国际前沿趋势\cite{jadeWhitzard2026}，中文场景对齐、国产算力防护、行业安全赛道与全球并跑，基础机理、可证明安全、AGI前沿风险研究整体跟跑。核心短板集中于三方面：一是原生攻击机理原创性不足，多为欧美现有攻击范式的本土化适配，缺乏奠基性成果；二是大模型智能体专属可证明安全体系空白，形式化验证、决策鲁棒性数学证明能力薄弱；三是超级对齐、多智能体链式风险等AGI级前沿研究体量不足，风险预判能力滞后\cite{zengyiCASIA2026,bigaiResearchDept2026}。
\end{itemize}


% \nopagebreak[4]
\enlargethispage{1cm}
\footnotesize
\setlength{\tabcolsep}{3pt}

\begin{longtblr}[
caption = {我国顶尖高校与科研研所AI智能体安全学术研究汇总},
label   = {tab:china_ai_agent_sec},
]{
colspec = {
|>{\rule{0pt}{2pt}}m{0.8cm}
|>{\rule{0pt}{2pt}}m{1.4cm}
|>{\rule{0pt}{2pt}}m{3cm}
|>{\rule{0pt}{2pt}}m{
    \dimexpr
    \linewidth
    - 0.8cm
    - 1.4cm
    - 3cm
    - 8\tabcolsep
    - 4\arrayrulewidth
    \relax
}|
},
rowhead = 1,
}

\hline
\textbf{研究主体}
& \textbf{高校 / 科研院所}
& \textbf{实验室/研究团队}
& \textbf{核心智能体安全研究方向} \\
\hline

% =========================================================
% 高校：共34行
% =========================================================

\SetCell[r=34]{c} 高校
& \SetCell[r=3]{c} {清华\\大学}
& 知识工程实验室与THUDM团队~\cite{tsinghuaKEG2026,tsinghuaTHUDM2026}
& \textbf{大模型智能体交互能力评测基准：}依托GLM系列与知识增强推理研究提出AgentBench，在8类交互环境下系统评测智能体推理、决策与长期执行能力~\cite{thuAgentBenchGithub2026}。 \\
\cline{3-4}

& &
OpenBMB与THUNLP-MT工具学习团队~\cite{openBMB2026}
& \textbf{工具调用与API调度安全评估：}构建ToolBench与StableToolBench两套开源基准，为大模型工具选择、API调度场景提供可复现安全评测数据集与实验框架~\cite{thuToolBenchGithub2026}。 \\
\cline{3-4}

% & &
% CoAI对话式AI团队~\cite{tsinghuaCoAI2026}
% & \textbf{中文大模型安全语料与检测器构建：}面向中文价值对齐与风险可解释评估，搭建SafetyPrompts安全语料库，研发可定制安全检测模型ShieldLM~\cite{thuSafetyPrompts2023,shieldLM2024}。 \\
% \cline{3-4}

& &
TSAIL~\cite{tsinghuaTSAIL2026}
& \textbf{具身智能体对抗鲁棒性提升：}基于贝叶斯可信机器学习理论提出Rein-EAD算法，强化三维对抗场景下具身智能体视觉感知鲁棒性~\cite{reinEAD2025}。 \\
\cline{2-4}

& \SetCell[r=3]{c} {北京\\大学}
& PKU-Alignment团队~\cite{pkuAlignment2026,yangyaodongHomepage2026}
& \textbf{大模型价值对齐与安全强化学习框架：}提出Safe RLHF、Safe Policy Optimization等安全强化学习算法，开源BeaverTails、PKU-SafeRLHF安全偏好数据集，相关成果入选ICLR顶会Spotlight。 \\
\cline{3-4}

& &
信息安全实验室~\cite{pkuInfosecLab2026}
& \textbf{软件安全测试与漏洞分析：}聚焦分布式系统、二进制漏洞挖掘与自动化安全测试，成果持续发表于ASE、ISSTA、ICDCS等软件工程与分布式系统顶会~\cite{pkuInfosecLab2026}。 \\
\cline{3-4}

% & &
% 视觉信息智能学习实验室（VILLA）~\cite{villaLab2026}
% &  \\
% \cline{3-4}

& &
多智能体研究中心~\cite{pkuMultiAgentCenter2026}
& \textbf{数据分析智能体交互可靠性评测：}联合伯克利提出IDA-Bench基准，实验证明当前最优大模型完成数据分析智能体任务仅能达到人类40\%成功率~\cite{idabench2025}。 \\
\cline{2-4}

& \SetCell[r=4]{c} {上海交通大学}
& MVIG实验室~\cite{luceluMVIG2026}
& \textbf{具身智能与机器人感知系统：}深耕机器人视觉感知与人体姿态建模，开源AlphaPose、HAKE工具成为行业通用基线~\cite{luceluMVIG2026}。 \\
\cline{3-4}

& &
密码与计算机安全实验室（LoCCS）~\cite{sjtuLoCCS2026}
& \textbf{密码理论与系统安全基础：}研究底层密码原语、软硬件隔离机制，为智能体运行环境底层安全评估提供理论与工具支撑~\cite{sjtuLoCCS2026}。 \\
\cline{3-4}

& &
AI安全实验室~\cite{sjtuAISecLab2026}
& \textbf{智能体交互风险识别基准R-Judge：}将智能体安全判定定义为多轮对话风险识别任务，构建R-Judge评测基准，最优模型判别准确率仅74.42\%。 \\
\cline{3-4}

& &
跨媒体语言智能实验室（X-LANCE）~\cite{sjtuXLANCE2026}
& \textbf{语音交互智能体安全边界研究：}面向口语对话与认知型智能体，搭建语音场景安全边界分析完整体系~\cite{sjtuXLANCE2026}。 \\
\cline{2-4}

& \SetCell[r=4]{c} {浙江\\大学}
& 任奎团队（区块链与数据安全全国重点实验室）~\cite{kuiRenZJU2026}
& \textbf{全栈智能体安全评测平台AIcert：}研发覆盖数据、算法、模型、框架、系统五层安全检测的AIcert平台，已落地直播风控、轨道交通工控智能体场景~\cite{aicertZJU2026}。 \\
\cline{3-4}

& &
浙江大学--阿里巴巴集团AI安全联合实验室~\cite{zjuAlibabaAISafetyLab2025}
& \textbf{企业级智能体红队测试与安全治理：}以大模型智能体红队攻防、全生命周期安全治理为核心方向，支撑产业落地级安全体系搭建~\cite{zjuAlibabaAISafetyLab2025}。 \\
\cline{3-4}

& &
NESA NSSP Lab~\cite{zjuNESA2026}
& \textbf{AI辅助漏洞挖掘与智能模糊测试：}基于对抗学习设计自动化漏洞挖掘方案，面向智能体插件、边缘调度模块开展智能模糊测试研究~\cite{zjuNESA2026}。 \\
\cline{3-4}

& &
网络空间安全学院大模型与智能体安全团队~\cite{yangZiqiZJU2026,zjuFaculty2026}
& \textbf{移动端与Web漏洞挖掘：}围绕大模型智能体、移动应用、Web前端构建一体化漏洞挖掘与风险管控方案~\cite{yangZiqiZJU2026,zjuFaculty2026}。 \\
\cline{2-4}

& \SetCell[r=4]{c} {南京\\大学}
& LAMDA实验室~\cite{lamdaOfficial2026,zhouzhihuaHomepage2026}
& \textbf{学件范式与模型复用基座系统：}提出通用可复用模型“学件”理论，搭建开源模型复用基座平台北冥坞~\cite{lamdaOfficial2026}。 \\
\cline{3-4}

& &
LAMDA-NeSy团队~\cite{lamdaOfficial2026}
& \textbf{中文场景智能体任务规划基准：}发布面向真实出行场景的ChinaTravel规划基准，定为IJCAI 2025/2026国际竞赛官方评测集~\cite{chinaTravelIJCAI2026}。 \\
\cline{3-4}

& &
推理与学习研究组~\cite{gaoyangRLGroup2026}
& \textbf{强化学习驱动的具身智能体导航：}融合大模型规划与强化学习，研究开放环境下机器人自主导航、物体操作安全控制方案~\cite{gaoyangRLGroup2026}。 \\
\cline{3-4}

& &
COSEC实验室~\cite{cosecNJU2026}
& \textbf{深度学习与联邦学习安全：}研究模型梯度泄露、联邦聚合攻击防护方案，相关隐私安全技术已落地头部互联网企业智能体业务~\cite{cosecNJU2026}。 \\
\cline{2-4}

& \SetCell[r=4]{c} {复旦\\大学}
& NLP实验室~\cite{fudanNLPLab2026,qiuxipengHomepage2026}
& \textbf{插件增强对话模型技术路线：}开源我国首个插件增强大模型MOSS，打通预训练至工具调用全链路；联合产业机构推进OpenMOSS多模态开源智能体生态建设~\cite{openMOSS2026}。 \\
\cline{3-4}

& &
知识工场实验室~\cite{fudanKWLab2026}
& \textbf{大规模中文知识图谱服务：}构建CN-DBpedia、CN-Probase通用中文知识图谱，对外提供标准化知识检索API，支撑检索增强智能体可信问答~\cite{fudanKWLab2026}。 \\
\cline{3-4}

& &
系统软件与安全实验室~\cite{whitzardAI2026}
& \textbf{系统安全与恶意代码检测：}深耕底层系统漏洞、恶意样本识别与AI系统攻防，成果长期发表S\&P、USENIX Security、CCS安全顶会~\cite{whitzardAI2026}。 \\
\cline{3-4}

& &
白泽智能团队~\cite{whitzardAI2026}
& \textbf{大模型靶向式安全评测平台JADE：}搭建JADE定向安全评测系统，实测模型违规检出率超70\%；扩展至智能体MCP协议风险检测，参与制定《智能体应用安全基本要求》国家标准~\cite{jadeWhitzard2026,whitzardAI2026}。 \\
\cline{2-4}

& \SetCell[r=4]{c} {哈尔滨工业大学}
& HIT-SCIR~\cite{hitSCIROfficial2026,liutingHomepage2026}
& \textbf{中文语言技术基础平台建设：}长期维护LTP中文自然语言处理工具栈，为对话、检索类智能体提供底层语言安全预处理能力。 \\
\cline{3-4}

& &
健康智能组~\cite{hitirHealthAgent2026}
& \textbf{多智能体协同临床诊断：}提出多智能体协同进化诊疗框架，诊断准确率提升10\%以上，已落地全国30余家医疗机构临床智能体系统~\cite{hitirHealthAgent2026}。 \\
\cline{3-4}

& &
自主智能无人系统全国重点实验室~\cite{aiusHIT2026}
& \textbf{分布式协同控制理论方法：}构建混合智能优化与稳定性证明体系，支撑无人集群多智能体协同安全调度~\cite{aiusHIT2026}。 \\
\cline{3-4}

& &
多智能体可靠与安全控制团队~\cite{aiusHIT2026}
& \textbf{多智能体系统可靠与安全控制：}面向无人集群、工业协同场景，研究分布式多智能体容错、安全约束控制理论与落地方案~\cite{aiusHIT2026}。 \\
\cline{2-4}

& \SetCell[r=4]{c} {中国科学技术大学}
& 张卫明教授团队~\cite{zhangweimingUSTC2026}
& \textbf{大模型越狱攻击防范与水印嵌入：}研究大模型提示越狱防御、生成文本溯源水印技术，斩获IJCAI对抗防御赛道冠军、DFDC伪造检测亚军~\cite{zhangweimingUSTC2026}。 \\
\cline{3-4}

& &
杨威副教授团队~\cite{yangweiUSTC2026}
& \textbf{量子AI与安全多方计算：}结合量子计算、隐私计算技术，设计面向分布式智能体的安全数据交互协议~\cite{yangweiUSTC2026}。 \\
\cline{3-4}

& &
黄刘生教授团队~\cite{huangliushengUSTC2026}
& \textbf{信息隐藏与量子信息安全：}研究多媒体隐写、量子保密通信，为多智能体跨节点隐蔽信息传输提供安全防护方案~\cite{huangliushengUSTC2026}。 \\
\cline{3-4}

& &
先进数据系统实验室（ADSL）~\cite{ustcADSL2026}
& \textbf{大规模存储与分布式机器学习：}面向分布式训练集群、海量智能体数据存储开展底层安全优化，成果发表FAST、OSDI、SOSP等系统顶会~\cite{ustcADSL2026}。 \\
\cline{2-4}

& \SetCell[r=4]{c} {北京航空航天大学}
& 韦星星教授团队~\cite{weixingxingBUAA2026}
& \textbf{多模态大模型可信度评估：}构建图文跨模态幻觉检测、可信度量化体系，面向多模态交互智能体开展风险量化评估~\cite{weixingxingBUAA2026}。 \\
\cline{3-4}

& &
可信智能软件实验室~\cite{zhengzhengBUAA2026}
& \textbf{软件可靠性工程与测试方法：}研究高可信智能体软件形式化测试、故障容错技术，负责人长期担任ISSRE、PRDC等国际可靠性会议主席~\cite{zhengzhengBUAA2026}。 \\
\cline{3-4}

& &
需求工程团队~\cite{zhangliBUAA2026}
& \textbf{软件建模与需求工程分析：}面向自主智能体构建安全需求建模、风险识别标准化流程~\cite{zhangliBUAA2026}。 \\
\cline{3-4}

& &
通用自主多模态智能体课题组（GAMMA-LAB）~\cite{gammaLabBUAA2026}
& \textbf{无人机多智能体协同设计验证：}研发ArtiAero平台，基于多智能体协同完成无人机从概念设计、模型生成到仿真安全验证全流程自动化~\cite{artiaeroBUAA2026,gammaLabBUAA2026}。 \\
\hline

% =========================================================
% 科研院所：共27行
% =========================================================

\SetCell[r=1]{c} 科研院所
& \SetCell[r=3]{c} {中国科学院自动化所}
& AI安全与超级对齐北京市重点实验室~\cite{zengyiCASIA2026}
& \textbf{大模型安全攻防评估平台PandaGuard：}搭建灵御PandaGuard评测平台，集成19类攻击、12套防御方案，完成49款主流大模型安全量化测评~\cite{pandaguardCASIA2026}。 \\
\cline{3-4}

& &
复杂系统管理与控制国家重点实验室~\cite{wangfeiyueCASIA2026}
& \textbf{平行智能与复杂系统建模：}基于平行智能理论构建社会、工业多智能体复杂系统仿真与风险推演框架~\cite{wangfeiyueCASIA2026}。 \\
\cline{3-4}

& &
多智能体决策与科研智能体平台（JidiAI/ScienceOne）~\cite{jidiAI2026,scienceOne2026}
& \textbf{多智能体博弈决策开放平台：}JidiAI接入40余种博弈仿真环境，服务海内外50余家科研机构；ScienceOne整合300+科学工具，实现科研智能体自主规划推理~\cite{jidiAI2026,scienceOne2026}。 \\
\cline{2-4}

& \SetCell[r=4]{c} {中国科学院计算所}
& 智能算法安全全国重点实验室~\cite{chenweiICT2026}
& \textbf{机器学习安全理论与可信应用：}围绕基础机器学习攻防理论，构建面向行业智能体的可信部署、风险管控整套技术体系~\cite{chenweiICT2026}。 \\
\cline{3-4}

& &
前瞻研究实验室/数字内容合成与伪造检测实验室~\cite{caojuanICT2026}
& \textbf{虚假新闻与伪造内容检测平台：}落地“AI识谣”“睿鉴图灵”多模态伪造检测系统，支撑国家级舆情、内容安全智能体任务部署~\cite{caojuanICT2026}。 \\
\cline{3-4}

& &
处理器芯片全国重点实验室~\cite{ictCASOfficial2026}
& \textbf{国产通用CPU与开源芯片生态：}自主研发龙芯通用处理器，开源香山RISC-V芯片架构，为端侧智能体提供底层硬件可信基座~\cite{ictCASOfficial2026}。 \\
\cline{3-4}

& &
智能处理器方向~\cite{ictCASOfficial2026}
& \textbf{深度学习处理器架构设计：}提出寒武纪DianNao系列深度学习加速器架构，成果入选ISCA年度标杆成果、《Science》开创性进展~\cite{ictCASOfficial2026}。 \\
\cline{2-4}

& \SetCell[r=4]{c} {中国科学院信工所}
& 陈恺研究员团队~\cite{kaiChenTeam2026}
& \textbf{系统安全与AI安全研究：}聚焦操作系统底层漏洞、大模型智能体攻防，相关成果获CCF-IEEE CS青年科学家奖~\cite{kaiChenTeam2026}。 \\
\cline{3-4}

& &
虎嵩林研究员团队~\cite{huSonglinUCAS2026}
& \textbf{大模型安全测评平台Galaxy：}研发星河Galaxy安全测评系统，牵头举办首届全国生成式AI安全攻防大赛。 \\
\cline{3-4}

& &
葛仕明/赵险峰多媒体安全团队~\cite{geShimingIMSG2026,zhaoXianfengUCAS2026}
& \textbf{AI视觉安全与多媒体隐蔽通信：}研究深度学习视觉伪造检测、多媒体信息隐藏技术，防范多模态智能体隐蔽信道泄露风险~\cite{geShimingIMSG2026,zhaoXianfengUCAS2026}。 \\
\cline{3-4}

& &
王伟平研究员团队~\cite{wangWeipingUCAS2026}
& \textbf{大数据与数据安全治理：}面向政企大数据智能体，构建数据分级、访问管控、全链路安全治理框架~\cite{wangWeipingUCAS2026}。 \\
\cline{2-4}

& \SetCell[r=4]{c} {北京通用AI研究院}
& 跨媒体通用AI全国重点实验室~\cite{bigaiResearchDept2026}
& \textbf{通用AI统一理论与认知架构：}搭建通用人工智能统一建模、认知推理理论体系，支撑通用自主智能体底层设计~\cite{bigaiResearchDept2026}。 \\
\cline{3-4}

& &
通用AI安全北京市重点实验室~\cite{bigaiResearchDept2026}
& \textbf{通用智能体全链条安全治理：}覆盖价值对齐、认知推理、物理执行三层空间，构建通用智能体端到端安全管控体系~\cite{bigaiResearchDept2026}。 \\
\cline{3-4}

& &
TongTest评测团队
& \textbf{通用AI评级测试标准：}发布TongTest分级评测标准与平台，为通用智能体分级监管、准入审查提供量化依据。 \\
\cline{3-4}

& &
多智能体实验室~\cite{bigaiResearchDept2026}
& \textbf{混合动机博弈与演化博弈论：}基于演化博弈分析人机混合多智能体交互中的合作、背叛与安全均衡机制~\cite{bigaiResearchDept2026}。 \\
\cline{2-4}

& \SetCell[r=4]{c} {北京智源AI研究院}
& FlagEval
& \textbf{大模型安全与价值观评测：}完成140余款开源/闭源大模型全维度测评，将安全价值对齐列为七大核心评测能力维度之一。 \\
\cline{3-4}

& &
FlagOpen/FlagOS~\cite{flagopenBAAI2026,baaiFlagOS2026}
& \textbf{国产算力智能体部署框架：}搭建多芯片统一软件栈FlagOS，提供原生智能体调度能力，适配国产服务器与边缘端部署场景~\cite{flagopenBAAI2026,baaiFlagOS2026}。 \\
\cline{3-4}

& &
智源世界模型创新中心~\cite{baaiWorldModel2026}
& \textbf{具身智能体安全规划基础模型：}开源悟界·Physis-v0.1世界模型，为机器人、物理交互智能体提供仿真安全规划底层基座~\cite{baaiWorldModel2026}。 \\
\cline{3-4}

& &
具身交互世界模型实验室~\cite{baaiEmbodiedAI2026,baaiConferenceWorldModel2026}
& \textbf{垂直领域智能体任务规划落地：}基于RoboBrain、RoboOS平台，面向工业、服务机器人智能体完成安全任务规划方案落地验证~\cite{baaiEmbodiedAI2026,baaiConferenceWorldModel2026}。 \\
\cline{2-4}

& \SetCell[r=4]{c} {上海AI实验室}
& 安全可信AI中心（AI45）~\cite{ai45SHLab2026}
& \textbf{模型内生安全推理能力提升：}提出SafeLadder分层安全框架，发布SafeWork-R1模型增强大模型原生风险识别与安全推理能力~\cite{safework2025}。 \\
\cline{3-4}

& &
书安（Intern-Shannon）团队~\cite{intershannon2026}
& \textbf{产业级安全智能体操作系统：}研发面向政企高安全场景的智能体OS，采用“安全即服务”架构统一管控多智能体权限与交互风险~\cite{intershannon2026}。 \\
\cline{3-4}

& &
Intern-Discovery/InternAgent团队~\cite{internAgentSHLab2026}
& \textbf{自主科学发现闭环智能体框架：}开源InternAgent 1.5多智能体闭环系统，实现科研数据采集、实验执行、结论推理全自主化安全流程~\cite{internAgent15SHLab2026}。 \\
\cline{3-4}

& &
OpenMMLab, OpenGVLab开源平台~\cite{openmmlabSHLab2026,opengvlabSHLab2026}
& \textbf{多模态开源工具链建设：}搭建从视觉感知、多模态理解到模型安全评测、智能体部署的完整开源工具链OpenCompass、InternVL等~\cite{opencompassSHLab2026}。 \\
\cline{2-4}

& \SetCell[r=4]{c} {中关村实验室}
& 网络空间安全体系化研究团队
& \textbf{身份认证与供应链可信框架：}构建覆盖身份鉴权、权限管控、数据可信、软件供应链安全的一体化智能体防护体系~\cite{cyberspaceTrustModel2026}。 \\
\cline{3-4}

& &
APT与未知攻击危害消解团队~\cite{ccfAgentSecurityForum2025}
& \textbf{安全运营智能体威胁溯源：}提出数字替身仿真消解思路，用于安全运营智能体的攻击行为还原、威胁根因定位~\cite{ccfAgentSecurityForum2025}。 \\
\cline{3-4}

& &
大模型网络安全知识评测合作团队~\cite{csebenchmark2025}
& \textbf{网络安全知识盲区评测基准：}构建CSEBenchmark安全知识库，包含345个知识点、11050道测评题，量化验证大模型网络安全推理缺陷~\cite{csebenchmark2025,csebenchmarkGithub2025}。 \\
\cline{3-4}

& &
智能体通信安全与标准协同平台~\cite{agentCommunicationWorkshop2026}
& \textbf{智能体通信安全标准协同：}主办智能体通信安全研讨会与AISCAI安全大赛，推动多智能体交互通信安全行业标准共建~\cite{agentCommunicationWorkshop2026}。 \\
\hline

\end{longtblr}




% \begin{table}[htbp]
% \centering
% \scriptsize
% \renewcommand{\arraystretch}{0.7}
% \caption*{表\ref{tab:china_ai_agent_sec_1}（续）--科研院所}
% % \label{tab:china_ai_agent_sec_2}
% \begin{tabular}{|m{1cm}|m{1cm}|m{2cm}|m{8cm}|}
% \hline
% \textbf{地区} & \textbf{单位名称} & \textbf{实验室/研究团队} & \textbf{核心智能体安全研究方向} \\
% \hline
% \multirow{5}{=}{科研院所}
% & \multirow{1}{=}{中国科学院自动化所}
% & \multirow{1}{=}{AI安全与超级对齐北京市重点实验室（曾毅团队）~\cite{zengyiCASIA2026}}
% & \begin{itemize}[tabitem]
% \item 发布PandaGuard（灵御）多智能体博弈评测平台，将越狱安全建模为多智能体攻防系统
% \item 平台集成19类攻击范式，覆盖大模型越狱风险检测~\cite{pandaguardCASIA2026}
% \item 依托王飞跃团队群体与博弈智能理论提供底层理论支撑~\cite{wangfeiyueCASIA2026}
% \end{itemize} \\
% \cline{2-4}
% & \multirow{1}{=}{中国科学院软件所}
% & \multirow{1}{=}{智能博弈重点实验室（原互联网软件技术实验室）~\cite{issSmartGame2026}}
% & \begin{itemize}[tabitem]
% \item 长期主攻可信软件过程管理、软件风险控制、程序形式化验证与自动推理
% \item 我国软件工程领域标杆实验室
% \item 公开资料暂无面向自主智能体安全的专项项目与细分成果
% \end{itemize} \\
% \cline{2-4}
% & \multirow{1}{=}{中国科学院信工所}
% & \multirow{1}{=}{虎嵩林团队~\cite{huSonglinUCAS2026}}
% & \begin{itemize}[tabitem]
% \item 研究大数据、NLP、知识图谱与大模型全链路安全
% \item 推出Galexy（星河）大模型安全测评平台，提出“攻、检、防”循环评测体系
% \item 覆盖不安全指令调用、具身智能体恶意行为等智能体典型风险~\cite{huSonglinGalexy2024}
% \end{itemize} \\
% \cline{2-4}
% & \multirow{1}{=}{上海AI实验室}
% & \multirow{1}{=}{安全可信AI团队}
% & \begin{itemize}[tabitem]
% \item 2026年发布产业级安全智能体操作系统“书安”Intern-Shannon，安全即服务产业落地~\cite{intershannon2026}
% \item 推出SafeWork-R1，探索智能能力与安全约束协同演化路径~\cite{safework2025}
% \item 依托OpenMMLab、InternLM开源体系提供底层工程支撑
% \end{itemize} \\
% \cline{2-4}
% & \multirow{1}{=}{之江实验室}
% & \multirow{1}{=}{——}
% & \begin{itemize}[tabitem]
% \item 核心方向：智能感知、通用AI、智能计算、智能网络、智能系统基础设施
% \item 代表性工程：南湖计算框架、三体计算星座、021基础大模型
% \item 公开资料暂无专门智能体安全团队与对应细分成果，待补充权威文献完善
% \end{itemize} \\
% \hline
% \end{tabular}
% \end{table}


% \begin{table}[htbp]
% \centering
% \scriptsize
% \renewcommand{\arraystretch}{0.7}
% \caption{我国高校/科研院所智能体安全相关样本占比（2023--2026，近似统计）}
% \label{tab:china_agent_sec_statistics}
% \begin{tabular}{|m{5cm}|m{1.5cm}|m{2cm}|m{5.5cm}|}
% \hline
% \textbf{高校/机构} & \textbf{样本数} & \textbf{占比} & \textbf{相关方向} \\
% \hline
% 清华大学 & 7 & 10.4\% & \begin{itemize}[tabitem]
% \item 可信AI
% \item 安全评测
% \item 治理规范
% \end{itemize} \\
% \hline
% 北京大学 & 6 & 9.0\% & \begin{itemize}[tabitem]
% \item 软件智能体
% \item 系统安全
% \item 安全对齐
% \end{itemize} \\
% \hline
% 中国科学院自动化所 & 5 & 7.5\% & \begin{itemize}[tabitem]
% \item 多模态智能体
% \item 具身智能体安全
% \item 多智能体博弈
% \end{itemize} \\
% \hline
% 中国科学院信工所 & 5 & 7.5\% & \begin{itemize}[tabitem]
% \item 网络安全
% \item 工具调用安全
% \item 大模型安全测评
% \end{itemize} \\
% \hline
% 中国科学院计算所 & 4 & 6.0\% & \begin{itemize}[tabitem]
% \item 系统安全
% \item 可信执行环境
% \item 智能体底层安全
% \end{itemize} \\
% \hline
% 上海交通大学 & 5 & 7.5\% & \begin{itemize}[tabitem]
% \item 大模型安全
% \item 安全评测
% \item 智能体攻防
% \end{itemize} \\
% \hline
% 浙江大学 & 5 & 7.5\% & \begin{itemize}[tabitem]
% \item 隐私计算
% \item 数据安全
% \item 端侧智能体安全
% \end{itemize} \\
% \hline
% 南京大学 & 4 & 6.0\% & \begin{itemize}[tabitem]
% \item 多智能体系统
% \item 软件工程安全
% \item 智能体协同安全
% \end{itemize} \\
% \hline
% 复旦大学 & 4 & 6.0\% & \begin{itemize}[tabitem]
% \item RAG安全
% \item 知识安全
% \item 智能体MCP协议安全
% \end{itemize} \\
% \hline
% 哈尔滨工业大学 & 4 & 6.0\% & \begin{itemize}[tabitem]
% \item NLP安全
% \item 多模态智能体安全
% \item 大模型对齐
% \end{itemize} \\
% \hline
% 中国科学技术大学 & 4 & 6.0\% & \begin{itemize}[tabitem]
% \item 网络安全
% \item 可信系统
% \item 智能体底层安全
% \end{itemize} \\
% \hline
% 北航、北邮、西电、华中科大 & 8 & 11.9\% & \begin{itemize}[tabitem]
% \item 工程化安全
% \item 通信网络安全
% \item 智能体落地安全
% \end{itemize} \\
% \hline
% 南开、同济、深大、南科大、国防科大等 & 6 & 9.0\% & \begin{itemize}[tabitem]
% \item 行业场景安全
% \item 高可靠智能体系统
% \item 垂直领域合规
% \end{itemize} \\
% \hline
% \end{tabular}
% \end{table}

% \begin{table}[htbp]
% \centering
% \caption{我国智能体安全研究方向样本占比}
% \begin{tabular}{p{5cm}p{2cm}p{7cm}}
% \hline
% 研究方向 & 占比 & 说明 \\
% \hline
% 大模型内容安全与测评 & 24\% & 价值观、幻觉、越狱、备案测评 \\
% 网络安全智能体与攻防自动化 & 20\% & 漏洞挖掘、威胁研判、安全运营 \\
% 工具/插件/工作流安全 & 16\% & 插件准入、MCP、工作流审计 \\
% RAG、知识库与记忆安全 & 14\% & 知识库污染、长期记忆、数据脱敏 \\
% 行业智能体治理 & 16\% & 政务、金融、工业、医疗、教育 \\
% 身份授权、沙箱与运行审计 & 10\% & 企业权限、日志、可信执行 \\
% \hline
% \end{tabular}
% \end{table}

% \begin{figure}[htbp]
% \centering
% \includegraphics[width=0.48\textwidth]{figures/china_academic_share.png}
% \includegraphics[width=0.48\textwidth]{figures/china_direction_share.png}
% \caption{我国高校/院所样本占比与研究方向占比}
% \end{figure}


% \subsection{全球学术研究格局小结：对比}
% 综合来看，全球学术研究已形成相对清晰的区域分工。北美研究更集中于前沿攻防、红队评测、工具协议安全、记忆攻击和真实部署风险；欧洲研究更强调合规、伦理、隐私、问责和人类监督；亚太与我国研究则更侧重政策牵引、行业场景、标准建设和工程应用。三类路径共同推动了智能体安全研究的发展，但整体仍呈现“风险点发现快、体系化整合慢；运行态研究强、全生命周期研究弱；技术护栏研究多、治理证据链研究少；产品能力披露多、第三方可验证评测少”的特征。

% 进一步比较可以发现，全球学术研究的差异并不只是地域差异，而是问题定义方式的差异。北美研究通常从“系统如何被攻破”出发，强调攻击样本、红队评测、真实部署漏洞和安全基准，因此更容易发现新型风险。欧洲研究通常从“系统如何被合法、透明、可问责地使用”出发，强调权利保护、责任分配、风险评估和人类监督，因此更容易形成治理框架。我国和亚太研究通常从“系统如何在真实行业场景中落地”出发，强调数据边界、业务权限、平台责任和标准化测评，因此更接近政企应用需求。这三类路径共同构成智能体安全研究的知识来源，但尚未在同一理论框架内完成整合。

% 从研究对象看，当前全球研究内容已经覆盖 prompt、工具、MCP、RAG、长期记忆、多智能体协同、Web Agent、代码执行和身份授权等关键对象，但大多数研究仍以单一对象为中心。例如，工具投毒研究能够揭示工具描述和返回值的污染风险，却往往较少同时讨论工具被调用前的业务授权、调用中的沙箱约束和调用后的审计证据；长期记忆研究能够揭示跨会话污染风险，却较少与数据生命周期、用户删除权、退役清理和合规留痕结合；多智能体安全研究能够说明恶意节点和错误上下文的传播问题，却尚未充分讨论子任务权限分配、协作责任归因和组织级治理。由此可见，现有研究的短板不是缺少风险点，而是缺少贯穿能力组件、工程流程和责任主体的系统化解释。

% 从方法论看，北美红队和基准研究的优势在于可复现攻击样本和可度量指标，但指标多集中于攻击成功率、泄露率、拒答率、任务成功率等即时结果。欧洲治理研究的优势在于规范性和责任分析，但很多要求尚未转化为可执行的技术控制。我国行业研究的优势在于场景复杂、业务链条长、合规约束强，但公开数据集、开放基准和跨机构复现实验仍不足。

\section{全球AI智能体安全主流产品}
全球AI智能体安全的产业落地层面侧重工程化风控、运行时防护、生态治理与合规适配，形成了学术机理牵引、产业产品落地、区域规则约束的递进式发展体系。
北美、欧洲、我国三大区域形成了差异化的产业落地路径，本节主要围绕北美、欧洲及我国的主流AI智能体安全产品进行调研与分析。
% 北美依托基础模型原生创新与云生态优势，主导智能体安全底层协议、开发者标准与全栈防护体系建设；欧洲立足数据主权与合规监管要求，聚焦高可信、可解释、私有化的主权级智能体安全方案；中国依托大模型产业规模化迭代与政企场景刚需，形成以云平台Agent开发底座、安全网关、行业风控为核心的工程化落地体系。

\subsection{北美主流AI智能体安全产品}
北美是全球AI智能体技术研发、产品落地与安全体系构建的核心策源地，依托头部科技企业、开源社区与云原生技术生态，形成覆盖基础模型、开发工具、算力底座、行业部署全链路的智能体产业矩阵\cite{openaichatgptagent2025,nvidiaagenttoolkit2026}。
北美头部机构围绕内生对齐、运行时隔离、权限治理、链路可观测、硬件可信底座方向，推出标准化安全组件、企业级管控套件与合规基准规范，形成成熟完备的工程化落地体系\cite{anthropicresponsiblescaling2025,googletoolgovernance2025}。
本节梳理北美主流AI智能体安全产品、技术文档与工程落地统计，主要发现如下（详见表\~ref{tab:north-america-agent-security-products}）：

\begin{itemize}[tabitem]
\item \textbf{产业形成全栈闭环格局，头部企业分层主导模型、云、算力、应用四级生态。}
% 北美AI智能体安全由头部商业科技企业主导。
OpenAI依托GPT-4o企业版、Stateful Runtime环境管控上层通用智能体安全，通过Enterprise连接器白名单、写入操作分级审批机制规范开发者与企业工作流权限，实现通用智能体工作流安全管控闭环\cite{openaichatgptagent2025,openaiagentssdk2026}；Anthropic基于Claude系列模型与MCP协议，标准化定义企业私有数据源、内网插件的访问边界，筑牢企业私有场景智能体交互安全基线\cite{mcpsafetyaudit2025}；Google依托Google Cloud Agent身份体系与Project Mariner云隔离架构，实现智能体运行环境与本地基础设施物理隔离，支撑云原生智能体规模化安全部署\cite{googleagentidentity2026,googletoolgovernance2025}；Meta开源Llama系列模型配套Llama Guard、Prompt Guard防护组件，主导社区级安全评测基准 CyberSecEval建设，构建开源生态智能体安全防护体系；NVIDIA基于GPU机密计算域与Isaac仿真平台，提供端侧可信执行环境与物理机器人智能体风险验证能力，补齐硬件层与具身智能体安全短板\cite{nvidiaconfidentialcomputing2026,nvidiaisaacsim2026}；Cursor、Perplexity等垂直厂商针对代码生成、检索增强场景，补充细粒度API调用风控与检索污染检测能力，安全能力自上而下覆盖算法、软件、部署、硬件全层级，完善细分场景安全防护体系\cite{perplexityenterprise2026,cursorsecurity2026}。

\item \textbf{安全范式全面前置，构建模型层预防护与运行时强管控的防护体系。}
% 北美工程化安全范式采用前置内生安全架构，构建模型层预防护与运行时强管控的防护体系。
模型层方面，Anthropic与OpenAI将禁止越权、数据最小披露、恶意指令拒绝等价值约束嵌入预训练、RLHF与推理全过程，搭载ASL-3实时安全分类器拦截高危请求，实现模型内生安全对齐\cite{openaichatgptagentsystemcard2025}；运行时层面，主流产品采用容器化沙箱、独立虚拟机隔离智能体代码执行与系统调用，如Anthropic Claude Cowork受限执行容器、OpenAI代码解释器隔离环境，从执行环境层面规避恶意操作风险\cite{anthropiccomputeruse2026,openaiagentkit2025}。
% 流程层面，通过SDK强制工作流编排、高危写入操作人工审批、会话级不可篡改日志审计，对工具调用参数、返回数据、访问时序进行全链路记录，日志独立存储且禁止智能体自身读取，从源头阻断上下文劫持、隐性提示注入与持久化漏洞驻留风险，实现智能体全生命周期安全可控\cite{openaiobservability2026,cursorhooks2025}。

\item \textbf{双轨体系差异化落地，以最小权限治理 + 软硬协同底座支撑规模化企业部署：}
北美形成分工明确的闭源 / 开源双轨安全体系。闭源阵营（OpenAI、Anthropic、Google）面向政企交付私有化部署套件、SOC 2/ISO 27001合规管控模块与闭环漏洞响应流程，提供租户级隔离、数据静态与传输加密、年度渗透测试服务，满足高合规场景落地要求\cite{anthropicclaudecode2025,googlevertexagentbuilder2026}；开源阵营（Meta、NVIDIA）开放安全护栏代码、评测数据集与NeMo Guardrails 工具链，降低中小开发者安全准入门槛，赋能开源生态智能体安全标准化建设\cite{nvidianemoguardrails2026,metallamafirewall2026}。企业治理层面统一以最小权限与可观测性为核心，基于IAM条件访问控制、细粒度角色权限、全链路tracing追踪跨服务、跨API调用行为，精准管控连接器外联、内网数据库访问等高风险入口\cite{googleagentidentity2026,openaiobservability2026}；硬件层面联动 NVIDIA 机密算力、Google 物理隔离云基座，结合科学智能体漏洞挖掘与机器人攻防仿真，构建纯软件防护无法实现的底层安全壁垒，支撑工业控制、无人系统、科学计算等高风险场景规模化落地\cite{nvidiaconfidentialcomputing2026,googlecoscientist2025}。
\end{itemize}

% 依赖项（必须放在主文档导言区，即 \begin{document} 之前）：
% \usepackage{tabularray}
% \UseTblrLibrary{varwidth}
% \usepackage{enumitem}
\begingroup
% \nopagebreak[4]
\enlargethispage{1cm}
\footnotesize
\centering

\DefTblrTemplate{conthead-text}{default}{（续表）}
\DefTblrTemplate{contfoot-text}{default}{续下页}

\begin{longtblr}[
  caption = {北美主流AI智能体安全产品与落地方向},
  label   = {tab:north-america-agent-security-products},
]{
  width   = \linewidth,
  colspec = {
    |Q[wd=0.8cm,l,m]
    |Q[wd=1.0cm,l,m]
    |Q[wd=2.cm,l,m]
    |Q[wd=0.57\linewidth,l,m]
    |X[1,l,m]|
  },
  cells   = {l,m},
  colsep  = 3pt,
  row{1}  = {font=\bfseries},
  rowhead = 1,
  % 每行内容与上下边框之间各保留 5pt 空间。
  rowsep  = 5pt,
  measure = vbox,
  stretch = -1,
}
\hline
\textbf{国家} & \textbf{企业} & \textbf{所属产业领域} & \textbf{安全研究方向} & \textbf{产品/成果} \\
\hline
美国
& OpenAI
&
\begin{itemize}[leftmargin=*,topsep=0pt,partopsep=0pt,itemsep=1pt,parsep=0pt]
    \item 通用基础大模型与ChatGPT终端应用
    \item 面向开发者的智能体API与软件开发套件
    \item 企业级智能体构建框架与标准化评测基准
    \item 智能体对齐技术与安全防护体系研究
\end{itemize}
& 
% \begin{itemize}[leftmargin=*,topsep=0pt,partopsep=0pt,itemsep=1pt,parsep=0pt]
%     \item \textbf{终端用户自主任务智能体的网页、文件、连接器与代码执行风险控制。}
%     未来趋势是从“是否拒答”转向“能否在真实行动链中控制风险”：浏览器、连接器、文件和终端会成为统一任务环境，提示注入、敏感数据外泄、误购买/误发送、越权访问会成为核心风险面。研究重点应放在上下文级权限最小化、敏感动作显式确认、连接器数据隔离、任务过程可中断/可接管，以及对网页与文件中恶意指令的持续检测。
%     \item \textbf{开发者智能体的工具编排、状态管理、可观测性与数据隔离。}
%     开发者智能体会走向并行、多代理、长周期后台执行，安全边界不能只靠模型本身，而要依赖沙箱、工作区隔离、网络策略、工具调用审计、测试证据和可回放轨迹。未来研究应关注工具仿真与部署仿真、代码仓库状态一致性、跨任务状态污染、供应链/README 注入、以及从终端日志和测试结果中构建可验证的智能体行为证据链。
%     \item \textbf{企业低代码智能体的流程治理、护栏配置、评测闭环与连接器管理。}
%     企业场景会更强调“可配置、可审计、可分级授权”的治理能力：不同部门、用户、数据源和业务流程需要不同的护栏策略，而不是统一的全局拒绝规则。未来趋势是把安全能力产品化为策略编排、连接器权限分层、客户自主管控、隐私保护检测、部署前仿真评测和上线后监测闭环，让低代码智能体既能接入真实业务系统，又能把高风险动作锁在可追责、可回滚、可验证的流程里。
%     \end{itemize}
    \begin{itemize}[leftmargin=*,topsep=0pt,partopsep=0pt,itemsep=1pt,parsep=0pt]
    \item 研究终端用户智能体在网页、文件、连接器、代码执行场景下的风险防控方法。
    \item 开展开发者智能体工具编排、状态管控、运行可观测性与数据隔离机制研究。
    \item 研究企业低代码智能体流程治理、安全护栏配置、评测闭环和连接器管理方案。\\
    \textbf{未来研究方向趋势：}
    \item 研究智能体安全评测与部署仿真技术，提前排查智能体越权、误操作与数据泄露等现实问题~\cite{openaideploymentsimulation2026}
    % \item 提示注入与跨上下文攻击防护：随着智能体频繁读取网页、邮件、文档和代码仓库，外部内容中的恶意指令会成为核心攻击面，需要研究跨来源内容隔离、指令可信度判定和上下文污染检测机制。~\cite{openaichatgptagent2025,ma2026autodojo}
    % \item 高风险能力的分级访问与持续监控：对于网络安全、生物安全、代码执行等高风险能力，未来重点不只是模型拒答，而是建立用户分级、场景分级、工具权限分级和上线后异常行为监测体系。~\cite{openaipreparedness2025,openairosalind2026}
    % \item 智能体行为可解释与责任追踪：长链路任务中，系统需要记录智能体为何调用某个工具、使用了哪些数据、产生了哪些外部影响，从而支持审计、复盘、合规和责任划分。~\cite{openaitoolsagents2025,openaiagentkit2025}
    % \item 多智能体协作中的安全边界：未来企业和开发场景会出现多个智能体并行协作，需要研究智能体之间的权限传递、任务拆分、共享状态污染、结果互信和冲突仲裁机制。~\cite{openaitoolsagents2025,xu2025trustparadox}
    % \item 面向连接器生态的数据治理：智能体连接企业 SaaS、知识库、代码库和本地文件后，连接器会成为新的安全边界，需要重点研究细粒度授权、最小可见数据、敏感字段脱敏和跨系统数据流追踪。~\cite{openaiagentkit2025,hou2025mcp,wang2025mcptox}
    \item 研究多智能体协作环境的安全边界，攻克智能体间权限传递、任务划分、状态污染、可信验证和冲突仲裁关键技术。~\cite {openaitoolsagents2025,xu2025trustparadox}
    \item 针对连接器生态开展数据治理研究，实现细粒度授权、最小数据访问、敏感字段脱敏和跨系统数据流追踪~\cite {openaiagentkit2025,hou2025mcp,wang2025mcptox}
    \end{itemize}
& 
    \begin{itemize}[leftmargin=*,topsep=0pt,partopsep=0pt,itemsep=1pt,parsep=0pt]
        \item ChatGPT‑Agent及其系统安全卡，落实用户确认、运行监督、连接器权限管控以及提示注入防御机制~\cite{openaichatgptagent2025,openaichatgptagentsystemcard2025}
        \item Responses API与Agents‑SDK，提供工具调用、运行追踪、智能体任务转交、安全约束以及人工审批能力~\cite{openaiagentssdk2026,openaiobservability2026}
        \item AgentKit套件，涵盖可视化编排Agent‑Builder、前端对话组件ChatKit、评测工具Evals和安全防护模块~\cite{openaiagentkit2025,openaiagentssdk2026}
    \end{itemize}\\
\cline{2-5}
& Anth-\newline ropic
& \begin{itemize}[leftmargin=*,topsep=0pt,partopsep=0pt,itemsep=1pt,parsep=0pt]
    \item Claude 基础模型以及企业级协作落地应用
    \item 工具调用、计算机使用与代码智能体
    \item AI 安全、对齐与开放协议生态
  \end{itemize}
& \begin{itemize}[leftmargin=*,topsep=0pt,partopsep=0pt,itemsep=1pt,parsep=0pt]
    \item 模型内生对齐、宪法式约束与前沿能力分级治理
    \item 工具型智能体的 schema约束、沙箱化计算机使用与人类授权管控
    \item 企业数据源接入方式、代码智能体插件化与连接器权限边界 \\ 
    \textbf{未来研究方向趋势：}
    \item 安全管控由模型静态对齐转向智能体运行时控制平面，依托沙箱隔离、MCP 私网隧道、全链路审计，分离推理模块与执行环境，由企业管控数据和资源边界~\cite {anthropicmanagedagentsinfra2026,anthropicsandboxing2025}
    \item 智能体启用独立身份体系，摒弃继承用户全部权限的模式，采用任务级临时授权与双重鉴权机制约束敏感操作~\cite {claudeagentidentity2026,debenedetti2025camel}
    \item 基于委派图和信任传播规范多智能体协作，落实最小权限原则，对跨‑Agent 消息、权限转交全程溯源，防范安全风险链式扩散~\cite {anthropicmultiagent2025,claudedreaming2026,xu2025trustparadox}
    %
    % \item 将安全重心从模型对齐转向智能体运行时控制平面。
    % 未来安全架构将围绕沙箱、出站网络策略、边界凭据注入、MCP 私网隧道和全链路审计展开，使智能体的“推理大脑”与“执行环境”解耦，并由企业掌握代码、文件、网络和数据边界。~\cite{anthropicmanagedagentsinfra2026,anthropicsandboxing2025}
    % \item 让智能体成为具有独立身份的非人主体。 权限模型将由“智能体继承用户全部权限”转向独立智能体身份、任务级临时授权和双重鉴权，即只有智能体自身权限与请求者权限同时满足时才允许执行敏感操作。~\cite{claudeagentidentity2026,debenedetti2025camel}
    % \item 多智能体安全将围绕委派图和信任传播展开。 每个子智能体应获得独立上下文、最小必要信息和任务限定工具权限，并对共享文件、跨智能体消息、权限转授和结果合并建立来源追踪，避免单个智能体受攻击后沿协作链横向扩散。~\cite{anthropicmultiagent2025,claudedreaming2026,xu2025trustparadox}
    % \item 将提示注入防护演化为信息流与能力流隔离。 单纯依靠模型识别恶意文本难以提供稳定安全保证，未来系统将显式区分可信指令与不可信数据，并在工具调用层实施数据来源标记、能力约束、敏感信息流检查和执行前策略验证。~\cite{anthropicpromptinjection2025,debenedetti2025camel}
    % \item 将持久记忆成为新的安全数据平面。 智能体记忆需要具备来源证明、分级读写权限、版本控制、回滚、删除和跨智能体隔离能力，同时防范恶意网页或文件植入休眠信息、在后续会话中触发敏感行为的延迟记忆投毒。~\cite{claudememory2026,claudedreaming2026,pulipaka2026sleepermemory}
    %  
    % \item 网络安全将进入自主攻防速度竞争。 漏洞发现、利用、分诊和修复周期会明显缩短，防御体系将逐步采用隔离运行的安全智能体持续扫描代码与资产、生成并验证补丁，但高风险利用验证和生产变更仍需要独立检查与人工决策点。~\cite{claudeaicyberoffense2026}
    \end{itemize}
& \begin{itemize}[leftmargin=*,topsep=0pt,partopsep=0pt,itemsep=1pt,parsep=0pt]
    \item Constitutional AI、Claude Constitution、Claude系统卡与 Responsible Scaling Policy ~\cite{anthropicclaudeconstitution2025,anthropicresponsiblescaling2025}
    \item Claude Tool Use、Computer Use、Messages API工具调用机制 ~\cite{anthropictooluse2026,anthropiccomputeruse2026}
    \item Model Context Protocol、Claude Code 与企业级 MCP治理实践 ~\cite{anthropicclaudecode2025,mcpsafetyaudit2025}
    \end{itemize} \\
\cline{2-5}
& Google \newline  Google DeepMind \newline  Google Cloud
& \begin{itemize}[leftmargin=*,topsep=0pt,partopsep=0pt,itemsep=1pt,parsep=0pt]
    \item 搜索、网页与云原生企业应用
    \item Gemini智能体平台与云安全治理
    \item 科学智能体与漏洞攻防研究
  \end{itemize}
& \begin{itemize}[leftmargin=*,topsep=0pt,partopsep=0pt,itemsep=1pt,parsep=0pt]
    \item 云原生智能体的构建、部署、观测、评测与生命周期治理
    \item 企业工具/API注册、Agent身份、最小权限与跨服务访问控制
    \item 科学推理与安全攻防智能体的受控实验、漏洞发现和结果可复现
    \textbf{未来研究方向趋势：}
    % \item 智能体控制平面与系统级防御：Google 的趋势不是只依赖模型对齐或拒答，而是把智能体当作可能失误甚至不完全可信的内部执行者，围绕沙箱、权限分级、行为监控、同步阻断和响应时延构建 AI Control 控制平面。~\cite{googleaicontrolroadmap2026,deepmindagentsecurityblog2026,xiang2026systemdefenses}
    %
    \item 智能体安全评测由静态问答范式转向真实任务轨迹、工具调用链与上线监控驱动，重点防控编码、检索、企业流程等高权限场景下的越权、误操作与异常执行风险~\cite {googleaicontrolroadmap2026,deepmindagentsecurityblog2026,debenedetti2024agentdojo}
    \item 智能体安全研究向多层生态体系演进，涵盖单体智能体、多智能体系统与产业级智能体网络，聚焦跨智能体权限传递、身份认证、协作协议与生态韧性治理~\cite {googlethreelayersagentsecurity2026,habler2025securea2a}
    \item 针对科研与编码类高阶智能体开展边界管控研究，约束自主实验、自动编程与科学推理能力，构建权限隔离、执行边界校验、结果复核与风险传播阻断机制~\cite {deepmindalphaevolve2025,gottweis2025coscientist,googleaicontrolroadmap2026}
    %
    % \item 提示注入与外部内容隔离：当智能体频繁读取网页、邮件、文档、代码仓库和检索结果时，安全重点会从“识别恶意文本”转向系统架构层面的来源隔离、可信指令判定、上下文污染检测和高风险工具调用前的强制校验。~\cite{googlegeminicomputeruse2025,debenedetti2024agentdojo,xiang2026systemdefenses}
    %
    % \item Agentic RAG 的可信检索与可审计知识治理：Google Research 的 Agentic RAG 方向显示，企业智能体安全不只是回答准确率问题，而是要判断检索上下文是否充分、来源是否可信、答案是否可追溯，并把知识库更新、遗忘、审计纳入治理闭环。~\cite{googleresearchagenticrag2026,googleresearchsufficientcontext2025,joren2024sufficientcontext,googleunlearningaudit2026}
    %
    % \item 高风险能力的实时监控与分级阻断：随着智能体具备更强的网络安全、代码修改和系统操作能力，未来防护会按危害等级从异步审计升级到实时预防，对高风险动作实施 supervisor 复核、阻断、回滚和人工升级。~\cite{googleaicontrolroadmap2026,deepmindagentsecurityblog2026,googlegeminicomputeruse2025}
    \end{itemize}
& \begin{itemize}[leftmargin=*,topsep=0pt,partopsep=0pt,itemsep=1pt,parsep=0pt]
    \item Vertex AI Agent Builder、Agent Development Kit、Agent Engine 与 Agent Search ~\cite{googlevertexagentbuilder2026,googleadk2026,googleagentsearch2026}
    \item Cloud API Registry 工具治理、Agent Identity、IAM 条件与 Agent Gateway ~\cite{googletoolgovernance2025,googleagentidentity2026}
    \item AI co-scientist、AlphaEvolve、Project Naptime、Big Sleep/CodeMender ~\cite{googlecoscientist2025,googlealphaevolve2025,googlecodemender2025}
    \end{itemize} \\
\cline{2-5}
& Meta
& \begin{itemize}[leftmargin=*,topsep=0pt,partopsep=0pt,itemsep=1pt,parsep=0pt]
    \item 开源/开放权重基础模型
    \item 开源安全护栏与安全评测
    \item 端侧部署与开发者生态
  \end{itemize}
& \begin{itemize}[leftmargin=*,topsep=0pt,partopsep=0pt,itemsep=1pt,parsep=0pt]
    \item 开源模型输入/输出内容安全分类与多语种安全策略适配
    \item 代码生成、网络安全能力与提示注入风险的开放评测
    \item 智能体运行时安全监控、间接提示注入和危险代码风险拦截 \\ 
% \end{itemize}
%
% \vspace{0.5em}
% 2026重点研究方向\cite{meta_risk_review_2026,meta_support_safety_ai_2026,meta_anti_scam_tools_2026}：
% \vspace{1.5em}
%
% \begin{itemize}[leftmargin=*,topsep=0pt,partopsep=0pt,itemsep=1pt,parsep=0pt]
    \textbf{未来研究方向趋势：}
    \item 依托 AI 实现全生命周期风险审查，前置排查合规安全隐患，并对模型运行状态开展持续性风险监测 ~\cite{meta_risk_review_2026,meta_support_safety_ai_2026,meta_anti_scam_tools_2026}
    \item 借助人工智能开展诈骗识别、身份核验与违规内容甄别，兼顾违规检出效率，规避过度管控问题 ~\cite{meta_risk_review_2026,meta_support_safety_ai_2026,meta_anti_scam_tools_2026}
    \item 落实私密对话隔离、敏感数据分隔与用户自主授权机制，厘清 AI 调用私有数据、跨服务交互过程中的隐私边界~\cite{meta_risk_review_2026,meta_support_safety_ai_2026,meta_anti_scam_tools_2026}
    % \item 将风险治理从单项隐私评审扩展为跨公司的统一风险审查，使安全标准、法律要求和专家判断能够一致应用
\end{itemize}

& \begin{itemize}[leftmargin=*,topsep=0pt,partopsep=0pt,itemsep=1pt,parsep=0pt]
    \item Llama Guard、Prompt Guard 与 Llama Responsible Use Guide ~\cite{metapurplellama2026}
    \item CyberSec- Eval、Code Shield（用于不安全代码过滤与网络安全风险测试）~\cite{metapurplellama2026}
    \item LlamaFire- wall、集成 PromptGuard 2、Agent Alignment Checks 与 CodeShield ~\cite{metallamaprotections2026,metallamafirewall2026}
    \end{itemize} \\
\cline{2-5}
& NVIDIA
& \begin{itemize}[leftmargin=*,topsep=0pt,partopsep=0pt,itemsep=1pt,parsep=0pt]
    \item AI算力基础设施与GPU软件栈
    \item 企业级智能体编排与护栏框架
    \item 工业、机器人与物理AI仿真
  \end{itemize}
& \begin{itemize}[leftmargin=*,topsep=0pt,partopsep=0pt,itemsep=1pt,parsep=0pt]
    \item 生产级多智能体工作流的编排、工具连接、评测与性能剖析
    \item 对话式AI与RAG智能体的输入/输出内容护栏和策略约束
    \item 可信算力、模型/数据在用保护与物理AI安全仿真 \\ 
% \end{itemize}
%
% \vspace{0.5em}
% 2026重点研究方向\cite{nvidia_trustworthy_ai_cybersecurity_2026}：
% \vspace{1.5em}

% \begin{itemize}[leftmargin=*,topsep=0pt,partopsep=0pt,itemsep=1pt,parsep=0pt]
    \textbf{未来研究方向趋势：}
    \item 面向长期运行、自主进化智能体，研究沙箱隔离、最小权限、网络与文件访问控制以及受策略约束的运行环境\cite{nvidia_sandboxing_agentic_workflows_2026}
    % \item 防御代码仓库、AGENTS.md、MCP 响应等外部内容引发的间接提示注入，建立不可信内容传播跟踪和执行前检查机制
    \item 将安全边界下沉至 AI 工厂基础设施，以硬件信任根、机密计算、零信任网络和 DPU 隔离保护模型、数据、软件供应链与智能体\cite{nvidia_doca_agentic_ai_security_2026}
    \item 推进智能体技能与工具的能力治理，包括来源证明、安全扫描、机器可读技能卡、加密签名及发布审核\cite{nvidia_verified_agent_skills_2026}
    % \item 研究 AI 生成代码和命令的安全执行，通过语法约束、执行隔离、敏感操作审批及可审计日志限制破坏范围
    % \item 面向机器人智能体，研究新技能和控制原语的安全提出、验证与纳入，以及仿真到现实迁移和长期记忆安全\cite{nvidia_aspire_2026}
\end{itemize}



& \begin{itemize}[leftmargin=*,topsep=0pt,partopsep=0pt,itemsep=1pt,parsep=0pt]
    \item NVIDIA NeMo Agent Toolkit，支持MCP/A2A、workflow、evaluation 与 profiler ~\cite{nvidiaagenttoolkit2026}
    \item NVIDIA NeMo Guardrails、NIM与安全RAG部署参考 ~\cite{nvidianemoguardrails2026}
    \item NVIDIA Confidential Computing、Isaac Sim/Isaac Lab、Cosmos World Foundation Models ~\cite{nvidiaconfidentialcomputing2026,nvidiaisaacsim2026,nvidiaisaaclab2026,nvidiacosmos2025}
    \end{itemize} \\
\cline{2-5}
& Perple-\newline xity
& \begin{itemize}[leftmargin=*,topsep=0pt,partopsep=0pt,itemsep=1pt,parsep=0pt]
    \item 搜索原生问答与实时联网检索
    \item 浏览器智能体与网页任务自动化
    \item 企业研究智能体与API平台
  \end{itemize}
& \begin{itemize}[leftmargin=*,topsep=0pt,partopsep=0pt,itemsep=1pt,parsep=0pt]
    \item 联网检索型智能体的来源可信、引用溯源与搜索范围控制
    \item 浏览器智能体的网页上下文处理、跨站任务代理与提示注入风险防护
    \item 企业级多步任务智能体的连接器、访问边界、数据留存和网络策略治理 \\
% \end{itemize}

% \vspace{0.5em}
% 2026重点研究方向\cite{perplexity_secure_intelligence_institute_announcement_2026,perplexity_enterprise_security_2026}：
% \vspace{1.5em}

% \begin{itemize}[leftmargin=*,topsep=0pt,partopsep=0pt,itemsep=1pt,parsep=0pt]
    \textbf{未来研究方向趋势：}
    \item 研究自主智能体认证授权与权限治理体系，明确工具调用、高危操作、数据访问及用户确认的安全边界规范。
    \item 深耕可信 AI、鲁棒机器学习、应用密码学与可用隐私安全技术，搭建校企长期联合科研体系。
    \item 探究大规模多模型智能系统安全机制，实现模型路由、上下文传输与多模型协作的隔离可控、全程溯源与策略一致执行~\cite {perplexity_secure_intelligence_institute_2026}
    % cite{perplexity_secure_intelligence_institute_2026}
\end{itemize}
% • 以 Secure Intelligence Institute 为核心，开展前沿 AI 系统的安全、隐私、可信度和实用防御研究，并把成果直接转化到 Perplexity 系统
% • 重点研究自主智能体的认证、授权和权限治理，确定工具使用、高风险操作、数据访问及用户确认的安全边界
% • 持续研究浏览器智能体在开放网页环境中的提示注入、恶意网页内容、跨站操作及敏感会话风险
% • 推进可用隐私与安全、应用密码学、鲁棒机器学习和可信 AI，并与高校建立长期联合研究网络
% • 加强智能体及开发者环境的软件供应链防御，包括只读终端扫描、依赖与本地文件检查以及快速事件响应
% • 研究大规模多模型智能系统的安全，使模型路由、上下文传递及多模型协作过程具备隔离、可追踪与一致策略执行能力
& \begin{itemize}[leftmargin=*,topsep=0pt,partopsep=0pt,itemsep=1pt,parsep=0pt]
    \item Perplexity Sonar API、Search API、Pro Search与引用型答案机制 ~\cite{perplexitysonar2026,perplexitysearchapi2026,perplexityprivacysecurity2026}
    \item Comet Browser/Comet Assistant（结合网页智能体安全评测研究）~\cite{perplexitycomet2025,bravecometpromptinjection2025}
    \item Perplexity Enterprise、Perplexity Computer、Enterprise connectors 与 Computer Network Firewall Policy ~\cite{perplexityenterprise2026,perplexitycomputerenterprise2026}
    \end{itemize} \\
\cline{2-5}
& Cursor / Anysphere
& \begin{itemize}[leftmargin=*,topsep=0pt,partopsep=0pt,itemsep=1pt,parsep=0pt]
    \item AI IDE与代码智能体
    \item 云端/本地软件工程自动化
    \item 企业研发安全与MCP工具接入
  \end{itemize}
& \begin{itemize}[leftmargin=*,topsep=0pt,partopsep=0pt,itemsep=1pt,parsep=0pt]
    \item IDE内代码智能体的本地工程访问、终端/浏览器操作与命令审查
    \item 云端代码智能体的隔离执行环境、PR产物验证与企业内网部署
    \item 企业级代码数据保护、模型/MCP/网络访问控制和审计合规
    \end{itemize}
& \begin{itemize}[leftmargin=*,topsep=0pt,partopsep=0pt,itemsep=1pt,parsep=0pt]
    \item Cursor Agent、Plan Mode、Browser Controls、Hooks（用于计划、审计、命令阻断与上下文控制）~\cite{cursor17changelog2025,cursorplanmode2025,cursorhooks2025}
    \item Cursor Cloud Agents、Bugbot、Self-hosted Cloud Agents ~\cite{cursorcloudagents2025,cursorbugbot2025,cursorselfhostedcloudagents2026}
    \item Cursor Enterprise、Privacy Mode、SSO/SAML、RBAC、SCIM、audit logs、repository/model/MCP controls ~\cite{cursorsecurity2026,cursorenterprise2026}
    \end{itemize} \\
\hline
\end{longtblr}
\endgroup


% \begin{center}
% % \nopagebreak[4]
% \enlargethispage{1cm}
% \footnotesize
% \begin{longtable}{

% |>{\rule{0pt}{2pt}}m{0.8cm}|
% >{\rule{0pt}{2pt}}m{1.2cm}|
% >{\rule{0pt}{2pt}}m{2.8cm}|
% >{\rule{0pt}{2pt}}m{0.45\linewidth}|
% >{\rule{0pt}{2pt}}m{\dimexpr\linewidth - 0.8cm - 1.2cm - 2.8cm - 0.3\linewidth - 10\tabcolsep\relax}|
% }
% \caption{北美主流AI智能体安全产品与落地方向}\label{tab:north-america-agent-security-products}\\
% \hline
% \textbf{国家} & \textbf{企业} & \textbf{所属产业领域} & \textbf{安全研究方向} & \textbf{产品/成果} \\
% \hline
% \endfirsthead
% \multicolumn{5}{l}{\tablename\ \thetable\ -- 续上页}\\
% \hline
% \textbf{国家} & \textbf{企业} & \textbf{所属产业领域} & \textbf{安全研究方向} & \textbf{产品/成果} \\
% \hline
% \endhead
% \hline
% \multicolumn{5}{r}{续下页}\\
% \endfoot
% \hline
% \endlastfoot


% \multirow{12}{=}{美国}
% & \multirow{3}{=}{OpenAI}
% &
% \begin{itemize}[tabitem]
%     \item 通用基础大模型与ChatGPT终端应用
%     \item 面向开发者的智能体API与软件开发套件
%     \item 企业级智能体构建框架与标准化评测基准
%     \item 智能体对齐技术与安全防护体系研究
% \end{itemize}
% & 
% % \begin{itemize}[tabitem]
% %     \item \textbf{终端用户自主任务智能体的网页、文件、连接器与代码执行风险控制。}
% %     未来趋势是从“是否拒答”转向“能否在真实行动链中控制风险”：浏览器、连接器、文件和终端会成为统一任务环境，提示注入、敏感数据外泄、误购买/误发送、越权访问会成为核心风险面。研究重点应放在上下文级权限最小化、敏感动作显式确认、连接器数据隔离、任务过程可中断/可接管，以及对网页与文件中恶意指令的持续检测。
% %     \item \textbf{开发者智能体的工具编排、状态管理、可观测性与数据隔离。}
% %     开发者智能体会走向并行、多代理、长周期后台执行，安全边界不能只靠模型本身，而要依赖沙箱、工作区隔离、网络策略、工具调用审计、测试证据和可回放轨迹。未来研究应关注工具仿真与部署仿真、代码仓库状态一致性、跨任务状态污染、供应链/README 注入、以及从终端日志和测试结果中构建可验证的智能体行为证据链。
% %     \item \textbf{企业低代码智能体的流程治理、护栏配置、评测闭环与连接器管理。}
% %     企业场景会更强调“可配置、可审计、可分级授权”的治理能力：不同部门、用户、数据源和业务流程需要不同的护栏策略，而不是统一的全局拒绝规则。未来趋势是把安全能力产品化为策略编排、连接器权限分层、客户自主管控、隐私保护检测、部署前仿真评测和上线后监测闭环，让低代码智能体既能接入真实业务系统，又能把高风险动作锁在可追责、可回滚、可验证的流程里。
% %     \end{itemize}
%     \begin{itemize}[tabitem]
%     \item 终端用户自主任务智能体的网页、文件、连接器与代码执行风险控制。

%     \item 开发者智能体的工具编排、状态管理、可观测性与数据隔离。
   
%     \item 企业低代码智能体的流程治理、护栏配置、评测闭环与连接器管理。

%     未来研究方向趋势：
%     \item 智能体安全评测与部署仿真：未来需要构建更接近真实业务流程的评测体系，用历史任务、真实工具调用轨迹和模拟部署环境提前发现智能体在上线后的越权、误操作与数据泄露风险。~\cite{openaideploymentsimulation2026}
%     % \item 提示注入与跨上下文攻击防护：随着智能体频繁读取网页、邮件、文档和代码仓库，外部内容中的恶意指令会成为核心攻击面，需要研究跨来源内容隔离、指令可信度判定和上下文污染检测机制。~\cite{openaichatgptagent2025,ma2026autodojo}
%     % \item 高风险能力的分级访问与持续监控：对于网络安全、生物安全、代码执行等高风险能力，未来重点不只是模型拒答，而是建立用户分级、场景分级、工具权限分级和上线后异常行为监测体系。~\cite{openaipreparedness2025,openairosalind2026}
%     % \item 智能体行为可解释与责任追踪：长链路任务中，系统需要记录智能体为何调用某个工具、使用了哪些数据、产生了哪些外部影响，从而支持审计、复盘、合规和责任划分。~\cite{openaitoolsagents2025,openaiagentkit2025}
%     \item 多智能体协作中的安全边界：未来企业和开发场景会出现多个智能体并行协作，需要研究智能体之间的权限传递、任务拆分、共享状态污染、结果互信和冲突仲裁机制。~\cite{openaitoolsagents2025,xu2025trustparadox}
%     \item 面向连接器生态的数据治理：智能体连接企业 SaaS、知识库、代码库和本地文件后，连接器会成为新的安全边界，需要重点研究细粒度授权、最小可见数据、敏感字段脱敏和跨系统数据流追踪。~\cite{openaiagentkit2025,hou2025mcp,wang2025mcptox}

%     \end{itemize}
% & 
%     \begin{itemize}[tabitem]
%         \item ChatGPT‑Agent及其系统安全白皮书，落实用户确认、运行监督、连接器权限管控以及提示注入防御机制~\cite{openaichatgptagent2025,openaichatgptagentsystemcard2025}
%         \item Responses API与Agents‑SDK，提供工具调用、运行追踪、智能体任务转交、安全约束以及人工审批能力~\cite{openaiagentssdk2026,openaiobservability2026}
%         \item AgentKit套件，涵盖可视化编排Agent‑Builder、前端对话组件ChatKit、评测工具Evals和安全防护模块~\cite{openaiagentkit2025,openaiagentssdk2026}
%     \end{itemize}\\
% \cline{2-5}
% & \multirow{3}{=}{Anth-\\ropic}
% & \begin{itemize}[tabitem]
%     \item Claude 基础模型与企业协作应用
%     \item 工具调用、计算机使用与代码智能体
%     \item AI安全、对齐与开放协议生态
%   \end{itemize}
% & \begin{itemize}[tabitem]
%     \item 模型内生对齐、宪法式约束与前沿能力分级治理
%     \item 工具型智能体的 schema约束、沙箱化计算机使用与人类授权
%     \item 企业数据源接入、代码智能体插件化与连接器权限边界
%     \item 将安全重心从模型对齐转向智能体运行时控制平面。 未来安全架构将围绕沙箱、出站网络策略、边界凭据注入、MCP 私网隧道和全链路审计展开，使智能体的“推理大脑”与“执行环境”解耦，并由企业掌握代码、文件、网络和数据边界。~\cite{anthropicmanagedagentsinfra2026,anthropicsandboxing2025}
%     \item 让智能体成为具有独立身份的非人主体。 权限模型将由“智能体继承用户全部权限”转向独立智能体身份、任务级临时授权和双重鉴权，即只有智能体自身权限与请求者权限同时满足时才允许执行敏感操作。~\cite{claudeagentidentity2026,debenedetti2025camel}、
%     % \item 将提示注入防护演化为信息流与能力流隔离。 单纯依靠模型识别恶意文本难以提供稳定安全保证，未来系统将显式区分可信指令与不可信数据，并在工具调用层实施数据来源标记、能力约束、敏感信息流检查和执行前策略验证。~\cite{anthropicpromptinjection2025,debenedetti2025camel}
%     % \item 将持久记忆成为新的安全数据平面。 智能体记忆需要具备来源证明、分级读写权限、版本控制、回滚、删除和跨智能体隔离能力，同时防范恶意网页或文件植入休眠信息、在后续会话中触发敏感行为的延迟记忆投毒。~\cite{claudememory2026,claudedreaming2026,pulipaka2026sleepermemory}
%     \item 多智能体安全将围绕委派图和信任传播展开。 每个子智能体应获得独立上下文、最小必要信息和任务限定工具权限，并对共享文件、跨智能体消息、权限转授和结果合并建立来源追踪，避免单个智能体受攻击后沿协作链横向扩散。~\cite{anthropicmultiagent2025,claudedreaming2026,xu2025trustparadox}
%     %  
%     % \item 网络安全将进入自主攻防速度竞争。 漏洞发现、利用、分诊和修复周期会明显缩短，防御体系将逐步采用隔离运行的安全智能体持续扫描代码与资产、生成并验证补丁，但高风险利用验证和生产变更仍需要独立检查与人工决策点。~\cite{claudeaicyberoffense2026}
%     \end{itemize}
% & \begin{itemize}[tabitem]
%     \item Constitutional AI、Claude Constitution、Claude系统卡与 Responsible Scaling Policy ~\cite{anthropicclaudeconstitution2025,anthropicresponsiblescaling2025}
%     \item Claude Tool Use、Computer Use、Messages API工具调用机制 ~\cite{anthropictooluse2026,anthropiccomputeruse2026}
%     \item Model Context Protocol、Claude Code 与企业级 MCP治理实践 ~\cite{anthropicclaudecode2025,mcpsafetyaudit2025}
%     \end{itemize} \\
% \cline{2-5}
% & \multirow{3}{=}{\makecell[c]{Google \\ Google DeepMind \\ Google Cloud}}
% & \begin{itemize}[tabitem]
%     \item 搜索、网页与云原生企业应用
%     \item Gemini智能体平台与云安全治理
%     \item 科学智能体与漏洞攻防研究
%   \end{itemize}
% & \begin{itemize}[tabitem]
%     \item 云原生智能体的构建、部署、观测、评测与生命周期治理
%     \item 企业工具/API注册、Agent身份、最小权限与跨服务访问控制
%     \item 科学推理与安全攻防智能体的受控实验、漏洞发现和结果可复现
%     未来研究方向趋势：
%     % \item 智能体控制平面与系统级防御：Google 的趋势不是只依赖模型对齐或拒答，而是把智能体当作可能失误甚至不完全可信的内部执行者，围绕沙箱、权限分级、行为监控、同步阻断和响应时延构建 AI Control 控制平面。~\cite{googleaicontrolroadmap2026,deepmindagentsecurityblog2026,xiang2026systemdefenses}

%     \item 面向真实部署轨迹的智能体安全评测：未来评测将从静态问答转向真实任务轨迹、工具调用链和上线监控数据，重点发现智能体在编码、检索、企业流程和高权限操作中的误删、越权、过度执行与异常行为。~\cite{googleaicontrolroadmap2026,deepmindagentsecurityblog2026,debenedetti2024agentdojo}

%     \item 多层智能体生态安全：Google 将 agent security 拆成单体智能体、多智能体系统和社会/产业级网络三层，说明未来研究会更关注智能体间权限传递、身份认证、跨组织协作协议、信誉机制和生态级韧性。~\cite{googlethreelayersagentsecurity2026,habler2025securea2a}

%     % \item 提示注入与外部内容隔离：当智能体频繁读取网页、邮件、文档、代码仓库和检索结果时，安全重点会从“识别恶意文本”转向系统架构层面的来源隔离、可信指令判定、上下文污染检测和高风险工具调用前的强制校验。~\cite{googlegeminicomputeruse2025,debenedetti2024agentdojo,xiang2026systemdefenses}

%     % \item Agentic RAG 的可信检索与可审计知识治理：Google Research 的 Agentic RAG 方向显示，企业智能体安全不只是回答准确率问题，而是要判断检索上下文是否充分、来源是否可信、答案是否可追溯，并把知识库更新、遗忘、审计纳入治理闭环。~\cite{googleresearchagenticrag2026,googleresearchsufficientcontext2025,joren2024sufficientcontext,googleunlearningaudit2026}

%     \item 高能力科研与编码智能体的边界控制：AlphaEvolve 和 AI co-scientist 代表了 Google 对自主发现、自动编程和科学假设生成的推进，未来安全研究需要关注这类智能体的实验权限、代码执行边界、结果验证、错误传播和人类专家复核机制。~\cite{deepmindalphaevolve2025,gottweis2025coscientist,googleaicontrolroadmap2026}

%     % \item 高风险能力的实时监控与分级阻断：随着智能体具备更强的网络安全、代码修改和系统操作能力，未来防护会按危害等级从异步审计升级到实时预防，对高风险动作实施 supervisor 复核、阻断、回滚和人工升级。~\cite{googleaicontrolroadmap2026,deepmindagentsecurityblog2026,googlegeminicomputeruse2025}
%     \end{itemize}
% & \begin{itemize}[tabitem]
%     \item Vertex AI Agent Builder、Agent Development Kit、Agent Engine 与 Agent Search ~\cite{googlevertexagentbuilder2026,googleadk2026,googleagentsearch2026}
%     \item Cloud API Registry 工具治理、Agent Identity、IAM 条件与 Agent Gateway ~\cite{googletoolgovernance2025,googleagentidentity2026}
%     \item AI co-scientist、AlphaEvolve、Project Naptime、Big Sleep/CodeMender ~\cite{googlecoscientist2025,googlealphaevolve2025,googlecodemender2025}
%     \end{itemize} \\
% \cline{2-5}
% & \multirow{3}{=}{Meta}
% & \begin{itemize}[tabitem]
%     \item 开源/开放权重基础模型
%     \item 开源安全护栏与安全评测
%     \item 端侧部署与开发者生态
%   \end{itemize}
% & \begin{itemize}[tabitem]
%     \item 开源模型输入/输出内容安全分类与多语种安全策略适配
%     \item 代码生成、网络安全能力与提示注入风险的开放评测
%     \item 智能体运行时安全监控、间接提示注入和危险代码风险拦截
% \end{itemize}

% \vspace{0.5em}
% 2026重点研究方向\cite{meta_risk_review_2026,meta_support_safety_ai_2026,meta_anti_scam_tools_2026}：
% \vspace{1.5em}


% \begin{itemize}[tabitem]
%     \item 以 AI 驱动的持续风险审查覆盖产品全生命周期，在上线前识别隐私、安全、合规风险，并在运行中持续监测
%     \item 研究 AI 辅助的内容执法、诈骗识别与身份真实性验证，在提高严重违规识别率的同时减少过度执法
%     \item 强化私密 AI 交互、敏感数据隔离和用户可控授权，重点处理 AI 接入私人数据及跨应用服务时的隐私边界
%     % \item 将风险治理从单项隐私评审扩展为跨公司的统一风险审查，使安全标准、法律要求和专家判断能够一致应用
% \end{itemize}

% & \begin{itemize}[tabitem]
%     \item Llama Guard、Prompt Guard 与 Llama Responsible Use Guide ~\cite{metapurplellama2026}
%     \item CyberSecEval、Code Shield（用于不安全代码过滤与网络安全风险测试）~\cite{metapurplellama2026}
%     \item LlamaFirewall、集成 PromptGuard 2、Agent Alignment Checks 与 CodeShield ~\cite{metallamaprotections2026,metallamafirewall2026}
%     \end{itemize} \\
% \cline{2-5}
% & \multirow{3}{=}{NVIDIA}
% & \begin{itemize}[tabitem]
%     \item AI算力基础设施与GPU软件栈
%     \item 企业级智能体编排与护栏框架
%     \item 工业、机器人与物理AI仿真
%   \end{itemize}
% & \begin{itemize}[tabitem]
%     \item 生产级多智能体工作流的编排、工具连接、评测与性能剖析
%     \item 对话式AI与RAG智能体的输入/输出内容护栏和策略约束
%     \item 可信算力、模型/数据在用保护与物理AI安全仿真
% \end{itemize}

% \vspace{0.5em}
% 2026重点研究方向\cite{nvidia_trustworthy_ai_cybersecurity_2026}：
% \vspace{1.5em}

% \begin{itemize}[tabitem]
%     \item 面向长期运行、自主进化智能体，研究沙箱隔离、最小权限、网络与文件访问控制以及受策略约束的运行环境\cite{nvidia_sandboxing_agentic_workflows_2026}
%     \item 防御代码仓库、AGENTS.md、MCP 响应等外部内容引发的间接提示注入，建立不可信内容传播跟踪和执行前检查机制
%     \item 将安全边界下沉至 AI 工厂基础设施，以硬件信任根、机密计算、零信任网络和 DPU 隔离保护模型、数据、软件供应链与智能体\cite{nvidia_doca_agentic_ai_security_2026}
%     \item 推进智能体技能与工具的能力治理，包括来源证明、安全扫描、机器可读技能卡、加密签名及发布审核\cite{nvidia_verified_agent_skills_2026}
%     \item 研究 AI 生成代码和命令的安全执行，通过语法约束、执行隔离、敏感操作审批及可审计日志限制破坏范围
%     \item 面向机器人智能体，研究新技能和控制原语的安全提出、验证与纳入，以及仿真到现实迁移和长期记忆安全\cite{nvidia_aspire_2026}
% \end{itemize}



% & \begin{itemize}[tabitem]
%     \item NVIDIA NeMo Agent Toolkit，支持MCP/A2A、workflow、evaluation 与 profiler ~\cite{nvidiaagenttoolkit2026}
%     \item NVIDIA NeMo Guardrails、NIM与安全RAG部署参考 ~\cite{nvidianemoguardrails2026}
%     \item NVIDIA Confidential Computing、Isaac Sim/Isaac Lab、Cosmos World Foundation Models ~\cite{nvidiaconfidentialcomputing2026,nvidiaisaacsim2026,nvidiaisaaclab2026,nvidiacosmos2025}
%     \end{itemize} \\
% \cline{2-5}
% & \multirow{3}{=}{Perple-\\xity}
% & \begin{itemize}[tabitem]
%     \item 搜索原生问答与实时联网检索
%     \item 浏览器智能体与网页任务自动化
%     \item 企业研究智能体与API平台
%   \end{itemize}
% & \begin{itemize}[tabitem]
%     \item 联网检索型智能体的来源可信、引用溯源与搜索范围控制
%     \item 浏览器智能体的网页上下文处理、跨站任务代理与提示注入风险防护
%     \item 企业级多步任务智能体的连接器、访问边界、数据留存和网络策略治理
% \end{itemize}

% \vspace{0.5em}
% 2026重点研究方向\cite{perplexity_secure_intelligence_institute_announcement_2026,perplexity_enterprise_security_2026}：
% \vspace{1.5em}

% \begin{itemize}[tabitem]
%     \item 重点研究自主智能体的认证、授权和权限治理，确定工具使用、高风险操作、数据访问及用户确认的安全边界
%     \item 推进可用隐私与安全、应用密码学、鲁棒机器学习和可信 AI，并与高校建立长期联合研究网络
%     \item 研究大规模多模型智能系统的安全，使模型路由、上下文传递及多模型协作过程具备隔离、可追踪与一致策略执行能力\cite{perplexity_secure_intelligence_institute_2026}
% \end{itemize}
% % • 以 Secure Intelligence Institute 为核心，开展前沿 AI 系统的安全、隐私、可信度和实用防御研究，并把成果直接转化到 Perplexity 系统
% % • 重点研究自主智能体的认证、授权和权限治理，确定工具使用、高风险操作、数据访问及用户确认的安全边界
% % • 持续研究浏览器智能体在开放网页环境中的提示注入、恶意网页内容、跨站操作及敏感会话风险
% % • 推进可用隐私与安全、应用密码学、鲁棒机器学习和可信 AI，并与高校建立长期联合研究网络
% % • 加强智能体及开发者环境的软件供应链防御，包括只读终端扫描、依赖与本地文件检查以及快速事件响应
% % • 研究大规模多模型智能系统的安全，使模型路由、上下文传递及多模型协作过程具备隔离、可追踪与一致策略执行能力
% & \begin{itemize}[tabitem]
%     \item Perplexity Sonar API、Search API、Pro Search与引用型答案机制 ~\cite{perplexitysonar2026,perplexitysearchapi2026,perplexityprivacysecurity2026}
%     \item Comet Browser/Comet Assistant（结合网页智能体安全评测研究）~\cite{perplexitycomet2025,bravecometpromptinjection2025}
%     \item Perplexity Enterprise、Perplexity Computer、Enterprise connectors 与 Computer Network Firewall Policy ~\cite{perplexityenterprise2026,perplexitycomputerenterprise2026}
%     \end{itemize} \\
% \cline{2-5}
% & \multirow{3}{=}{Cursor / Anysphere}
% & \begin{itemize}[tabitem]
%     \item AI IDE与代码智能体
%     \item 云端/本地软件工程自动化
%     \item 企业研发安全与MCP工具接入
%   \end{itemize}
% & \begin{itemize}[tabitem]
%     \item IDE内代码智能体的本地工程访问、终端/浏览器操作与命令审查
%     \item 云端代码智能体的隔离执行环境、PR产物验证与企业内网部署
%     \item 企业级代码数据保护、模型/MCP/网络访问控制和审计合规
%     \end{itemize}
% & \begin{itemize}[tabitem]
%     \item Cursor Agent、Plan Mode、Browser Controls、Hooks（用于计划、审计、命令阻断与上下文控制）~\cite{cursor17changelog2025,cursorplanmode2025,cursorhooks2025}
%     \item Cursor Cloud Agents、Bugbot、Self-hosted Cloud Agents ~\cite{cursorcloudagents2025,cursorbugbot2025,cursorselfhostedcloudagents2026}
%     \item Cursor Enterprise、Privacy Mode、SSO/SAML、RBAC、SCIM、audit logs、repository/model/MCP controls ~\cite{cursorsecurity2026,cursorenterprise2026}
%     \end{itemize} 
% \end{longtable}
% \end{center}


% \renewcommand{\arraystretch}{0.7}
% \begin{table*}[htbp]
% \centering
% \caption{北美主流 AI 智能体安全产品与落地方向}
% \label{tab:north-america-agent-security-products}
% \fontsize{7pt}{7.8pt}\selectfont
% \setlength{\tabcolsep}{3pt}
% \renewcommand{\arraystretch}{1.08}

% \begin{tabularx}{\textwidth}{|>{\centering\arraybackslash}m{0.07\textwidth}|>{\raggedright\arraybackslash}m{0.12\textwidth}|>{\hspace{0.35em}\raggedright\arraybackslash}m{0.22\textwidth}|>{\hspace{0.35em}\raggedright\arraybackslash}X|}

% \hline
% \multicolumn{1}{|c|}{\rule[-0.75ex]{0pt}{3.0ex}所属区域} &
% \multicolumn{1}{c|}{\rule[-0.75ex]{0pt}{3.0ex}企业名称} &
% \multicolumn{1}{c|}{\rule[-0.75ex]{0pt}{3.0ex}所属产业领域} &
% \multicolumn{1}{c|}{\rule[-0.75ex]{0pt}{3.0ex}安全研究方向与产品/成果} \\

% \hline


% & OpenAI
% &\begin{itemize}[leftmargin=1.35em,labelsep=0.25em,itemsep=0.05em,topsep=0pt,parsep=0pt,partopsep=0pt,label=\textbullet]
%     \item 通用基础模型与 ChatGPT 应用
%     \item 开发者智能体 API 与 SDK
%     \item 企业级智能体构建与评测工具
%   \end{itemize}
% & \begin{tabular}{@{}>{\raggedright\arraybackslash}p{0.40\linewidth}@{\hspace{0.8em}}>{\raggedright\arraybackslash}p{0.54\linewidth}@{}}
%     \rule[-0.35em]{0pt}{1.35em}\textbullet\hspace{0.25em}终端用户自主任务智能体的网页、文件、连接器与代码执行风险控制
%     & \rule[-0.35em]{0pt}{1.35em}\textbullet\hspace{0.25em}ChatGPT agent 及其系统卡，强调用户确认、主动监督、连接器访问限制与提示注入防御 ~\cite{openaichatgptagent2025,openaichatgptagentsystemcard2025} \\[0.15em]
%     \textbullet\hspace{0.25em}开发者智能体的工具编排、状态管理、可观测性与数据隔离
%     & \textbullet\hspace{0.25em}Responses API、Agents SDK，支持工具调用、tracing、handoff、guardrails 与人工审批 ~\cite{openaiagentssdk2026,openaiobservability2026} \\[0.15em]
%     \textbullet\hspace{0.25em}企业低代码智能体的流程治理、护栏配置、评测闭环与连接器管理
%     & \textbullet\hspace{0.25em}AgentKit、Agent Builder、ChatKit、Evals 与 Guardrails ~\cite{openaiagentkit2025,openaiagentssdk2026}\rule[-0.45em]{0pt}{1.25em}
%   \end{tabular} \\
% \cline{2-4}

% & Anthropic
% & \begin{itemize}[leftmargin=1.35em,labelsep=0.25em,itemsep=0.05em,topsep=0pt,parsep=0pt,partopsep=0pt,label=\textbullet]
%     \item Claude 基础模型与企业协作应用
%     \item 工具调用、计算机使用与代码智能体
%     \item AI 安全、对齐与开放协议生态
%   \end{itemize}
% & \begin{tabular}{@{}>{\raggedright\arraybackslash}p{0.40\linewidth}@{\hspace{0.8em}}>{\raggedright\arraybackslash}p{0.54\linewidth}@{}}
%     \textbullet\hspace{0.25em}模型内生对齐、宪法式约束与前沿能力分级治理
%     & \textbullet\hspace{0.25em}Constitutional AI、Claude Constitution、Claude 系统卡与 Responsible Scaling Policy ~\cite{bai2022constitutionalai,anthropicclaudeconstitution2025,anthropicresponsiblescaling2025} \\[0.15em]
%     \textbullet\hspace{0.25em}工具型智能体的 schema 约束、沙箱化计算机使用与人类授权
%     & \textbullet\hspace{0.25em}Claude Tool Use、Computer Use、Messages API 工具调用机制 ~\cite{anthropictooluse2026,anthropiccomputeruse2026} \\[0.15em]
%     \textbullet\hspace{0.25em}企业数据源接入、代码智能体插件化与连接器权限边界
%     & \textbullet\hspace{0.25em}Model Context Protocol、Claude Code 与企业级 MCP 治理实践 ~\cite{anthropicmcp2024,anthropicclaudecode2025,mcpsafetyaudit2025}
%   \end{tabular} \\

% \cline{2-4}

% & Google / Google DeepMind / Google Cloud
% & \begin{itemize}[leftmargin=1.35em,labelsep=0.25em,itemsep=0.05em,topsep=0pt,parsep=0pt,partopsep=0pt,label=\textbullet]
%     \item 搜索、网页与云原生企业应用
%     \item Gemini智能体平台与云安全治理
%     \item 科学智能体与漏洞攻防研究
%   \end{itemize}
% & \begin{tabular}{@{}>{\raggedright\arraybackslash}p{0.40\linewidth}@{\hspace{0.8em}}>{\raggedright\arraybackslash}p{0.54\linewidth}@{}}
%     \textbullet\hspace{0.25em}云原生智能体的构建、部署、观测、评测与生命周期治理
%     & \textbullet\hspace{0.25em}Vertex AI Agent Builder、Agent Development Kit、Agent Engine 与 Agent Search ~\cite{googlevertexagentbuilder2026,googleadk2026,googleagentsearch2026} \\[0.15em]
%     \textbullet\hspace{0.25em}企业工具/API 注册、Agent 身份、最小权限与跨服务访问控制
%     & \textbullet\hspace{0.25em}Cloud API Registry 工具治理、Agent Identity、IAM 条件与 Agent Gateway ~\cite{googletoolgovernance2025,googleagentidentity2026} \\[0.15em]
%     \textbullet\hspace{0.25em}科学推理与安全攻防 Agent 的受控实验、漏洞发现和结果可复现
%     & \textbullet\hspace{0.25em}AI co-scientist、AlphaEvolve、Project Naptime、Big Sleep / CodeMender ~\cite{googlecoscientist2025,googlealphaevolve2025,googlenaptime2024,googlecodemender2025}
%   \end{tabular} \\

% \cline{2-4}

% 美国
% & Meta
% & \begin{itemize}[leftmargin=1.35em,labelsep=0.25em,itemsep=0.05em,topsep=0pt,parsep=0pt,partopsep=0pt,label=\textbullet]
%     \item 开源/开放权重基础模型
%     \item 开源安全护栏与安全评测
%     \item 端侧部署与开发者生态
%   \end{itemize}
% & \begin{tabular}{@{}>{\raggedright\arraybackslash}p{0.40\linewidth}@{\hspace{0.8em}}>{\raggedright\arraybackslash}p{0.54\linewidth}@{}}
%     \textbullet\hspace{0.25em}开源模型输入/输出内容安全分类与多语种安全策略适配
%     & \textbullet\hspace{0.25em}Llama Guard、Prompt Guard 与 Llama Responsible Use Guide ~\cite{metapurplellama2026,metallamaguard2024,metallamaresponsibleuse2024} \\[0.15em]
%     \textbullet\hspace{0.25em}代码生成、网络安全能力与提示注入风险的开放评测
%     & \textbullet\hspace{0.25em}CyberSecEval、Code Shield，用于不安全代码过滤与网络安全风险测试 ~\cite{metapurplellama2026,metacyberseceval2023} \\[0.15em]
%     \textbullet\hspace{0.25em}Agent 运行时安全监控、间接提示注入和危险代码风险拦截
%     & \textbullet\hspace{0.25em}LlamaFirewall，集成 PromptGuard 2、Agent Alignment Checks 与 CodeShield ~\cite{metallamaprotections2026,metallamafirewall2026}
%   \end{tabular} \\

% \cline{2-4}

% & NVIDIA
% & \begin{itemize}[leftmargin=1.35em,labelsep=0.25em,itemsep=0.05em,topsep=0pt,parsep=0pt,partopsep=0pt,label=\textbullet]
%     \item AI 算力基础设施与 GPU 软件栈
%     \item 企业级 Agent 编排与护栏框架
%     \item 工业、机器人与物理 AI 仿真
%   \end{itemize}
% & \begin{tabular}{@{}>{\raggedright\arraybackslash}p{0.40\linewidth}@{\hspace{0.8em}}>{\raggedright\arraybackslash}p{0.54\linewidth}@{}}
%     \textbullet\hspace{0.25em}生产级多智能体工作流的编排、工具连接、评测与性能剖析
%     & \textbullet\hspace{0.25em}NVIDIA NeMo Agent Toolkit，支持 MCP/A2A、workflow、evaluation 与 profiler ~\cite{nvidiaagenttoolkit2026} \\[0.15em]
%     \textbullet\hspace{0.25em}对话式 AI 与 RAG Agent 的输入/输出内容护栏和策略约束
%     & \textbullet\hspace{0.25em}NVIDIA NeMo Guardrails、NIM 与安全 RAG 部署参考 ~\cite{nvidianemoguardrails2026,rebedea2023nemo} \\[0.15em]
%     \textbullet\hspace{0.25em}可信算力、模型/数据在用保护与物理 AI 安全仿真
%     & \textbullet\hspace{0.25em}NVIDIA Confidential Computing、Isaac Sim / Isaac Lab、Cosmos World Foundation Models ~\cite{nvidiaconfidentialcomputing2026,nvidiaisaacsim2026,nvidiaisaaclab2026,nvidiacosmos2025}
%   \end{tabular} \\

% \cline{2-4}

% & Perplexity
% & \begin{itemize}[leftmargin=1.35em,labelsep=0.25em,itemsep=0.05em,topsep=0pt,parsep=0pt,partopsep=0pt,label=\textbullet]
%     \item 搜索原生问答与实时联网检索
%     \item 浏览器智能体与网页任务自动化
%     \item 企业研究智能体与 API 平台
%   \end{itemize}
% & \begin{tabular}{@{}>{\raggedright\arraybackslash}p{0.40\linewidth}@{\hspace{0.8em}}>{\raggedright\arraybackslash}p{0.54\linewidth}@{}}
%     \textbullet\hspace{0.25em}联网检索型 Agent 的来源可信、引用溯源与搜索范围控制
%     & \textbullet\hspace{0.25em}Perplexity Sonar API、Search API、Pro Search 与引用型答案机制 ~\cite{perplexitysonar2026,perplexitysearchapi2026,perplexityprivacysecurity2026} \\[0.15em]
%     \textbullet\hspace{0.25em}浏览器 Agent 的网页上下文处理、跨站任务代理与提示注入风险防护
%     & \textbullet\hspace{0.25em}Comet Browser / Comet Assistant，并结合网页 Agent 安全评测研究 ~\cite{perplexitycomet2025,bravecometpromptinjection2025,greshake2023indirectprompt,agentdojo2024} \\[0.15em]
%     \textbullet\hspace{0.25em}企业级多步任务 Agent 的连接器、访问边界、数据留存和网络策略治理
%     & \textbullet\hspace{0.25em}Perplexity Enterprise、Perplexity Computer、Enterprise connectors 与 Computer Network Firewall Policy ~\cite{perplexityenterprise2026,perplexitycomputerenterprise2026}
%   \end{tabular} \\

% \cline{2-4}

% & Cursor / Anysphere
% & \begin{itemize}[leftmargin=1.35em,labelsep=0.25em,itemsep=0.05em,topsep=0pt,parsep=0pt,partopsep=0pt,label=\textbullet]
%     \item AI IDE 与代码智能体
%     \item 云端/本地软件工程自动化
%     \item 企业研发安全与 MCP 工具接入
%   \end{itemize}
% & \begin{tabular}{@{}>{\raggedright\arraybackslash}p{0.40\linewidth}@{\hspace{0.8em}}>{\raggedright\arraybackslash}p{0.54\linewidth}@{}}
%     \textbullet\hspace{0.25em}IDE 内代码智能体的本地工程访问、终端/浏览器操作与命令审查
%     & \textbullet\hspace{0.25em}Cursor Agent、Plan Mode、Browser Controls、Hooks，用于计划、审计、命令阻断与上下文控制 ~\cite{cursor17changelog2025,cursorplanmode2025,cursorhooks2025} \\[0.15em]
%     \textbullet\hspace{0.25em}云端代码智能体的隔离执行环境、PR 产物验证与企业内网部署
%     & \textbullet\hspace{0.25em}Cursor Cloud Agents、Bugbot、Self-hosted Cloud Agents ~\cite{cursorcloudagents2025,cursorbugbot2025,cursorselfhostedcloudagents2026} \\[0.15em]
%     \textbullet\hspace{0.25em}企业级代码数据保护、模型/MCP/网络访问控制和审计合规
%     & \textbullet\hspace{0.25em}Cursor Enterprise、Privacy Mode、SSO/SAML、RBAC、SCIM、audit logs、repository/model/MCP controls ~\cite{cursorsecurity2026,cursorenterprise2026}
%   \end{tabular} \\

% \hline
% \end{tabularx}
% \end{table*}


\subsection{欧洲主流AI智能体安全产品}
欧洲AI智能体安全产业深度绑定数据主权、区域监管合规、高可信可解释性三大底层逻辑，形成\textbf{合规驱动迭代、主权约束扩张}的差异化产业体系\cite{alephphariaaipdf2025,mistralaistudio2025}。欧盟依托《通用数据保护条例》（GDPR）与《AI法案》（EU AI Act）构建严苛的AI治理框架，从法律层面对智能体数据处理、高风险应用分级、权限边界、审计追溯提出强制性约束，成为区域企业产品研发与迭代的核心准则\cite{alephintelligencelayer2026}。
%
本节梳理了欧洲主要国家的主流AI智能体产品和技术成果，主要发现如下：（详见表~\ref{tab:europe-agent-security-products}）
\begin{itemize}[tabitem]
    \item \textbf{驱动逻辑监管导向，形成“合规牵引技术、主权约束生态”的发展范式。}
    欧盟全域研发以GDPR与EU AI Act为法定前置条件，所有智能体产品设计优先适配数据驻留、高风险分级、可审计追溯等强制性条款\cite{alephphariaaipdf2025}。Mistral AI、Aleph Alpha的安全护栏、私有化能力均直接对标法案要求，技术研发完全围绕合规需求展开，实现法规标准与产品能力的深度耦合\cite{mistralforge2026,alephphariamodels2026}；俄罗斯以地缘数据自主、规避域外生态依赖为核心诉求，形成监管合规与地缘安全双重驱动逻辑，与欧盟合规导向形成互补，二者底层均以\textbf{主权可控}为核心目标，构建非扩张型、安全优先的产业发展路径\cite{yandexaistudiosecurity2026}。

    \item \textbf{技术聚焦高可信研究，构建“可解释+数据闭环+私有部署”研究体系：}
    欧盟基础层企业错位发力、赛道分工明确：Mistral AI侧重通用智能体运行时安全，覆盖提示注入检测、越狱拦截、工具调用权限管控等通用护栏能力，兼顾开发者生态灵活度与企业合规硬性需求，形成通用场景标准化安全能力体系\cite{mistralagentsapi2025,mistralstudio2026}；Aleph Alpha聚焦政务、工业等高监管高风险场景，主打可解释推理、决策溯源、全流程人工复核机制，构建可验证、可追溯、可审计的高可信智能体体系，精准适配欧洲高等级合规场景\cite{alephphariaai2026,alephintelligencelayer2026}。英国Darktrace走垂直差异化路线，不布局基础模型研发，深耕网络安全域智能体，依托无监督基线建模实现跨域威胁识别与自主响应，成为全球SOC可信自治智能体标杆产品\cite{darktraceactiveai2026,darktracecyberanalyst2026}。俄罗斯Yandex核心聚焦生态隔离与数据自主，通过无状态推理、本地基础设施部署、训练数据回流禁用等机制，实现俄语生态智能体全链路属地化管控，从底层规避域外数据泄露、模型劫持与生态依赖风险\cite{yandexaistudioonprem2026,yandexaisearch2026}。

    \item \textbf{生态治理保守闭环，依托窄域联动构筑合规可信的全球差异化壁垒：}
    相较于北美全域拥抱MCP协议、打造无边界全球互联生态，欧俄区域对智能体跨平台、跨域互联持保守管控态度，拒绝无边界生态开放，以安全合规、数据主权为优先准则\cite{mistralagentsapi2025,yandexaistudio2026}。区域主流产品普遍强化工具调用白名单边界、文档库精细化访问权限、智能体交接（handoff）强约束机制，严格限制跨主体未授权上下文流转，从架构源头阻断数据泄露、上下文劫持与横向渗透风险\cite{mistralforge2026,yandexaistudiosecurity2026}；Yandex MCP Hub、Mistral工具管理体系均采用严格白名单准入机制，区别于北美自由插件接入模式，以强准入、强审计、强管控筑牢闭环安全生态\cite{yandexaistudio2026,mistralstudio2026}。同时，区域产学研高度耦合，欧洲学界可证明安全、形式化验证等前沿学术成果高效转化为企业合规产品，实现理论研究与产业落地双向赋能\cite{alephphariaaipdf2025}。
\end{itemize}


% 依赖项（必须放在主文档导言区，即 \begin{document} 之前）：
% \usepackage{tabularray}
% \UseTblrLibrary{varwidth}
% \usepackage{enumitem}
\begingroup
% \nopagebreak[4]
\enlargethispage{1cm}
\footnotesize
\centering

\DefTblrTemplate{conthead-text}{default}{（续表）}
\DefTblrTemplate{contfoot-text}{default}{续下页}

\begin{longtblr}[
  caption = {欧洲主流AI智能体安全产品与落地方向},
  label   = {tab:europe-agent-security-products},
]{
  width   = \linewidth,
  colspec = {
    |Q[wd=0.8cm,l,m]
    |Q[wd=1.2cm,l,m]
    |Q[wd=2.8cm,l,m]
    |Q[wd=0.3\linewidth,l,m]
    |X[1,l,m]|
  },
  cells   = {l,m},
  colsep  = 3pt,
  row{1}  = {font=\bfseries},
  rowhead = 1,
  % 每行内容与上下边框之间各保留 5pt 空间。
  rowsep  = 5pt,
  measure = vbox,
  stretch = -1,
}
\hline
\textbf{国家} & \textbf{企业} & \textbf{所属产业领域} & \textbf{安全研究方向} & \textbf{产品/成果} \\
\hline
法国
& Mistral AI
& \begin{itemize}[leftmargin=*,topsep=0pt,partopsep=0pt,itemsep=1pt,parsep=0pt]
    \item 欧洲基础模型与企业AI平台
    \item 终端智能体、企业助手与开发者API
    \item 私有化部署、模型定制与主权AI
  \end{itemize}
& \begin{itemize}[leftmargin=*,topsep=0pt,partopsep=0pt,itemsep=1pt,parsep=0pt]
    \item 终端与企业智能体的数据使用、连接器访问、组织级隐私和权限控制
    \item 开发者智能体的工具调用、持久记忆、多工具连接与编排控制
    \item 生产级智能体的可观测性、持久运行、资产治理与自托管部署
  \end{itemize}
& \begin{itemize}[leftmargin=*,topsep=0pt,partopsep=0pt,itemsep=1pt,parsep=0pt]
    \item Le Chat Enterprise与企业级连接器/知识库应用 ~\cite{mistrallachatenterprise2025}
    \item Mistral Agents API、支持持久状态、多Agent、内置工具、handoff与MCP连接 ~\cite{mistralagentsapi2025}
    \item Mistral AI Studio、Agent Runtime、AI Registry、Forge 与 self-hosted/hybrid deployment ~\cite{mistralaistudio2025,mistralstudio2026,mistralforge2026}
  \end{itemize} \\
\hline
德国
& Aleph Alpha
& \begin{itemize}[leftmargin=*,topsep=0pt,partopsep=0pt,itemsep=1pt,parsep=0pt]
    \item 专用化、主权型大语言模型
    \item 高可信、可解释企业/政务AI平台
    \item 欧洲基础设施与私有化部署
  \end{itemize}
& \begin{itemize}[leftmargin=*,topsep=0pt,partopsep=0pt,itemsep=1pt,parsep=0pt]
    \item 高监管场景下的可解释、可追溯、数据主权与组织级AI应用治理
    \item 德语区及欧洲多语种、工程/工业领域专用模型的安全对齐与合规训练
    \item 企业和公共部门的本地化部署、生产评测、无用户数据训练与审计可控
  \end{itemize}
& \begin{itemize}[leftmargin=*,topsep=0pt,partopsep=0pt,itemsep=1pt,parsep=0pt]
    \item PhariaAI、PhariaAssistant、PhariaStudio、PhariaOS ~\cite{alephphariaai2026,alephphariaaipdf2025}
    \item Pharia-1-LLM-7B-control/control-aligned ~\cite{alephphariamodels2026}
    \item Intelligence Layer SDK、on-premise installation、PhariaOS ~\cite{alephintelligencelayer2026,alephphariaaipdf2025}
  \end{itemize} \\
\hline
英国
& Dark-\newline trace
& \begin{itemize}[leftmargin=*,topsep=0pt,partopsep=0pt,itemsep=1pt,parsep=0pt]
    \item AI原生网络安全平台
    \item SOC自动化、威胁检测与自主响应
    \item 邮件、网络、云、端点、身份与OT安全
  \end{itemize}
& \begin{itemize}[leftmargin=*,topsep=0pt,partopsep=0pt,itemsep=1pt,parsep=0pt]
    \item 基于自学习AI的组织行为基线建模、已知/未知威胁检测与实时风险闭环
    \item SOC告警自动研判、事件关联调查、上下文综合与自然语言报告生成
    \item 精准自主响应与跨域安全控制，降低横向移动、邮件攻击和云端异常扩散风险
  \end{itemize}
& \begin{itemize}[leftmargin=*,topsep=0pt,partopsep=0pt,itemsep=1pt,parsep=0pt]
    \item Darktrace ActiveAI Security Platform ~\cite{darktraceactiveai2026}
    \item Cyber AI Analyst ~\cite{darktraceactiveai2026,darktracecyberanalyst2026}
    \item Autonomous Response、NETWORK、EMAIL、CLOUD、ENDPOINT、IDENTITY、OT、SECURE AI ~\cite{darktraceautonomousresponse2026,darktraceactiveai2026}
  \end{itemize} \\
\hline
俄罗斯
& Yandex
& \begin{itemize}[leftmargin=*,topsep=0pt,partopsep=0pt,itemsep=1pt,parsep=0pt]
    \item 搜索、云平台与俄语大模型服务
    \item AI Studio、Agent Atelier与企业RAG/AISearch
    \item Alice 助手生态与云安全合规工具
  \end{itemize}
& \begin{itemize}[leftmargin=*,topsep=0pt,partopsep=0pt,itemsep=1pt,parsep=0pt]
    \item 企业智能体构建、RAG、网页搜索、MCP工具接入与业务流程集成
    \item 本地化或私有基础设施内运行智能体，控制数据、请求、处理结果和访问日志
    \item 云安全配置、合规建议、告警解释与安全运营辅助
  \end{itemize}
& \begin{itemize}[leftmargin=*,topsep=0pt,partopsep=0pt,itemsep=1pt,parsep=0pt]
    \item Yandex AI Studio、Agent Atelier、AI Search、MCP Hub与Yandex Workflows~\cite{yandexaistudio2026,yandexaisearch2026}
    \item Yandex AI Studio on-premises与请求日志禁用、访问管理、Audit Trails等安全设置~\cite{yandexaistudioonprem2026,yandexaistudiosecurity2026}
    \item Yandex Security Deck AI assistant ~\cite{yandexsecuritydeckassistant2026}
  \end{itemize} \\
\hline
\end{longtblr}
\endgroup





% \renewcommand{\arraystretch}{0.7}
% \begin{table*}[htbp]
% \centering
% \caption{欧洲及周边主流 AI 智能体安全产品与落地方向}
% \label{tab:europe-agent-security-products}
% \fontsize{6.7pt}{7.4pt}\selectfont
% \setlength{\tabcolsep}{3pt}
% \renewcommand{\arraystretch}{1.06}

% \begin{tabularx}{\textwidth}{|>{\centering\arraybackslash}m{0.065\textwidth}|>{\raggedright\arraybackslash}m{0.13\textwidth}|>{\hspace{0.45em}\raggedright\arraybackslash}m{0.20\textwidth}|>{\hspace{0.35em}\raggedright\arraybackslash}X|}
% \hline
% \multicolumn{1}{|c|}{所属区域} & \multicolumn{1}{c|}{企业名称} & \multicolumn{1}{c|}{所属产业领域} & \multicolumn{1}{c|}{安全研究方向与产品/成果} \\
% \hline

% 法国
% & Mistral AI
% & \begin{itemize}[leftmargin=1.15em,labelsep=0.18em,itemsep=0.05em,topsep=0pt,parsep=0pt,partopsep=0pt,label=\textbullet]
%     \item 欧洲基础模型与企业 AI 平台
%     \item 终端智能体、企业助手与开发者 API
%     \item 私有化部署、模型定制与主权 AI
%   \end{itemize}
% & \begin{tabular}{@{}>{\raggedright\arraybackslash}p{0.39\linewidth}@{\hspace{0.7em}}>{\raggedright\arraybackslash}p{0.56\linewidth}@{}}
%     \textbullet\hspace{0.18em}终端与企业智能体的数据使用、连接器访问、组织级隐私和权限控制
%     & \textbullet\hspace{0.18em}Le Chat Enterprise 与企业级连接器/知识库应用 ~\cite{mistrallachatenterprise2025} \\[0.18em]
%     \textbullet\hspace{0.18em}开发者 Agent 的工具调用、持久记忆、多工具连接与编排控制
%     & \textbullet\hspace{0.18em}Mistral Agents API，支持持久状态、多 Agent、内置工具、handoff 与 MCP 连接 ~\cite{mistralagentsapi2025} \\[0.18em]
%     \textbullet\hspace{0.18em}生产级智能体的可观测性、持久运行、资产治理与自托管部署
%     & \textbullet\hspace{0.18em}Mistral AI Studio、Agent Runtime、AI Registry、Forge 与 self-hosted / hybrid deployment ~\cite{mistralaistudio2025,mistralstudio2026,mistralforge2026}
%   \end{tabular} \\

% \hline

% 德国
% & Aleph Alpha
% & \begin{itemize}[leftmargin=1.15em,labelsep=0.18em,itemsep=0.05em,topsep=0pt,parsep=0pt,partopsep=0pt,label=\textbullet]
%     \item 专用化、主权型大语言模型
%     \item 高可信、可解释企业/政务 AI 平台
%     \item 欧洲基础设施与私有化部署
%   \end{itemize}
% & \begin{tabular}{@{}>{\raggedright\arraybackslash}p{0.39\linewidth}@{\hspace{0.7em}}>{\raggedright\arraybackslash}p{0.56\linewidth}@{}}
%     \textbullet\hspace{0.18em}高监管场景下的可解释、可追溯、数据主权与组织级 AI 应用治理
%     & \textbullet\hspace{0.18em}PhariaAI、PhariaAssistant、PhariaStudio、PhariaOS ~\cite{alephphariaai2026,alephphariaaipdf2025} \\[0.18em]
%     \textbullet\hspace{0.18em}德语区及欧洲多语种、工程/工业领域专用模型的安全对齐与合规训练
%     & \textbullet\hspace{0.18em}Pharia-1-LLM-7B-control / control-aligned ~\cite{alephphariamodels2026} \\[0.18em]
%     \textbullet\hspace{0.18em}企业和公共部门的本地化部署、生产评测、无用户数据训练与审计可控
%     & \textbullet\hspace{0.18em}Intelligence Layer SDK、on-premise installation、PhariaOS ~\cite{alephintelligencelayer2026,alephphariaaipdf2025}
%   \end{tabular} \\

% \hline

% 英国
% & Darktrace
% & \begin{itemize}[leftmargin=1.15em,labelsep=0.18em,itemsep=0.05em,topsep=0pt,parsep=0pt,partopsep=0pt,label=\textbullet]
%     \item AI 原生网络安全平台
%     \item SOC 自动化、威胁检测与自主响应
%     \item 邮件、网络、云、端点、身份与 OT 安全
%   \end{itemize}
% & \begin{tabular}{@{}>{\raggedright\arraybackslash}p{0.39\linewidth}@{\hspace{0.7em}}>{\raggedright\arraybackslash}p{0.56\linewidth}@{}}
%     \textbullet\hspace{0.18em}基于自学习 AI 的组织行为基线建模、已知/未知威胁检测与实时风险闭环
%     & \textbullet\hspace{0.18em}Darktrace ActiveAI Security Platform ~\cite{darktraceactiveai2026} \\[0.18em]
%     \textbullet\hspace{0.18em}SOC 告警自动研判、事件关联调查、上下文综合与自然语言报告生成
%     & \textbullet\hspace{0.18em}Cyber AI Analyst ~\cite{darktraceactiveai2026,darktracecyberanalyst2026} \\[0.18em]
%     \textbullet\hspace{0.18em}精准自主响应与跨域安全控制，降低横向移动、邮件攻击和云端异常扩散风险
%     & \textbullet\hspace{0.18em}Autonomous Response、NETWORK、EMAIL、CLOUD、ENDPOINT、IDENTITY、OT、SECURE AI ~\cite{darktraceautonomousresponse2026,darktraceactiveai2026}
%   \end{tabular} \\

% \hline

% 俄罗斯
% & Yandex
% & \begin{itemize}[leftmargin=1.15em,labelsep=0.18em,itemsep=0.05em,topsep=0pt,parsep=0pt,partopsep=0pt,label=\textbullet]
%     \item 搜索、云平台与俄语大模型服务
%     \item AI Studio、Agent Atelier 与企业 RAG / AI Search
%     \item Alice 助手生态与云安全合规工具
%   \end{itemize}
% & \begin{tabular}{@{}>{\raggedright\arraybackslash}p{0.39\linewidth}@{\hspace{0.7em}}>{\raggedright\arraybackslash}p{0.56\linewidth}@{}}
%     \textbullet\hspace{0.18em}企业智能体构建、RAG、网页搜索、MCP 工具接入与业务流程集成
%     & \textbullet\hspace{0.18em}Yandex AI Studio、Agent Atelier、AI Search、MCP Hub 与 Yandex Workflows ~\cite{yandexaistudio2026,yandexaisearch2026} \\[0.18em]
%     \textbullet\hspace{0.18em}本地化或私有基础设施内运行智能体，控制数据、请求、处理结果和访问日志
%     & \textbullet\hspace{0.18em}Yandex AI Studio on-premises 与请求日志禁用、访问管理、Audit Trails 等安全设置 ~\cite{yandexaistudioonprem2026,yandexaistudiosecurity2026} \\[0.18em]
%     \textbullet\hspace{0.18em}云安全配置、合规建议、告警解释与安全运营辅助
%     & \textbullet\hspace{0.18em}Yandex Security Deck AI assistant ~\cite{yandexsecuritydeckassistant2026}
%   \end{tabular} \\

% \hline
% \end{tabularx}
% \end{table*}


% \renewcommand{\arraystretch}{0.7}
% \begin{table*}[htbp]
% \centering
% \caption{欧洲/俄罗斯主流 AI 智能体安全产品与落地方向}
% \label{tab:europe-russia-agent-security-products}
% \footnotesize
% \begin{tabularx}{\textwidth}{|p{0.08\textwidth} |p{0.14\textwidth}| p{0.19\textwidth}| X|}
% \hline
% 所属区域 & 企业名称 & 所属产业赛道 & 智能体安全核心研究与落地方向 \\
% \hline
% 法国/欧盟 & Mistral AI & 欧洲基础模型、开发者 Agent API、企业私有化与主权 AI 平台 & 方向：LLM 注入安全、Agent 安全框架与欧洲主权合规。重点：Agent 级 guardrails 与 moderation；PII、有害内容和 jailbreak 检测；多 Agent handoff、工具调用、代码执行和文档库工具权限边界；企业定制模型与私有部署中的数据驻留、GDPR/欧盟 AI Act 合规和审计 ~\cite{mistralvibe2026,mistralagents2026,mistralguardrails2026,mistralforge2026,gdpr2016,euaiact2024}。 \\
% \hline
% 德国/欧盟 & Aleph Alpha & 德国主权 AI、可解释企业模型与高可信行业智能体 & 方向：可解释 Agent、主权 AI 与高风险行业合规治理。重点：欧洲基础设施上的数据主权和客户域内训练；企业/政府场景中的证据溯源和可验证洞察；高风险决策中的可解释性、人类审查和合规评估；行业知识接入时的权限继承、数据隔离和审计 ~\cite{alephalpha2026,phariaai2026,atman2023,gdpr2016,euaiact2024}。 \\
% \hline
% 英国 & Darktrace & 网络安全运营智能体、自主检测响应与 AI 安全防护平台 & 方向：SOC 智能体、自主检测响应与企业 AI 安全防护。重点：无监督业务基线建模和异常行为识别；告警关联、攻击链解释和自动化调查；自主响应/隔离动作的人机监督和策略约束；网络、终端、云、邮件、OT 跨域遥测融合 ~\cite{darktraceplatform2026,darktracenetwork2026,darktracesecureai2026,agentsecurity2024}。 \\
% \hline
% 俄罗斯 & Yandex & 俄语大模型、AI Studio 企业 Agent 平台、属地化云与私有化部署 & 方向：俄语生态 Agent、企业工具编排与属地化/私有化部署安全。重点：知识库、工具连接、互联网/文件搜索、代码解释器和 MCP Hub 管理；AI Studio 的访问控制、逻辑隔离、加密、审计和 guardrails；客户数据不用于训练、Stateless 推理和私有端点；on-premises 场景下数据与处理结果留存在客户基础设施内 ~\cite{yandexaliceai2026,yandexmodelgallery2026,yandexaistudiosecurity2026,yandexaistudioonprem2026}。 \\
% \hline
% \end{tabularx}
% \end{table*}


\subsection{我国主流AI智能体安全产品}
我国AI智能体安全产业以\textbf{场景倒逼、工程落地}为核心底层逻辑，形成业务优先于理论、工程适配原生防护的差异化发展体系\cite{aliyunopenclawdocs2026,tencentagentgateway2026,huaweiagentarts2026}。我国智能体安全研发聚焦政企办公、工业流程、知识库检索、代码开发、安全运营等高频落地场景，依托国产化软硬件生态与大规模政企交付需求，构建以外挂式边界防护为主体、私有化部署为载体、结构化生态分工为支撑的工程化安全体系，在垂直场景治理与信创适配领域形成独有产业特征，在底层原生技术与国际标准层面存在差异化短板\cite{baiduqianfanagent2026,iflyteksparkenterpriseagent2026}。本节梳理了我国头部厂商的主要落地产品与技术成果，主要发现如下（详见表~\ref{tab:china-agent-security-products-verified-cited}）：
\begin{itemize}[tabitem]
    \item \textbf{产业驱动呈现场景倒逼特征，形成“场景牵引工程防护”的主导发展范式。}
    头部企业安全能力均针对性适配垂直场景痛点：科大讯飞聚焦语音交互隐写与误触发风险构建多模态安全规则\cite{iflytekxingchenagent2026}；华为针对工业智能体高危操作设计熔断机制，约束执行层越权行为\cite{huaweicodearts2026}；360围绕智能体资产测绘与插件供应链风险补齐攻防能力，重点解决智能体外部生态带来的安全隐患\cite{threesixtyagentskillsreport2026}，所有技术研发均由具体场景问题直接牵引。同时监管侧以落地合规为核心导向，重点围绕算法备案、数据本地留存、等级保护提出硬性要求，不侧重可证明安全、宪法式对齐等前沿理论研究，进一步强化了产业工程化、落地化的核心导向\cite{baiduqianfanagent2026}。

    \item \textbf{技术架构以外挂边界防护为核心，构建“网关护栏+权限审计+租户隔离”三层工程体系。}
    与北美主打模型内生对齐、欧洲深耕可解释溯源的技术路线不同，我国主流采用轻量化、易部署的外挂式、运行时旁路防护架构\cite{aliyunaiguardrails2026}。以腾讯 AI 智能体网关\cite{tencentagentgateway2026}、阿里云OpenClaw\cite{aliyunopenclawdocs2026}、火山引擎 LLM 防火墙\cite{volcenginellmfirewall2026}为代表，在模型出入口、工具调用链路、上下文流转节点设置统一管控枢纽，实现提示注入拦截、敏感数据脱敏、恶意插件检测等基础防护能力。全域形成标准化三层治理架构：边界层依托安全网关、流量检测、MCP协议防护构筑外部屏障；运行层通过最小权限管控、租户隔离、全链路操作审计规范运行行为；资产层完成Agent资产盘点、影子实例测绘、Skill供应链风控。该架构部署成本低、通用性强，高度适配我国大规模政企批量交付的核心需求\cite{tencentadp2026}。

    \item \textbf{生态分工结构化且深度绑定信息技术应用创新产业，形成四级生态与国产化闭环治理格局：}
    云底座厂商（百度、阿里、腾讯、华为、字节等）主导开发平台、运行时内核与身份权限体系，定义全域安全治理底座\cite{baiduqianfanagent2026,huaweiagentidentity2026}；行业垂直厂商（科大讯飞）依托多模态壁垒，深耕政务、医疗等合规行业，构建私有化数据闭环\cite{iflyteksparkenterpriseagent2026}；专业安全厂商（360）补齐漏洞溯源、威胁狩猎等攻防能力，填补大模型原生安全短板\cite{threesixtymodelguard2026}；开源厂商（DeepSeek）聚焦开发者生态与沙箱隔离，服务中小开发群体\cite{deepseektoolcalls2026}。同时产业深度绑定信息技术应用创新产业（信创）体系，基于昇腾/鲲鹏算力、飞桨框架构建国产化可控闭环，在加密机制、审计粒度上差异化适配，形成全球独有的信创导向智能体安全体系\cite{paddlepaddle2026,huaweiagentarts2026}。
    % 整体产业在工程落地、场景治理领域全球领跑，在底层协议、内生对齐等原生技术领域跟跑北美与欧盟，依托细分壁垒构建差异化全球竞争定位\cite{aliyunopenclawassistant2026}。
\end{itemize}

% 依赖项（必须放在主文档导言区，即 \begin{document} 之前）：
% \usepackage{tabularray}
% \UseTblrLibrary{varwidth}
% \usepackage{enumitem}
\begingroup
% \nopagebreak[4]
\enlargethispage{1cm}
\footnotesize
\centering

\DefTblrTemplate{conthead-text}{default}{（续表）}
\DefTblrTemplate{contfoot-text}{default}{续下页}

\begin{longtblr}[
  caption = {我国主流 AI 智能体安全产品与落地方向},
  label   = {tab:china-agent-security-products-verified-cited},
]{
  width   = \linewidth,
  colspec = {
    |Q[wd=0.8cm,l,m]
    |Q[wd=1cm,l,m]
    |Q[wd=3cm,l,m]
    |Q[wd=0.4\linewidth,l,m]
    |X[1,l,m]|
  },
  cells   = {l,m},
  colsep  = 3pt,
  row{1}  = {font=\bfseries},
  rowhead = 1,
  % 每行内容与上下边框之间各保留 5pt 空间。
  rowsep  = 5pt,
  measure = vbox,
  stretch = -1,
}
\hline
\textbf{国家} & \textbf{企业} & \textbf{所属产业领域} & \textbf{安全研究方向} & \textbf{产品/成果} \\
\hline
中国大陆
& 百度
& \begin{itemize}[leftmargin=*,topsep=0pt,partopsep=0pt,itemsep=1pt,parsep=0pt]
    \item 搜索增强大模型
    \item 智能云智能体开发平台
    \item 飞桨与国产AI生态
  \end{itemize}
& \begin{itemize}[leftmargin=*,topsep=0pt,partopsep=0pt,itemsep=1pt,parsep=0pt]
    \item 企业智能体构建、RAG、工具/MCP接入与模型服务治理
    % \item 企业智能体任务执行、Skill生态与消息平台集成
    \item 国产AI框架生态、模型开发与工程化部署基础
  \end{itemize}

\vspace{0.5em}
2026重点研究方向：
\vspace{1.5em}

\begin{itemize}[leftmargin=*,topsep=0pt,partopsep=0pt,itemsep=1pt,parsep=0pt]
    \item 面向智能体复杂任务推理，研究更高效的强化学习、长上下文奖励分配与推理能力优化\cite{baidu_acl_ai_safety_2026}
    \item 构建覆盖数据、身份权限、记忆、工具调用和外部系统连接的智能体全链路安全基础设施\cite{baidu_agent_security_center_2026}
    \item 加强 Skill/MCP 供应链安全，研究恶意指令、过度权限、凭证泄露与高风险工具调用的准入检测和持续巡检\cite{baidu_ai_security_solution_2026,baidu_qianfan_platform_2026}
    \item 推进策略即代码、输入输出安全策略和语义干预，使企业能够按业务场景定义、执行并审计智能体行为规则\cite{baidu_qianfan_security_policy_2026,baidu_qianfan_intervention_policy_2026}
  \end{itemize}

% • 面向智能体复杂任务推理，研究更高效的强化学习、长上下文奖励分配与推理能力优化
% • 推进精细化模型安全对齐，通过参数级风险诊断和局部微调，降低安全对齐对模型通用能力的损耗
% • 构建覆盖数据、身份权限、记忆、工具调用和外部系统连接的智能体全链路安全基础设施
% • 加强 Skill/MCP 供应链安全，研究恶意指令、过度权限、凭证泄露与高风险工具调用的准入检测和持续巡检
% • 推进策略即代码、输入输出安全策略和语义干预，使企业能够按业务场景定义、执行并审计智能体行为规则



& \begin{itemize}[leftmargin=*,topsep=0pt,partopsep=0pt,itemsep=1pt,parsep=0pt]
    \item 百度千帆智能体开发平台：Agent引擎、工具/MCP、模型服务与企业服务能力 ~\cite{baiduqianfanagent2026}
    \item 百度DuClaw：AI智能体服务、Skill 系统、SkillMarket与消息平台配置 ~\cite{baiduduclaw2026,baiduduclawskillmarket2026}
    \item 飞桨PaddlePaddle与百度智能云千帆工程生态 ~\cite{paddlepaddle2026,baiduqianfanagent2026}
  \end{itemize} \\
\cline{2-5}
& 阿里巴巴/阿里云
& \begin{itemize}[leftmargin=*,topsep=0pt,partopsep=0pt,itemsep=1pt,parsep=0pt]
    \item 通义/Qwen模型与百炼智能体应用
    \item 云原生Workflow、RAG与插件编排
    \item OpenClaw运行时与AI应用安全护栏
  \end{itemize}
& \begin{itemize}[leftmargin=*,topsep=0pt,partopsep=0pt,itemsep=1pt,parsep=0pt]
    \item 输入输出安全护栏、Prompt注入检测、敏感信息和恶意链接识别
    \item 智能体运行时防护、工具调用审计、Skill检测与密钥托管
    \item RAG/插件/代码解释器边界治理与企业应用编排
  \end{itemize}

\vspace{0.5em}
2026重点研究方向：
\vspace{1.5em}

\begin{itemize}[leftmargin=*,topsep=0pt,partopsep=0pt,itemsep=1pt,parsep=0pt]
    \item 构建覆盖智能体开发期与运行期的全生命周期防护，统一管理模型、知识库、身份、工具、MCP及Skill资产关系\cite{alibaba_agent_security_center_2026}
    \item 加强智能体身份和动态权限治理，研究人机身份映射、细粒度授权、凭证托管以及关键操作二次确认\cite{alibaba_hitl_agent_security_2026}
    \item 面向Skill投毒、间接提示注入、工具投毒和恶意命令执行，发展运行时行为检测、阻断和AI红队测试\cite{alibaba_agent_risk_detection_2026}
    \item 研究Agent原生API和网络安全，通过语义化敏感数据识别、动态脱敏、外联访问控制与全链路行为审计保护数据流\cite{alibaba_agentic_api_security_2026,alibaba_agent_firewall_2026}
    \item 推进多智能体协同安全运营，以专业Agent集群完成威胁感知、攻击链推理、协同调查及闭环响应\cite{alibaba_agentic_soc_2026}
\end{itemize}
  
& \begin{itemize}[leftmargin=*,topsep=0pt,partopsep=0pt,itemsep=1pt,parsep=0pt]
    \item 阿里云AI安全护栏/OpenClaw运行时防护方案 ~\cite{aliyunaiguardrails2026,aliyunopenclawdocs2026}
    \item openclaw-security-assistant 与 OpenClaw runtime security 插件 ~\cite{aliyunopenclawassistant2026,aliyunopenclawdocs2026}
    \item 阿里云百炼Model Studio智能体应用/工作流、Assistant API代码解释器 ~\cite{aliyunbailian2026,aliyunopenclawdocs2026}
  \end{itemize} \\
\cline{2-5}
& 腾讯
& \begin{itemize}[leftmargin=*,topsep=0pt,partopsep=0pt,itemsep=1pt,parsep=0pt]
    \item 混元大模型与腾讯云ADP
    \item LLM+RAG、Workflow与Multi-Agent 应用
    \item CodeBuddy代码智能体与MCP工具生态
  \end{itemize}
& \begin{itemize}[leftmargin=*,topsep=0pt,partopsep=0pt,itemsep=1pt,parsep=0pt]
    \item MCP/智能体接入网关防护、检测规则、限流、鉴权与脱敏
    \item 租户隔离、内容风控、权限审计、监控告警与日志分析
    \item 代码智能体的本地工程访问、MCP工具调用和命令执行安全
  \end{itemize}

\vspace{0.5em}
2026重点研究方向：
\vspace{1.5em}

\begin{itemize}[leftmargin=*,topsep=0pt,partopsep=0pt,itemsep=1pt,parsep=0pt]
    \item 加强Agent Skills与MCP供应链安全研究，重点评估恶意指令、硬编码凭证、危险依赖、权限滥用及命令注入\cite{tencent_agent_skills_security_2026}
    \item 建设面向真实Skill生态的标准化安全评测基准，同时衡量Skill自身可信度和外部扫描器的检测效果、误报及漏报\cite{tencent_skilltrustbench_2026}
    \item 持续研究价值对齐、隐私保护、深度伪造检测、智能体安全和AI生成内容检测等从模型内生安全到系统防御的综合安全体系\cite{tencent_esg_report_2025_published_2026}
\end{itemize}

% • 加强Agent Skills与MCP供应链安全研究，重点评估恶意指令、硬编码凭证、危险依赖、权限滥用及命令注入
% • 建设面向真实Skill生态的标准化安全评测基准，同时衡量Skill自身可信度和外部扫描器的检测效果、误报及漏报
% • 推进“Agent扫描Agent”的自动化安全研究，由安全智能体理解技能语义、比对功能描述与实际行为，并验证风险路径
% • 建立覆盖基础设施、代码、模型、RAG、工作流和运行时行为的智能体红队演习与上线前安全基线
% • 发展以AI Agent安全网关为核心的统一控制面，对模型、API和MCP实施身份认证、限流、注入检测、敏感数据脱敏及行为审计
% • 持续研究价值对齐、隐私保护、深度伪造检测、智能体安全和AI生成内容检测等从模型内生安全到系统防御的综合安全体系
  
& \begin{itemize}[leftmargin=*,topsep=0pt,partopsep=0pt,itemsep=1pt,parsep=0pt]
    \item 腾讯云AI智能体安全网关：MCP Server封装、检测规则配置、Prompt注入/脱敏/访问鉴权等防护 ~\cite{tencentagentgateway2026,tencentagentgatewayquickstart2026}
    \item 腾讯云智能体开发平台（Tencent Cloud ADP）多租户隔离、内容安全、权限与操作审计、监控告警 ~\cite{tencentadp2026,tencenthunyuan2026}
    \item 腾讯云代码助手CodeBuddy：默认只读、命令/修改审批、MCP本地服务器与安全配置 ~\cite{tencentcodebuddysecurity2026,tencentcodebuddymcp2026}
  \end{itemize} \\
\cline{2-5}
& 智谱
& \begin{itemize}[leftmargin=*,topsep=0pt,partopsep=0pt,itemsep=1pt,parsep=0pt]
    \item GLM大模型
    \item 长时程Coding Agent与应用级智能体
    \item GUI/手机智能体与多模态工具调用
    \item 企业大模型API、模型精调、推理与评测平台
  \end{itemize}
& \begin{itemize}[leftmargin=*,topsep=0pt,partopsep=0pt,itemsep=1pt,parsep=0pt]
    \item 长时程代码智能体训练与评测中的反奖励作弊，检测并阻断智能体读取隐藏测试数据、复制参考答案及绕过任务的工具调
    \item 长上下文、高并发智能体服务的稳定性与可靠性研究，重点解决乱码、重复、罕见字符生成及长任务运行异常\cite{zhipu_coding_agent_serving_reliability_2026,zhipu_glm51_process_quality_2026}
    \item GUI/手机智能体的敏感操作确认、人工接管和受控执行，避免智能体在登录、验证码、支付等高风险场景中自主越权\cite{zhipu_open_autoglm_2026}
    \item 模型输入输出及多模态内容安全审核，实现风险分级、违规内容定位、输入拦截、输出限制和生成终止\cite{zhipu_security_audit_2026}
    \item 面向开放式机器学习系统的稳健性、公平性、隐私与可追溯治理，建设可审核、可监督、可追溯的模型服务\cite{zhipu_security_risk_notice_2026}
  \end{itemize}
& \begin{itemize}[leftmargin=*,topsep=0pt,partopsep=0pt,itemsep=1pt,parsep=0pt]
    \item GLM-5.2、GLM-5.1：面向长时程任务、工具调用和Agentic Engineering的旗舰模型\cite{zhipu_autoglm_phone_documentation_2026,zhipu_glm52_2026,glm5team2026glm5,liu2024autoglm} %最后一篇是有关GUI自主智能体的技术基础
    \item Z.ai Agent Mode、ZCode：支持内置Skills、多轮协作、长任务、远程开发及移动设备控制
    \item Open-AutoGLM、AutoGLM-Phone-9B：开源手机GUI智能体模型与框架，内置敏感操作确认和人工接管\cite{zhipu_open_autoglm_2026}
    \item 智谱大模型开放平台：模型API、智能体开发、模型精调、推理、评测及多层安全防护\cite{zhipu_bigmodel_platform_2026}
    \item Content Moderation API：支持文本、图片、音频和视频内容安全检测\cite{zhipu_content_moderation_api_2026,zhipu_security_audit_2026}
  \end{itemize} \\
\cline{2-5}
% & \multirow{3}{=}{华为}
% & \begin{itemize}[leftmargin=*,topsep=0pt,partopsep=0pt,itemsep=1pt,parsep=0pt]
%     \item AgentArts企业智能体平台与OfficeAce
%     \item 盘古大模型、ModelArts Studio与云端智能体运行时
%     \item CodeArts代码智能体与昇腾/鲲鹏工程生态
%   \end{itemize}
% & \begin{itemize}[leftmargin=*,topsep=0pt,partopsep=0pt,itemsep=1pt,parsep=0pt]
%     \item 身份鉴权、API Key 认证、沙箱工具与MCP网关接入控制
%     \item 高危动作审批护栏与PC端安全审批管理
%     \item 代码数据边界、文件保护、权限控制和隔离执行
%   \end{itemize}
% & \begin{itemize}[leftmargin=*,topsep=0pt,partopsep=0pt,itemsep=1pt,parsep=0pt]
%     \item 华为云AgentArts/AgentIdentity/Sand b\allowbreak -ox Tool：API Key认证、代码执行与文件管理沙箱、MCP Gateway   %\allowbreak允许断行 否则华为云这几个字的间距会被拉的很大
%     ~\cite{huaweiagentarts2026,huaweiagentidentity2026,huaweiagentartssandbox2026}
%     \item OfficeAce/AgentArts PC安全审批管理、对敏感操作进行审批拦截和授权确认 ~\cite{huaweiofficeaceapproval2026}
%     \item 华为云码道（CodeArts）：dotfile保护、读写边界、沙箱执行、高危命令阻断与风险提示 ~\cite{huaweicodearts2026,huaweicodeartssandbox2026}
%   \end{itemize} \\
% \cline{2-5}
& 字节跳动
& \begin{itemize}[leftmargin=*,topsep=0pt,partopsep=0pt,itemsep=1pt,parsep=0pt]
    \item 豆包大模型与火山方舟AI云
    \item 扣子Coze智能体、工作流与Skill生态
    \item 飞书aily/智能伙伴与组织协同Agent
  \end{itemize}
& \begin{itemize}[leftmargin=*,topsep=0pt,partopsep=0pt,itemsep=1pt,parsep=0pt]
    \item 大模型应用防火墙、Prompt攻击拦截、模型滥用和敏感数据泄露防护
    \item Skill最小授权、凭证保护、OAuth PKCE与企业成员权限管理
    \item 权限继承、数据保护、模型输入控制策略与组织协同治理
    未来研究方向趋势：
    % \item 智能体安全评测与部署仿真：未来需要面向真实长时任务、工具调用轨迹和模拟部署环境，提前评估智能体上线后的越权、误操作和数据泄露风险。~\cite{bytedanceseedpublicpapers2026,zhu2026edgebench}
    % \item 工具调用与非可信内容隔离：未来需要在系统层面隔离任务指令、外部网页/文档内容和工具权限，降低搜索、浏览器、代码仓库等工具型 Agent 的间接提示注入风险。~\cite{he2025pasa}
    % \item 长期记忆的安全治理：未来需要研究 Agent 记忆的选择、遗忘、权限控制和隐私保护，避免多模态长期记忆积累敏感信息或被跨任务滥用。~\cite{zou2026taskmem,zhang2025functiontokens}
    \item 多模态/具身 Agent 的误操作风险控制：未来需要从离线成功率评测转向在线干预、异常检测、动作约束和人类接管机制，降低 VLA 与机器人 Agent 的真实环境风险。~\cite{zhang2026uam,li2026handitloop}
    \item 面向科研与专业知识工作的可信 Agent：未来需要强化引用可验证性、证据链追踪、幻觉检测和抗操纵排序，避免科研 Agent 在高知识密度场景中放大错误信息。~\cite{he2025pasa}
    \item 从 Responsible AI 到 Agent Security 的机制化转化：未来可将评测基准、红队样例、运行时监控、权限策略和记忆治理整合为统一的 Agent 安全工程体系。~\cite{bytedanceseedresponsibleai2026,he2025pasa,zou2026taskmem,zhang2025functiontokens}


  \end{itemize}
& \begin{itemize}[leftmargin=*,topsep=0pt,partopsep=0pt,itemsep=1pt,parsep=0pt]
    \item 火山引擎大模型应用防火墙（覆盖输入/输出防护、Prompt攻击、恶意链接、敏感数据和模型滥用风险）~\cite{volcenginellmfirewall2026,volcenginesecuritywhitepaper2026}
    \item 扣子Coze技能安全指南、OAuth PKCE授权流程和企业版权限管理 ~\cite{cozeskillsecurity2026,cozeoauthpkce2026,volcenginecozepro2026}
    \item 飞书智能伙伴/aily数据保护与策略管理，支持敏感数据分级、模型输入控制与组织级权限一致性 ~\cite{feishuaily2026,feishuailydataprotection2026}
  \end{itemize} \\
\cline{2-5}

& 月之暗面 Kimi
& \begin{itemize}[leftmargin=*,topsep=0pt,partopsep=0pt,itemsep=1pt,parsep=0pt]
    \item Kimi K2/K2.7 Code 的工具调用、长上下文与 Agentic Coding
    \item Kimi Code、Claude Code、Cline、RooCode 等开发者工作流接入
    \item Web Search、官方工具、MCP Server 与多智能体/长时任务执行
  \end{itemize}
& \begin{itemize}[leftmargin=*,topsep=0pt,partopsep=0pt,itemsep=1pt,parsep=0pt]
    \item 工具调用参数、外部网页内容与用户指令边界
    \item MCP Server、官方工具、代码执行器与文件/网络权限配置
    \item 长时编码、多智能体协作与上下文压缩过程中的越权和误操作
    未来研究方向趋势：
    \item 工具调用与 MCP 最小权限治理：未来 Kimi 需要围绕 Tool Calls、官方工具和 MCP Server 建立细粒度授权、参数校验、调用审计与可回滚机制。~\cite{kimiapi_toolcalls2026,kimiapi_officialtools2026,kimiapi_mcp2026}
    \item 长时编码 Agent 的沙箱化执行：面向 Kimi Code 和 K2.7 Code 的长程编程任务，应强化代码执行、文件修改、依赖安装和网络访问的隔离验证。~\cite{kimiapi_k27code2026,kimik26release2026}
    \item Web Search 与外部内容的间接提示注入防御：Kimi 的搜索、抓取和网页阅读能力需要区分可信指令与非可信内容，并配套动态攻击评测。~\cite{kimiapi_websearch2026}
    \item 多智能体协作的监督与责任归因：面向 Swarm 和长时任务执行，应研究计划级监控、子智能体权限分层、结果归并审计和异常中止策略。~\cite{kimik26release2026,kimik2thinkingblog2025}
    \item 记忆、上下文压缩与安全可解释性：未来 Kimi 可进一步研究长上下文压缩、持久记忆和工具反馈在跨任务复用时的隐私保护与可解释审计。~\cite{kimiapi_officialtools2026,kimiapi_k27code2026}
  \end{itemize}
& \begin{itemize}[leftmargin=*,topsep=0pt,partopsep=0pt,itemsep=1pt,parsep=0pt]
    \item Kimi K2 面向工具使用、推理和自主问题求解优化，并在技术报告中强调 agentic data synthesis 与环境交互训练~\cite{kimik2technicalreport2025,moonshotkimik2github2025}
    \item Kimi K2.7 Code 支持 256K 上下文、长程编码、多步工具调用和主流 Agent 编程工具接入~\cite{kimiapi_k27code2026}
    \item Kimi API 提供 tool calls、Web Search、官方工具和 ModelScope MCP Server 接入能力，涉及搜索、抓取、代码执行、记忆和外部服务同步等安全边界~\cite{kimiapi_toolcalls2026,kimiapi_websearch2026,kimiapi_officialtools2026,kimiapi_mcp2026}
  \end{itemize} \\
\cline{2-5}

& Mini-\allowbreak Max
& \begin{itemize}[leftmargin=*,topsep=0pt,partopsep=0pt,itemsep=1pt,parsep=0pt]
    \item 通用大模型、Agentic Coding 与长上下文生产力智能体
    \item 多模态生成、语音/音乐/视频与角色陪伴式 Agent
    \item 开放平台 API、工具调用、MCP 与企业级智能体应用
  \end{itemize}
& \begin{itemize}[leftmargin=*,topsep=0pt,partopsep=0pt,itemsep=1pt,parsep=0pt]
    \item 长程 Agent 任务中的工具调用、浏览搜索、代码执行与权限边界
    \item 多智能体协作、长上下文记忆压缩与任务状态继承的可审计性
    \item 语音复刻、视频生成、角色扮演 Agent 的身份滥用、内容安全与隐私治理
    未来研究方向趋势：
    \item 长程任务 Agent 的可审计执行：MiniMax Agent/Mavis 和 M 系列模型面向长时间、多步骤任务，未来需要强化任务计划、工具调用、状态继承和最终交付物的全链路审计。~\cite{minimaxagentteam2026,minimaxm25blog2026}

    % \item 工具调用与 MCP 权限治理：随着 Function Call、MCP 和 Server-side Tools 接入增多，未来需要细粒度授权、参数校验、调用回放和最小权限沙箱。~\cite{minimaxfunctioncall2026,minimaxmcp2026,minimaxservertools2026}

    \item Agent 强化学习的安全奖励建模：Forge 将 RL 推向真实世界 Agent，未来需要把越权调用、错误执行、隐私泄露和不可逆操作纳入 reward/verifier 体系。~\cite{minimaxforge2026}

    % \item 编程 Agent 的代码执行安全：MiniMax Code 与 M 系列编码模型需要围绕仓库修改、依赖安装、测试执行和终端操作建立隔离环境与风险分级审批。~\cite{minimaxm2blog2025,minimaxm3blog2026,jimenez2024swebench}

    \item 多模态与角色 Agent 的身份安全：MiniMax 在语音、视频、音乐和 Role-Play Agent 上布局较深，未来需要加强语音复刻授权、合成内容标识、人格越界和未成年人保护。~\cite{minimaxvoiceclone2026,minimaxm2her2026}
  \end{itemize}
& \begin{itemize}[leftmargin=*,topsep=0pt,partopsep=0pt,itemsep=1pt,parsep=0pt]
    \item MiniMax M2/M2.5/M3 面向代码、工具调用、搜索、办公和智能体任务优化，覆盖 SWE-Bench、BrowseComp 等真实任务基准~\cite{minimaxm2blog2025,minimaxm25blog2026,minimaxm3blog2026}
    \item MiniMax Agent Team/Mavis、Forge Agent RL 框架面向长程任务、持续进化和真实世界 Agent 强化学习~\cite{minimaxagentteam2026,minimaxforge2026}
    \item MiniMax 开放平台提供 Function Call、MCP、Server-side Tools、Mini-Agent 与语音复刻等能力，形成工具接入和多模态 Agent 产品底座~\cite{minimaxfunctioncall2026,minimaxmcp2026,minimaxservertools2026,minimaxminiagent2026,minimaxvoiceclone2026}
  \end{itemize} \\
\cline{1-4}


& Deep- Seek
& \begin{itemize}[leftmargin=*,topsep=0pt,partopsep=0pt,itemsep=1pt,parsep=0pt]
    \item 开放权重基础模型与推理模型
    \item API 服务与开发者工具调用生态
    \item 开源模型二次开发、本地部署与安全评测
  \end{itemize}
& \begin{itemize}[leftmargin=*,topsep=0pt,partopsep=0pt,itemsep=1pt,parsep=0pt]
    \item API / 模型输入输出治理、JSON 输出与兼容式调用
    \item 工具调用边界、JSON Schema 校验与严格函数调用
    \item 开源权重供应链、推理模型透明度与安全评测基础
    未来研究方向趋势：
    % \item 推理型智能体的安全评测与过程监督：DeepSeek-R1 展示了强化学习驱动的长链路推理能力，但未来安全研究需要评测推理轨迹中的目标偏移、隐性越权、错误自洽和高风险任务放大效应，而不只是看最终回答是否安全。~\cite{deepseekr12025,zhou2025hiddenrisksr1,marjanovic2025thoughtology}

    % \item 工具调用智能体的权限边界与调用审计：随着 DeepSeek API 支持 tool calls，未来重点会转向函数调用前的意图校验、参数约束、最小权限授权、调用日志审计和失败回滚，防止模型把普通推理错误转化为真实外部副作用。~\cite{deepseekToolCalls2026,debenedetti2024agentdojo,xiang2026systemdefenses}

    % \item 代码智能体的执行隔离与供应链安全：DeepSeek-Coder 系列提升了代码生成和修复能力，也意味着未来需要重点研究代码执行沙箱、依赖安装约束、仓库权限隔离、自动补丁验证和恶意代码/后门检测。~\cite{deepseekcoderv22024,deepseekToolCalls2026,xiang2026systemdefenses}

    % \item 长上下文与缓存机制下的数据隔离：DeepSeek 的 context caching 和高效推理基础设施会推动企业级长任务智能体落地，但也需要研究缓存命中隔离、跨会话数据残留、敏感上下文污染和租户级访问控制。~\cite{deepseekContextCaching2026,deepseekv32024,deepseekfyiFlashMLA2026}

    \item 开源高能力模型的下游安全治理：DeepSeek 的开放模型路线降低了构建本地智能体和垂直智能体的门槛，未来安全重点会从模型发布前评测扩展到微调后评测、蒸馏模型监管、下游工具接入审计和行业场景风险分级。~\cite{deepseekr12025,deepseekv32024,zhang2025chisafetydeepseek}

    \item 数学与形式化推理智能体的可验证安全：DeepSeekMath 和 DeepSeek-Prover-V2 表明形式化推理会成为智能体可靠性的核心路径，未来需要把形式证明、子目标分解验证和可执行检查器用于高风险工具调用、合约代码、策略配置和安全规则验证。~\cite{deepseekproverv22025,marjanovic2025thoughtology}

    \item 高性能推理基础设施的安全可观测性：DeepSeek-V3、DualPipe 和 FlashMLA 代表了大规模低成本推理趋势，未来 agent 安全需要在高吞吐系统中实现细粒度日志、异常行为检测、速率隔离、资源滥用防护和多租户安全监控。~\cite{deepseekfyiDualPipe2026,deepseekRateIsolation2026}
  \end{itemize}
& \begin{itemize}[leftmargin=*,topsep=0pt,partopsep=0pt,itemsep=1pt,parsep=0pt]
    \item DeepSeek API、JSON Output、OpenAI-compatible API 调用格式 ~\cite{deepseekapi2026,deepseekjsonoutput2026}
    \item DeepSeek Tool Calls、Strict Mode / Strict Function Calling ~\cite{deepseektoolcalls2026}
    \item DeepSeek-R1、DeepSeek-V3 系列技术报告和开源模型仓库 ~\cite{deepseekr12025,deepseekv3github2025}
    
  \end{itemize} \\
  \cline{2-5}

& 科大讯飞
& \begin{itemize}[leftmargin=*,topsep=0pt,partopsep=0pt,itemsep=1pt,parsep=0pt]
    \item 星火大模型与语音多模态AI
    \item 星火企业智能体平台与星辰智能体开发
    \item 行业垂直智能体与智能编程助手
  \end{itemize}
& \begin{itemize}[leftmargin=*,topsep=0pt,partopsep=0pt,itemsep=1pt,parsep=0pt]
    \item 企业信源、知识库、工具调用与任务规划边界
    \item 插件、MCP Server、知识库与工作流编排配置
    \item 私域知识检索、企业级知识管理与行业问答治理
  \end{itemize}
& \begin{itemize}[leftmargin=*,topsep=0pt,partopsep=0pt,itemsep=1pt,parsep=0pt]
    \item 星火企业智能体平台：企业信源、知识库、工具/技能接入和智能任务规划 ~\cite{iflyteksparkenterpriseagent2026}
    \item 讯飞星辰智能体开发平台及开发指南（支持插件、MCP Server、知识库和工作流节点配置）~\cite{iflytekxingchenagent2026,iflytekagentguide2026}
    \item 星火知识库与星火智能体场景文档（可用于企业私域知识问答和低代码智能体构建）~\cite{iflyteksparkknowledgebase2026,iflyteksceneagent2026}
  \end{itemize} \\
\cline{2-5}

& 360
& \begin{itemize}[leftmargin=*,topsep=0pt,partopsep=0pt,itemsep=1pt,parsep=0pt]
    \item 大模型/Agent 安全防护与风险治理
    \item 安全运营智能体与安全大模型
    \item 企业智能体工厂与 Skill 供应链治理
  \end{itemize}
& \begin{itemize}[leftmargin=*,topsep=0pt,partopsep=0pt,itemsep=1pt,parsep=0pt]
    \item 大模型内容护栏、内容评测、模型漏洞检测与智能体行为管控
    \item 智能体生命周期风险、Skill 入口、框架层、生态层和沙箱层安全分析
    \item 安全运营自动化、告警研判、威胁狩猎与响应闭环
  \end{itemize}
& \begin{itemize}[leftmargin=*,topsep=0pt,partopsep=0pt,itemsep=1pt,parsep=0pt]
    \item 360 大模型卫士防护系统/大模型安全卫士：内容安全、模型漏洞检测、敏感问题代答、AI智能体安全智能体 ~\cite{threesixtymodelguard2026,threesixtywhitepaper2025}
    \item 360《智能体安全实践报告》与《智能体安全报告》：攻击面、Skill 风险、行为审计与智能防御路线 ~\cite{threesixtyagentsafetyreport2025,threesixtyagentskillsreport2026}
    \item 360 安全运营智能体：9大AI安全智能体矩阵、告警研判、威胁狩猎和响应处置 ~\cite{threesixtysecurityagent2026}
  \end{itemize} \\
\hline
\end{longtblr}
\endgroup


% \begin{center}
% % \nopagebreak[4]
% \enlargethispage{1cm}
% \footnotesize
% \setlength{\tabcolsep}{3pt}
% % \renewcommand{\arraystretch}{0.7}
% \begin{longtable}{
% |>{\rule{0pt}{2pt}}m{0.8cm}
% |>{\rule{0pt}{2pt}}m{1.cm}
% |>{\rule{0pt}{2pt}}m{3.0cm}
% |>{\rule{0pt}{2pt}}m{0.25\linewidth}
% |>{\rule{0pt}{2pt}}m{\dimexpr\linewidth - 0.8cm - 1.cm - 3.0cm - 0.25\linewidth - 10\tabcolsep\relax}|}
% \caption{我国主流 AI 智能体安全产品与落地方向}\label{tab:china-agent-security-products-verified-cited}\\
% \hline
% \textbf{国家} & \textbf{企业} & \textbf{所属产业领域} & \textbf{安全研究方向} & \textbf{产品/成果} \\
% \hline
% \endfirsthead
% \multicolumn{5}{l}{\tablename\ \thetable\ -- 续上页}\\
% \hline
% \textbf{国家} & \textbf{企业} & \textbf{所属产业领域} & \textbf{安全研究方向} & \textbf{产品/成果} \\
% \hline
% \endhead
% \hline
% \multicolumn{5}{r}{续下页}\\
% \endfoot
% \hline
% \endlastfoot

% \multirow{12}{=}{中国大陆}
% & \multirow{3}{=}{百度}
% & \begin{itemize}[tabitem]
%     \item 搜索增强大模型
%     \item 智能云智能体开发平台
%     \item 飞桨与国产AI生态
%   \end{itemize}
% & \begin{itemize}[tabitem]
%     \item 企业智能体构建、RAG、工具/MCP接入与模型服务治理
%     % \item 企业智能体任务执行、Skill生态与消息平台集成
%     \item 国产AI框架生态、模型开发与工程化部署基础
%   \end{itemize}

% \vspace{0.5em}
% 2026重点研究方向：
% \vspace{1.5em}

% \begin{itemize}[tabitem]
%     \item 面向智能体复杂任务推理，研究更高效的强化学习、长上下文奖励分配与推理能力优化\cite{baidu_acl_ai_safety_2026}
%     \item 构建覆盖数据、身份权限、记忆、工具调用和外部系统连接的智能体全链路安全基础设施\cite{baidu_agent_security_center_2026}
%     \item 加强 Skill/MCP 供应链安全，研究恶意指令、过度权限、凭证泄露与高风险工具调用的准入检测和持续巡检\cite{baidu_ai_security_solution_2026,baidu_qianfan_platform_2026}
%     \item 推进策略即代码、输入输出安全策略和语义干预，使企业能够按业务场景定义、执行并审计智能体行为规则\cite{baidu_qianfan_security_policy_2026,baidu_qianfan_intervention_policy_2026}
%   \end{itemize}

% % • 面向智能体复杂任务推理，研究更高效的强化学习、长上下文奖励分配与推理能力优化
% % • 推进精细化模型安全对齐，通过参数级风险诊断和局部微调，降低安全对齐对模型通用能力的损耗
% % • 构建覆盖数据、身份权限、记忆、工具调用和外部系统连接的智能体全链路安全基础设施
% % • 加强 Skill/MCP 供应链安全，研究恶意指令、过度权限、凭证泄露与高风险工具调用的准入检测和持续巡检
% % • 推进策略即代码、输入输出安全策略和语义干预，使企业能够按业务场景定义、执行并审计智能体行为规则



% & \begin{itemize}[tabitem]
%     \item 百度千帆智能体开发平台：Agent引擎、工具/MCP、模型服务与企业服务能力 ~\cite{baiduqianfanagent2026}
%     \item 百度DuClaw：AI智能体服务、Skill 系统、SkillMarket与消息平台配置 ~\cite{baiduduclaw2026,baiduduclawskillmarket2026}
%     \item 飞桨PaddlePaddle与百度智能云千帆工程生态 ~\cite{paddlepaddle2026,baiduqianfanagent2026}
%   \end{itemize} \\
% \cline{2-5}
% & \multirow{3}{=}{阿里巴巴/阿里云}
% & \begin{itemize}[tabitem]
%     \item 通义/Qwen模型与百炼智能体应用
%     \item 云原生Workflow、RAG与插件编排
%     \item OpenClaw运行时与AI应用安全护栏
%   \end{itemize}
% & \begin{itemize}[tabitem]
%     \item 输入输出安全护栏、Prompt注入检测、敏感信息和恶意链接识别
%     \item 智能体运行时防护、工具调用审计、Skill检测与密钥托管
%     \item RAG/插件/代码解释器边界治理与企业应用编排
%   \end{itemize}

% \vspace{0.5em}
% 2026重点研究方向：
% \vspace{1.5em}

% \begin{itemize}[tabitem]
%     \item 构建覆盖智能体开发期与运行期的全生命周期防护，统一管理模型、知识库、身份、工具、MCP及Skill资产关系\cite{alibaba_agent_security_center_2026}
%     \item 加强智能体身份和动态权限治理，研究人机身份映射、细粒度授权、凭证托管以及关键操作二次确认\cite{alibaba_hitl_agent_security_2026}
%     \item 面向Skill投毒、间接提示注入、工具投毒和恶意命令执行，发展运行时行为检测、阻断和AI红队测试\cite{alibaba_agent_risk_detection_2026}
%     \item 研究Agent原生API和网络安全，通过语义化敏感数据识别、动态脱敏、外联访问控制与全链路行为审计保护数据流\cite{alibaba_agentic_api_security_2026,alibaba_agent_firewall_2026}
%     \item 推进多智能体协同安全运营，以专业Agent集群完成威胁感知、攻击链推理、协同调查及闭环响应\cite{alibaba_agentic_soc_2026}
% \end{itemize}
  
% & \begin{itemize}[tabitem]
%     \item 阿里云AI安全护栏/OpenClaw运行时防护方案 ~\cite{aliyunaiguardrails2026,aliyunopenclawdocs2026}
%     \item openclaw-security-assistant 与 OpenClaw runtime security 插件 ~\cite{aliyunopenclawassistant2026,aliyunopenclawdocs2026}
%     \item 阿里云百炼Model Studio智能体应用/工作流、Assistant API代码解释器 ~\cite{aliyunbailian2026,aliyunopenclawdocs2026}
%   \end{itemize} \\
% \cline{2-5}
% & \multirow{3}{=}{腾讯}
% & \begin{itemize}[tabitem]
%     \item 混元大模型与腾讯云ADP
%     \item LLM+RAG、Workflow与Multi-Agent 应用
%     \item CodeBuddy代码智能体与MCP工具生态
%   \end{itemize}
% & \begin{itemize}[tabitem]
%     \item MCP/智能体接入网关防护、检测规则、限流、鉴权与脱敏
%     \item 租户隔离、内容风控、权限审计、监控告警与日志分析
%     \item 代码智能体的本地工程访问、MCP工具调用和命令执行安全
%   \end{itemize}

% \vspace{0.5em}
% 2026重点研究方向：
% \vspace{1.5em}

% \begin{itemize}[tabitem]
%     \item 加强Agent Skills与MCP供应链安全研究，重点评估恶意指令、硬编码凭证、危险依赖、权限滥用及命令注入\cite{tencent_agent_skills_security_2026}
%     \item 建设面向真实Skill生态的标准化安全评测基准，同时衡量Skill自身可信度和外部扫描器的检测效果、误报及漏报\cite{tencent_skilltrustbench_2026}
%     \item 持续研究价值对齐、隐私保护、深度伪造检测、智能体安全和AI生成内容检测等从模型内生安全到系统防御的综合安全体系\cite{tencent_esg_report_2025_published_2026}
% \end{itemize}

% % • 加强Agent Skills与MCP供应链安全研究，重点评估恶意指令、硬编码凭证、危险依赖、权限滥用及命令注入
% % • 建设面向真实Skill生态的标准化安全评测基准，同时衡量Skill自身可信度和外部扫描器的检测效果、误报及漏报
% % • 推进“Agent扫描Agent”的自动化安全研究，由安全智能体理解技能语义、比对功能描述与实际行为，并验证风险路径
% % • 建立覆盖基础设施、代码、模型、RAG、工作流和运行时行为的智能体红队演习与上线前安全基线
% % • 发展以AI Agent安全网关为核心的统一控制面，对模型、API和MCP实施身份认证、限流、注入检测、敏感数据脱敏及行为审计
% % • 持续研究价值对齐、隐私保护、深度伪造检测、智能体安全和AI生成内容检测等从模型内生安全到系统防御的综合安全体系
  
% & \begin{itemize}[tabitem]
%     \item 腾讯云AI智能体安全网关：MCP Server封装、检测规则配置、Prompt注入/脱敏/访问鉴权等防护 ~\cite{tencentagentgateway2026,tencentagentgatewayquickstart2026}
%     \item 腾讯云智能体开发平台（Tencent Cloud ADP）多租户隔离、内容安全、权限与操作审计、监控告警 ~\cite{tencentadp2026,tencenthunyuan2026}
%     \item 腾讯云代码助手CodeBuddy：默认只读、命令/修改审批、MCP本地服务器与安全配置 ~\cite{tencentcodebuddysecurity2026,tencentcodebuddymcp2026}
%   \end{itemize} \\
% \cline{2-5}
% & \multirow{3}{=}{智谱}
% & \begin{itemize}[tabitem]
%     \item GLM大模型
%     \item 长时程Coding Agent与应用级智能体
%     \item GUI/手机智能体与多模态工具调用
%     \item 企业大模型API、模型精调、推理与评测平台
%   \end{itemize}
% & \begin{itemize}[tabitem]
%     \item 长时程代码智能体训练与评测中的反奖励作弊，检测并阻断智能体读取隐藏测试数据、复制参考答案及绕过任务的工具调
%     \item 长上下文、高并发智能体服务的稳定性与可靠性研究，重点解决乱码、重复、罕见字符生成及长任务运行异常\cite{zhipu_coding_agent_serving_reliability_2026,zhipu_glm51_process_quality_2026}
%     \item GUI/手机智能体的敏感操作确认、人工接管和受控执行，避免智能体在登录、验证码、支付等高风险场景中自主越权\cite{zhipu_open_autoglm_2026}
%     \item 模型输入输出及多模态内容安全审核，实现风险分级、违规内容定位、输入拦截、输出限制和生成终止\cite{zhipu_security_audit_2026}
%     \item 面向开放式机器学习系统的稳健性、公平性、隐私与可追溯治理，建设可审核、可监督、可追溯的模型服务\cite{zhipu_security_risk_notice_2026}
%   \end{itemize}
% & \begin{itemize}[tabitem]
%     \item GLM-5.2、GLM-5.1：面向长时程任务、工具调用和Agentic Engineering的旗舰模型\cite{zhipu_autoglm_phone_documentation_2026,zhipu_glm52_2026,glm5team2026glm5,liu2024autoglm} %最后一篇是有关GUI自主智能体的技术基础
%     \item Z.ai Agent Mode、ZCode：支持内置Skills、多轮协作、长任务、远程开发及移动设备控制
%     \item Open-AutoGLM、AutoGLM-Phone-9B：开源手机GUI智能体模型与框架，内置敏感操作确认和人工接管\cite{zhipu_open_autoglm_2026}
%     \item 智谱大模型开放平台：模型API、智能体开发、模型精调、推理、评测及多层安全防护\cite{zhipu_bigmodel_platform_2026}
%     \item Content Moderation API：支持文本、图片、音频和视频内容安全检测\cite{zhipu_content_moderation_api_2026,zhipu_security_audit_2026}
%   \end{itemize} \\
% \cline{2-5}
% % & \multirow{3}{=}{华为}
% % & \begin{itemize}[tabitem]
% %     \item AgentArts企业智能体平台与OfficeAce
% %     \item 盘古大模型、ModelArts Studio与云端智能体运行时
% %     \item CodeArts代码智能体与昇腾/鲲鹏工程生态
% %   \end{itemize}
% % & \begin{itemize}[tabitem]
% %     \item 身份鉴权、API Key 认证、沙箱工具与MCP网关接入控制
% %     \item 高危动作审批护栏与PC端安全审批管理
% %     \item 代码数据边界、文件保护、权限控制和隔离执行
% %   \end{itemize}
% % & \begin{itemize}[tabitem]
% %     \item 华为云AgentArts/AgentIdentity/Sand b\allowbreak -ox Tool：API Key认证、代码执行与文件管理沙箱、MCP Gateway   %\allowbreak允许断行 否则华为云这几个字的间距会被拉的很大
% %     ~\cite{huaweiagentarts2026,huaweiagentidentity2026,huaweiagentartssandbox2026}
% %     \item OfficeAce/AgentArts PC安全审批管理、对敏感操作进行审批拦截和授权确认 ~\cite{huaweiofficeaceapproval2026}
% %     \item 华为云码道（CodeArts）：dotfile保护、读写边界、沙箱执行、高危命令阻断与风险提示 ~\cite{huaweicodearts2026,huaweicodeartssandbox2026}
% %   \end{itemize} \\
% % \cline{2-5}
% & \multirow{3}{=}{字节跳动}
% & \begin{itemize}[tabitem]
%     \item 豆包大模型与火山方舟AI云
%     \item 扣子Coze智能体、工作流与Skill生态
%     \item 飞书aily/智能伙伴与组织协同Agent
%   \end{itemize}
% & \begin{itemize}[tabitem]
%     \item 大模型应用防火墙、Prompt攻击拦截、模型滥用和敏感数据泄露防护
%     \item Skill最小授权、凭证保护、OAuth PKCE与企业成员权限管理
%     \item 权限继承、数据保护、模型输入控制策略与组织协同治理
%     未来研究方向趋势：
%     % \item 智能体安全评测与部署仿真：未来需要面向真实长时任务、工具调用轨迹和模拟部署环境，提前评估智能体上线后的越权、误操作和数据泄露风险。~\cite{bytedanceseedpublicpapers2026,zhu2026edgebench}
%     % \item 工具调用与非可信内容隔离：未来需要在系统层面隔离任务指令、外部网页/文档内容和工具权限，降低搜索、浏览器、代码仓库等工具型 Agent 的间接提示注入风险。~\cite{he2025pasa}
%     % \item 长期记忆的安全治理：未来需要研究 Agent 记忆的选择、遗忘、权限控制和隐私保护，避免多模态长期记忆积累敏感信息或被跨任务滥用。~\cite{zou2026taskmem,zhang2025functiontokens}
%     \item 多模态/具身 Agent 的误操作风险控制：未来需要从离线成功率评测转向在线干预、异常检测、动作约束和人类接管机制，降低 VLA 与机器人 Agent 的真实环境风险。~\cite{zhang2026uam,li2026handitloop}
%     \item 面向科研与专业知识工作的可信 Agent：未来需要强化引用可验证性、证据链追踪、幻觉检测和抗操纵排序，避免科研 Agent 在高知识密度场景中放大错误信息。~\cite{he2025pasa}
%     \item 从 Responsible AI 到 Agent Security 的机制化转化：未来可将评测基准、红队样例、运行时监控、权限策略和记忆治理整合为统一的 Agent 安全工程体系。~\cite{bytedanceseedresponsibleai2026,he2025pasa,zou2026taskmem,zhang2025functiontokens}


%   \end{itemize}
% & \begin{itemize}[tabitem]
%     \item 火山引擎大模型应用防火墙（覆盖输入/输出防护、Prompt攻击、恶意链接、敏感数据和模型滥用风险）~\cite{volcenginellmfirewall2026,volcenginesecuritywhitepaper2026}
%     \item 扣子Coze技能安全指南、OAuth PKCE授权流程和企业版权限管理 ~\cite{cozeskillsecurity2026,cozeoauthpkce2026,volcenginecozepro2026}
%     \item 飞书智能伙伴/aily数据保护与策略管理，支持敏感数据分级、模型输入控制与组织级权限一致性 ~\cite{feishuaily2026,feishuailydataprotection2026}
%   \end{itemize} \\
% \cline{2-5}

% & \multirow{3}{=}{月之暗面 Kimi}
% & \begin{itemize}[tabitem]
%     \item Kimi K2/K2.7 Code 的工具调用、长上下文与 Agentic Coding
%     \item Kimi Code、Claude Code、Cline、RooCode 等开发者工作流接入
%     \item Web Search、官方工具、MCP Server 与多智能体/长时任务执行
%   \end{itemize}
% & \begin{itemize}[tabitem]
%     \item 工具调用参数、外部网页内容与用户指令边界
%     \item MCP Server、官方工具、代码执行器与文件/网络权限配置
%     \item 长时编码、多智能体协作与上下文压缩过程中的越权和误操作
%     未来研究方向趋势：
%     \item 工具调用与 MCP 最小权限治理：未来 Kimi 需要围绕 Tool Calls、官方工具和 MCP Server 建立细粒度授权、参数校验、调用审计与可回滚机制。~\cite{kimiapi_toolcalls2026,kimiapi_officialtools2026,kimiapi_mcp2026}
%     \item 长时编码 Agent 的沙箱化执行：面向 Kimi Code 和 K2.7 Code 的长程编程任务，应强化代码执行、文件修改、依赖安装和网络访问的隔离验证。~\cite{kimiapi_k27code2026,kimik26release2026}
%     \item Web Search 与外部内容的间接提示注入防御：Kimi 的搜索、抓取和网页阅读能力需要区分可信指令与非可信内容，并配套动态攻击评测。~\cite{kimiapi_websearch2026}
%     \item 多智能体协作的监督与责任归因：面向 Swarm 和长时任务执行，应研究计划级监控、子智能体权限分层、结果归并审计和异常中止策略。~\cite{kimik26release2026,kimik2thinkingblog2025}
%     \item 记忆、上下文压缩与安全可解释性：未来 Kimi 可进一步研究长上下文压缩、持久记忆和工具反馈在跨任务复用时的隐私保护与可解释审计。~\cite{kimiapi_officialtools2026,kimiapi_k27code2026}
%   \end{itemize}
% & \begin{itemize}[tabitem]
%     \item Kimi K2 面向工具使用、推理和自主问题求解优化，并在技术报告中强调 agentic data synthesis 与环境交互训练~\cite{kimik2technicalreport2025,moonshotkimik2github2025}
%     \item Kimi K2.7 Code 支持 256K 上下文、长程编码、多步工具调用和主流 Agent 编程工具接入~\cite{kimiapi_k27code2026}
%     \item Kimi API 提供 tool calls、Web Search、官方工具和 ModelScope MCP Server 接入能力，涉及搜索、抓取、代码执行、记忆和外部服务同步等安全边界~\cite{kimiapi_toolcalls2026,kimiapi_websearch2026,kimiapi_officialtools2026,kimiapi_mcp2026}
%   \end{itemize} \\
% \cline{2-5}

% & \multirow{3}{=}{MiniMax}
% & \begin{itemize}[tabitem]
%     \item 通用大模型、Agentic Coding 与长上下文生产力智能体
%     \item 多模态生成、语音/音乐/视频与角色陪伴式 Agent
%     \item 开放平台 API、工具调用、MCP 与企业级智能体应用
%   \end{itemize}
% & \begin{itemize}[tabitem]
%     \item 长程 Agent 任务中的工具调用、浏览搜索、代码执行与权限边界
%     \item 多智能体协作、长上下文记忆压缩与任务状态继承的可审计性
%     \item 语音复刻、视频生成、角色扮演 Agent 的身份滥用、内容安全与隐私治理
%     未来研究方向趋势：
%     \item 长程任务 Agent 的可审计执行：MiniMax Agent/Mavis 和 M 系列模型面向长时间、多步骤任务，未来需要强化任务计划、工具调用、状态继承和最终交付物的全链路审计。~\cite{minimaxagentteam2026,minimaxm25blog2026}

%     % \item 工具调用与 MCP 权限治理：随着 Function Call、MCP 和 Server-side Tools 接入增多，未来需要细粒度授权、参数校验、调用回放和最小权限沙箱。~\cite{minimaxfunctioncall2026,minimaxmcp2026,minimaxservertools2026}

%     \item Agent 强化学习的安全奖励建模：Forge 将 RL 推向真实世界 Agent，未来需要把越权调用、错误执行、隐私泄露和不可逆操作纳入 reward/verifier 体系。~\cite{minimaxforge2026}

%     % \item 编程 Agent 的代码执行安全：MiniMax Code 与 M 系列编码模型需要围绕仓库修改、依赖安装、测试执行和终端操作建立隔离环境与风险分级审批。~\cite{minimaxm2blog2025,minimaxm3blog2026,jimenez2024swebench}

%     \item 多模态与角色 Agent 的身份安全：MiniMax 在语音、视频、音乐和 Role-Play Agent 上布局较深，未来需要加强语音复刻授权、合成内容标识、人格越界和未成年人保护。~\cite{minimaxvoiceclone2026,minimaxm2her2026}
%   \end{itemize}
% & \begin{itemize}[tabitem]
%     \item MiniMax M2/M2.5/M3 面向代码、工具调用、搜索、办公和智能体任务优化，覆盖 SWE-Bench、BrowseComp 等真实任务基准~\cite{minimaxm2blog2025,minimaxm25blog2026,minimaxm3blog2026}
%     \item MiniMax Agent Team/Mavis、Forge Agent RL 框架面向长程任务、持续进化和真实世界 Agent 强化学习~\cite{minimaxagentteam2026,minimaxforge2026}
%     \item MiniMax 开放平台提供 Function Call、MCP、Server-side Tools、Mini-Agent 与语音复刻等能力，形成工具接入和多模态 Agent 产品底座~\cite{minimaxfunctioncall2026,minimaxmcp2026,minimaxservertools2026,minimaxminiagent2026,minimaxvoiceclone2026}
%   \end{itemize} \\
% \cline{1-4}


% & \multirow{3}{=}{Deep- Seek}
% & \begin{itemize}[tabitem]
%     \item 开放权重基础模型与推理模型
%     \item API 服务与开发者工具调用生态
%     \item 开源模型二次开发、本地部署与安全评测
%   \end{itemize}
% & \begin{itemize}[tabitem]
%     \item API / 模型输入输出治理、JSON 输出与兼容式调用
%     \item 工具调用边界、JSON Schema 校验与严格函数调用
%     \item 开源权重供应链、推理模型透明度与安全评测基础
%     未来研究方向趋势：
%     % \item 推理型智能体的安全评测与过程监督：DeepSeek-R1 展示了强化学习驱动的长链路推理能力，但未来安全研究需要评测推理轨迹中的目标偏移、隐性越权、错误自洽和高风险任务放大效应，而不只是看最终回答是否安全。~\cite{deepseekr12025,zhou2025hiddenrisksr1,marjanovic2025thoughtology}

%     % \item 工具调用智能体的权限边界与调用审计：随着 DeepSeek API 支持 tool calls，未来重点会转向函数调用前的意图校验、参数约束、最小权限授权、调用日志审计和失败回滚，防止模型把普通推理错误转化为真实外部副作用。~\cite{deepseekToolCalls2026,debenedetti2024agentdojo,xiang2026systemdefenses}

%     % \item 代码智能体的执行隔离与供应链安全：DeepSeek-Coder 系列提升了代码生成和修复能力，也意味着未来需要重点研究代码执行沙箱、依赖安装约束、仓库权限隔离、自动补丁验证和恶意代码/后门检测。~\cite{deepseekcoderv22024,deepseekToolCalls2026,xiang2026systemdefenses}

%     % \item 长上下文与缓存机制下的数据隔离：DeepSeek 的 context caching 和高效推理基础设施会推动企业级长任务智能体落地，但也需要研究缓存命中隔离、跨会话数据残留、敏感上下文污染和租户级访问控制。~\cite{deepseekContextCaching2026,deepseekv32024,deepseekfyiFlashMLA2026}

%     \item 开源高能力模型的下游安全治理：DeepSeek 的开放模型路线降低了构建本地智能体和垂直智能体的门槛，未来安全重点会从模型发布前评测扩展到微调后评测、蒸馏模型监管、下游工具接入审计和行业场景风险分级。~\cite{deepseekr12025,deepseekv32024,zhang2025chisafetydeepseek}

%     \item 数学与形式化推理智能体的可验证安全：DeepSeekMath 和 DeepSeek-Prover-V2 表明形式化推理会成为智能体可靠性的核心路径，未来需要把形式证明、子目标分解验证和可执行检查器用于高风险工具调用、合约代码、策略配置和安全规则验证。~\cite{deepseekproverv22025,marjanovic2025thoughtology}

%     \item 高性能推理基础设施的安全可观测性：DeepSeek-V3、DualPipe 和 FlashMLA 代表了大规模低成本推理趋势，未来 agent 安全需要在高吞吐系统中实现细粒度日志、异常行为检测、速率隔离、资源滥用防护和多租户安全监控。~\cite{deepseekfyiDualPipe2026,deepseekRateIsolation2026}
%   \end{itemize}
% & \begin{itemize}[tabitem]
%     \item DeepSeek API、JSON Output、OpenAI-compatible API 调用格式 ~\cite{deepseekapi2026,deepseekjsonoutput2026}
%     \item DeepSeek Tool Calls、Strict Mode / Strict Function Calling ~\cite{deepseektoolcalls2026}
%     \item DeepSeek-R1、DeepSeek-V3 系列技术报告和开源模型仓库 ~\cite{deepseekr12025,deepseekv3github2025}
    
%   \end{itemize} \\
%   \cline{2-5}

% & \multirow{3}{=}{科大讯飞}
% & \begin{itemize}[tabitem]
%     \item 星火大模型与语音多模态AI
%     \item 星火企业智能体平台与星辰智能体开发
%     \item 行业垂直智能体与智能编程助手
%   \end{itemize}
% & \begin{itemize}[tabitem]
%     \item 企业信源、知识库、工具调用与任务规划边界
%     \item 插件、MCP Server、知识库与工作流编排配置
%     \item 私域知识检索、企业级知识管理与行业问答治理
%   \end{itemize}
% & \begin{itemize}[tabitem]
%     \item 星火企业智能体平台：企业信源、知识库、工具/技能接入和智能任务规划 ~\cite{iflyteksparkenterpriseagent2026}
%     \item 讯飞星辰智能体开发平台及开发指南（支持插件、MCP Server、知识库和工作流节点配置）~\cite{iflytekxingchenagent2026,iflytekagentguide2026}
%     \item 星火知识库与星火智能体场景文档（可用于企业私域知识问答和低代码智能体构建）~\cite{iflyteksparkknowledgebase2026,iflyteksceneagent2026}
%   \end{itemize} \\
% \cline{2-5}

% & \multirow{3}{=}{360}
% & \begin{itemize}[tabitem]
%     \item 大模型/Agent 安全防护与风险治理
%     \item 安全运营智能体与安全大模型
%     \item 企业智能体工厂与 Skill 供应链治理
%   \end{itemize}
% & \begin{itemize}[tabitem]
%     \item 大模型内容护栏、内容评测、模型漏洞检测与智能体行为管控
%     \item 智能体生命周期风险、Skill 入口、框架层、生态层和沙箱层安全分析
%     \item 安全运营自动化、告警研判、威胁狩猎与响应闭环
%   \end{itemize}
% & \begin{itemize}[tabitem]
%     \item 360 大模型卫士防护系统/大模型安全卫士：内容安全、模型漏洞检测、敏感问题代答、AI智能体安全智能体 ~\cite{threesixtymodelguard2026,threesixtywhitepaper2025}
%     \item 360《智能体安全实践报告》与《智能体安全报告》：攻击面、Skill 风险、行为审计与智能防御路线 ~\cite{threesixtyagentsafetyreport2025,threesixtyagentskillsreport2026}
%     \item 360 安全运营智能体：9大AI安全智能体矩阵、告警研判、威胁狩猎和响应处置 ~\cite{threesixtysecurityagent2026}
%   \end{itemize} \\
% \cline{2-5}
% \end{longtable}
% \end{center}


% ==========================================
\section{产业界与学术界联动}
自主智能体安全的技术演进与体系化治理，本质是学术界基础理论突破与产业界工程化落地双向赋能、闭环迭代的结果。
本节聚焦梳理产业界与学术界联动的协作模式、技术代差、风险痛点与区域特征，主要方向包括：开源基础大模型与原生智能体安全底座、通用商用智能体安全评估与治理规范与国产中文智能体评测基准与本土产学研安全体系。
% 为智能体安全领域的技术溯源、趋势研判与本土化建设提供实证依据。

% \subsection{开源基础大模型与原生智能体安全底座}
% 开源基础大模型是原生智能体安全的底层技术载体，其核心架构、安全护栏体系与研发主体更迭，直接决定智能体感知、决策、行动全链路的安全边界~\cite{nyuLeCun2026,googleGeminiEnterpriseAgentPlatform2026}。
% % 当前全球形成以Meta体系世界模型安全、Google体系云端企业智能体工程化、ETH‑CMU自动化攻防评测为核心的三类产学协同范式~\cite{metaLlamaFirewall2025,grayswanSeriesA2026}。
% 本小节梳理了头部机构核心学者、自研项目与商用产品，聚焦原生底座架构演进、安全护栏部署、攻防基准构建，厘清基础层智能体安全从学术理论到产业工程化的转化逻辑，主要发现如下(详见表~\ref{tab:industry_academia_agent_security})：
% \begin{itemize}[tabitem]
%     \item \textbf{全球原生智能体安全底座形成二元技术路线分化。}Meta依托JEPA世界模型聚焦真实世界表征与预测安全，借助多模态世界模型预判动作执行后果，从感知底层降低智能体越界风险~\cite{nyuLeCun2026,vjepa2025}；Google基于Gemini生态侧重多智能体协同与企业运行时风控，依托IAM权限体系、全链路追踪实现企业环境下智能体权限管控，底层技术范式相互独立~\cite{googleADK2025,deepmindCoScientist2025}。

%     \item \textbf{核心人员流动重构产学研发分工。}FAIR核心学者人事变动之后，Meta内部团队承接Llama系列安全护栏工程化迭代工作，持续优化Llama‑Firewall适配Agent场景~\cite{pineauDeparture2025,metaLlamaFirewall2025}；AMI Labs依托Yann LeCun团队独立开展JEPA架构商业化落地，面向医疗机器人、工业智能体开展世界模型安全的产业化验证，形成企业工程落地与学术技术转化的分工格局~\cite{amiLabsFunding2026,lejepaIdentifiability2026}。

%     \item \textbf{攻防评测形成完整产学闭环。}学术界苏黎世联邦理工SPY‑Lab开源AgentDojo平台，把间接提示注入、工具劫持问题构建成标准化可复现评测环境，夯实攻防实验的研究基础~\cite{agentdojoGithub2026}；产业端Gray Swan AI基于卡内基梅隆大学团队理论成果落地Cygnal、Shade工业级红队平台，成果被Anthropic采纳用于Claude安全评估；核心学者Zico Kolter进入OpenAI董事会，将自动化红队研究成果上升为顶层治理规则，完成学术成果向企业安全规范的闭环落地~\cite{latentspaceGrayswan2026,kolterWikipedia2026}。
% \end{itemize}


% \begin{center}
% % \nopagebreak[4]
% \enlargethispage{1cm}
% \footnotesize
% \setlength{\tabcolsep}{3pt}
% \begin{longtable}{|>{\rule{0pt}{2pt}}m{0.07\textwidth}
% |>{\rule{0pt}{2pt}}m{0.12\textwidth}
% |>{\rule{0pt}{2pt}}m{\dimexpr\linewidth - 0.07\textwidth - 0.12\textwidth - 6\tabcolsep - 3\arrayrulewidth\relax}|}
% \caption{开源基础底座、云端企业智能体与自动化攻防评测体系}
% \label{tab:industry_academia_agent_security} \\
% \hline
% \textbf{研究方向} & \textbf{企业/机构 + 核心学者} & \textbf{自研项目、产学合作、商用安全产品} \\
% \hline
% \endfirsthead
% \multicolumn{3}{l}{\tablename\ \thetable\ -- 续上页}\\
% \hline
% \textbf{研究方向} & \textbf{企业/机构 + 核心学者} & \textbf{自研项目、产学合作、商用安全产品} \\
% \hline
% \endhead
% \hline
% \multicolumn{3}{r}{续下页}\\
% \endfoot
% \hline
% \endlastfoot

% \multirow{6}{=}{开源基础大模型原生智能体安全底座}
% & Meta FAIR；Yann LeCun, Joelle Pineau
% & \begin{itemize}[tabitem]
% \item \textbf{世界模型安全底座：}以JEPA、Image-JEPA、Video-JEPA、V-JEPA 2构成学术路线演进，支撑多模态智能体的状态表征、感知预测和行动后果建模~\cite{nyuLeCun2026,vjepa2025}；
% \item \textbf{开源安全护栏体系：}公开Llama Guard、Prompt Guard和LlamaFirewall，其中LlamaFirewall定位为面向安全AI Agent的开源护栏系统~\cite{metaLlamaFirewall2025}；
% \item \textbf{团队组织演进：}Joelle Pineau卸任FAIR负责人、LeCun离开Meta后，相关研发由公司新组建的Meta Superintelligence Labs承接~\cite{pineauDeparture2025,lecunDeparture2025}。
% \end{itemize} \\
% \cline{2-3}
% & 纽约大学、AMI Labs；Yann LeCun
% & \begin{itemize}[tabitem]
% \item \textbf{JEPA架构产业化：}LeCun离职后创办AMI Labs推动JEPA架构商业化，首批目标行业为医疗、机器人和工业自动化~\cite{amiLabsFunding2026,lecunMITTechReview2026}；
% \item \textbf{理论与基准验证：}2026年发表的JEPA相关论文证明该架构在特定条件下能够恢复真实世界底层结构，同时用配套基准指出当前模型面对细微视觉变化时仍较脆弱~\cite{lejepaIdentifiability2026}。
% \end{itemize} \\
% \hline

% \multirow{9}{=}{云端原生企业级大模型智能体研发与风控}
% & Google DeepMind；Demis Hassabis
% & \begin{itemize}[tabitem]
% \item \textbf{多智能体协同基础理论：}AlphaStar、RoboCat等研究成果支撑多智能体协同博弈和机器人任务规划的基础理论；
% \item \textbf{企业级智能体产品落地：}相关研究能力被工程化为Gemini Enterprise Agent Platform和Managed Agents，落实到API注册、IAM权限管控和运行日志审计~\cite{googleGeminiEnterpriseAgentPlatform2026,googleADK2025}；
% \item \textbf{科研智能体协作平台：}AI Co-Scientist基于Gemini构建，由一组专业化Agent协同完成科学假设的生成、辩论、排序和迭代~\cite{deepmindCoScientist2025,gottweisAICoScientist2025}。
% \end{itemize} \\
% \cline{2-3}
% & Google Research；Jeff Dean
% & \begin{itemize}[tabitem]
% \item \textbf{云端Agent风控机制：}分布式大规模机器学习与系统工程能力支撑Gemini Enterprise Agent Platform的API注册、IAM权限管控和运行日志审计~\cite{googleADK2025,googleGeminiEnterpriseAgentPlatform2026}；
% \item \textbf{托管式智能体治理：}Managed Agents将全链路Tracing审计、IAM权限管控等系统工程理论嵌入托管式智能体服务，对应企业Agent运行时风控需求~\cite{googleADK2025}。
% \end{itemize} \\
% \cline{2-3}
% & 斯坦福大学医学院、MIT；Gary Peltz, Ritu Raman
% & \begin{itemize}[tabitem]
% \item \textbf{产学研安全验证：}AI Co-Scientist于2025年2月首发时联合全球100多个科研机构共同开发和测试，Gary Peltz、Ritu Raman公开验证其科研任务能力~\cite{deepmindCoScientist2025}；
% \item \textbf{科研落地与企业试用：}Co-Scientist已产出抗菌素耐药性、植物免疫等湿实验验证成果，并被第一三共、拜耳作物科学等企业内部试用~\cite{labcriticsCoScientist2026}。
% \end{itemize} \\
% \hline

% \multirow{5}{=}{智能体攻防评测基准与自动化红队体系}
% & 苏黎世联邦理工SPY Lab；Florian Tramèr
% & \begin{itemize}[tabitem]
% \item \textbf{自动化攻防基准：}AgentDojo将间接提示注入、工具劫持等标准化攻防任务集统一到动态评测环境中，已由ETH SPY Lab开源部署、~\cite{agentdojoGithub2026}；
% \item \textbf{基准复现实验：}AgentDojo使邮件、日历、文件和工具调用场景中的提示注入攻击具备可复现、可比较的实验对象。
% \end{itemize} \\*
% \nobreakcline{2-3}
% & CMU, Gray Swan AI；Zico Kolter, Matt Fredrikson, Andy Zou, Dawn Song
% & \begin{itemize}[tabitem]
% \item \textbf{工业级红队平台：}Gray Swan AI将模型鲁棒性、神经网络形式化验证理论转化为检测能力，代表性产品包括红队测试平台Cygnal和对抗性红队工具Shade~\cite{grayswanSeriesA2026}；
% \item \textbf{提示注入鲁棒性评测：}Shade已被Anthropic用于评估模型在代码环境中对提示注入攻击的鲁棒性，Gray Swan的间接提示注入研究也被Anthropic Claude Mythos的system card引用为权威研究依据~\cite{latentspaceGrayswan2026}；
% \item \textbf{模型风险治理嵌入：}Kolter已出任OpenAI董事会成员并兼任安全与安保委员会主席，使自动化红队研究与模型风险治理进入更高层级的产业治理流程~\cite{kolterWikipedia2026}。
% \end{itemize} 
% \end{longtable}
% \end{center}


% \subsection{通用商用智能体安全评估与治理规范}
% 随着云端商用智能体与工具生态规模化落地，安全治理从模型鲁棒性延伸至协议交互、运行时管控、隐私计算、外部审计的全生命周期。海外形成“企业研发‑学界定标‑第三方核验”的成熟产学体系，覆盖通用智能体评估、MCP工具协议、端侧硬件隐私三大方向\cite{openaiRedTeamNetwork2025,mcpFormalFramework2026,nvidiaConfidentialComputingProduct2026}。本小节结合头部厂商、国家级安全机构与顶尖高校成果，归纳商用智能体的治理范式、协议漏洞矛盾与硬件隐私安全落地规律，主要发现如下（详见表~\ref{tab:commercial_mcp_edge_agent_security}）：
% \begin{itemize}[tabitem]
%     \item \textbf{头部厂商建立标准化商用治理范式。}OpenAI搭建外部专家红队测试网络、在Agents‑SDK内置多层安全护栏，配套全链路Tracing追踪与人在回路审批机制，并发布System‑Card公开安全评估报告，构建起完整闭环体系，成为全球商用智能体安全治理标杆范式\cite{openaiChatGPTAgentSystemCard2025,openaiAgentsSDK2025,openaiGuardrailsApprovalsDoc2026}；英国UK‑AISI、FAR.AI开展第三方独立安全审计，斯坦福CRFM提出HELM评测理论，从第三方与学术层面完善评估方法论，形成产学核验的治理格局\cite{venturebeatChatGPTAgent2025}。

%     \item \textbf{工具协议存在结构性权责矛盾。}Anthropic推出的MCP协议生态规模快速扩张，但OX‑Security发现协议原生存在RCE高危漏洞；出于生态灵活性考量，厂商拒绝在架构层面完成彻底修复，把输入校验、风险净化的安全责任转嫁至下游开发者。伯克利CHAI团队从目标对齐理论剖析该问题本质，学术界也意识到借助ProVerif、Tamarin对MCP开展形式化验证难度很高，充分暴露出协议生态扩张与安全约束之间的固有冲突\cite{oxSecurityMCP2026,berkeleyCHAI2026,mcpFormalFramework2026}。

%     \item \textbf{智能体隐私安全构建垂直产学链条。}产业端Apple依靠PCC端云分层架构实现数据最小化处理与可验证云端推理，NVIDIA依托GPU硬件信任根搭建机密计算环境；学术层面EPFL‑SPRING‑Lab的隐私工程理论、NYU针对GPU的侧信道分析成果，为硬件可信基座、端侧隐私智能体提供理论支撑，形成理论研究‑硬件实现‑产品落地的完整链条\cite{springLab2026,karriNYU2026}。
% \end{itemize}


% \begin{center}
% % \nopagebreak[4]
% \enlargethispage{1cm}
% \footnotesize
% \setlength{\tabcolsep}{3pt}
% \begin{longtable}{|>{\rule{0pt}{2pt}}m{0.07\textwidth}
% |>{\rule{0pt}{2pt}}m{0.12\textwidth}
% |>{\rule{0pt}{2pt}}m{\dimexpr\linewidth - 0.07\textwidth - 0.12\textwidth - 6\tabcolsep - 3\arrayrulewidth\relax}|}
% \caption{通用商用智能体、MCP工具调用协议与端侧硬件隐私计算智能体}
% \label{tab:commercial_mcp_edge_agent_security} \\
% \hline
% \textbf{研究方向} & \textbf{企业/机构 + 核心学者} & \textbf{自研项目、产学合作、商用安全产品} \\
% \hline
% \endfirsthead
% \multicolumn{3}{l}{\tablename\ \thetable\ -- 续上页}\\
% \hline
% \textbf{研究方向} & \textbf{企业/机构 + 核心学者} & \textbf{自研项目、产学合作、商用安全产品} \\
% \hline
% \endhead
% \hline
% \multicolumn{3}{r}{续下页}\\
% \endfoot
% \hline
% \endlastfoot

% \multirow{6}{=}{通用商用智能体安全评估与治理规范}
% & OpenAI：Sam Altman、Greg Brockman
% & \begin{itemize}[tabitem]
% \item \textbf{ChatGPT Agent交互智能体：}搭建外部多专家红队测试网络，Red Teaming Network由16位具备生物安全背景的博士研究员组成，共提交110次攻击尝试，其中16次超过内部风险阈值~\cite{openaiChatGPTAgentSystemCard2025,openaiRedTeamNetwork2025}。
% \item \textbf{开源安全护栏体系：}OpenAI Agents SDK内置三层输入/输出/工具护栏，支持检测风险时立即中止运行的tripwire机制，配合运行时Tracing审计与高风险动作的human-in-the-loop审批机制~\cite{openaiAgentsSDK2025,openaiGuardrailsDoc2026,openaiTracingDoc2026,openaiGuardrailsApprovalsDoc2026}。
% \item \textbf{全球统一风险披露体系：}System Card持续迭代智能体能力、风险与透明度评估指标，全产品线强制落地，并以公开白皮书系统阐述外部专家参与风险评估的方法论~\cite{openaiChatGPTAgentSystemCard2025,openaiExternalRedTeaming2025}。
% \end{itemize} \\
% \cline{2-3}
% & 英国AI安全研究所UK AISI、FAR.AI独立安全实验室
% & \begin{itemize}[tabitem]
% \item \textbf{深度访问链路漏洞审计：}UK AISI获得对ChatGPT Agent内部推理链和策略文本的深度访问权限，完成官方漏洞审计测试，测出7个通用漏洞~\cite{venturebeatChatGPTAgent2025}。
% \item \textbf{商用智能体安全缺陷批评：}FAR.AI经40小时独立测试发现3个部分漏洞，提出商用智能体过度依赖运行时监控的结构性安全缺陷~\cite{venturebeatChatGPTAgent2025}。
% \end{itemize} \\
% \cline{2-3}
% & 斯坦福CRFM：Percy Liang
% & \begin{itemize}[tabitem]
% \item \textbf{全链路可比较评测理论对齐：}HELM全链路可比较评测理论与OpenAI Agents SDK的Tracing架构在方法论上形成理论对应，共同强调把系统行为拆解为可追溯、可比较的具体环节~\cite{openaiTracingDoc2026}。
% \end{itemize} \\
% \hline

% \multirow{7}{=}{智能体工具调用MCP协议与目标对齐安全}
% & Anthropic：Dario Amodei、Jared Kaplan
% & \begin{itemize}[tabitem]
% \item \textbf{Agent全生态工具交互协议：}MCP工具通信协议基于JSON-RPC 2.0构建，支撑Claude Agent全生态工具交互，截至2026年初已形成超1万活跃服务器、17.7万注册工具、月均9700万次SDK下载的生态规模~\cite{mcpFormalFramework2026}。
% \item \textbf{高危分级智能体系统：}Claude Code代码智能体、Mythos高危分级智能体系统均采用分层行为约束护栏~\cite{anthropicClaudeCodeSecurity2026}。
% \item \textbf{协议架构安全责任取舍：}面对OX Security披露的架构级RCE漏洞，Anthropic确认该行为符合设计预期，以维持协议灵活性为由拒绝修复，将输入净化责任划给下游开发者~\cite{oxSecurityMCP2026}。
% \end{itemize} \\
% \cline{2-3}
% & 加州伯克利CHAI：Stuart Russell
% & \begin{itemize}[tabitem]
% \item \textbf{人类兼容AI理论与目标对齐：}Russell作为Berkeley CHAI创始主任，其人类兼容AI理论关注目标不确定性和系统可控性，与高自治Agent的目标漂移、价值对齐问题存在直接理论关联~\cite{berkeleyCHAI2026}。
% \end{itemize} \\
% \cline{2-3}
% & OX Security、Check Point安全厂商；密码学形式化验证学界
% & \begin{itemize}[tabitem]
% \item \textbf{MCP协议远程代码执行漏洞披露：}OX Security披露MCP协议"设计如此"的远程代码执行漏洞，影响超1.5亿次下载和20万台服务器，攻陷11个主流MCP市场中的9个，产生至少10个高危或严重级别CVE编号~\cite{oxSecurityMCP2026,thehackernewsMCP2026}。
% \item \textbf{Claude Code代码执行漏洞修复：}Check Point于2025年9至10月披露Claude Code三项漏洞（CVE-2025-59536，CVSS 8.7；CVE-2026-21852，CVSS 5.3），由厂商协同完整修复，分别在1.0.111版本和2.0.65版本发布安全补丁~\cite{checkpointClaudeCode2026,githubClaudeCodeAdvisory2026}。
% \item \textbf{MCP威胁分类与形式化验证：}学术界已提出MCP的威胁分类模型，并指出借鉴Tamarin、ProVerif等密码协议验证工具对MCP做形式化验证仍是开放且困难的研究方向~\cite{mcpFormalFramework2026,mcpSoK2026}。
% \end{itemize} \\
% \hline

% \multirow{5}{=}{端侧与硬件机密计算隐私智能体安全技术}
% & Apple：Craig Federighi 隐私安全实验室
% & \begin{itemize}[tabitem]
% \item \textbf{端云分层隐私架构：}Apple Intelligence移动端本地智能体采用端云分层数据最小化架构~\cite{applePCCSecurityGuide2026}。
% \item \textbf{可验证云端推理平台：}Private Cloud Compute（PCC）通过Virtual Research Environment向全球研究者开放本地复现能力，配合开源形式化验证工具、函数库与安全悬赏计划，构建透明可信推理体系~\cite{appleSecurityResearch2026}。
% \end{itemize} \\
% \cline{2-3}
% & EPFL SPRING Lab：Carmela Troncoso
% & \begin{itemize}[tabitem]
% \item \textbf{隐私工程理论支撑端侧架构：}Troncoso团队长期研究隐私工程、数据最小化和匿名化技术，相关理论与端侧/云端隐私智能体架构问题高度相关~\cite{springLab2026}。
% \end{itemize} \\
% \cline{2-3}
% & NVIDIA：Jensen Huang，NYU；Ramesh Karri
% & \begin{itemize}[tabitem]
% \item \textbf{GPU硬件机密计算平台：}Confidential Agent GPU硬件机密计算平台自H100起首次支持机密计算，通过片上硬件信任根实现安全启动和远程认证，在Blackwell架构上进一步延伸~\cite{nvidiaConfidentialComputingProduct2026}。
% \item \textbf{硬件侧信道安全研究支撑：}Karri教授长期研究硬件木马、芯片供应链安全和侧信道分析，与GPU侧信道信息泄露防护议题直接相关，学术界也已对TEE模式下大模型推理性能开销做独立评测~\cite{karriNYU2026}。
% \end{itemize} 

% \end{longtable}
% \end{center}




% \subsection{国产中文智能体评测基准与本土产学研安全体系}
% 中文语义特征、合规要求与国产化算力生态，决定我国智能体安全无法直接复用海外技术体系。我国聚焦中文场景评测、安全对齐训练、垂直行业多智能体风控\cite{pandaguardCASIA2026}。本小节基于我国核心机构成果，系统归纳本土安全体系的差异化特征、技术能力与产业落地路径，主要发现如下（详见表~\ref{tab:chinese_agent_benchmark_security}）：
% \begin{itemize}[tabitem]
%     \item \textbf{我国高校建成差异化中文评测体系。}清华大学依托ToolBench、AgentBench搭建工具调用和综合任务评测基准，适配海量真实API与中文环境；上海交通大学提出R‑Judge评测体系，聚焦对话交互场景下智能体行为风险判定，二者形成互补的本土化基准，填补海外评测集对中文语境适配不足的短板。

%     \item \textbf{国产安全对齐形成开源可复现技术栈。}北京大学PKU‑Alignment团队提出Safe RLHF与PKU‑SafeRLHF算法，配套BeaverTails高质量中文偏好数据集和align‑anything全模态开源框架，为我国大模型智能体后训练安全对齐提供可落地的基础组件\cite{pkuAlignAnything2026}；浙江大学依托AIcert平台完成大量国产模型安全评估，进一步完善对齐效果的检验手段，共同构筑国产智能体安全对齐训练的核心开源底座\cite{aicertZJU2026}。

%     \item \textbf{国产安全具备强行业导向特征：}中科院自动化所打造PandaGuard攻防评测平台与JidiAI博弈环境，从科研层面开展多智能体风险推演\cite{pandaguardCASIA2026,jidiAI2026}；上海AI实验室推出Intern‑Shannon安全智能体操作系统和SafeWork‑R1模型，面向政企场景完成产业级落地\cite{intershannon2026,safework2025}；华为联合上交IPADS团队搭建SMARTS自动驾驶仿真平台并开展芯片可信环境研究，面向交通行业完成落地验证\cite{noahArkSMARTS2026,ipadsSJTU2026}。
%     % 整体摒弃海外通用智能体研发路线，形成科研验证、平台研发、行业场景落地的分层发展模式。
% \end{itemize}



% \begin{center}
% % \nopagebreak[4]
% \enlargethispage{1cm}
% \footnotesize
% \setlength{\tabcolsep}{3pt}
% \begin{longtable}{|>{\rule{0pt}{2pt}}m{0.07\textwidth}
% |>{\rule{0pt}{2pt}}m{0.14\textwidth}
% |>{\rule{0pt}{2pt}}m{\dimexpr\linewidth - 0.07\textwidth - 0.14\textwidth - 6\tabcolsep - 3\arrayrulewidth\relax}|}
% \caption{国产中文智能体评测基准与本土产学研安全体系}
% \label{tab:chinese_agent_benchmark_security} \\
% \hline
% \textbf{研究方向} & \textbf{企业/机构 + 核心学者} & \textbf{自研项目、产学合作、商用安全产品} \\
% \hline
% \endfirsthead
% \multicolumn{3}{l}{\tablename\ \thetable\ -- 续上页}\\
% \hline
% \textbf{研究方向} & \textbf{企业/机构 + 核心学者} & \textbf{自研项目、产学合作、商用安全产品} \\
% \hline
% \endhead
% \hline
% \multicolumn{3}{r}{续下页}\\
% \endfoot
% \hline
% \endlastfoot

% \multirow{3}{=}{中文本土智能体评测与国产化标准体系}
% & 清华大学THUNLP/AIR；朱军、唐杰、刘知远、孙茂松
% & \begin{itemize}[tabitem]
% \item \textbf{工具调用评测基准：}ToolBench工具调用评测基准，覆盖16000多个真实API中文工具场景~\cite{tsinghuaAIInstitute2026,tsinghuaAIR2026}。
% \item \textbf{全栈智能体综合评测：}AgentBench八场景全栈智能体综合评测基准，联合俄亥俄州立大学、加州大学伯克利分校共建。
% \end{itemize} \\
% \cline{2-3}
% & \mbox{北京大学}PKU-Alignment；杨耀东、刘健
% & \begin{itemize}[tabitem]
% \item \textbf{安全强化学习训练框架：}Safe RLHF、PKU-SafeRLHF训练框架，BeaverTails中文安全偏好数据集；
% \item \textbf{全模态对齐训练框架：}Aligner、ProgressGym等后训练对齐系列成果，开源align-anything全模态对齐训练、数据和测评框架~\cite{pkuAlignAnything2026}。
% \end{itemize} \\
% \cline{2-3}
% & 上海交通大学；谷大武（LoCCS）、张倬胜（AI安全实验室）
% & \begin{itemize}[tabitem]
% \item \textbf{密码学与系统安全底座：}谷大武团队围绕密码理论、软硬件与系统安全展开研究，为智能体运行环境提供底层安全基础~\cite{sjtuLoCCS2026}。
% \item \textbf{智能体行为安全评测基准：}张倬胜团队提出R-Judge智能体行为安全评测基准，覆盖5类应用场景、27个风险场景与10类风险类型~\cite{sjtuAISecLab2026}。
% \end{itemize} \\
% \hline

% \multirow{3}{=}{我国顶尖高校开源智能体对齐安全体系}
% & 浙江大学区块链与数据安全全国重点实验室；任奎
% & \begin{itemize}[tabitem]
% \item \textbf{通用AI安全评测平台：}AIcertAI安全评测平台，覆盖数据质量、算法安全、模型安全、框架安全与系统安全检测，已评测LLaMA、Gemma、Qwen等35个开源大模型，并在淘宝直播风控、中车株洲所等场景形成示范应用~\cite{aicertZJU2026}。
% \item \textbf{企业级安全联合实验室：}浙江大学--阿里巴巴集团AI安全联合实验室将智能体安全评测列为首期核心研究方向，支撑企业级大模型Agent红队测试与安全治理~\cite{zjuAlibabaAISafetyLab2025}。
% \end{itemize} \\
% \cline{2-3}
% & 北京大学PKU-Alignment；杨耀东、刘健
% & \begin{itemize}[tabitem]
% \item \textbf{后训练安全对齐生态：}维护Safe RLHF、PKU-SafeRLHF、BeaverTails等安全对齐训练与数据体系，构成我国较早的可复现安全强化学习框架~\cite{pkuAlignAnything2026}。
% \end{itemize} \\
% \cline{2-3}
% & 上海交通大学AI安全实验室；张倬胜
% & \begin{itemize}[tabitem]
% \item \textbf{自然语言处理与自主智能体安全：}专注自然语言处理、预训练语言模型与自主智能体安全，R-Judge基准支撑智能体行为风险识别研究~\cite{sjtuAISecLab2026}。
% \end{itemize} \\
% \hline

% \multirow{3}{=}{我国科研院所、产业实验室行业智能体安全落地}
% & 中国科学院自动化所；曾毅、王飞跃团队
% & \begin{itemize}[tabitem]
% \item \textbf{多智能体博弈攻防评测平台：}PandaGuard（灵御）多智能体博弈攻防评测平台，集成19种攻击算法与12种防御机制~\cite{pandaguardCASIA2026}。
% \item \textbf{多智能体决策开放平台：}JidiAI多智能体开源开放平台面向智能体博弈决策领域，吸引我国外50余所高校机构注册使用~\cite{jidiAI2026}。
% \end{itemize} \\
% \cline{2-3}
% & 上海AI实验室；安全可信AI团队
% & \begin{itemize}[tabitem]
% \item \textbf{产业级安全智能体操作系统：}Intern-Shannon（书安）产业级安全智能体操作系统，以“安全即服务”模式支撑高安全场景部署~\cite{intershannon2026}。
% \item \textbf{安全能力协同演化与科研智能体：}SafeWork-R1通过SafeLadder框架实现安全与智能协同进化，依托InternLM/OpenMMLab开源底座；InternAgent构建面向自主科学发现的闭环多智能体框架~\cite{safework2025,internAgentSHLab2026}。
% \end{itemize} \\
% \hline
% & 华为诺亚方舟实验室；郝建业（决策推理实验室），上交陈海波（IPADS实验室）
% & \begin{itemize}[tabitem]
% \item \textbf{自动驾驶多智能体仿真平台：}SMARTS大规模多智能体强化学习仿真平台，支持数百至数千个智能体在真实交通场景中的交互学习~\cite{noahArkSMARTS2026}。
% \item \textbf{可信执行环境与操作系统级支撑：}sNPU可信执行环境研究被ISCA 2024接收，陈海波担任OpenHarmony技术指导委员会创始主席，IPADS与华为在鸿蒙操作系统上保持长期合作关系~\cite{ipadsSJTU2026}。
% \end{itemize} 

% \end{longtable}
% \end{center}
\subsection{开源基础大模型与原生智能体安全底座}
开源基础大模型是原生智能体安全的底层技术载体，其核心架构、安全护栏体系与研发主体更迭，直接决定智能体感知、决策、行动全链路的安全边界~\cite{nyuLeCun2026,googleGeminiEnterpriseAgentPlatform2026}。
% 当前全球形成以Meta体系世界模型安全、Google体系云端企业智能体工程化、ETH‑CMU自动化攻防评测为核心的三类产学协同范式~\cite{metaLlamaFirewall2025,grayswanSeriesA2026}。
本小节梳理了头部机构核心学者、自研项目与商用产品，聚焦原生底座架构演进、安全护栏部署、攻防基准构建，厘清基础层智能体安全从学术理论到产业工程化的转化逻辑，主要发现如下(详见表~\ref{tab:industry_academia_agent_security})：
\begin{itemize}[tabitem]
    \item \textbf{全球原生智能体安全底座形成二元技术路线分化。}Meta依托JEPA世界模型聚焦真实世界表征与预测安全，借助多模态世界模型预判动作执行后果，从感知底层降低智能体越界风险~\cite{nyuLeCun2026,vjepa2025}；Google基于Gemini生态侧重多智能体协同与企业运行时风控，依托IAM权限体系、全链路追踪实现企业环境下智能体权限管控，底层技术范式相互独立~\cite{googleADK2025,deepmindCoScientist2025}。

    \item \textbf{核心人员流动重构产学研发分工。}FAIR核心学者人事变动之后，Meta内部团队承接Llama系列安全护栏工程化迭代工作，持续优化Llama‑Firewall适配Agent场景~\cite{pineauDeparture2025,metaLlamaFirewall2025}；AMI Labs依托Yann LeCun团队独立开展JEPA架构商业化落地，面向医疗机器人、工业智能体开展世界模型安全的产业化验证，形成企业工程落地与学术技术转化的分工格局~\cite{amiLabsFunding2026,lejepaIdentifiability2026}。

    \item \textbf{攻防评测形成完整产学闭环。}学术界苏黎世联邦理工SPY‑Lab开源AgentDojo平台，把间接提示注入、工具劫持问题构建成标准化可复现评测环境，夯实攻防实验的研究基础~\cite{agentdojoGithub2026}；产业端Gray Swan AI基于卡内基梅隆大学团队理论成果落地Cygnal、Shade工业级红队平台，成果被Anthropic采纳用于Claude安全评估；核心学者Zico Kolter进入OpenAI董事会，将自动化红队研究成果上升为顶层治理规则，完成学术成果向企业安全规范的闭环落地~\cite{latentspaceGrayswan2026,kolterWikipedia2026}。
\end{itemize}

% 依赖项（必须放在主文档导言区，即 \begin{document} 之前）：
% \usepackage{tabularray}
% \UseTblrLibrary{varwidth}
% \usepackage{enumitem}
\begingroup
% \nopagebreak[4]
\enlargethispage{1cm}
\footnotesize
\centering

\DefTblrTemplate{conthead-text}{default}{（续表）}
\DefTblrTemplate{contfoot-text}{default}{续下页}

\begin{longtblr}[
  caption = {开源基础底座、云端企业智能体与自动化攻防评测体系},
  label   = {tab:industry_academia_agent_security},
]{
  width   = \linewidth,
  colspec = {
    |Q[wd=0.07\linewidth,l,m]
    |Q[wd=0.12\linewidth,l,m]
    |X[1.00,l,m]|
  },
  cells   = {l,m},
  colsep  = 3pt,
  row{1}  = {font=\bfseries},
  rowhead = 1,
  rowsep  = 5pt,
  %measure = vbox,
  stretch = -1,
}
\hline
研究方向 & 企业/机构 + 核心学者 & 自研项目、产学合作、商用安全产品 \\
\hline

\SetCell[r=2]{l,m} 开源基础大模型原生智能体安全底座
& Meta FAIR；Yann LeCun, Joelle Pineau
& \begin{itemize}[leftmargin=*,topsep=0pt,partopsep=0pt,itemsep=1pt,parsep=0pt]
\item \textbf{世界模型安全底座：}以JEPA、Image-JEPA、Video-JEPA、V-JEPA 2构成学术路线演进，支撑多模态智能体的状态表征、感知预测和行动后果建模~\cite{nyuLeCun2026,vjepa2025}；
\item \textbf{开源安全护栏体系：}公开Llama Guard、Prompt Guard和LlamaFirewall，其中LlamaFirewall定位为面向安全AI Agent的开源护栏系统~\cite{metaLlamaFirewall2025}；
\item \textbf{团队组织演进：}Joelle Pineau卸任FAIR负责人、LeCun离开Meta后，相关研发由公司新组建的Meta Superintelligence Labs承接~\cite{pineauDeparture2025,lecunDeparture2025}。
\end{itemize} \\
\cline{2-3}
& 纽约大学、AMI Labs；Yann LeCun
& \begin{itemize}[leftmargin=*,topsep=0pt,partopsep=0pt,itemsep=1pt,parsep=0pt]
\item \textbf{JEPA架构产业化：}LeCun离职后创办AMI Labs推动JEPA架构商业化，首批目标行业为医疗、机器人和工业自动化~\cite{amiLabsFunding2026,lecunMITTechReview2026}；
\item \textbf{理论与基准验证：}2026年发表的JEPA相关论文证明该架构在特定条件下能够恢复真实世界底层结构，同时用配套基准指出当前模型面对细微视觉变化时仍较脆弱~\cite{lejepaIdentifiability2026}。
\end{itemize} \\
\hline

\SetCell[r=3]{l,m} 云端原生企业级大模型智能体研发与风控
& Google DeepMind；Demis Hassabis
& \begin{itemize}[leftmargin=*,topsep=0pt,partopsep=0pt,itemsep=1pt,parsep=0pt]
\item \textbf{多智能体协同基础理论：}AlphaStar、RoboCat等研究成果支撑多智能体协同博弈和机器人任务规划的基础理论；
\item \textbf{企业级智能体产品落地：}相关研究能力被工程化为Gemini Enterprise Agent Platform和Managed Agents，落实到API注册、IAM权限管控和运行日志审计~\cite{googleGeminiEnterpriseAgentPlatform2026,googleADK2025}；
\item \textbf{科研智能体协作平台：}AI Co-Scientist基于Gemini构建，由一组专业化Agent协同完成科学假设的生成、辩论、排序和迭代~\cite{deepmindCoScientist2025,gottweisAICoScientist2025}。
\end{itemize} \\
\cline{2-3}
& Google Research；Jeff Dean
& \begin{itemize}[leftmargin=*,topsep=0pt,partopsep=0pt,itemsep=1pt,parsep=0pt]
\item \textbf{云端Agent风控机制：}分布式大规模机器学习与系统工程能力支撑Gemini Enterprise Agent Platform的API注册、IAM权限管控和运行日志审计~\cite{googleADK2025,googleGeminiEnterpriseAgentPlatform2026}；
\item \textbf{托管式智能体治理：}Managed Agents将全链路Tracing审计、IAM权限管控等系统工程理论嵌入托管式智能体服务，对应企业Agent运行时风控需求~\cite{googleADK2025}。
\end{itemize} \\
\cline{2-3}
& 斯坦福大学医学院、MIT；Gary Peltz, Ritu Raman
& \begin{itemize}[leftmargin=*,topsep=0pt,partopsep=0pt,itemsep=1pt,parsep=0pt]
\item \textbf{产学研安全验证：}AI Co-Scientist于2025年2月首发时联合全球100多个科研机构共同开发和测试，Gary Peltz、Ritu Raman公开验证其科研任务能力~\cite{deepmindCoScientist2025}；
\item \textbf{科研落地与企业试用：}Co-Scientist已产出抗菌素耐药性、植物免疫等湿实验验证成果，并被第一三共、拜耳作物科学等企业内部试用~\cite{labcriticsCoScientist2026}。
\end{itemize} \\
\hline

\SetCell[r=2]{l,m} 智能体攻防评测基准与自动化红队体系
& 苏黎世联邦理工SPY Lab；Florian Tramèr
& \begin{itemize}[leftmargin=*,topsep=0pt,partopsep=0pt,itemsep=1pt,parsep=0pt]
\item \textbf{自动化攻防基准：}AgentDojo将间接提示注入、工具劫持等标准化攻防任务集统一到动态评测环境中，已由ETH SPY Lab开源部署、~\cite{agentdojoGithub2026}；
\item \textbf{基准复现实验：}AgentDojo使邮件、日历、文件和工具调用场景中的提示注入攻击具备可复现、可比较的实验对象。
\end{itemize} \\
\cline{2-3}
& CMU, Gray Swan AI；Zico Kolter, Matt Fredrikson, Andy Zou, Dawn Song
& \begin{itemize}[leftmargin=*,topsep=0pt,partopsep=0pt,itemsep=1pt,parsep=0pt]
\item \textbf{工业级红队平台：}Gray Swan AI将模型鲁棒性、神经网络形式化验证理论转化为检测能力，代表性产品包括红队测试平台Cygnal和对抗性红队工具Shade~\cite{grayswanSeriesA2026}；
\item \textbf{提示注入鲁棒性评测：}Shade已被Anthropic用于评估模型在代码环境中对提示注入攻击的鲁棒性，Gray Swan的间接提示注入研究也被Anthropic Claude Mythos的system card引用为权威研究依据~\cite{latentspaceGrayswan2026}；
\item \textbf{模型风险治理嵌入：}Kolter已出任OpenAI董事会成员并兼任安全与安保委员会主席，使自动化红队研究与模型风险治理进入更高层级的产业治理流程~\cite{kolterWikipedia2026}。
\end{itemize} \\
\hline
\end{longtblr}
\endgroup

% \begingroup

% \begin{center}
% % \nopagebreak[4]
% \enlargethispage{1cm}
% \footnotesize
% \setlength{\tabcolsep}{3pt}
% \begin{longtblr}[
% caption = {开源基础底座、云端企业智能体与自动化攻防评测体系},
% label   = {tab:industry_academia_agent_security},
% ]
% {
% colspec = {|m{0.07\textwidth}
% |m{0.12\textwidth}
% |m{\dimexpr\linewidth - 0.07\textwidth - 0.12\textwidth - 6\tabcolsep - 3\arrayrulewidth\relax}|},
% rowhead = 1,
% }
% \hline
% \textbf{研究方向} & \textbf{企业/机构 + 核心学者} & \textbf{自研项目、产学合作、商用安全产品} \\
% \hline

% \SetCell[r=6]{c} 开源基础大模型原生智能体安全底座
% & Meta FAIR；Yann LeCun, Joelle Pineau
% & \begin{itemize}[tabitem]
% \item \textbf{世界模型安全底座：}以JEPA、Image-JEPA、Video-JEPA、V-JEPA 2构成学术路线演进，支撑多模态智能体的状态表征、感知预测和行动后果建模~\cite{nyuLeCun2026,vjepa2025}；
% \item \textbf{开源安全护栏体系：}公开Llama Guard、Prompt Guard和LlamaFirewall，其中LlamaFirewall定位为面向安全AI Agent的开源护栏系统~\cite{metaLlamaFirewall2025}；
% \item \textbf{团队组织演进：}Joelle Pineau卸任FAIR负责人、LeCun离开Meta后，相关研发由公司新组建的Meta Superintelligence Labs承接~\cite{pineauDeparture2025,lecunDeparture2025}。
% \end{itemize} \\
% \cline{2-3}
% & 纽约大学、AMI Labs；Yann LeCun
% & \begin{itemize}[tabitem]
% \item \textbf{JEPA架构产业化：}LeCun离职后创办AMI Labs推动JEPA架构商业化，首批目标行业为医疗、机器人和工业自动化~\cite{amiLabsFunding2026,lecunMITTechReview2026}；
% \item \textbf{理论与基准验证：}2026年发表的JEPA相关论文证明该架构在特定条件下能够恢复真实世界底层结构，同时用配套基准指出当前模型面对细微视觉变化时仍较脆弱~\cite{lejepaIdentifiability2026}。
% \end{itemize} \\
% \hline

% \SetCell[r=9]{c} 云端原生企业级大模型智能体研发与风控
% & Google DeepMind；Demis Hassabis
% & \begin{itemize}[tabitem]
% \item \textbf{多智能体协同基础理论：}AlphaStar、RoboCat等研究成果支撑多智能体协同博弈和机器人任务规划的基础理论；
% \item \textbf{企业级智能体产品落地：}相关研究能力被工程化为Gemini Enterprise Agent Platform和Managed Agents，落实到API注册、IAM权限管控和运行日志审计~\cite{googleGeminiEnterpriseAgentPlatform2026,googleADK2025}；
% \item \textbf{科研智能体协作平台：}AI Co-Scientist基于Gemini构建，由一组专业化Agent协同完成科学假设的生成、辩论、排序和迭代~\cite{deepmindCoScientist2025,gottweisAICoScientist2025}。
% \end{itemize} \\
% \cline{2-3}
% & Google Research；Jeff Dean
% & \begin{itemize}[tabitem]
% \item \textbf{云端Agent风控机制：}分布式大规模机器学习与系统工程能力支撑Gemini Enterprise Agent Platform的API注册、IAM权限管控和运行日志审计~\cite{googleADK2025,googleGeminiEnterpriseAgentPlatform2026}；
% \item \textbf{托管式智能体治理：}Managed Agents将全链路Tracing审计、IAM权限管控等系统工程理论嵌入托管式智能体服务，对应企业Agent运行时风控需求~\cite{googleADK2025}。
% \end{itemize} \\
% \cline{2-3}
% & 斯坦福大学医学院、MIT；Gary Peltz, Ritu Raman
% & \begin{itemize}[tabitem]
% \item \textbf{产学研安全验证：}AI Co-Scientist于2025年2月首发时联合全球100多个科研机构共同开发和测试，Gary Peltz、Ritu Raman公开验证其科研任务能力~\cite{deepmindCoScientist2025}；
% \item \textbf{科研落地与企业试用：}Co-Scientist已产出抗菌素耐药性、植物免疫等湿实验验证成果，并被第一三共、拜耳作物科学等企业内部试用~\cite{labcriticsCoScientist2026}。
% \end{itemize} \\
% \hline

% \SetCell[r=5]{c} 智能体攻防评测基准与自动化红队体系
% & 苏黎世联邦理工SPY Lab；Florian Tramèr
% & \begin{itemize}[tabitem]
% \item \textbf{自动化攻防基准：}AgentDojo将间接提示注入、工具劫持等标准化攻防任务集统一到动态评测环境中，已由ETH SPY Lab开源部署、~\cite{agentdojoGithub2026}；
% \item \textbf{基准复现实验：}AgentDojo使邮件、日历、文件和工具调用场景中的提示注入攻击具备可复现、可比较的实验对象。
% \end{itemize} \\
% \hline
% & CMU, Gray Swan AI；Zico Kolter, Matt Fredrikson, Andy Zou, Dawn Song
% & \begin{itemize}[tabitem]
% \item \textbf{工业级红队平台：}Gray Swan AI将模型鲁棒性、神经网络形式化验证理论转化为检测能力，代表性产品包括红队测试平台Cygnal和对抗性红队工具Shade~\cite{grayswanSeriesA2026}；
% \item \textbf{提示注入鲁棒性评测：}Shade已被Anthropic用于评估模型在代码环境中对提示注入攻击的鲁棒性，Gray Swan的间接提示注入研究也被Anthropic Claude Mythos的system card引用为权威研究依据~\cite{latentspaceGrayswan2026}；
% \item \textbf{模型风险治理嵌入：}Kolter已出任OpenAI董事会成员并兼任安全与安保委员会主席，使自动化红队研究与模型风险治理进入更高层级的产业治理流程~\cite{kolterWikipedia2026}。
% \end{itemize} \\
% \hline
% \end{longtblr}
% \end{center}
% \endgroup


\subsection{通用商用智能体安全评估与治理规范}
随着云端商用智能体与工具生态规模化落地，安全治理从模型鲁棒性延伸至协议交互、运行时管控、隐私计算、外部审计的全生命周期。海外形成“企业研发‑学界定标‑第三方核验”的成熟产学体系，覆盖通用智能体评估、MCP工具协议、端侧硬件隐私三大方向\cite{openaiRedTeamNetwork2025,mcpFormalFramework2026,nvidiaConfidentialComputingProduct2026}。本小节结合头部厂商、国家级安全机构与顶尖高校成果，归纳商用智能体的治理范式、协议漏洞矛盾与硬件隐私安全落地规律，主要发现如下（详见表~\ref{tab:commercial_mcp_edge_agent_security}）：
\begin{itemize}[tabitem]
    \item \textbf{头部厂商建立标准化商用治理范式。}OpenAI搭建外部专家红队测试网络、在Agents‑SDK内置多层安全护栏，配套全链路Tracing追踪与人在回路审批机制，并发布System‑Card公开安全评估报告，构建起完整闭环体系，成为全球商用智能体安全治理标杆范式\cite{openaiChatGPTAgentSystemCard2025,openaiAgentsSDK2025,openaiGuardrailsApprovalsDoc2026}；英国UK‑AISI、FAR.AI开展第三方独立安全审计，斯坦福CRFM提出HELM评测理论，从第三方与学术层面完善评估方法论，形成产学核验的治理格局\cite{venturebeatChatGPTAgent2025}。

    \item \textbf{工具协议存在结构性权责矛盾。}Anthropic推出的MCP协议生态规模快速扩张，但OX‑Security发现协议原生存在RCE高危漏洞；出于生态灵活性考量，厂商拒绝在架构层面完成彻底修复，把输入校验、风险净化的安全责任转嫁至下游开发者。伯克利CHAI团队从目标对齐理论剖析该问题本质，学术界也意识到借助ProVerif、Tamarin对MCP开展形式化验证难度很高，充分暴露出协议生态扩张与安全约束之间的固有冲突\cite{oxSecurityMCP2026,berkeleyCHAI2026,mcpFormalFramework2026}。

    \item \textbf{智能体隐私安全构建垂直产学链条。}产业端Apple依靠PCC端云分层架构实现数据最小化处理与可验证云端推理，NVIDIA依托GPU硬件信任根搭建机密计算环境；学术层面EPFL‑SPRING‑Lab的隐私工程理论、NYU针对GPU的侧信道分析成果，为硬件可信基座、端侧隐私智能体提供理论支撑，形成理论研究‑硬件实现‑产品落地的完整链条\cite{springLab2026,karriNYU2026}。
\end{itemize}


% \begingroup

% \begin{center}
% % \nopagebreak[4]
% \enlargethispage{1cm}
% \footnotesize
% \setlength{\tabcolsep}{3pt}
% \begin{longtblr}[
% caption = {通用商用智能体、MCP工具调用协议与端侧硬件隐私计算智能体},
% label   = {tab:commercial_mcp_edge_agent_security},
% ]
% {
% colspec = {|m{0.07\textwidth}
% |m{0.12\textwidth}
% |m{\dimexpr\linewidth - 0.07\textwidth - 0.12\textwidth - 6\tabcolsep - 3\arrayrulewidth\relax}|},
% rowhead = 1,
% }
% \hline
% \textbf{研究方向} & \textbf{企业/机构 + 核心学者} & \textbf{自研项目、产学合作、商用安全产品} \\
% \hline

% \SetCell[r=6]{c} 通用商用智能体安全评估与治理规范
% & OpenAI：Sam Altman、Greg Brockman
% & \begin{itemize}[tabitem]
% \item \textbf{ChatGPT Agent交互智能体：}搭建外部多专家红队测试网络，Red Teaming Network由16位具备生物安全背景的博士研究员组成，共提交110次攻击尝试，其中16次超过内部风险阈值~\cite{openaiChatGPTAgentSystemCard2025,openaiRedTeamNetwork2025}。
% \item \textbf{开源安全护栏体系：}OpenAI Agents SDK内置三层输入/输出/工具护栏，支持检测风险时立即中止运行的tripwire机制，配合运行时Tracing审计与高风险动作的human-in-the-loop审批机制~\cite{openaiAgentsSDK2025,openaiGuardrailsDoc2026,openaiTracingDoc2026,openaiGuardrailsApprovalsDoc2026}。
% \item \textbf{全球统一风险披露体系：}System Card持续迭代智能体能力、风险与透明度评估指标，全产品线强制落地，并以公开白皮书系统阐述外部专家参与风险评估的方法论~\cite{openaiChatGPTAgentSystemCard2025,openaiExternalRedTeaming2025}。
% \end{itemize} \\
% \cline{2-3}
% & 英国AI安全研究所UK AISI、FAR.AI独立安全实验室
% & \begin{itemize}[tabitem]
% \item \textbf{深度访问链路漏洞审计：}UK AISI获得对ChatGPT Agent内部推理链和策略文本的深度访问权限，完成官方漏洞审计测试，测出7个通用漏洞~\cite{venturebeatChatGPTAgent2025}。
% \item \textbf{商用智能体安全缺陷批评：}FAR.AI经40小时独立测试发现3个部分漏洞，提出商用智能体过度依赖运行时监控的结构性安全缺陷~\cite{venturebeatChatGPTAgent2025}。
% \end{itemize} \\
% \cline{2-3}
% & 斯坦福CRFM：Percy Liang
% & \begin{itemize}[tabitem]
% \item \textbf{全链路可比较评测理论对齐：}HELM全链路可比较评测理论与OpenAI Agents SDK的Tracing架构在方法论上形成理论对应，共同强调把系统行为拆解为可追溯、可比较的具体环节~\cite{openaiTracingDoc2026}。
% \end{itemize} \\
% \hline

% \SetCell[r=7]{c} 智能体工具调用MCP协议与目标对齐安全
% & Anthropic：Dario Amodei、Jared Kaplan
% & \begin{itemize}[tabitem]
% \item \textbf{Agent全生态工具交互协议：}MCP工具通信协议基于JSON-RPC 2.0构建，支撑Claude Agent全生态工具交互，截至2026年初已形成超1万活跃服务器、17.7万注册工具、月均9700万次SDK下载的生态规模~\cite{mcpFormalFramework2026}。
% \item \textbf{高危分级智能体系统：}Claude Code代码智能体、Mythos高危分级智能体系统均采用分层行为约束护栏~\cite{anthropicClaudeCodeSecurity2026}。
% \item \textbf{协议架构安全责任取舍：}面对OX Security披露的架构级RCE漏洞，Anthropic确认该行为符合设计预期，以维持协议灵活性为由拒绝修复，将输入净化责任划给下游开发者~\cite{oxSecurityMCP2026}。
% \end{itemize} \\
% \cline{2-3}
% & 加州伯克利CHAI：Stuart Russell
% & \begin{itemize}[tabitem]
% \item \textbf{人类兼容AI理论与目标对齐：}Russell作为Berkeley CHAI创始主任，其人类兼容AI理论关注目标不确定性和系统可控性，与高自治Agent的目标漂移、价值对齐问题存在直接理论关联~\cite{berkeleyCHAI2026}。
% \end{itemize} \\
% \cline{2-3}
% & OX Security、Check Point安全厂商；密码学形式化验证学界
% & \begin{itemize}[tabitem]
% \item \textbf{MCP协议远程代码执行漏洞披露：}OX Security披露MCP协议"设计如此"的远程代码执行漏洞，影响超1.5亿次下载和20万台服务器，攻陷11个主流MCP市场中的9个，产生至少10个高危或严重级别CVE编号~\cite{oxSecurityMCP2026,thehackernewsMCP2026}。
% \item \textbf{Claude Code代码执行漏洞修复：}Check Point于2025年9至10月披露Claude Code三项漏洞（CVE-2025-59536，CVSS 8.7；CVE-2026-21852，CVSS 5.3），由厂商协同完整修复，分别在1.0.111版本和2.0.65版本发布安全补丁~\cite{checkpointClaudeCode2026,githubClaudeCodeAdvisory2026}。
% \item \textbf{MCP威胁分类与形式化验证：}学术界已提出MCP的威胁分类模型，并指出借鉴Tamarin、ProVerif等密码协议验证工具对MCP做形式化验证仍是开放且困难的研究方向~\cite{mcpFormalFramework2026,mcpSoK2026}。
% \end{itemize} \\
% \hline

% \SetCell[r=5]{c} 端侧与硬件机密计算隐私智能体安全技术
% & Apple：Craig Federighi 隐私安全实验室
% & \begin{itemize}[tabitem]
% \item \textbf{端云分层隐私架构：}Apple Intelligence移动端本地智能体采用端云分层数据最小化架构~\cite{applePCCSecurityGuide2026}。
% \item \textbf{可验证云端推理平台：}Private Cloud Compute（PCC）通过Virtual Research Environment向全球研究者开放本地复现能力，配合开源形式化验证工具、函数库与安全悬赏计划，构建透明可信推理体系~\cite{appleSecurityResearch2026}。
% \end{itemize} \\
% \cline{2-3}
% & EPFL SPRING Lab：Carmela Troncoso
% & \begin{itemize}[tabitem]
% \item \textbf{隐私工程理论支撑端侧架构：}Troncoso团队长期研究隐私工程、数据最小化和匿名化技术，相关理论与端侧/云端隐私智能体架构问题高度相关~\cite{springLab2026}。
% \end{itemize} \\
% \cline{2-3}
% & NVIDIA：Jensen Huang，NYU；Ramesh Karri
% & \begin{itemize}[tabitem]
% \item \textbf{GPU硬件机密计算平台：}Confidential Agent GPU硬件机密计算平台自H100起首次支持机密计算，通过片上硬件信任根实现安全启动和远程认证，在Blackwell架构上进一步延伸~\cite{nvidiaConfidentialComputingProduct2026}。
% \item \textbf{硬件侧信道安全研究支撑：}Karri教授长期研究硬件木马、芯片供应链安全和侧信道分析，与GPU侧信道信息泄露防护议题直接相关，学术界也已对TEE模式下大模型推理性能开销做独立评测~\cite{karriNYU2026}。
% \end{itemize} \\
% \hline

% \end{longtblr}
% \end{center}
% \endgroup

% 依赖项（必须放在主文档导言区，即 \begin{document} 之前）：
% \usepackage{tabularray}
% \UseTblrLibrary{varwidth}
% \usepackage{enumitem}
\begingroup
% \nopagebreak[4]
\enlargethispage{1cm}
\footnotesize
\centering

\DefTblrTemplate{conthead-text}{default}{（续表）}
\DefTblrTemplate{contfoot-text}{default}{续下页}

\begin{longtblr}[
  caption = {通用商用智能体、MCP工具调用协议与端侧硬件隐私计算智能体},
  label   = {tab:commercial_mcp_edge_agent_security},
]{
  width   = \linewidth,
  colspec = {
    |Q[wd=0.07\linewidth,l,m]
    |Q[wd=0.12\linewidth,l,m]
    |X[1.00,l,m]|
  },
  cells   = {l,m},
  colsep  = 3pt,
  row{1}  = {font=\bfseries},
  rowhead = 1,
  rowsep  = 5pt,
  measure = vbox,
  stretch = -1,
}
\hline
研究方向 & 企业/机构 + 核心学者 & 自研项目、产学合作、商用安全产品 \\
\hline

\SetCell[r=3]{l,m} 通用商用智能体安全评估与治理规范
& OpenAI：Sam Altman、Greg Brockman
& \begin{itemize}[leftmargin=*,topsep=0pt,partopsep=0pt,itemsep=1pt,parsep=0pt]
\item \textbf{ChatGPT Agent交互智能体：}搭建外部多专家红队测试网络，Red Teaming Network由16位具备生物安全背景的博士研究员组成，共提交110次攻击尝试，其中16次超过内部风险阈值~\cite{openaiChatGPTAgentSystemCard2025,openaiRedTeamNetwork2025}。
\item \textbf{开源安全护栏体系：}OpenAI Agents SDK内置三层输入/输出/工具护栏，支持检测风险时立即中止运行的tripwire机制，配合运行时Tracing审计与高风险动作的human-in-the-loop审批机制~\cite{openaiAgentsSDK2025,openaiGuardrailsDoc2026,openaiTracingDoc2026,openaiGuardrailsApprovalsDoc2026}。
\item \textbf{全球统一风险披露体系：}System Card持续迭代智能体能力、风险与透明度评估指标，全产品线强制落地，并以公开白皮书系统阐述外部专家参与风险评估的方法论~\cite{openaiChatGPTAgentSystemCard2025,openaiExternalRedTeaming2025}。
\end{itemize} \\
\cline{2-3}
& 英国AI安全研究所UK AISI、FAR.AI独立安全实验室
& \begin{itemize}[leftmargin=*,topsep=0pt,partopsep=0pt,itemsep=1pt,parsep=0pt]
\item \textbf{深度访问链路漏洞审计：}UK AISI获得对ChatGPT Agent内部推理链和策略文本的深度访问权限，完成官方漏洞审计测试，测出7个通用漏洞~\cite{venturebeatChatGPTAgent2025}。
\item \textbf{商用智能体安全缺陷批评：}FAR.AI经40小时独立测试发现3个部分漏洞，提出商用智能体过度依赖运行时监控的结构性安全缺陷~\cite{venturebeatChatGPTAgent2025}。
\end{itemize} \\
\cline{2-3}
& 斯坦福CRFM：Percy Liang
& \begin{itemize}[leftmargin=*,topsep=0pt,partopsep=0pt,itemsep=1pt,parsep=0pt]
\item \textbf{全链路可比较评测理论对齐：}HELM全链路可比较评测理论与OpenAI Agents SDK的Tracing架构在方法论上形成理论对应，共同强调把系统行为拆解为可追溯、可比较的具体环节~\cite{openaiTracingDoc2026}。
\end{itemize} \\
\hline

\SetCell[r=3]{l,m} 智能体工具调用MCP协议与目标对齐安全
& Anthropic：Dario Amodei、Jared Kaplan
& \begin{itemize}[leftmargin=*,topsep=0pt,partopsep=0pt,itemsep=1pt,parsep=0pt]
\item \textbf{Agent全生态工具交互协议：}MCP工具通信协议基于JSON-RPC 2.0构建，支撑Claude Agent全生态工具交互，截至2026年初已形成超1万活跃服务器、17.7万注册工具、月均9700万次SDK下载的生态规模~\cite{mcpFormalFramework2026}。
\item \textbf{高危分级智能体系统：}Claude Code代码智能体、Mythos高危分级智能体系统均采用分层行为约束护栏~\cite{anthropicClaudeCodeSecurity2026}。
\item \textbf{协议架构安全责任取舍：}面对OX Security披露的架构级RCE漏洞，Anthropic确认该行为符合设计预期，以维持协议灵活性为由拒绝修复，将输入净化责任划给下游开发者~\cite{oxSecurityMCP2026}。
\end{itemize} \\
\cline{2-3}
& 加州伯克利CHAI：Stuart Russell
& \begin{itemize}[leftmargin=*,topsep=0pt,partopsep=0pt,itemsep=1pt,parsep=0pt]
\item \textbf{人类兼容AI理论与目标对齐：}Russell作为Berkeley CHAI创始主任，其人类兼容AI理论关注目标不确定性和系统可控性，与高自治Agent的目标漂移、价值对齐问题存在直接理论关联~\cite{berkeleyCHAI2026}。
\end{itemize} \\
\cline{2-3}
& OX Security、Check Point安全厂商；密码学形式化验证学界
& \begin{itemize}[leftmargin=*,topsep=0pt,partopsep=0pt,itemsep=1pt,parsep=0pt]
\item \textbf{MCP协议远程代码执行漏洞披露：}OX Security披露MCP协议"设计如此"的远程代码执行漏洞，影响超1.5亿次下载和20万台服务器，攻陷11个主流MCP市场中的9个，产生至少10个高危或严重级别CVE编号~\cite{oxSecurityMCP2026,thehackernewsMCP2026}。
\item \textbf{Claude Code代码执行漏洞修复：}Check Point于2025年9至10月披露Claude Code三项漏洞（CVE-2025-59536，CVSS 8.7；CVE-2026-21852，CVSS 5.3），由厂商协同完整修复，分别在1.0.111版本和2.0.65版本发布安全补丁~\cite{checkpointClaudeCode2026,githubClaudeCodeAdvisory2026}。
\item \textbf{MCP威胁分类与形式化验证：}学术界已提出MCP的威胁分类模型，并指出借鉴Tamarin、ProVerif等密码协议验证工具对MCP做形式化验证仍是开放且困难的研究方向~\cite{mcpFormalFramework2026,mcpSoK2026}。
\end{itemize} \\
\hline

\SetCell[r=3]{l,m} 端侧与硬件机密计算隐私智能体安全技术
& Apple：Craig Federighi 隐私安全实验室
& \begin{itemize}[leftmargin=*,topsep=0pt,partopsep=0pt,itemsep=1pt,parsep=0pt]
\item \textbf{端云分层隐私架构：}Apple Intelligence移动端本地智能体采用端云分层数据最小化架构~\cite{applePCCSecurityGuide2026}。
\item \textbf{可验证云端推理平台：}Private Cloud Compute（PCC）通过Virtual Research Environment向全球研究者开放本地复现能力，配合开源形式化验证工具、函数库与安全悬赏计划，构建透明可信推理体系~\cite{appleSecurityResearch2026}。
\end{itemize} \\
\cline{2-3}
& EPFL SPRING Lab：Carmela Troncoso
& \begin{itemize}[leftmargin=*,topsep=0pt,partopsep=0pt,itemsep=1pt,parsep=0pt]
\item \textbf{隐私工程理论支撑端侧架构：}Troncoso团队长期研究隐私工程、数据最小化和匿名化技术，相关理论与端侧/云端隐私智能体架构问题高度相关~\cite{springLab2026}。
\end{itemize} \\
\cline{2-3}
& NVIDIA：Jensen Huang，NYU；Ramesh Karri
& \begin{itemize}[leftmargin=*,topsep=0pt,partopsep=0pt,itemsep=1pt,parsep=0pt]
\item \textbf{GPU硬件机密计算平台：}Confidential Agent GPU硬件机密计算平台自H100起首次支持机密计算，通过片上硬件信任根实现安全启动和远程认证，在Blackwell架构上进一步延伸~\cite{nvidiaConfidentialComputingProduct2026}。
\item \textbf{硬件侧信道安全研究支撑：}Karri教授长期研究硬件木马、芯片供应链安全和侧信道分析，与GPU侧信道信息泄露防护议题直接相关，学术界也已对TEE模式下大模型推理性能开销做独立评测~\cite{karriNYU2026}。
\end{itemize} \\
\hline
\end{longtblr}
\endgroup



\subsection{国产中文智能体评测基准与本土产学研安全体系}
中文语义特征、合规要求与国产化算力生态，决定我国智能体安全无法直接复用海外技术体系。我国聚焦中文场景评测、安全对齐训练、垂直行业多智能体风控\cite{pandaguardCASIA2026}。本小节基于我国核心机构成果，系统归纳本土安全体系的差异化特征、技术能力与产业落地路径，主要发现如下（详见表~\ref{tab:chinese_agent_benchmark_security}）：
\begin{itemize}[tabitem]
    \item \textbf{我国高校建成差异化中文评测体系。}清华大学依托ToolBench、AgentBench搭建工具调用和综合任务评测基准，适配海量真实API与中文环境；上海交通大学提出R‑Judge评测体系，聚焦对话交互场景下智能体行为风险判定，二者形成互补的本土化基准，填补海外评测集对中文语境适配不足的短板。

    \item \textbf{国产安全对齐形成开源可复现技术栈。}北京大学PKU‑Alignment团队提出Safe RLHF与PKU‑SafeRLHF算法，配套BeaverTails高质量中文偏好数据集和align‑anything全模态开源框架，为我国大模型智能体后训练安全对齐提供可落地的基础组件\cite{pkuAlignAnything2026}；浙江大学依托AIcert平台完成大量国产模型安全评估，进一步完善对齐效果的检验手段，共同构筑国产智能体安全对齐训练的核心开源底座\cite{aicertZJU2026}。

    \item \textbf{国产安全具备强行业导向特征：}中科院自动化所打造PandaGuard攻防评测平台与JidiAI博弈环境，从科研层面开展多智能体风险推演\cite{pandaguardCASIA2026,jidiAI2026}；上海AI实验室推出Intern‑Shannon安全智能体操作系统和SafeWork‑R1模型，面向政企场景完成产业级落地\cite{intershannon2026,safework2025}；华为联合上交IPADS团队搭建SMARTS自动驾驶仿真平台并开展芯片可信环境研究，面向交通行业完成落地验证\cite{noahArkSMARTS2026,ipadsSJTU2026}。
    % 整体摒弃海外通用智能体研发路线，形成科研验证、平台研发、行业场景落地的分层发展模式。
\end{itemize}



\begingroup

% \nopagebreak[4]
\enlargethispage{1cm}
\footnotesize
\setlength{\tabcolsep}{3pt}
\begin{longtblr}[
caption = {国产中文智能体评测基准与本土产学研安全体系},
label   = {tab:chinese_agent_benchmark_security},
]
{
colspec = {|m{0.07\textwidth}
|m{0.14\textwidth}
|m{\dimexpr\linewidth - 0.07\textwidth - 0.14\textwidth - 6\tabcolsep - 3\arrayrulewidth\relax}|},
rowhead = 1,
}
\hline
\textbf{研究方向} & \textbf{企业/机构 + 核心学者} & \textbf{自研项目、产学合作、商用安全产品} \\
\hline

\SetCell[r=3]{c} 中文本土智能体评测与国产化标准体系
& 清华大学THUNLP/AIR；朱军、唐杰、刘知远、孙茂松
& \begin{itemize}[tabitem]
\item \textbf{工具调用评测基准：}ToolBench工具调用评测基准，覆盖16000多个真实API中文工具场景~\cite{tsinghuaAIInstitute2026,tsinghuaAIR2026}。
\item \textbf{全栈智能体综合评测：}AgentBench八场景全栈智能体综合评测基准，联合俄亥俄州立大学、加州大学伯克利分校共建。
\end{itemize} \\
\cline{2-3}
& \mbox{北京大学}PKU-Alignment；杨耀东、刘健
& \begin{itemize}[tabitem]
\item \textbf{安全强化学习训练框架：}Safe RLHF、PKU-SafeRLHF训练框架，BeaverTails中文安全偏好数据集；
\item \textbf{全模态对齐训练框架：}Aligner、ProgressGym等后训练对齐系列成果，开源align-anything全模态对齐训练、数据和测评框架~\cite{pkuAlignAnything2026}。
\end{itemize} \\
\cline{2-3}
& 上海交通大学；谷大武（LoCCS）、张倬胜（AI安全实验室）
& \begin{itemize}[tabitem]
\item \textbf{密码学与系统安全底座：}谷大武团队围绕密码理论、软硬件与系统安全展开研究，为智能体运行环境提供底层安全基础~\cite{sjtuLoCCS2026}。
\item \textbf{智能体行为安全评测基准：}张倬胜团队提出R-Judge智能体行为安全评测基准，覆盖5类应用场景、27个风险场景与10类风险类型~\cite{sjtuAISecLab2026}。
\end{itemize} \\
\hline

\SetCell[r=3]{c} 我国顶尖高校开源智能体对齐安全体系
& 浙江大学区块链与数据安全全国重点实验室；任奎
& \begin{itemize}[tabitem]
\item \textbf{通用AI安全评测平台：}AIcertAI安全评测平台，覆盖数据质量、算法安全、模型安全、框架安全与系统安全检测，已评测LLaMA、Gemma、Qwen等35个开源大模型，并在淘宝直播风控、中车株洲所等场景形成示范应用~\cite{aicertZJU2026}。
\item \textbf{企业级安全联合实验室：}浙江大学--阿里巴巴集团AI安全联合实验室将智能体安全评测列为首期核心研究方向，支撑企业级大模型Agent红队测试与安全治理~\cite{zjuAlibabaAISafetyLab2025}。
\end{itemize} \\
\cline{2-3}
& 北京大学PKU-Alignment；杨耀东、刘健
& \begin{itemize}[tabitem]
\item \textbf{后训练安全对齐生态：}维护Safe RLHF、PKU-SafeRLHF、BeaverTails等安全对齐训练与数据体系，构成我国较早的可复现安全强化学习框架~\cite{pkuAlignAnything2026}。
\end{itemize} \\
\cline{2-3}
& 上海交通大学AI安全实验室；张倬胜
& \begin{itemize}[tabitem]
\item \textbf{自然语言处理与自主智能体安全：}专注自然语言处理、预训练语言模型与自主智能体安全，R-Judge基准支撑智能体行为风险识别研究~\cite{sjtuAISecLab2026}。
\end{itemize} \\
\hline

\SetCell[r=3]{c} 我国科研院所、产业实验室行业智能体安全落地
& 中国科学院自动化所；曾毅、王飞跃团队
& \begin{itemize}[tabitem]
\item \textbf{多智能体博弈攻防评测平台：}PandaGuard（灵御）多智能体博弈攻防评测平台，集成19种攻击算法与12种防御机制~\cite{pandaguardCASIA2026}。
\item \textbf{多智能体决策开放平台：}JidiAI多智能体开源开放平台面向智能体博弈决策领域，吸引我国外50余所高校机构注册使用~\cite{jidiAI2026}。
\end{itemize} \\
\cline{2-3}
& 上海AI实验室；安全可信AI团队
& \begin{itemize}[tabitem]
\item \textbf{产业级安全智能体操作系统：}Intern-Shannon（书安）产业级安全智能体操作系统，以“安全即服务”模式支撑高安全场景部署~\cite{intershannon2026}。
\item \textbf{安全能力协同演化与科研智能体：}SafeWork-R1通过SafeLadder框架实现安全与智能协同进化，依托InternLM/OpenMMLab开源底座；InternAgent构建面向自主科学发现的闭环多智能体框架~\cite{safework2025,internAgentSHLab2026}。
\end{itemize} \\
\hline
& 华为诺亚方舟实验室；郝建业（决策推理实验室），上交陈海波（IPADS实验室）
& \begin{itemize}[tabitem]
\item \textbf{自动驾驶多智能体仿真平台：}SMARTS大规模多智能体强化学习仿真平台，支持数百至数千个智能体在真实交通场景中的交互学习~\cite{noahArkSMARTS2026}。
\item \textbf{可信执行环境与操作系统级支撑：}sNPU可信执行环境研究被ISCA 2024接收，陈海波担任OpenHarmony技术指导委员会创始主席，IPADS与华为在鸿蒙操作系统上保持长期合作关系~\cite{ipadsSJTU2026}。
\end{itemize} \\
\hline
\end{longtblr}
\endgroup
% 依赖项（必须放在主文档导言区，即 \begin{document} 之前）：
% \usepackage{tabularray}
% \UseTblrLibrary{varwidth}
% \usepackage{enumitem}
% \begingroup
% % \nopagebreak[4]
% \enlargethispage{1cm}
% \footnotesize
% \centering

% \DefTblrTemplate{conthead-text}{default}{（续表）}
% \DefTblrTemplate{contfoot-text}{default}{续下页}

% \begin{longtblr}[
%   caption = {国产中文智能体评测基准与本土产学研安全体系},
%   label   = {tab:chinese_agent_benchmark_security},
% ]{
%   width   = \linewidth,
%   colspec = {
%     |Q[wd=0.07\linewidth,l,m]
%     |Q[wd=0.14\linewidth,l,m]
%     |X[1.00,l,m]|
%   },
%   cells   = {l,m},
%   colsep  = 3pt,
%   row{1}  = {font=\bfseries},
%   rowhead = 1,
%   rowsep  = 5pt,
%   measure = vbox,
%   vspan = even,
%   stretch = -1,

  
% }
% \hline
% 研究方向 & 企业/机构 + 核心学者 & 自研项目、产学合作、商用安全产品 \\
% \hline

% \SetCell[r=3]{l,m} 中文本土智能体评测与国产化标准体系
% & 清华大学THUNLP/AIR；朱军、唐杰、刘知远、孙茂松
% & \begin{itemize}[leftmargin=*,topsep=0pt,partopsep=0pt,itemsep=1pt,parsep=0pt]
% \item \textbf{工具调用评测基准：}ToolBench工具调用评测基准，覆盖16000多个真实API中文工具场景~\cite{tsinghuaAIInstitute2026,tsinghuaAIR2026}。
% \item \textbf{全栈智能体综合评测：}AgentBench八场景全栈智能体综合评测基准，联合俄亥俄州立大学、加州大学伯克利分校共建。
% \end{itemize} \\
% \cline{2-3}
% & \mbox{北京大学}PKU-Alignment；杨耀东、刘健
% & \begin{itemize}[leftmargin=*,topsep=0pt,partopsep=0pt,itemsep=1pt,parsep=0pt]
% \item \textbf{安全强化学习训练框架：}Safe RLHF、PKU-SafeRLHF训练框架，BeaverTails中文安全偏好数据集；
% \item \textbf{全模态对齐训练框架：}Aligner、ProgressGym等后训练对齐系列成果，开源align-anything全模态对齐训练、数据和测评框架~\cite{pkuAlignAnything2026}。
% \end{itemize} \\
% \cline{2-3}
% & 上海交通大学；谷大武（LoCCS）、张倬胜（AI安全实验室）
% & \begin{itemize}[leftmargin=*,topsep=0pt,partopsep=0pt,itemsep=1pt,parsep=0pt]
% \item \textbf{密码学与系统安全底座：}谷大武团队围绕密码理论、软硬件与系统安全展开研究，为智能体运行环境提供底层安全基础~\cite{sjtuLoCCS2026}。
% \item \textbf{智能体行为安全评测基准：}张倬胜团队提出R-Judge智能体行为安全评测基准，覆盖5类应用场景、27个风险场景与10类风险类型~\cite{sjtuAISecLab2026}。
% \end{itemize} \\
% \hline

% \SetCell[r=3]{l,m} 我国顶尖高校开源智能体对齐安全体系
% & 浙江大学区块链与数据安全全国重点实验室；任奎
% & \begin{itemize}[leftmargin=*,topsep=0pt,partopsep=0pt,itemsep=1pt,parsep=0pt]
% \item \textbf{通用AI安全评测平台：}AIcertAI安全评测平台，覆盖数据质量、算法安全、模型安全、框架安全与系统安全检测，已评测LLaMA、Gemma、Qwen等35个开源大模型，并在淘宝直播风控、中车株洲所等场景形成示范应用~\cite{aicertZJU2026}。
% \item \textbf{企业级安全联合实验室：}浙江大学--阿里巴巴集团AI安全联合实验室将智能体安全评测列为首期核心研究方向，支撑企业级大模型Agent红队测试与安全治理~\cite{zjuAlibabaAISafetyLab2025}。
% \end{itemize} \\
% \cline{2-3}
% & 北京大学PKU-Alignment；杨耀东、刘健
% & \begin{itemize}[leftmargin=*,topsep=0pt,partopsep=0pt,itemsep=1pt,parsep=0pt]
% \item \textbf{后训练安全对齐生态：}维护Safe RLHF、PKU-SafeRLHF、BeaverTails等安全对齐训练与数据体系，构成我国较早的可复现安全强化学习框架~\cite{pkuAlignAnything2026}。
% \end{itemize} \\
% \cline{2-3}
% & 上海交通大学AI安全实验室；张倬胜
% & \begin{itemize}[leftmargin=*,topsep=0pt,partopsep=0pt,itemsep=1pt,parsep=0pt]
% \item \textbf{自然语言处理与自主智能体安全：}专注自然语言处理、预训练语言模型与自主智能体安全，R-Judge基准支撑智能体行为风险识别研究~\cite{sjtuAISecLab2026}。
% \end{itemize} \\
% \hline

% \SetCell[r=3]{l,m} 我国科研院所、产业实验室行业智能体安全落地
% & 中国科学院自动化所；曾毅、王飞跃团队
% & \begin{itemize}[leftmargin=*,topsep=0pt,partopsep=0pt,itemsep=1pt,parsep=0pt]
% \item \textbf{多智能体博弈攻防评测平台：}PandaGuard（灵御）多智能体博弈攻防评测平台，集成19种攻击算法与12种防御机制~\cite{pandaguardCASIA2026}。
% \item \textbf{多智能体决策开放平台：}JidiAI多智能体开源开放平台面向智能体博弈决策领域，吸引我国外50余所高校机构注册使用~\cite{jidiAI2026}。
% \end{itemize} \\
% \cline{2-3}
% & 上海AI实验室；安全可信AI团队
% & \begin{itemize}[leftmargin=*,topsep=0pt,partopsep=0pt,itemsep=1pt,parsep=0pt]
% \item \textbf{产业级安全智能体操作系统：}Intern-Shannon（书安）产业级安全智能体操作系统，以“安全即服务”模式支撑高安全场景部署~\cite{intershannon2026}。
% \item \textbf{安全能力协同演化与科研智能体：}SafeWork-R1通过SafeLadder框架实现安全与智能协同进化，依托InternLM/OpenMMLab开源底座；InternAgent构建面向自主科学发现的闭环多智能体框架~\cite{safework2025,internAgentSHLab2026}。
% \end{itemize} \\
% \cline{2-3}
% & 华为诺亚方舟实验室；郝建业（决策推理实验室），上交陈海波（IPADS实验室）
% & \begin{itemize}[leftmargin=*,topsep=0pt,partopsep=0pt,itemsep=1pt,parsep=0pt]
% \item \textbf{自动驾驶多智能体仿真平台：}SMARTS大规模多智能体强化学习仿真平台，支持数百至数千个智能体在真实交通场景中的交互学习~\cite{noahArkSMARTS2026}。
% \item \textbf{可信执行环境与操作系统级支撑：}sNPU可信执行环境研究被ISCA 2024接收，陈海波担任OpenHarmony技术指导委员会创始主席，IPADS与华为在鸿蒙操作系统上保持长期合作关系~\cite{ipadsSJTU2026}。
% \end{itemize} \\
% \hline
% \end{longtblr}
% \endgroup


% ========== 全球政策法规与标准治理现状（国家 / 区域制度==========
\section{全球AI智能体安全政策与法规}
当前全球AI治理已从通用大模型合规管控，细化聚焦至AI智能体专属安全治理层面。各国及地区基于自身技术优势、产业格局、数据主权诉求与地缘安全考量，形成了差异化的政策法规体系。
% 北美侧重技术安全基线与市场化合规问责，欧洲坚守风险分级与强监管域外管辖，我国立足全生命周期治理与产业合规协同。上述治理路径的分化，既反映了AI智能体安全风险的跨域共性，也体现了区域治理价值取向的底层差异。
%
本节立足全球治理格局，分区域梳理北美、欧洲与我国的AI智能体安全政策法规体系，系统归集各区域主管机构、核心文件、生效规则与管控要点。


% \begin{itemize}[leftmargin=*,topsep=0pt,partopsep=0pt]
%     \item 全球AI治理重心呈现\textbf{从通用大模型向专用智能体下沉}的明确趋势，政策条款针对性覆盖智能体工具调用、自主交互、多主体协同等特有风险，告别了普惠式AI监管阶段。
%     \item 全球形成三大差异化治理阵营，北美以行政令+州级立法构建分层问责体系，欧洲依托AI ACT实现风险分级与强执法管辖，中国构建多部门协同的全生命周期治理架构。
%     \item AI智能体政策普遍绑定\textbf{安全合规+产业发展}双重目标，所有区域均在风险管控基础上，配套算力、生态、标准等产业支撑条款，规避单一监管带来的产业抑制问题。
% \end{itemize}


\subsection{北美AI智能体安全政策与法规}
北美作为全球AI技术研发与产业落地的核心区域，其AI智能体安全治理呈现联邦统筹、州级细化、多部门分权协同的双层架构~\cite{Executive_Order_14409, M‑25‑21}。
% 美国依托白宫行政令、OMB管理备忘录、州级专项立法形成三级管控体系，聚焦前沿大模型驱动的智能体灾难性风险、网络安全滥用、透明度缺失三大核心问题~\cite{SB‑53,2025‑A6453A‑Bill}；加拿大立足主权AI生态建设，以国家战略引导产业发展，通过修订PIPEDA法案规范智能体全链路个人信息处理行为，形成国家战略与隐私法规相互配合的治理模式~\cite{AI_for_ALL,PIPEDA}。
本节基于美国、加拿大现行及已公示待生效政策文件，汇总北美地区AI智能体安全治理的主管主体、法规谱系与核心管控条款，客观呈现北美分权式治理模式的规则特征与覆盖边界，主要发现如下（详见表~\ref{tab:north-america-ai-agent-policy}）：

\begin{itemize}[tabitem]
    \item \textbf{美国形成\textbf{联邦顶层规制+州级场景强约束}的双层治理结构。}联邦层面依靠EO 14409行政令、OMB的M‑25‑21备忘录从国家层面制定通用安全基线，重点统筹前沿模型安全评估、关键基础设施防护与联邦机构AI应用管控~\cite{Executive_Order_14409,M‑25‑21}；加州SB‑53法案、纽约州RAISE法案进行场景化加码，分别针对通用前沿模型透明度、金融行业智能体灾难性风险评估提出更加严苛的落地要求，形成联邦底线约束、地方根据行业特征收紧合规标准的格局~\cite{SB‑53,2025‑A6453A‑Bill}。

    \item \textbf{北美政策高度重视AI智能体的灾难性风险前置管控。}美国联邦行政令强制界定前沿模型判定标准，要求开发者完成安全测试、风险缓解措施备案；加州与纽约州进一步把灾难性风险评估设置为部署前强制性门槛，配套安全事件上报机制和民事处罚手段，将事后监管转变为部署前风险管控，从法律层面对智能体自主链式攻击、越权访问等潜在重大风险提前设防~\cite{Executive_Order_14409,SB‑53,2025‑A6453A‑Bill}。

    \item \textbf{加拿大治理体系以\textbf{隐私合规+主权生态}为核心。}加拿大《AI for All》国家战略布局本土算力、人才储备与科研成果转化，打造独立于美国的主权AI生态~\cite{AI_for_ALL}；修订之后的PIPEDA法案将智能体场景纳入隐私监管范围，对数据收集范围、知情同意、跨境流转、访问审计作出硬性约束，通过隐私合规划定智能体数据安全边界，实现生态建设和安全约束双向并行~\cite{PIPEDA}。
\end{itemize}

\begin{center}
% \nopagebreak[4]
\enlargethispage{1cm}
\footnotesize
\setlength{\tabcolsep}{3pt}

% \begin{longtblr}[
% caption = {北美AI智能体安全政策和法规汇总},
% label   = {tab:north-america-ai-agent-policy},
% ]
% {
% colspec = {
% |m{0.8cm}
% |m{2cm}
% |m{2.2cm}
% |m{1.5cm}
% |m{
%     \dimexpr
%     \linewidth
%     - 0.8cm
%     - 2cm
%     - 2.2cm
%     - 1.5cm
%     - 10\tabcolsep
%     - 5\arrayrulewidth
%     \relax
% }|
% },
% rowhead = 1,
% }

\begin{center}
% \nopagebreak[4]
\enlargethispage{1cm}
\footnotesize

\begin{longtblr}[
  caption = {北美AI智能体安全政策和法规汇总},
  label   = {tab:north-america-ai-agent-policy},
]{
  width   = \linewidth,
  colspec = {
    |Q[wd=0.8cm,l,m]
    |Q[wd=2cm,l,m]
    |Q[wd=2.2cm,l,m]
    |Q[wd=1.5cm,l,m]
    |X[1,l,m]|
  },
  cells   = {l,m},
  colsep  = 3pt,
  row{1}  = {font=\bfseries},
  rowhead = 1,
  rowsep  = 5pt,
  measure = vbox,
  stretch = -1,
}

\hline
\textbf{国家}
& \textbf{主管部门}
& \textbf{政策法规}
& \textbf{发布日期}
& \textbf{核心内容} \\
\hline

\SetCell[r=4]{c,m} 美国
& \begin{itemize}[tabitem]
    \item 白宫（牵头）
    \item 财政部（深度参与）
    \item NSA（深度参与）
\end{itemize}
& EO 14409《促进先进人工智能创新与安全行政令》~\cite{Executive_Order_14409}
& 2026.06.02
& \begin{itemize}[tabitem]
    \item 建立covered frontier model的联邦基准、评估和自愿安全框架；
    \item 建立AI网络安全信息共享和协调机制，服务关键基础设施与网络防护；
    \item 要求执法部门优先打击AI辅助的非法计算机访问、网络入侵等犯罪行为。
\end{itemize} \\
\cline{2-5}

% & \begin{itemize}[tabitem]
%     \item OMB
% \end{itemize}

& 白宫管理与预算办公室（OMB）
& M-25-21联邦高影响AI系统管控备忘录~\cite{M-25-21}
& 2025.04.03
& \begin{itemize}[tabitem]
    \item 要求联邦机构建立AI治理机制，并对高影响AI用例实施风险识别、测试、监测和问责；
    \item 要求维护 AI 用例清单和相关合规信息，提高公众可见性和机构内部治理能力；
    \item 强调采购、隐私、公平、安全、人工监督和公共信任要求，推动负责任的联邦AI部署。
\end{itemize} \\
\cline{2-5}

% & \begin{itemize}[tabitem]
%     \item CPPA
% \end{itemize}

& 加州隐私保护局（CPPA）
& SB 53《前沿AI透明度法案》~\cite{SB-53}
& 2025.09.29（2026.01.01生效）
& \begin{itemize}[tabitem]
    \item 要求大型前沿模型开发者公开或提交前沿模型安全框架、风险评估和缓解措施；
    \item 建立关键安全事件报告机制，覆盖可能造成重大危害的前沿模型安全事故；
    \item 设置举报人保护和州总检察长执法机制，对违反透明度与安全义务的主体追责。
\end{itemize} \\
\cline{2-5}

% & \begin{itemize}[tabitem]
%     \item 纽约州金融服务部（NYDFS）
% \end{itemize}

& 纽约州金融服务部（NYDFS）
& RAISE负责任AI安全法案~\cite{2025-A6453A-Bill}
& 2025.12.16（2027.01.01生效）
& \begin{itemize}[tabitem]
    \item 面向大型前沿AI模型开发者设定安全与安保协议披露义务；
    \item 要求在部署或重大更新前评估灾难性风险，并说明测试、缓解和安全控制措施；
    \item 授权州总检察长提起民事诉讼，并对违反安全协议、信息披露和员工保护义务的主体实施民事处罚。
\end{itemize} \\
\hline

\SetCell[r=2]{c,m} 加拿大

% & \begin{itemize}[tabitem]
%     \item 创新、科学与经济发展部
% \end{itemize}

& 创新、科学与经济发展部
& 《AI for All》全国人工智能战~\cite{AI_for_ALL}
& 2026.06.08
& \begin{itemize}[tabitem]
    \item 推进加拿大AI生态、人才培养、科研转化和公共部门/产业应用；
    \item 强化可信AI、安全治理、标准协作和公众信任建设；
    \item 支持主权AI算力、数据与基础设施能力，提升加拿大本土AI发展支撑。
\end{itemize} \\
\cline{2-5}

% & \begin{itemize}[tabitem]
%     \item 联邦隐私专员办公室
% \end{itemize}

& 联邦隐私专员办公室
& PIPEDA个人信息保护法案（AI修订版~\cite{PIPEDA}
& 长期生效，2026修订
& \begin{itemize}[tabitem]
    \item 规范加拿大私营部门组织在商业活动中收集、使用和披露个人信息的行为，并适用于联邦监管企业的员工个人信息及跨省或跨境流动的个人信息；
    \item 要求组织事先明确个人信息处理目的并取得个人知情同意，同时限制个人信息的收集、使用、披露和保存范围；
    \item 要求组织落实问责、信息准确性、安全保护和公开透明义务，并保障个人访问、更正个人信息及对组织合规情况提出异议的权利。
\end{itemize} \\
\hline

\end{longtblr}
\end{center}

% \begin{center}
% % \nopagebreak[4]
% \enlargethispage{1cm}
% \footnotesize
% \setlength{\tabcolsep}{3pt}
% \begin{longtable}{|>{\rule{0pt}{2pt}}m{0.8cm}
% |>{\rule{0pt}{2pt}}m{2cm}
% |>{\rule{0pt}{2pt}}m{2.2cm}
% |>{\rule{0pt}{2pt}}m{1.5cm}
% |>{\rule{0pt}{2pt}}m{\dimexpr\linewidth - 0.8cm - 2cm - 2.2cm - 1.5cm - 10\tabcolsep\relax}|}
% \caption{北美AI智能体安全政策和法规汇总}\label{tab:north-america-ai-agent-policy}\\
% \hline
% \textbf{国家} & \textbf{主管部门} & \textbf{政策法规} & \textbf{发布日期} & \textbf{核心内容}\\
% \hline
% \endfirsthead
% \multicolumn{5}{l}{\tablename\ \thetable\ -- 续上页}\\
% \hline
% \textbf{国家} & \textbf{主管部门} & \textbf{政策法规} & \textbf{发布日期} & \textbf{核心内容}\\
% \hline
% \endhead
% \hline
% \multicolumn{5}{r}{续下页}\\
% \endfoot
% \hline
% \endlastfoot

% \multirow{4}{=}{美国}
% & \begin{itemize}[tabitem]
%     \item 白宫（牵头）
%     \item 财政部（深度参与）
%     \item NSA（深度参与）
% \end{itemize}
% & EO 14409《促进先进人工智能创新与安全行政令》~\cite{Executive_Order_14409}
% & 2026.06.02
% & \begin{itemize}[tabitem]
%     \item 建立covered frontier model的联邦基准、评估和自愿安全框架；
%     \item 建立AI网络安全信息共享和协调机制，服务关键基础设施与网络防护；
%     \item 要求执法部门优先打击AI辅助的非法计算机访问、网络入侵等犯罪行为。
% \end{itemize} \\
% \cline{2-5}
% % & \begin{itemize}[tabitem]
% %     \item OMB
% % \end{itemize}
% &白宫管理与预算办公室（OMB）
% & M-25-21联邦高影响AI系统管控备忘录~\cite{M-25-21}
% & 2025.04.03
% & \begin{itemize}[tabitem]
%     \item 要求联邦机构建立AI治理机制，并对高影响AI用例实施风险识别、测试、监测和问责；
%     \item 要求维护 AI 用例清单和相关合规信息，提高公众可见性和机构内部治理能力；
%     \item 强调采购、隐私、公平、安全、人工监督和公共信任要求，推动负责任的联邦AI部署。
% \end{itemize} \\
% \cline{2-5}
% % & \begin{itemize}[tabitem]
% %     \item CPPA
% % \end{itemize}
% &加州隐私保护局（CPPA）
% & SB 53《前沿AI透明度法案》~\cite{SB-53}
% & 2025.09.29（2026.01.01生效）
% & \begin{itemize}[tabitem]
%     \item 要求大型前沿模型开发者公开或提交前沿模型安全框架、风险评估和缓解措施；
%     \item 建立关键安全事件报告机制，覆盖可能造成重大危害的前沿模型安全事故；
%     \item 设置举报人保护和州总检察长执法机制，对违反透明度与安全义务的主体追责。
% \end{itemize} \\
% \cline{2-5}
% % & \begin{itemize}[tabitem]
% %     \item 纽约州金融服务部（NYDFS）
% % \end{itemize}
% &纽约州金融服务部（NYDFS）
% & RAISE负责任AI安全法案~\cite{2025-A6453A-Bill}
% & 2025.12.16（2027.01.01生效）
% & \begin{itemize}[tabitem]
%     \item 面向大型前沿AI模型开发者设定安全与安保协议披露义务；
%     \item 要求在部署或重大更新前评估灾难性风险，并说明测试、缓解和安全控制措施；
%     \item 授权州总检察长提起民事诉讼，并对违反安全协议、信息披露和员工保护义务的主体实施民事处罚。
% \end{itemize} \\
% \hline
% \multirow{2}{=}{加拿大}
% % & \begin{itemize}[tabitem]
% %     \item 创新、科学与经济发展部
% % \end{itemize}
% & 创新、科学与经济发展部
% & 《AI for All》全国人工智能战~\cite{AI_for_ALL}
% & 2026.06.08
% & \begin{itemize}[tabitem]
%     \item 推进加拿大AI生态、人才培养、科研转化和公共部门/产业应用；
%     \item 强化可信AI、安全治理、标准协作和公众信任建设；
%     \item 支持主权AI算力、数据与基础设施能力，提升加拿大本土AI发展支撑。
% \end{itemize} \\
% \cline{2-5}
% % & \begin{itemize}[tabitem]
% %     \item 联邦隐私专员办公室
% % \end{itemize}
% & 联邦隐私专员办公室
% & PIPEDA个人信息保护法案（AI修订版~\cite{PIPEDA}
% & 长期生效，2026修订
% & \begin{itemize}[tabitem]
%     \item 规范加拿大私营部门组织在商业活动中收集、使用和披露个人信息的行为，并适用于联邦监管企业的员工个人信息及跨省或跨境流动的个人信息；
%     \item 要求组织事先明确个人信息处理目的并取得个人知情同意，同时限制个人信息的收集、使用、披露和保存范围；
%     \item 要求组织落实问责、信息准确性、安全保护和公开透明义务，并保障个人访问、更正个人信息及对组织合规情况提出异议的权利。
% \end{itemize} 
% \end{longtable}
% \end{center}


\paragraph{北美AI领域出口管制体系:}
美国与加拿大共同构成北美AI战略技术出口管制体系，二者均依托基础立法+动态管制清单搭建两用技术管控框架，但在执法逻辑、管控工具、政策导向层面存在显著差异。
% 美国形成多层级行政规则、实体清单与定向芯片审查组合管控模式，同时配套面向可信主体的AI技术对外输出计划；加拿大以统一进出口许可法案为根基，通过年度清单修订、敏感技术目录划定AI相关技术风险边界。
本节客观呈现北美区域AI算力、模型、算法技术跨境流转的合规约束标准，为跨区域AI智能体技术出口风险对比提供统一实证依据，核心发现如下（详见表~\ref{tab:north_america_ai_export_control}）：
\begin{itemize}[tabitem]
    \item \textbf{北美两国均采用“基础立法+动态清单”双层管制架构，规则持续高频迭代。}美国以ECRA+EAR、加拿大以EIPA+ECL作为长期稳定法定底座，每年通过行政令、清单修订调整AI与半导体管控参数，美国还出现专项AI管制规则发布后废止的政策波动现象。
    \item \textbf{美国管控工具维度更丰富，形成限制输出与定向开放并行双路径。}除清单、实体清单、专项芯片审查外，美国增设《美国人工智能出口计划》面向可信买方推广全栈AI技术；加拿大仅依托清单修订与敏感技术目录完成风险识别，无专项技术输出扶持机制。
    \item \textbf{美加管制判定均以技术参数、物项类别为核心依据，不针对AI概念全域禁限。}两国均不笼统约束所有AI框架、数据集，仅硬件、模型、技术落入清单明确参数范围，结合目的地、最终用户、转出口风险综合判定出口许可义务。
\end{itemize}



\begingroup
\enlargethispage{1cm}
\footnotesize
\centering
\begin{longtblr}[
  caption = {北美AI智能体领域出口管制政策汇总},
  label   = {tab:north_america_ai_export_control},
]{
  width   = \linewidth,
  colspec = {
    |Q[wd=0.8cm,l,m]
    |Q[wd=1.5cm,l,m]
    |Q[wd=2.5cm,l,m]
    |Q[wd=2cm,l,m]
    |X[1.00,l,m]|
  },
  cells   = {l,m},
  colsep  = 3pt,
  row{1}  = {font=\bfseries},
  rowhead = 1,
  rowsep  = 5pt,
  measure = vbox,
  stretch = -1,
}
\hline
国家 & 主管部门 & 政策法规/管制工具 & 发布/生效日期 & 核心内容 \\
\hline
美国
% & \begin{itemize}[tabitem]
%     \item 国会
%     \item 工业和安全局(BIS)
% \end{itemize}
% & 《出口管制改革法案》ECRA+《出口管理条例》EAR~\cite{ECRA,EAR}
% & ECRA 2018；EAR持续修订
% & \begin{itemize}[tabitem]
%     \item ECRA确立美国对军民两用、敏感和新兴基础技术实施出口管制的法定基础；
%     \item EAR通过CCL/ECCN、许可要求、目的地、最终用户和最终用途规则实施具体管制；
%     \item BIS可通过实体清单、未经验证清单、军事最终用户规则和特定 FDPR 扩展执法覆盖范围。
% \end{itemize} \\
% \cline{2-5}
% & \begin{itemize}[tabitem]
%     \item 商务部BIS
% \end{itemize}
& \SetCell[r=3]{l,m} 工业和安全局(BIS)
& 2025先进AI技术全域扩散管制最终规则~\cite{AI_Diffusion_Rule}
& 2025.01.15生效（2025.05.13被废除）~\cite{Rescission_of_Biden-Era_AI_Diffusion_Rule}
& \begin{itemize}[leftmargin=*,topsep=0pt,partopsep=0pt,itemsep=1pt,parsep=0pt]
    \item 建立面向先进AI芯片和部分闭源AI模型权重的全球许可框架；
    \item 通过国家分组、受信终端用户和许可例外/授权机制区分盟友、受信主体和受限目的地；
    \item 旨在降低先进计算芯片和部分闭源模型权重被转移、规避管制或用于危害美国国家安全与外交政策利益用途的风险。
\end{itemize} \\
\cline{3-5}
% & \begin{itemize}[tabitem]
%     \item 商务部BIS
% \end{itemize}
&
& 2026对华先进芯片审查修订规则~\cite{Revises_License_Review_Policy_for_Semiconductors_Exported_to_China}
& 2026.01.15生效
& \begin{itemize}[leftmargin=*,topsep=0pt,partopsep=0pt,itemsep=1pt,parsep=0pt]
    \item 对NVIDIA H200、AMD MI325X及类似芯片的对华出口许可申请，在满足规定安全条件后实行逐案审查；
    \item 申请人须证明相关出口不会减少目前可供美国客户使用的全球半导体产能；
    \item 中国购买方须建立出口合规和客户筛查程序，相关产品还须在美国接受独立第三方性能与安全测试。
\end{itemize} \\
\cline{3-5}
% & \begin{itemize}[tabitem]
%     \item 商务部BIS
% \end{itemize}
&
& 实体清单/UVL（AI专项增补）~\cite{Entity_List}
& 2024--2026动态扩容
& \begin{itemize}[leftmargin=*,topsep=0pt,partopsep=0pt,itemsep=1pt,parsep=0pt]
    \item 实体清单对被列名主体设置额外许可要求和许可审查政策；
    \item 未经验证清单用于标记 BIS 无法核验最终用途/最终用户的交易对象，并增加出口商尽调义务；
    \item AI 半导体、模型或技术相关限制应回到具体 EAR 条款、实体清单条目和最终用途规则逐条判断。
\end{itemize} \\
\cline{2-5}
& \begin{itemize}[leftmargin=*,topsep=0pt,partopsep=0pt,itemsep=1pt,parsep=0pt]
    \item 国际贸易管理局(ITA)
    \item BIS
\end{itemize}
& 《美国人工智能出口计划》配套规则~\cite{USA_AI_Exports_Program}
& 2026.03落地
& \begin{itemize}[leftmargin=*,topsep=0pt,partopsep=0pt,itemsep=1pt,parsep=0pt]
    \item 面向海外推广由AI优化计算硬件、数据中心存储、模型、网络安全措施和行业应用组成的全栈AI技术包；
    \item 设置预设联合体和按需联合体两种模式，分别服务持续性全球部署和具体项目机会；
    \item 由商务部长会同国务院及相关部门遴选项目，获批的全栈AI技术包向可信外国买方提供。
\end{itemize} \\
\hline
加拿大
% & \begin{itemize}[tabitem]
%     \item 全球事务部 GAC（出口管制司）
% \end{itemize}
& 全球事务部(GAC)
& 《进出口许可法》EIPA+《出口管制清单》ECL（第 5506 项两用战略技术）~\cite{EIPA,ECL_Guide}
& EIPA长期生效；ECL持续年度修订
& \begin{itemize}[leftmargin=*,topsep=0pt,partopsep=0pt,itemsep=1pt,parsep=0pt]
    \item EIPA授权加拿大对特定货物、技术和目的地实施出口/进口许可管理；
    \item ECL列明受控货物与技术类别，出口商需按清单、目的地、最终用户和最终用途判断许可义务；
    \item AI相关软硬件只有在落入清单条目或其他制裁/最终用途规则时，才构成受控出口。
\end{itemize} \\
\cline{3-5}
% % & \begin{itemize}[tabitem]
% %     \item 全球事务部GAC
% % \end{itemize}
% &
% & 2024先进半导体与AI算力管制修订令SOR/2024-112 ~\cite{SOR2024-112}
% & 2024.07.20生效
% & \begin{itemize}[tabitem]
%     \item 修订加拿大Export Control List中的受控物项，以对接国际出口管制义务和战略技术变化；
%     \item 涉及半导体、先进制造或相关技术的控制应以清单条目和技术参数为准；
%     \item 许可审查聚焦物项分类、目的地、最终用户、最终用途和转出口风险，而非单独的“智能体训练”标签。
% \end{itemize} \\
% \cline{3-5}
% % & \begin{itemize}[tabitem]
% %     \item 全球事务部GAC
% % \end{itemize}
&
& 2025 战略技术单边管制修正案DORS/2025-89 ~\cite{DORS2025-89}
& 2025.04.25 生效
& \begin{itemize}[leftmargin=*,topsep=0pt,partopsep=0pt,itemsep=1pt,parsep=0pt]
    \item 对Export Control List进行修订，新增或调整特定战略货物、软件和技术控制；
    \item 出口许可义务取决于清单条目、技术参数、目的地和最终用途，而不是“AI数据集/框架”的泛化名称；
    \item 企业应进行物项分类、最终用户尽调和转出口风险审查。
\end{itemize} \\
\cline{3-5}
% & \begin{itemize}[tabitem]
%     \item 全球事务部GAC
% \end{itemize}
&
& 2025 新版《出口管制清单操作指南》~\cite{ECL_Guide_2025}
& 2025.07.01生效
& \begin{itemize}[leftmargin=*,topsep=0pt,partopsep=0pt,itemsep=1pt,parsep=0pt]
    \item 指导出口商理解Export Control List的分类结构、条目说明和技术参数；
    \item 许可判断应基于物项分类、目的地、最终用户、最终用途和适用豁免；
    \item 对AI或云访问的限制应另查具体清单条目、制裁措施或许可政策，不能由指南本身推导。
\end{itemize} \\
\cline{3-5}
% & \begin{itemize}[tabitem]
%     \item 全球事务部GAC
% \end{itemize}
&
& 《敏感技术研究领域清单》——人工智能与大数据技术~\cite{STL}
& 2024--2026持续扩容更新
& \begin{itemize}[leftmargin=*,topsep=0pt,partopsep=0pt,itemsep=1pt,parsep=0pt]
    \item 将人工智能与大数据技术列为加拿大十一类敏感技术领域之一，视为支撑多领域发展的跨领域基础技术；
    \item AI敏感技术覆盖数据科学、数字孪生、生成式AI、大语言模型、全谱系机器学习技术；
    \item 清单用于外资审查、出口管制、科研安全风险识别，最终风险判定需结合技术转移场景综合研判。
\end{itemize} \\
\hline
\end{longtblr}
\endgroup

% \begin{center}
% % \nopagebreak[4]
% \enlargethispage{1cm}
% \footnotesize
% \setlength{\tabcolsep}{3pt}
% \begin{longtable}{|>{\rule{0pt}{2pt}}m{0.8cm}
% |>{\rule{0pt}{2pt}}m{1.5cm}
% |>{\rule{0pt}{2pt}}m{2.5cm}
% |>{\rule{0pt}{2pt}}m{2cm}
% |>{\rule{0pt}{2pt}}m{\dimexpr\linewidth - 0.8cm - 1.5cm - 2.5cm - 2cm - 10\tabcolsep\relax}|}
% \caption{北美AI智能体领域出口管制政策汇总}\label{tab:north_america_ai_export_control}\\
% \hline
% \textbf{国家} & \textbf{主管部门} & \textbf{政策法规/管制工具} & \textbf{发布/生效日期} & \textbf{核心内容} \\
% \hline
% \endfirsthead
% \multicolumn{5}{l}{\tablename\ \thetable\ -- 续上页}\\
% \hline
% \textbf{国家} & \textbf{主管部门} & \textbf{政策法规/管制工具} & \textbf{发布/生效日期} & \textbf{核心内容} \\
% \hline
% \endhead
% \hline
% \multicolumn{5}{r}{续下页}\\
% \endfoot
% \hline
% \endlastfoot

% \multirow{5}{=}{美国}
% % & \begin{itemize}[tabitem]
% %     \item 国会
% %     \item 工业和安全局(BIS)
% % \end{itemize}
% % & 《出口管制改革法案》ECRA+《出口管理条例》EAR~\cite{ECRA,EAR}
% % & ECRA 2018；EAR持续修订
% % & \begin{itemize}[tabitem]
% %     \item ECRA确立美国对军民两用、敏感和新兴基础技术实施出口管制的法定基础；
% %     \item EAR通过CCL/ECCN、许可要求、目的地、最终用户和最终用途规则实施具体管制；
% %     \item BIS可通过实体清单、未经验证清单、军事最终用户规则和特定 FDPR 扩展执法覆盖范围。
% % \end{itemize} \\
% % \cline{2-5}
% % & \begin{itemize}[tabitem]
% %     \item 商务部BIS
% % \end{itemize}
% & \multirow{9}{=}{工业和安全局(BIS)}
% & 2025先进AI技术全域扩散管制最终规则~\cite{AI_Diffusion_Rule}
% & 2025.01.15生效（2025.05.13被废除）~\cite{Rescission_of_Biden-Era_AI_Diffusion_Rule}
% & \begin{itemize}[tabitem]
%     \item 建立面向先进AI芯片和部分闭源AI模型权重的全球许可框架；
%     \item 通过国家分组、受信终端用户和许可例外/授权机制区分盟友、受信主体和受限目的地；
%     \item 旨在降低先进计算芯片和部分闭源模型权重被转移、规避管制或用于危害美国国家安全与外交政策利益用途的风险。
% \end{itemize} \\
% \cline{3-5}
% % & \begin{itemize}[tabitem]
% %     \item 商务部BIS
% % \end{itemize}
% & 
% & 2026对华先进芯片审查修订规则~\cite{Revises_License_Review_Policy_for_Semiconductors_Exported_to_China}
% & 2026.01.15生效
% & \begin{itemize}[tabitem]
%     \item 对NVIDIA H200、AMD MI325X及类似芯片的对华出口许可申请，在满足规定安全条件后实行逐案审查；
%     \item 申请人须证明相关出口不会减少目前可供美国客户使用的全球半导体产能；
%     \item 中国购买方须建立出口合规和客户筛查程序，相关产品还须在美国接受独立第三方性能与安全测试。
% \end{itemize} \\
% \cline{3-5}
% % & \begin{itemize}[tabitem]
% %     \item 商务部BIS
% % \end{itemize}
% & 
% & 实体清单/UVL（AI专项增补）~\cite{Entity_List}
% & 2024--2026动态扩容
% & \begin{itemize}[tabitem]
%     \item 实体清单对被列名主体设置额外许可要求和许可审查政策；
%     \item 未经验证清单用于标记 BIS 无法核验最终用途/最终用户的交易对象，并增加出口商尽调义务；
%     \item AI 半导体、模型或技术相关限制应回到具体 EAR 条款、实体清单条目和最终用途规则逐条判断。
% \end{itemize} \\
% \cline{2-5}
% & \begin{itemize}[tabitem]
%     \item 国际贸易管理局(ITA)
%     \item BIS
% \end{itemize}

% & 《美国人工智能出口计划》配套规则~\cite{USA_AI_Exports_Program}
% & 2026.03落地
% & \begin{itemize}[tabitem]
%     \item 面向海外推广由AI优化计算硬件、数据中心存储、模型、网络安全措施和行业应用组成的全栈AI技术包；
%     \item 设置预设联合体和按需联合体两种模式，分别服务持续性全球部署和具体项目机会；
%     \item 由商务部长会同国务院及相关部门遴选项目，获批的全栈AI技术包向可信外国买方提供。
% \end{itemize} \\
% \hline
% \multirow{5}{=}{加拿大}
% % & \begin{itemize}[tabitem]
% %     \item 全球事务部 GAC（出口管制司）
% % \end{itemize}
% & \multirow{4}{=}{全球事务部(GAC)}
% & 《进出口许可法》EIPA+《出口管制清单》ECL（第 5506 项两用战略技术）~\cite{EIPA,ECL_Guide}
% & EIPA长期生效；ECL持续年度修订
% & \begin{itemize}[tabitem]
%     \item EIPA授权加拿大对特定货物、技术和目的地实施出口/进口许可管理；
%     \item ECL列明受控货物与技术类别，出口商需按清单、目的地、最终用户和最终用途判断许可义务；
%     \item AI相关软硬件只有在落入清单条目或其他制裁/最终用途规则时，才构成受控出口。
% \end{itemize} \\
% \cline{3-5}
% % % & \begin{itemize}[tabitem]
% % %     \item 全球事务部GAC
% % % \end{itemize}
% % & 
% % & 2024先进半导体与AI算力管制修订令SOR/2024-112 ~\cite{SOR2024-112}
% % & 2024.07.20生效
% % & \begin{itemize}[tabitem]
% %     \item 修订加拿大Export Control List中的受控物项，以对接国际出口管制义务和战略技术变化；
% %     \item 涉及半导体、先进制造或相关技术的控制应以清单条目和技术参数为准；
% %     \item 许可审查聚焦物项分类、目的地、最终用户、最终用途和转出口风险，而非单独的“智能体训练”标签。
% % \end{itemize} \\
% % \cline{3-5}
% % % & \begin{itemize}[tabitem]
% % %     \item 全球事务部GAC
% % % \end{itemize}
% & 
% & 2025 战略技术单边管制修正案DORS/2025-89 ~\cite{DORS2025-89}
% & 2025.04.25 生效
% & \begin{itemize}[tabitem]
%     \item 对Export Control List进行修订，新增或调整特定战略货物、软件和技术控制；
%     \item 出口许可义务取决于清单条目、技术参数、目的地和最终用途，而不是“AI数据集/框架”的泛化名称；
%     \item 企业应进行物项分类、最终用户尽调和转出口风险审查。
% \end{itemize} \\
% \cline{3-5}
% % & \begin{itemize}[tabitem]
% %     \item 全球事务部GAC
% % \end{itemize}
% & 
% & 2025 新版《出口管制清单操作指南》~\cite{ECL_Guide_2025}
% & 2025.07.01生效
% & \begin{itemize}[tabitem]
%     \item 指导出口商理解Export Control List的分类结构、条目说明和技术参数；
%     \item 许可判断应基于物项分类、目的地、最终用户、最终用途和适用豁免；
%     \item 对AI或云访问的限制应另查具体清单条目、制裁措施或许可政策，不能由指南本身推导。
% \end{itemize} \\
% \cline{3-5}
% % & \begin{itemize}[tabitem]
% %     \item 全球事务部GAC
% % \end{itemize}
% & 
% & 《敏感技术研究领域清单》——人工智能与大数据技术~\cite{STL}
% & 2024--2026持续扩容更新
% & \begin{itemize}[tabitem]
%     \item 将人工智能与大数据技术列为加拿大十一类敏感技术领域之一，视为支撑多领域发展的跨领域基础技术；
%     \item AI敏感技术覆盖数据科学、数字孪生、生成式AI、大语言模型、全谱系机器学习技术；
%     \item 清单用于外资审查、出口管制、科研安全风险识别，最终风险判定需结合技术转移场景综合研判。
% \end{itemize} 
% \end{longtable}
% \end{center}

\subsection{欧洲AI智能体安全政策与法规}
欧洲是全球首个建立统一AI法律体系的区域，依托《人工智能法案》（AI‑ACT）为核心底座，叠加GDPR数据合规规则、各国专项操作守则，构建了层级清晰、风险导向、执法刚性最强的AI智能体治理架构~\cite{eu_ai_act_2024}。
% 欧盟依靠统一立法实现全部成员国强制适用；英国脱欧之后脱离欧盟法案约束，形成行业自愿安全标准搭配数据合规指引的双轨治理模式；瑞士、俄罗斯立足自身国家安全诉求，分别出台数据本地驻留、本土模型优先部署的管控条例，形成欧盟统一规则和国别差异化约束并行的格局~\cite{uk_ai_cyber_security_code_2025,Russia_AI_Draft}。
% 本节整合欧盟、英国、瑞士、俄罗斯的政策文件，梳理区域内统一立法、国别细则、行业标准的层级关系，拆解
本小节梳理欧洲主要国家在智能体人工监督、数据跨境、漏洞治理、主权管控等维度的差异化规则，主要发现如下（详见表~\ref{tab:europe_agent_policy_summary}）:

\begin{itemize}[tabitem]
    \item \textbf{欧盟以AI ACT为核心实现全成员国统一强制监管。}依据Reg.2024/1689法案，欧盟对高风险智能体划定统一合规底线，硬性规定产品上线前开展标准化红队安全测试、运行期间全链路操作审计和全程人工监督，并且法案具备域外管辖权，违规企业最高可处以全球年度营收7\%的巨额罚款，依靠严苛惩罚倒逼企业落实安全责任~\cite{eu_ai_act_2024}。

    \item \textbf{欧洲形成AI安全与数据合规深度绑定的治理逻辑。}EDPB出台AI场景配套GDPR实施细则，将数据最小化原则落实到智能体工具调用、知识库检索、联网搜索环节，用户可以随时撤销授权；同时确立72小时数据泄露强制上报制度，智能体所有访问行为留存完整审计日志，把隐私合规作为智能体准入的前置条件~\cite{eu_gpai_code_of_practice_2025}。

    \item \textbf{非欧盟欧洲国家侧重主权安全与场景管控。}俄罗斯联邦AI法案明确关键基础设施仅允许本土智能体运行，对政务、能源领域的智能体设置互联网访问白名单，防范域外模型带来安全隐患~\cite{Russia_AI_Draft}；瑞士修订FADP法案，要求商用智能体的数据必须在境内存储，高‑风险场景优先采用TEE硬件隔离方案，借助硬件层面筑牢隐私安全屏障，形成对欧盟法规的补充性区域治理规则。
\end{itemize}



% \begin{center}
% % \nopagebreak[4]
% \enlargethispage{0.8cm}
% \footnotesize
% \setlength{\tabcolsep}{3pt}

% \begin{longtblr}[
% caption = {欧洲AI智能体安全政策和法规汇总},
% label   = {tab:europe_agent_policy_summary},
% ]
% {
% colspec = {
% |m{1cm}
% |m{2cm}
% |m{2.2cm}
% |m{1.5cm}
% |m{
%     \dimexpr
%     \linewidth
%     - 1cm
%     - 2cm
%     - 2.2cm
%     - 1.5cm
%     - 10\tabcolsep
%     - 5\arrayrulewidth
%     \relax
% }|
% },
% rowhead = 1,
% }

\begin{center}
% \nopagebreak[4]
\enlargethispage{0.8cm}
\footnotesize

\begin{longtblr}[
  caption = {欧洲AI智能体安全政策和法规汇总},
  label   = {tab:europe_agent_policy_summary},
]{
  width   = \linewidth,
  colspec = {
    |Q[wd=1cm,l,m]
    |Q[wd=2cm,l,m]
    |Q[wd=2.2cm,l,m]
    |Q[wd=1.5cm,l,m]
    |X[1,l,m]|
  },
  cells   = {l,m},
  colsep  = 3pt,
  row{1}  = {font=\bfseries},
  rowhead = 1,
  rowsep  = 5pt,
  measure = vbox,
  stretch = -1,
}

\hline
\textbf{国家}
& \textbf{主管部门}
& \textbf{政策法规名称}
& \textbf{发布/生效日期}
& \textbf{核心内容} \\
\hline

\SetCell[r=2]{c} 欧盟
& \begin{itemize}[tabitem]
    \item 欧委会（牵头）
    \item 欧洲议会
\end{itemize}
& AI ACT\newline(Reg.2024/\newline1689)~\cite{eu_ai_act_2024}
& 2024.03通过，2026.08生效
& \begin{itemize}[tabitem]
    \item 按风险分级管控高风险智能体，强制全程人工监督机制；
    \item 大模型智能体上线前必须完成标准化红队测试、全链路操作审计；
    \item 具备域外管辖权，违规企业最高处以全球年度营收7\%罚款。
\end{itemize} \\
\cline{2-5}

% \hline
% 欧盟
% & \begin{itemize}[tabitem]
% \item EDPB
% \end{itemize}

& 欧洲数据保护委员会(EDBP)
& GDPR（AI专项指引）~\cite{eu_gpai_code_of_practice_2025}
& 2018正式生效，2025增补AI配套细则
& \begin{itemize}[tabitem]
    \item 智能体严格遵循数据最小化，禁止自主无差别抓取用户敏感信息；
    \item 用户可随时撤回联网检索、第三方插件调用、知识库访问授权；
    \item 发生数据泄露需72小时内强制上报监管机构，完整留存操作存证。
\end{itemize} \\
\hline

\SetCell[r=2]{c} 英国

% & \begin{itemize}[tabitem]
% \item DSIT
% \end{itemize}

& 科学、创新和技术部(DIST)
& 《AI网络安全操作守则》~\cite{uk_ai_cyber_security_code_2025}
& 2025.01.31（行业自愿落地标准）
& \begin{itemize}[tabitem]
    \item 制定智能体专属安全防护基线，抵御提示注入、供应链投毒漏洞；
    \item 高自主智能体必须配置分级人工复核与审批流程；
    \item 搭建行业统一智能体漏洞情报共享、攻击溯源协同机制。
\end{itemize} \\
\cline{2-5}

% & \begin{itemize}[tabitem]
% \item ICO
% \end{itemize}

& 信息专员办公室(ICO)
& 智能体数据合规指引
& 2026.03正式发布
& \begin{itemize}[tabitem]
    \item 企业对旗下智能体全部操作行为承担完整法律责任；
    \item 严格限制营销类Agent批量用户触达、无授权深度伪造内容生成；
    \item 金融、医疗等高风险场景智能体强制部署权限隔离安全沙箱。
\end{itemize} \\
\hline

% 瑞士
% % & \begin{itemize}[tabitem]
% % \item FDPIC
% % \end{itemize}
% & 联邦数据保护和信息专员(FDPIC)
% & revFADP修订数据保护法
% & 2023.09.01正式生效
% & \begin{itemize}[tabitem]
% \item 商用智能体处理数据强制本地驻留，跨境数据传输需签署专项合规协议；
% \item 智能体自动化决策全链路可追溯，违规可追究企业负责人连带处罚；
% \item 金融、医疗等高风险行业推荐基于TEE硬件隔离部署智能体。
% \end{itemize} \\
% \hline

俄罗斯
& \begin{itemize}[tabitem]
    \item 数字发展部（牵头）
    \item 联邦安全局（深度参与）
\end{itemize}
& 联邦AI法（草案）~\cite{Russia_AI_Draft}
& 2026.03对外公示，2027.09拟生效
& \begin{itemize}[tabitem]
    \item 关键基础设施仅允许本土主权智能体部署运行；
    \item 政务、能源类智能体执行互联网外联白名单严格管控；
    \item AI产品必须添加可识别透明标识，违规企业限制境内市场准入。
\end{itemize} \\
\hline

\end{longtblr}
\end{center}
% \begin{center}
% % \nopagebreak[4]
% \enlargethispage{0.8cm}
% \footnotesize
% \setlength{\tabcolsep}{3pt}
% \begin{longtable}{|>{\rule{0pt}{2pt}}m{1cm}
% |>{\rule{0pt}{2pt}}m{2cm}
% |>{\rule{0pt}{2pt}}m{2.2cm}
% |>{\rule{0pt}{2pt}}m{1.5cm}
% |>{\rule{0pt}{2pt}}m{\dimexpr\linewidth - 0.8cm - 2cm - 2.2cm - 1.5cm - 10\tabcolsep\relax}|}
% \caption{欧洲AI智能体安全政策和法规汇总}\label{tab:europe_agent_policy_summary}\\
% \hline
% \textbf{国家} & \textbf{主管部门} & \textbf{政策法规名称} & \textbf{发布/生效日期} & \textbf{核心内容} \\
% \hline
% \endfirsthead
% \multicolumn{5}{l}{\tablename\ \thetable\ -- 续上页}\\
% \hline
% \textbf{国家} & \textbf{主管部门} & \textbf{政策法规名称} & \textbf{发布/生效日期} & \textbf{核心内容} \\
% \hline
% \endhead
% \hline
% \multicolumn{5}{r}{续下页}\\
% \endfoot
% \hline
% \endlastfoot

% \multirow{4}{=}{欧盟}
% & \begin{itemize}[tabitem]
% \item 欧委会（牵头）
% \item 欧洲议会
% \end{itemize}
% & AI ACT\newline(Reg.2024/\newline1689)~\cite{eu_ai_act_2024}
% & 2024.03通过，2026.08生效
% & \begin{itemize}[tabitem]
% \item 按风险分级管控高风险智能体，强制全程人工监督机制；
% \item 大模型智能体上线前必须完成标准化红队测试、全链路操作审计；
% \item 具备域外管辖权，违规企业最高处以全球年度营收7\%罚款。
% \end{itemize} \\
% \cline{2-5}
% % \hline
% % 欧盟
% % & \begin{itemize}[tabitem]
% % \item EDPB
% % \end{itemize}
% & 欧洲数据保护委员会(EDBP)
% & GDPR（AI专项指引）~\cite{eu_gpai_code_of_practice_2025}
% & 2018正式生效，2025增补AI配套细则
% & \begin{itemize}[tabitem]
% \item 智能体严格遵循数据最小化，禁止自主无差别抓取用户敏感信息；
% \item 用户可随时撤回联网检索、第三方插件调用、知识库访问授权；
% \item 发生数据泄露需72小时内强制上报监管机构，完整留存操作存证。
% \end{itemize} \\
% \hline
% \multirow{4}{=}{英国}
% % & \begin{itemize}[tabitem]
% % \item DSIT
% % \end{itemize}
% & 科学、创新和技术部(DIST)
% & 《AI网络安全操作守则》~\cite{uk_ai_cyber_security_code_2025}
% & 2025.01.31（行业自愿落地标准）
% & \begin{itemize}[tabitem]
% \item 制定智能体专属安全防护基线，抵御提示注入、供应链投毒漏洞；
% \item 高自主智能体必须配置分级人工复核与审批流程；
% \item 搭建行业统一智能体漏洞情报共享、攻击溯源协同机制。
% \end{itemize} \\
% \cline{2-5}
% % & \begin{itemize}[tabitem]
% % \item ICO
% % \end{itemize}
% & 信息专员办公室(ICO)
% & 智能体数据合规指引
% & 2026.03正式发布
% & \begin{itemize}[tabitem]
% \item 企业对旗下智能体全部操作行为承担完整法律责任；
% \item 严格限制营销类Agent批量用户触达、无授权深度伪造内容生成；
% \item 金融、医疗等高风险场景智能体强制部署权限隔离安全沙箱。
% \end{itemize} \\
% \hline
% % 瑞士
% % % & \begin{itemize}[tabitem]
% % % \item FDPIC
% % % \end{itemize}
% % & 联邦数据保护和信息专员(FDPIC)
% % & revFADP修订数据保护法
% % & 2023.09.01正式生效
% % & \begin{itemize}[tabitem]
% % \item 商用智能体处理数据强制本地驻留，跨境数据传输需签署专项合规协议；
% % \item 智能体自动化决策全链路可追溯，违规可追究企业负责人连带处罚；
% % \item 金融、医疗等高风险行业推荐基于TEE硬件隔离部署智能体。
% % \end{itemize} \\
% % \hline
% 俄罗斯
% & \begin{itemize}[tabitem]
% \item 数字发展部（牵头）
% \item 联邦安全局（深度参与）
% \end{itemize}
% & 联邦AI法（草案）~\cite{Russia_AI_Draft}
% & 2026.03对外公示，2027.09拟生效
% & \begin{itemize}[tabitem]
% \item 关键基础设施仅允许本土主权智能体部署运行；
% \item 政务、能源类智能体执行互联网外联白名单严格管控；
% \item AI产品必须添加可识别透明标识，违规企业限制境内市场准入。
% \end{itemize} 
% \end{longtable}
% \end{center}



\paragraph{欧洲AI领域出口管制体系:}
欧盟、英国、荷兰构成欧洲AI两用技术出口管制核心主体，区域内形成欧盟统一两用物项框架、英国独立管制立法、荷兰半导体专项限制三层并行管控体系。
% 欧系管制兼顾区域统一规则与各国差异化产业约束，既依托通用两用物项条例搭建基础许可机制，又结合AI法案、半导体设备专项政令、国别制裁清单形成分层约束。
本节通过结构化台账归集欧盟、荷兰、英国的主管机构、核心法规、生效时间与管控边界，客观呈现欧洲AI出口管制的立法渊源、许可模式与差异化产业管控规则，为跨境AI智能体技术流转合规判定提供欧洲区域政策依据, 主要发现如下（详见表~\ref{tab:europe_ai_export_control}):
\begin{itemize}[tabitem]
    \item \textbf{欧盟形成统一基础管制框架，AI监管与出口管制双向联动。}以Reg.(EU) 2021/821两用物项条例作为全域统一管制底座，同步配套AI Act分级监管指引，高风险AI系统、通用大模型合规要求同步纳入出口授权审查范畴。
    \item \textbf{欧洲各国依托本土优势产业增设专项管制条款。}荷兰针对ASML先进光刻设备叠加配套AI技术单独出台逐案许可机制；英国独立立法将大模型权重、智能体源代码、高性能AI加速器直接纳入受控技术清单。
    \item \textbf{欧系许可与制裁体系具备明确国别差异化约束。}欧盟统一授权体系区分多类通用许可；英国OGEL通用许可仅对五眼联盟开放，对特定国家交易执行从严审查，同时依托动态制裁名单实现交易对手准入管控~\cite{The_Export_Control_2024}。
\end{itemize}


% \begin{center}
% % \nopagebreak[4]
% \enlargethispage{0.8cm}
% \footnotesize
% \setlength{\tabcolsep}{3pt}
% % \renewcommand{\arraystretch}{0.7}

% \begin{longtblr}[
% caption = {欧洲AI智能体领域出口管制政策汇总},
% label   = {tab:europe_ai_export_control},
% ]
% {
% colspec = {
% |>{\rule{0pt}{2pt}}m{1cm}
% |>{\rule{0pt}{2pt}}m{2.5cm}
% |>{\rule{0pt}{2pt}}m{2.5cm}
% |>{\rule{0pt}{2pt}}m{2cm}
% |>{\rule{0pt}{2pt}}m{
%     \dimexpr
%     \linewidth
%     - 0.8cm
%     - 2.5cm
%     - 2.5cm
%     - 2cm
%     - 10\tabcolsep
%     \relax
% }|
% },
% rowhead = 1,
% }

\begin{longtblr}[
  caption = {欧洲AI智能体领域出口管制政策汇总},
  label   = {tab:europe_ai_export_control},
]{
  width   = \linewidth,
  colspec = {
    |Q[wd=1cm,l,m]
    |Q[wd=2.5cm,l,m]
    |Q[wd=2.5cm,l,m]
    |Q[wd=2cm,l,m]
    |X[1,l,m]|
  },
  cells   = {l,m},
  colsep  = 3pt,
  row{1}  = {font=\bfseries},
  rowhead = 1,
  rowsep  = 5pt,
  measure = vbox,
  stretch = -1,
}
\hline
\textbf{国家}
& \textbf{主管部门}
& \textbf{管制法规/工具}
& \textbf{发布/生效日期}
& \textbf{核心内容} \\
\hline

\SetCell[r=2]{c} 欧盟
& \begin{itemize}[tabitem]
    \item 欧盟委员会（牵头）
    \item 欧洲议会（深度参与）
    \item 欧洲理事会（深度参与）
\end{itemize}
& 两用物项条例Reg.(EU) 2021/821
& 2021.09，持续增补AI管制条目，2025.11.15发布最新整理版本
& \begin{itemize}[tabitem]
    \item 建立欧盟对两用物项出口、过境、经纪活动和技术援助的统一管制制度；
    % \item 附件I所列货物、软件和技术出口原则上须取得授权，部分未列项物项在特定最终用途或网络监控风险下也可能受控
    \item 设置单项、全球和欧盟通用出口授权等许可形式，并结合物项、目的地、最终用户和最终用途实施审查。
\end{itemize} \\
\cline{2-5}

& \begin{itemize}[tabitem]
    \item 欧盟委员会（牵头）
\end{itemize}
& AI Act 配套出口管制指引~\cite{EU_AI_act}
& 2024.08.01生效，2026.08落地覆盖，部分独立高风险系统推迟
& \begin{itemize}[tabitem]
    \item AI Act按风险等级对AI系统实施分级监管，并禁止不可接受风险的AI用途；
    \item 高风险AI系统须满足风险管理、技术文档、记录保存、透明度、人类监督、准确性和网络安全等要求；
    \item 通用人工智能模型承担透明度义务，具有系统性风险的模型还需开展模型评估和严重事件报告。
\end{itemize} \\
\hline

荷兰
& \begin{itemize}[tabitem]
    \item 经济事务与气候部
\end{itemize}
& ASML设备配套AI管制新规~\cite{ASML}
& 2023–2026动态迭代更新
& \begin{itemize}[tabitem]
    \item 对特定先进半导体制造设备及相关技术实施国家出口许可要求；
    \item 2025年进一步扩大受控设备与技术范围，包括特定测量和检测设备，具体范围以相关部长令为准；
    \item 荷兰政府认为相关技术可能用于生产先进半导体，并在先进军事应用中发挥关键作用，因此实行逐案许可审查。
\end{itemize} \\
\hline

英国
% \multirow{2}{=}{英国}
% & \begin{itemize}[tabitem]
%     \item 英国商业贸易部（DBT）（牵头）
%     \item 出口管制联合局 (ECJU)（深度参与）
% \end{itemize}
% & 《出口管制法2002》+两用物项清单
% & 2002年立法，2024增补AI专项管控条目~\cite{The_Export_Control_2024}
% & \begin{itemize}[tabitem]
%     \item 大模型权重、智能体完整源代码列为受控技术，跨境云端传输同等要求出口许可；
%     \item 依据算力性能阈值对AI加速器实施分级审批管控；
%     \item OGEL通用许可仅面向五眼联盟，对华、俄相关交易执行从严推定审查。
% \end{itemize} \\
% \cline{2-5}

& \begin{itemize}[tabitem]
    \item 英国外交、联邦和发展办公室(FCDO)（牵头）
    \item OFSI（深度参与）
\end{itemize}
& 英国制裁名单（The UK Sanctions List）~\cite{sanction_list}
& 常年动态更新调整
& \begin{itemize}[tabitem]
    \item 英国制裁名单统一列明受英国各项制裁制度约束的个人、实体和船舶；涉及AI企业、研究机构、技术供应商或相关人员时，只有其被正式指定，才可据名单确认其制裁身份；
    \item 开展AI芯片、模型、软件、算力、技术服务或研发合作交易时，应核查交易对手是否列入名单，并进一步识别其是否由受制裁主体直接或间接拥有或控制。
\end{itemize} \\
\hline

\end{longtblr}
\end{center}
```


% \begin{center}
% % \nopagebreak[4]
% \enlargethispage{0.8cm}
% \footnotesize
% \setlength{\tabcolsep}{3pt}
% % \renewcommand{\arraystretch}{0.7}
% \begin{longtable}{|>{\rule{0pt}{2pt}}m{1cm}
% |>{\rule{0pt}{2pt}}m{2.5cm}
% |>{\rule{0pt}{2pt}}m{2.5cm}
% |>{\rule{0pt}{2pt}}m{2cm}
% |>{\rule{0pt}{2pt}}m{\dimexpr\linewidth - 0.8cm - 2.5cm - 2.5cm - 2cm - 10\tabcolsep\relax}|}
% \caption{欧洲AI智能体领域出口管制政策汇总}\label{tab:europe_ai_export_control}\\
% \hline
% \textbf{国家} & \textbf{主管部门} & \textbf{管制法规/工具} & \textbf{发布/生效日期} & \textbf{核心内容} \\
% \hline
% \endfirsthead

% \multicolumn{5}{l}{\tablename\ \thetable\ -- 续上页}\\
% \hline
% \textbf{地区/国家} & \textbf{主管部门} & \textbf{管制法规/工具} & \textbf{发布/生效日期} & \textbf{核心内容} \\
% \hline
% \endhead
% \hline
% \multicolumn{5}{r}{续下页}\\
% \endfoot
% \hline
% \endlastfoot

% \multirow{2}{=}{欧盟}
% & \begin{itemize}[tabitem]
%     \item 欧盟委员会（牵头）
%     \item 欧洲议会（深度参与）
%     \item 欧洲理事会（深度参与）
% \end{itemize}
% & 两用物项条例Reg.(EU) 2021/821
% & 2021.09，持续增补AI管制条目，2025.11.15发布最新整理版本
% & \begin{itemize}[tabitem]
%     \item 建立欧盟对两用物项出口、过境、经纪活动和技术援助的统一管制制度；
%     % \item 附件I所列货物、软件和技术出口原则上须取得授权，部分未列项物项在特定最终用途或网络监控风险下也可能受控
%     \item 设置单项、全球和欧盟通用出口授权等许可形式，并结合物项、目的地、最终用户和最终用途实施审查。
% \end{itemize} \\
% \cline{2-5}
% & \begin{itemize}[tabitem]
%     \item 欧盟委员会（牵头）
% \end{itemize}
% & AI Act 配套出口管制指引~\cite{EU_AI_act}
% & 2024.08.01生效，2026.08落地覆盖，部分独立高风险系统推迟
% & \begin{itemize}[tabitem]
%     \item AI Act按风险等级对AI系统实施分级监管，并禁止不可接受风险的AI用途；
%     \item 高风险AI系统须满足风险管理、技术文档、记录保存、透明度、人类监督、准确性和网络安全等要求；
%     \item 通用人工智能模型承担透明度义务，具有系统性风险的模型还需开展模型评估和严重事件报告。
% \end{itemize} \\
% \hline
% 荷兰
% & \begin{itemize}[tabitem]
%     \item 经济事务与气候部
% \end{itemize}
% & ASML设备配套AI管制新规~\cite{ASML}
% & 2023–2026动态迭代更新
% & \begin{itemize}[tabitem]
%     \item 对特定先进半导体制造设备及相关技术实施国家出口许可要求；
%     \item 2025年进一步扩大受控设备与技术范围，包括特定测量和检测设备，具体范围以相关部长令为准；
%     \item 荷兰政府认为相关技术可能用于生产先进半导体，并在先进军事应用中发挥关键作用，因此实行逐案许可审查。
% \end{itemize} \\
% \hline
% 英国
% % \multirow{2}{=}{英国}
% % & \begin{itemize}[tabitem]
% %     \item 英国商业贸易部（DBT）（牵头）
% %     \item 出口管制联合局 (ECJU)（深度参与）
% % \end{itemize}
% % & 《出口管制法2002》+两用物项清单
% % & 2002年立法，2024增补AI专项管控条目~\cite{The_Export_Control_2024}
% % & \begin{itemize}[tabitem]
% %     \item 大模型权重、智能体完整源代码列为受控技术，跨境云端传输同等要求出口许可；
% %     \item 依据算力性能阈值对AI加速器实施分级审批管控；
% %     \item OGEL通用许可仅面向五眼联盟，对华、俄相关交易执行从严推定审查。
% % \end{itemize} \\
% % \cline{2-5}
% & \begin{itemize}[tabitem]
%     \item 英国外交、联邦和发展办公室(FCDO)（牵头）
%     \item OFSI（深度参与）
% \end{itemize}
% & 英国制裁名单（The UK Sanctions List）~\cite{sanction_list}
% & 常年动态更新调整
% & \begin{itemize}[tabitem]
%     \item 英国制裁名单统一列明受英国各项制裁制度约束的个人、实体和船舶；涉及AI企业、研究机构、技术供应商或相关人员时，只有其被正式指定，才可据名单确认其制裁身份；
%     \item 开展AI芯片、模型、软件、算力、技术服务或研发合作交易时，应核查交易对手是否列入名单，并进一步识别其是否由受制裁主体直接或间接拥有或控制。
% \end{itemize} 
% \end{longtable}
% \end{center}


\subsection{我国AI智能体安全政策与法规}
我国AI智能体安全治理采用\textbf{多部门协同、合规与发展并重}的顶层架构，由网信办、发改委、工信部等部委分领域牵头，形成实施意见、治理框架、伦理办法、国家标准四层政策谱系~\cite{智能体规范应用与创新发展实施意见,人工智能安全治理框架2.0}。相较于海外区域治理模式，我国政策更加重视智能体从研发、训练、部署、运行直至退役的全流程风险管控，统筹推进行业规模化落地应用和安全风险闭环治理，平衡产业创新发展与安全约束的关系~\cite{“人工智能+”行动实施意见,AI科技伦理审查与服务办法（试行）}。
本节梳理我国国家级智能体专项政策、AI融合发展文件、安全治理框架、伦理规则与互联国家标准，客观呈现我国AI智能体安全治理的政策与法规，主要发现如下（详见表~\ref{tab:china_agent_policy_summary}）：

\begin{itemize}[tabitem]
    \item \textbf{构建四层递进式治理体系。}以《智能体规范应用与创新发展实施意见》划定监管红线，《“人工智能+”行动实施意见》夯实产业发展根基，《人工智能安全治理框架2.0》提供风险识别理论依据，《人工智能‑智能体互联》国家标准落地技术约束，分别落实监管约束、产业赋能、风险识别、技术合规四大治理目标，形成自上而下的完备治理链条~\cite{智能体规范应用与创新发展实施意见,“人工智能+”行动实施意见,人工智能安全治理框架2.0,《人工智能智能体互联》系列国标}。

    \item \textbf{治理重心聚焦AI智能体运行层专属风险。}依托《人工智能安全治理框架2.0》和专项实施意见，区别于通用大模型内容安全管控，重点针对身份伪造、权限越权、插件非法接入、知识库污染、多智能体链式协同失控等智能体特有安全隐患；强制落实工具调用审计、全周期日志留存、风险动态监测、事后应急处置机制，补齐智能体运行阶段的监管短板~\cite{人工智能安全治理框架2.0,智能体规范应用与创新发展实施意见}。

    \item \textbf{制定智能体互联专属国家标准。}由中国电子技术标准化研究院牵头制定智能体互联系列国标，统一身份编码、MCP‑类交互协议、接口规范、权限管控体系，从底层技术层面约束跨智能体未授权访问与数据泄露问题；将技术标准与行政监管配套结合，打造区别于欧美依靠商业厂商自主定义协议的模式，形成中国在多智能体可信协同领域的差异化治理优势~\cite{《人工智能智能体互联》系列国标}。
\end{itemize}



% \begin{center}
% % \nopagebreak[4]
% \enlargethispage{0.8cm}
% \footnotesize
% \setlength{\tabcolsep}{3pt}
% % \renewcommand{\arraystretch}{0.7}

% \begin{longtblr}[
% caption = {我国AI智能体安全政策和法规汇总},
% label   = {tab:china_agent_policy_summary},
% ]
% {
% colspec = {
% |>{\rule{0pt}{2pt}}m{0.8cm}
% |>{\rule{0pt}{2pt}}m{3.5cm}
% |>{\rule{0pt}{2pt}}m{2.cm}
% |>{\rule{0pt}{2pt}}m{1.5cm}
% |>{\rule{0pt}{2pt}}m{
%     \dimexpr
%     \linewidth
%     - 0.8cm
%     - 3.5cm
%     - 2.cm
%     - 1.5cm
%     - 10\tabcolsep
%     \relax
% }|
% },
% rowhead = 1,
% }

\begin{center}
\begin{longtblr}[
  caption = {我国AI智能体安全政策和法规汇总},
  label   = {tab:china_agent_policy_summary},
]{
  width   = \linewidth,
  colspec = {
    |Q[wd=0.8cm,l,m]
    |Q[wd=3.5cm,l,m]
    |Q[wd=2cm,l,m]
    |Q[wd=1.5cm,l,m]
    |X[1,l,m]|
  },
  cells   = {l,m},
  colsep  = 3pt,
  row{1}  = {font=\bfseries},
  rowhead = 1,
  rowsep  = 5pt,
  measure = vbox,
  stretch = -1,
}
\hline
\textbf{国家}
& \textbf{主管部门}
& \textbf{政策法规}
& \textbf{发布日期}
& \textbf{核心内容} \\
\hline

\SetCell[r=5]{c} 中国
& \begin{itemize}[tabitem]
    \item 网信办（牵头）
    \item 发改委（深度参与）
    \item 工信部（深度参与）
\end{itemize}
& 《智能体规范应用与创新发展实施意见》~\cite{智能体规范应用与创新发展实施意见}
& 2026.05.08发布
& \begin{itemize}[tabitem]
    \item 规范智能体应用创新，推动在重点行业和公共服务场景合规、安全、有序落地。
    \item 强化智能体身份、权限、工具调用、插件接入和数据使用等关键环节安全管理。
    \item 建立日志留存、风险监测、事件处置和责任闭环机制，提升可追溯、可监管能力。
\end{itemize} \\
\cline{2-5}

& \begin{itemize}[tabitem]
    \item 国家发改委（牵头）
    \item 工信部（深度参与）
    \item 商务部（深度参与）
    \item 教育部（深度参与）
\end{itemize}
& 《“人工智能+”行动实施意见》~\cite{“人工智能+”行动实施意见}
& 2025.08.21发布
& \begin{itemize}[tabitem]
    \item 推动人工智能与科技、产业、消费、民生、治理等领域深度融合，形成重点应用场景；
    \item 完善算力、数据、模型、开源生态、标准和产业协同等基础支撑；
    \item 强化安全治理、标准规范、风险防控和可信应用，促进AI规模化、规范化发展。
\end{itemize} \\
\cline{2-5}

& \begin{itemize}[tabitem]
    \item 国家互联网应急中心（牵头）
    \item 网信办（深度参与）
    \item 全国网安标委（参与）
\end{itemize}
& 《人工智能安全治理框架2.0》~\cite{人工智能安全治理框架2.0}
& 2025.09.01发布
& \begin{itemize}[tabitem]
    \item 从模型、数据、系统、应用和使用等维度识别AI安全风险，覆盖滥用、失控、隐私、内容和网络安全等问题；
    \item 提出AI全生命周期治理思路，包括风险评估、测试验证、监测预警、应急处置和责任管理；
    \item 可作为智能体提示注入、工具调用越权、知识库污染、多智能体协同等风险治理的参考框架。
\end{itemize} \\
\cline{2-5}

& \begin{itemize}[tabitem]
    \item 工信部（牵头）
    \item 发改委（深度参与）等10个部门
    % \item 教育部（深度参与）
    % \item 科技部（深度参与）
    % \item 农业农村部（深度参与）
    % \item 卫健委（深度参与）
    % \item 中国人民银行（深度参与）
    % \item 网信办（深度参与）
    % \item 中国科学院（深度参与）
    % \item 中国科协（深度参与）
\end{itemize}
& 《人工智能科技伦理审查与服务办法（试行）》~\cite{AI科技伦理审查与服务办法（试行）}
& 2026.04发布
& \begin{itemize}[tabitem]
    \item 对人工智能科技活动建立伦理审查、复核、备案和服务管理机制；
    \item 对可能带来较高伦理风险的研究、开发和应用活动，要求开展伦理风险评估和审查；
    \item 重大调整或风险变化时应进行持续跟踪、复核和合规管理，确保科技活动符合伦理规范。
\end{itemize} \\
\cline{2-5}

& \begin{itemize}[tabitem]
    \item 中国电子技术标准化研究院（牵头）
\end{itemize}
& 《人工智能 智能体互联》系列国家标准 ~\cite{《人工智能智能体互联》系列国标}
& 2026.06.26发布
& \begin{itemize}[tabitem]
    \item 构建智能体互联互通标准体系，统一身份识别、交互协议、接口和互操作要求；
    \item 规范智能体间协作、工具调用、权限管理和可信连接，降低未授权访问和滥用风险；
    \item 系列标准覆盖总体架构、身份码、身份管理、智能体描述、智能体发现、智能体交互和工具调用，为智能体互联互通提供统一技术依据。
\end{itemize} \\
\hline

\end{longtblr}
\end{center}


% \begin{center}
% % \nopagebreak[4]
% \enlargethispage{0.8cm}
% \footnotesize
% \setlength{\tabcolsep}{3pt}
% % \renewcommand{\arraystretch}{0.7}
% \begin{longtable}{|>{\rule{0pt}{2pt}}m{0.8cm}
% |>{\rule{0pt}{2pt}}m{3.5cm}
% |>{\rule{0pt}{2pt}}m{2.cm}
% |>{\rule{0pt}{2pt}}m{1.5cm}
% |>{\rule{0pt}{2pt}}m{\dimexpr\linewidth - 0.8cm - 3.5cm - 2.cm - 1.5cm - 10\tabcolsep\relax}|}
% \caption{我国AI智能体安全政策和法规汇总}\label{tab:china_agent_policy_summary}\\
% \hline
% \textbf{国家} & \textbf{主管部门} & \textbf{政策法规} & \textbf{发布日期} & \textbf{核心内容} \\
% \hline
% \endfirsthead

% \multicolumn{5}{l}{\tablename\ \thetable\ -- 续上页}\\
% \hline
% \textbf{国家} & \textbf{主管部门} & \textbf{政策法规} & \textbf{发布日期} & \textbf{核心内容} \\
% \hline
% \endhead
% \hline
% \multicolumn{5}{r}{续下页}\\
% \endfoot
% \hline
% \endlastfoot

% \multirow{5}{=}{我国}
% & \begin{itemize}[tabitem]
%     \item 网信办（牵头）
%     \item 发改委（深度参与）
%     \item 工信部（深度参与）
% \end{itemize}
% & 《智能体规范应用与创新发展实施意见》~\cite{智能体规范应用与创新发展实施意见}
% & 2026.05.08发布
% & \begin{itemize}[tabitem]
%     \item 规范智能体应用创新，推动在重点行业和公共服务场景合规、安全、有序落地。
%     \item 强化智能体身份、权限、工具调用、插件接入和数据使用等关键环节安全管理。
%     \item 建立日志留存、风险监测、事件处置和责任闭环机制，提升可追溯、可监管能力。
% \end{itemize} \\
% \cline{2-5}
% & \begin{itemize}[tabitem]
%     \item 国家发改委（牵头）
%     \item 工信部（深度参与）
%     \item 商务部（深度参与）
%     \item 教育部（深度参与）
% \end{itemize}
% & 《“人工智能+”行动实施意见》~\cite{“人工智能+”行动实施意见}
% & 2025.08.21发布
% & \begin{itemize}[tabitem]
%     \item 推动人工智能与科技、产业、消费、民生、治理等领域深度融合，形成重点应用场景；
%     \item 完善算力、数据、模型、开源生态、标准和产业协同等基础支撑；
%     \item 强化安全治理、标准规范、风险防控和可信应用，促进AI规模化、规范化发展。
% \end{itemize} \\
% \cline{2-5}
% & \begin{itemize}[tabitem]
%     \item 国家互联网应急中心（牵头）
%     \item 网信办（深度参与）
%     \item 全国网安标委（参与）
% \end{itemize}
% & 《人工智能安全治理框架2.0》~\cite{人工智能安全治理框架2.0}
% & 2025.09.01发布
% & \begin{itemize}[tabitem]
%     \item 从模型、数据、系统、应用和使用等维度识别AI安全风险，覆盖滥用、失控、隐私、内容和网络安全等问题；
%     \item 提出AI全生命周期治理思路，包括风险评估、测试验证、监测预警、应急处置和责任管理；
%     \item 可作为智能体提示注入、工具调用越权、知识库污染、多智能体协同等风险治理的参考框架。
% \end{itemize} \\
% \cline{2-5}
% & \begin{itemize}[tabitem]
%     \item 工信部（牵头）
%     \item 发改委（深度参与）等10个部门
%     % \item 教育部（深度参与）
%     % \item 科技部（深度参与）
%     % \item 农业农村部（深度参与）
%     % \item 卫健委（深度参与）
%     % \item 中国人民银行（深度参与）
%     % \item 网信办（深度参与）
%     % \item 中国科学院（深度参与）
%     % \item 中国科协（深度参与）
% \end{itemize}
% & 《人工智能科技伦理审查与服务办法（试行）》~\cite{AI科技伦理审查与服务办法（试行）}
% & 2026.04发布
% & \begin{itemize}[tabitem]
%     \item 对人工智能科技活动建立伦理审查、复核、备案和服务管理机制；
%     \item 对可能带来较高伦理风险的研究、开发和应用活动，要求开展伦理风险评估和审查；
%     \item 重大调整或风险变化时应进行持续跟踪、复核和合规管理，确保科技活动符合伦理规范。
% \end{itemize} \\
% \cline{2-5}
% & \begin{itemize}[tabitem]
%     \item 中国电子技术标准化研究院（牵头）
% \end{itemize}
% & 《人工智能 智能体互联》系列国家标准 ~\cite{《人工智能智能体互联》系列国标}
% & 2026.06.26发布
% & \begin{itemize}[tabitem]
%     \item 构建智能体互联互通标准体系，统一身份识别、交互协议、接口和互操作要求；
%     \item 规范智能体间协作、工具调用、权限管理和可信连接，降低未授权访问和滥用风险；
%     \item 系列标准覆盖总体架构、身份码、身份管理、智能体描述、智能体发现、智能体交互和工具调用，为智能体互联互通提供统一技术依据。
% \end{itemize} 
% \end{longtable}
% \end{center}


\paragraph{我国AI领域出口管制体系:}
依托《出口管制法》搭建顶层法律框架，我国形成以两用物项管制条例、限制出口技术目录为核心的AI相关技术跨境管控体系，实现算力硬件、AI软件、算法技术分层规范。
% 区别于欧美单边限制性管控逻辑，我国管制以清单制为核心判定标准，配套临时管制、最终用户核查、许可审批等全流程合规机制，兼顾产业正常国际合作与国家安全风险防控。
本节梳理我国管制主管权责、法规谱系、生效节点与管控判定规则，客观呈现我国AI两用物项出口管制的制度框架与实操边界，完成全球主要经济体AI出口管制体系的完整对照，具体信息详见表~\ref{tab:china_ai_export_control}：

\begin{itemize}[tabitem]
    \item \textbf{双层法规构成完整管制底座，多部门协同分工实施监管。}以《出口管制法》为上位法，配套《两用物项出口管制条例》形成基础制度；商务部牵头统筹审批，海关总署、科技部分环节参与核查，形成跨部门协同监管机制；
    \item \textbf{采用清单驱动判定模式，仅目录列明技术触发许可约束。}AI相关软硬件、模型、算法不实施全域一刀切管制，仅《禁止出口限制出口技术目录》及两用物项清单内明确标注的技术，才需要履行出口许可流程；
    \item \textbf{管制体系具备动态调整能力，适配AI技术迭代节奏。}限制出口技术目录分别于2023年、2025年完成修订调整，可根据AI产业发展、安全形势动态增减受控技术条目，同步衔接军民两用技术统一管理规则。
\end{itemize}


% 依赖项（必须放在主文档导言区，即 \begin{document} 之前）：
% \usepackage{tabularray}
% \UseTblrLibrary{varwidth}
% \usepackage{enumitem}
\begingroup
% \nopagebreak[4]
\enlargethispage{1cm}
\footnotesize
\centering

\DefTblrTemplate{conthead-text}{default}{（续表）}
\DefTblrTemplate{contfoot-text}{default}{续下页}

\begin{longtblr}[
  caption = {我国AI智能体领域出口管制政策汇总},
  label   = {tab:china_ai_export_control},
]{
  width   = \linewidth,
  colspec = {
    |Q[wd=0.8cm,l,m]
    |Q[wd=2.5cm,l,m]
    |Q[wd=2.5cm,l,m]
    |Q[wd=2cm,l,m]
    |X[1.00,l,m]|
  },
  cells   = {l,m},
  colsep  = 3pt,
  row{1}  = {font=\bfseries},
  rowhead = 1,
  rowsep  = 5pt,
  measure = vbox,
  stretch = -1,
}
\hline
国家 & 主管部门 & 政策法规/管制工具 & 发布/生效日期 & 核心内容 \\
\hline
\SetCell[r=2]{l,m} 中国
& \begin{itemize}[leftmargin=*,topsep=0pt,partopsep=0pt,itemsep=1pt,parsep=0pt]
    \item 商务部（牵头）
    \item 海关总署（深度参与）
\end{itemize}
& 《出口管制法》《两用物项出口管制条例》 ~\cite{中国出口管制法,两用物项出口管制条例}
& 2020年12月1日起施行；2024年12月1日起施行
& \begin{itemize}[leftmargin=*,topsep=0pt,partopsep=0pt,itemsep=1pt,parsep=0pt]
    \item 建立管制清单、临时管制、许可审批、最终用户和最终用途管理等出口管制基本制度；
    \item 两用物项出口需按清单、技术参数、目的地、最终用户、最终用途和转出口风险判断许可义务；
    \item AI相关硬件、软件、技术或模型只有在落入清单、临时管制或相关最终用途规则时，才纳入出口管制。
\end{itemize} \\
\cline{2-5}
& \begin{itemize}[leftmargin=*,topsep=0pt,partopsep=0pt,itemsep=1pt,parsep=0pt]
    \item 商务部（牵头）
    \item 科技部（深度参与）
\end{itemize}
& 《中国禁止出口限制出口技术目录》~\cite{中国禁止出口限制出口技术目录}
& 2023.12发布，2025.07调整
& \begin{itemize}[leftmargin=*,topsep=0pt,partopsep=0pt,itemsep=1pt,parsep=0pt]
    \item 将人工智能交互界面、语音合成、个性化信息推送等信息处理技术纳入限制出口技术范围；
    \item 属于限制出口目录的技术，应当依照技术出口许可程序办理相关手续；
    \item 目录所列技术同时属于军民两用技术的，纳入两用物项出口管制体系管理。
\end{itemize} \\
\hline
\end{longtblr}
\endgroup

% 依赖项（必须放在主文档导言区，即 \begin{document} 之前）：
% \usepackage{tabularray}
% \UseTblrLibrary{varwidth}
% \usepackage{enumitem}
\begingroup
% \nopagebreak[4]
\enlargethispage{1cm}
\footnotesize
\centering

\DefTblrTemplate{conthead-text}{default}{（续表）}
\DefTblrTemplate{contfoot-text}{default}{续下页}

\begin{longtblr}[
  caption = {我国AI智能体领域出口管制政策汇总},
  label   = {tab:china_ai_export_control},
]{
  width   = \linewidth,
  colspec = {
    |Q[wd=0.8cm,l,m]
    |Q[wd=2.5cm,l,m]
    |Q[wd=2.5cm,l,m]
    |Q[wd=2cm,l,m]
    |X[1.00,l,m]|
  },
  cells   = {l,m},
  colsep  = 3pt,
  row{1}  = {font=\bfseries},
  rowhead = 1,
  rowsep  = 5pt,
  measure = vbox,
  stretch = -1,
}
\hline
国家 & 主管部门 & 政策法规/管制工具 & 发布/生效日期 & 核心内容 \\
\hline
\SetCell[r=2]{l,m} 中国
& \begin{itemize}[leftmargin=*,topsep=0pt,partopsep=0pt,itemsep=1pt,parsep=0pt]
    \item 商务部（牵头）
    \item 海关总署（深度参与）
\end{itemize}
& 《出口管制法》《两用物项出口管制条例》 ~\cite{中国出口管制法,两用物项出口管制条例}
& 2020年12月1日起施行；2024年12月1日起施行
& \begin{itemize}[leftmargin=*,topsep=0pt,partopsep=0pt,itemsep=1pt,parsep=0pt]
    \item 建立管制清单、临时管制、许可审批、最终用户和最终用途管理等出口管制基本制度；
    \item 两用物项出口需按清单、技术参数、目的地、最终用户、最终用途和转出口风险判断许可义务；
    \item AI相关硬件、软件、技术或模型只有在落入清单、临时管制或相关最终用途规则时，才纳入出口管制。
\end{itemize} \\
\cline{2-5}
& \begin{itemize}[leftmargin=*,topsep=0pt,partopsep=0pt,itemsep=1pt,parsep=0pt]
    \item 商务部（牵头）
    \item 科技部（深度参与）
\end{itemize}
& 《中国禁止出口限制出口技术目录》~\cite{中国禁止出口限制出口技术目录}
& 2023.12发布，2025.07调整
& \begin{itemize}[leftmargin=*,topsep=0pt,partopsep=0pt,itemsep=1pt,parsep=0pt]
    \item 将人工智能交互界面、语音合成、个性化信息推送等信息处理技术纳入限制出口技术范围；
    \item 属于限制出口目录的技术，应当依照技术出口许可程序办理相关手续；
    \item 目录所列技术同时属于军民两用技术的，纳入两用物项出口管制体系管理。
\end{itemize} \\
\hline
\end{longtblr}
\endgroup


\begin{center}
% \nopagebreak[4]
\enlargethispage{1cm}
\footnotesize
\setlength{\tabcolsep}{3pt}
% \renewcommand{\arraystretch}{0.7}
\begin{longtable}{|>{\rule{0pt}{2pt}}m{0.8cm}
|>{\rule{0pt}{2pt}}m{2.5cm}
|>{\rule{0pt}{2pt}}m{2.5cm}
|>{\rule{0pt}{2pt}}m{2cm}
|>{\rule{0pt}{2pt}}m{\dimexpr\linewidth - 0.8cm - 2.5cm - 2.5cm - 2cm - 10\tabcolsep\relax}|}
\caption{我国AI智能体领域出口管制政策汇总}\label{tab:china_ai_export_control}\\
\hline
\textbf{国家} & \textbf{主管部门} & \textbf{政策法规/管制工具} & \textbf{发布/生效日期} & \textbf{核心内容} \\
\hline
\endfirsthead

\multicolumn{5}{l}{\tablename\ \thetable\ -- 续上页}\\
\hline
\textbf{国家} & \textbf{主管部门} & \textbf{政策法规/管制工具} & \textbf{发布/生效日期} & \textbf{核心内容} \\
\hline
\endhead
\hline
\multicolumn{5}{r}{续下页}\\
\endfoot
\hline
\endlastfoot

\multirow{2}{=}{中国}
& \begin{itemize}[tabitem]
    \item 商务部（牵头）
    \item 海关总署（深度参与）
\end{itemize}
& 《出口管制法》《两用物项出口管制条例》 ~\cite{中国出口管制法,两用物项出口管制条例}
& 2020年12月1日起施行；2024年12月1日起施行
& \begin{itemize}[tabitem]
    \item 建立管制清单、临时管制、许可审批、最终用户和最终用途管理等出口管制基本制度；
    \item 两用物项出口需按清单、技术参数、目的地、最终用户、最终用途和转出口风险判断许可义务；
    \item AI相关硬件、软件、技术或模型只有在落入清单、临时管制或相关最终用途规则时，才纳入出口管制。
\end{itemize} \\
\cline{2-5}
& \begin{itemize}[tabitem]
    \item 商务部（牵头）
    \item 科技部（深度参与）
\end{itemize}
& 《中国禁止出口限制出口技术目录》~\cite{中国禁止出口限制出口技术目录}
& 2023.12发布，2025.07调整
& \begin{itemize}[tabitem]
    \item 将人工智能交互界面、语音合成、个性化信息推送等信息处理技术纳入限制出口技术范围；
    \item 属于限制出口目录的技术，应当依照技术出口许可程序办理相关手续；
    \item 目录所列技术同时属于军民两用技术的，纳入两用物项出口管制体系管理。
\end{itemize} 
\end{longtable}
\end{center}


% ==========================================
\section{全球AI硬件与芯片底座安全研究现状}
相较于软件层风险可通过策略迭代、护栏优化快速闭环，硬件与芯片层风险具备强旁路性、底层不可见性与全局穿透性，侧信道信息泄露、多租户算力隔离失效、固件隐性后门、硬件木马与机密计算边界破损等缺陷，可直接旁路上层全部软件防护机制，从根源破坏智能体数据机密性、推理完整性与运行可信性，形成“软件防护有效、底层攻破即失效”的层级安全坍塌效应。
与此同时，随着多租户公有云Agent、私有化企业级智能体、端边协同嵌入式Agent规模化落地，安全信任锚点正从软件策略向硬件信任根下沉，硬件底座可信性已成为约束智能体越权行为、固化隐私数据边界、支撑合规举证的核心底层依托。
%
% 全球各区域基于自身芯片产业禀赋、机密计算技术路线、监管合规诉求与算力主权战略，形成了差异化的AI硬件与芯片安全发展范式：北美依托通用GPU生态主导底层攻防与机密计算原生创新，欧洲结合工业场景与监管体系深耕形式化验证与边缘硬件可信治理，中国立足信创自主可控推进国产算力栈适配与合规型安全补强，三者技术取向、驱动逻辑与攻防能力形成显著分野，构成全球AI底层安全的三区竞争与协作格局。
%
本节聚焦分析北美、欧洲与我国AI加速芯片、机密计算架构、固件驱动治理、芯片供应链、边缘可信底座的研究现状、产学研分工与产业落地形态。

\subsection{北美AI硬件与芯片安全研究现状}
北美作为全球AI算力硬件、芯片架构与机密计算技术的核心策源地，汇聚了顶尖高校安全团队、头部芯片厂商、云服务商与权威安全响应机构。其研究范畴覆盖端侧浏览器GPU微架构、云端多GPU算力集群、芯片设计供应链、端云协同机密推理四大场景，打通了从硬件侧信道机理分析到商业化机密计算产品落地的全链条。为系统梳理北美在AI硬件与芯片领域的安全研究脉络、主体分工、技术方案与落地成效，本节从学术侧与产业侧双维度汇总结构化成果，覆盖浏览器GPU侧信道、多租户显存隔离、硬件供应链、云端机密加速、驱动漏洞治理、端云私有推理六大核心方向，明确各主体的研究贡献、技术路径、应用价值与现存局限, 主要发现如下（详见表~\ref{tab:north-america-ai-hardware-security-nicematrix}）：

\begin{itemize}[tabitem]
\item \textbf{威胁全域向下穿透，软硬件防护存在不可逆层级差。}AI智能体攻击面从软件层延伸至硬件微架构，浏览器GPU侧信道、云端显存残留、芯片硬件木马均可旁路上层软件护栏，证明仅靠应用层与模型层防护无法实现全栈可信，软硬件协同防护是底层刚需。
\item \textbf{产学研闭环固化，形成标准化漏洞治理链路。}高校主导未知底层漏洞机理挖掘，NVIDIA/Intel等产业侧实现机密计算产品化，CERT/CC统筹漏洞披露与固件迭代，构建“挖掘—验证—披露—修复—迭代”的成熟闭环，主导全球硬件安全攻防范式。
\item \textbf{形成三元约束与技术垄断，构筑底层竞争壁垒。}硬件防护普遍存在安全强度、算力性能、生态开放性的固有权衡；北美垄断GPU架构、机密指令集、WebGPU标准与EDA工具链，在底层安全领域形成技术代差与话语权壁垒~\cite{darpaAISS}。

\end{itemize}

\begin{center}
% \nopagebreak[4]
\enlargethispage{1cm}
\footnotesize
\setlength{\tabcolsep}{3pt}

% \begin{longtblr}[
%   caption = {北美AI硬件与芯片安全研究汇总},
%   label   = {tab:north-america-ai-hardware-security-nicematrix},
% ]
% {
%   colspec = {
%     |m{0.8cm}|
%      m{1.2cm}|
%      m{1.2cm}|
%      m{0.22\linewidth}|
%      m{0.17\linewidth}|
%      m{0.17\linewidth}|
%      m{\dimexpr
%        \linewidth
%        - 0.8cm
%        - 1.2cm
%        - 1.2cm
%        - 0.22\linewidth
%        - 0.17\linewidth*2
%        - 14\tabcolsep
%        \relax}|
%   },
%   rowhead = 1,
%   cells   = {valign=m},
% }

\begin{longtblr}[
  caption = {北美AI硬件与芯片安全研究汇总},
  label   = {tab:north-america-ai-hardware-security-nicematrix},
]{
  width   = \linewidth,
  colspec = {
    |Q[wd=0.8cm,l,m]
    |Q[wd=1.2cm,l,m]
    |Q[wd=1.2cm,l,m]
    |X[1.30,l,m]
    |X[1.00,l,m]
    |X[1.00,l,m]
    |X[1.15,l,m]|
  },
  cells   = {l,m},
  colsep  = 3pt,
  row{1}  = {font=\bfseries},
  rowhead = 1,
  rowsep  = 5pt,
  measure = vbox,
  stretch = -1,
}

% \begin{longtblr}[
%   caption = {典型AI智能体安全事件汇总},
%   label   = {tab:agent_security_incident},
% ]{
%   width   = \linewidth,
%   colspec = {
%     |Q[wd=1.5cm,l,m]
%     |Q[wd=2cm,l,m]
%     |Q[wd=1.5cm,l,m]
%     |X[1.15,l,m]
%     |X[1,l,m]|
%   },
%   cells   = {l,m},
%   colsep  = 3pt,
%   row{1}  = {font=\bfseries},
%   rowhead = 1,
%   % 每行内容与上下边框之间各保留少量空间；由该行最高单元格决定整行高度。
%   rowsep  = 5pt,
%   measure = vbox,
%   stretch = -1,
% }
\hline
\textbf{核心维度}
& \textbf{细分落地}
& \textbf{核心主体}
& \textbf{代表性成果}
& \textbf{技术方案}
& \textbf{核心价值}
& \textbf{不足} \\
\hline

\centering 学术侧
& GPU渲染、浏览器侧信道
&
UT- \allowbreak Austin、CMU、UW、GT、UM、RUB、UCR 等GPU侧信道研究团队
&
\begin{itemize}[tabitem]
\item GPU.zip证实软件透明的图形压缩可造成跨源像素信息泄露。
\item Hot Pixels利用设备功耗、温度等混合信号实现浏览器像素窃取与历史嗅探。
\item WebGPU-SPY 证实 WebGPU 沙箱存在缓存侧信道，可用于用户指纹识别。
\end{itemize}
&
\begin{itemize}[tabitem]
\item 依托GPU渲染延迟、像素布局差异构造跨源观测通道；
\item 基于DVFS、温度与功耗传感器反向推断前端图形内容；
\item 利用WebGPU高并行特性与缓存争用构建浏览器侧信道。
\end{itemize}
&
\begin{itemize}[tabitem]
\item 证明浏览器智能体存在底层GPU硬件攻击面，不止应用层风险；
\item 推动浏览器、GPU驱动与图形API的隔离与精度安全优化；
\item 为端侧多模态智能体的底层硬件安全威胁建模提供依据。
\end{itemize}
&
\begin{itemize}[tabitem]
\item 攻击有效性依赖GPU型号、驱动版本与浏览器环境；
\item 安全缓解机制易牺牲图形性能与WebGPU可用性；
\item 现有研究缺乏企业生产环境与推理容器的真实场景验证。
\end{itemize} \\
\hline

\centering 学术侧
& GPU显存与多租户隔离
&
Trail of Bits、UCSC、UCR、PNNL等GPU内存安全研究团队
&
\begin{itemize}[tabitem]
\item LeftoverLocals 证实GPU局部内存清理不彻底，存在跨容器、跨用户数据泄露风险，影响AMD、Apple、Qualcomm 等主流厂商。
\item Spy in the GPU-box 发现多 GPU 集群可通过 L2 缓存争用构建隐蔽侧信道。
\item NVBleed、Beyond the Bridge 进一步证实 NVLink、PCIe 互连链路存在多租户泄露风险~\cite{nvbleed2025}。
\end{itemize}
&
\begin{itemize}[tabitem]
\item 读取GPU内核残留内存数据，恢复历史推理中间信息。
\item 通过缓存争用、性能计数器观测远端GPU负载状态。
\item 基于互连链路拥塞特征实现多租户应用指纹与数据窃取。
\end{itemize}
&
\begin{itemize}[tabitem]
\item 明确大模型推理中Prompt、张量、模型参数的GPU内存泄露威胁；
\item 证实共享GPU算力存在独立安全风险，需专属隔离策略；
\item 为云GPU调度、容器隔离与内存清零机制提供理论支撑。
\end{itemize}
&
\begin{itemize}[tabitem]
\item 攻击复现高度依赖设备、驱动与调度策略，通用性有限；
\item 严格内存清零与并发限制会降低GPU算力利用率；
\item 云厂商多租户底层策略未完全公开，第三方验证难度大。
\end{itemize} \\
\hline

\centering 学术侧
& 硬件木马与供应链安全
& UConn、Rice、UF、DARPA AISS及安全EDA/IP研究团队
&
\begin{itemize}[tabitem]
\item 硬件木马体系研究将威胁覆盖至IP复用、设计外包与测试缺陷等全供应链环节 ~\cite{tehranipoor2010trojan}。
\item 基于信息流与形式化验证方法，可在RTL/IP阶段检测恶意硬件逻辑。
\item DARPA AISS 项目将侧信道、逆向工程与硬件木马检测融入芯片自动化设计流程~\cite{darpaAISS}。
\end{itemize}
&
\begin{itemize}[tabitem]
\item 构建硬件木马触发机制、载荷特征与检测技术体系；
\item 在设计阶段通过形式化验证筛查违规硬件逻辑；
\item 将安全评估、IP校验与风险权衡融入SoC设计流程。
\end{itemize}
&
\begin{itemize}[tabitem]
\item 揭示AI加速芯片底层供应链可信性对智能体安全的核心影响；
\item 将硬件安全防护从事后修补前移至设计验证阶段；
\item 为AI加速器、NPU、定制SoC的安全认证提供技术框架。
\end{itemize}
&
\begin{itemize}[tabitem]
\item 现有研究与大模型专属加速硬件的适配性不足。
\item 大规模SoC的形式化验证成本高、覆盖范围有限。
\item 供应链安全依赖多主体协同，落地流程复杂、成本高。
\end{itemize} \\
\hline

\centering 产业侧
& CPU- \allowbreak GPU机密协同计算
&
NVIDIA、AMD、Intel等云端机密计算研发平台
&
\begin{itemize}[tabitem]
\item NVIDIA H100构建GPU硬件可信根，实现固件、设备证明与CPU-GPU加密传输。
\item AMD SEV-SNP、Intel TDX为机密虚拟机提供内存加密与远程证明，可与GPU机密计算协同适配 ~\cite{amdsev2026,inteltdx2026}。
\item NVIDIA Secure AI依托硬件证明链路，实现模型与敏感数据的可控释放~\cite{nvidiasecureai2025}。
\end{itemize}
&
\begin{itemize}[tabitem]
\item 基于硬件可信根与安全启动，构建GPU全链路可信认证体系；
\item 通过机密虚拟机保护主机内存完整性，联动GPU硬件可信通道；
\item 基于设备固件与镜像校验，实现AI工作流的可信准入控制。
\end{itemize}
&
\begin{itemize}[tabitem]
\item 保障大模型训练、推理过程中模型权重与业务数据的使用态机密性；
\item 支撑金融、政务等高敏感场景智能体在混合云环境安全部署；
\item 将底层硬件安全能力封装为可落地的企业AI安全服务。
\end{itemize}
&
\begin{itemize}[tabitem]
\item 可信链路跨软硬件多层级，可信计算基复杂、运维难度大。
\item 加密传输与远程证明会带来一定性能损耗与部署成本。
\item 核心校验逻辑依赖厂商闭源固件，透明度与可控性有限。
\end{itemize} \\
\hline

\centering 产业侧
& GPU漏洞治理与修复
&
Trail of Bits、CERT/CC、AMD、Apple、Google、Qualcomm 安全团队
&
\begin{itemize}[tabitem]
\item CERT/CC公开GPU局部内存泄露漏洞，统筹多厂商协同修复。
\item AMD发布安全公告，明确内存泄露漏洞编号、成因与固件修复方案。
\item 主流厂商针对不同设备与驱动版本迭代漏洞缓解与安全加固策略。
\end{itemize}
&
\begin{itemize}[tabitem]
\item 通过固件与驱动更新清理残留内存，限制非法并发访问；
\item 依托 CVE 与厂商公告体系建立标准化漏洞响应机制；
\item 对图形与计算接口进行版本化安全治理与权限管控。
\end{itemize}
&
\begin{itemize}[tabitem]
\item 实现学术漏洞研究向工业级安全修复的快速落地；
\item 完善 GPU 内存、寄存器与内核调度的安全基线规范；
\item 推动厂商将大模型多租户场景纳入常规安全评估范围。
\end{itemize}
&
\begin{itemize}[tabitem]
\item 漏洞修复链条长，终端用户升级迭代滞后性明显。
\item 安全隔离模式会限制GPU并发算力，降低资源利用率。
\item 老旧设备与闭源驱动长期安全维护无法保障。
\end{itemize} \\
\hline

\centering 产业侧
& 端云协同可信推理
& Apple PCC、Apple Silicon安全生态团队
&
\begin{itemize}[tabitem]
\item Apple PCC依托自研芯片与安全固件，构建可信云端推理节点。
\item PCC架构限制远程权限、日志留存与后台访问，杜绝数据非法留存。
\item 通过透明日志与设备证明机制，实现云端推理节点可验证、可追溯。
\end{itemize}
&
\begin{itemize}[tabitem]
\item 端侧校验云端节点证书与系统度量值，加密传输推理请求；
\item 依托安全启动与密钥隔离，最小化云端攻击面；
\item 通过权限受限架构与数据即时清除，降低运维滥用风险。
\end{itemize}
&
\begin{itemize}[tabitem]
\item 构建端云一体的隐私推理底座，适配个人智能体隐私场景；
\item 以硬件强制约束替代云端服务合规承诺，提升可信度；
\item 为隐私敏感型移动端智能体提供标准化落地范式。
\end{itemize}
&
\begin{itemize}[tabitem]
\item 高度依赖Apple封闭生态，跨平台适配性差、通用性弱；
\item 核心硬件与固件逻辑闭源，第三方安全审计能力有限；
\item 仅解决推理隐私问题，无法防护提示注入、越权调用等应用层风险。
\end{itemize} \\
\hline

\end{longtblr}
\end{center}

% \begin{center}
% % \nopagebreak[4]
% \enlargethispage{1cm}
% \footnotesize
% \setlength{\tabcolsep}{3pt}
% \begin{longtable}{
% |>{\rule{0pt}{2pt}}m{0.8cm}|
% >{\rule{0pt}{2pt}}m{1.2cm}|
% >{\rule{0pt}{2pt}}m{1.2cm}|
% >{\rule{0pt}{2pt}}m{0.22\linewidth}|
% >{\rule{0pt}{2pt}}m{0.17\linewidth}|
% >{\rule{0pt}{2pt}}m{0.17\linewidth}|
% >{\rule{0pt}{2pt}}m{\dimexpr\linewidth - 0.8cm - 1.2cm - 1.2cm - 0.22\linewidth - 0.17\linewidth*2 - 14\tabcolsep\relax}|
% }
% \caption{北美AI硬件与芯片安全研究汇总}\label{tab:north-america-ai-hardware-security-nicematrix}\\
% \hline
% \textbf{核心维度} & \textbf{细分落地} & \textbf{核心主体} & \textbf{代表性成果} & \textbf{技术方案} & \textbf{核心价值} & \textbf{不足} \\
% \hline
% \endfirsthead
% \multicolumn{7}{l}{\tablename\ \thetable\ -- 续上页}\\
% \hline
% \textbf{核心维度} & \textbf{细分落地} & \textbf{核心主体} & \textbf{代表性成果} & \textbf{技术方案} & \textbf{核心价值} & \textbf{不足} \\
% \hline
% \endhead
% \hline
% \multicolumn{7}{r}{续下页}\\
% \endfoot
% \hline
% \endlastfoot

% \centering 学术侧
% & GPU渲染、浏览器侧信道
% & 
% UT‑ \allowbreak Austin、CMU、UW、GT、UM、RUB、UCR 等GPU侧信道研究团队
% & \begin{itemize}[tabitem]
% \item GPU.zip证实软件透明的图形压缩可造成跨源像素信息泄露。
% \item Hot Pixels利用设备功耗、温度等混合信号实现浏览器像素窃取与历史嗅探。
% \item WebGPU‑SPY 证实 WebGPU 沙箱存在缓存侧信道，可用于用户指纹识别。
% \end{itemize}
% & \begin{itemize}[tabitem]
% \item 依托GPU渲染延迟、像素布局差异构造跨源观测通道；
% \item 基于DVFS、温度与功耗传感器反向推断前端图形内容；
% \item 利用WebGPU高并行特性与缓存争用构建浏览器侧信道。
% \end{itemize}
% & \begin{itemize}[tabitem]
% \item 证明浏览器智能体存在底层GPU硬件攻击面，不止应用层风险；
% \item 推动浏览器、GPU驱动与图形API的隔离与精度安全优化；
% \item 为端侧多模态智能体的底层硬件安全威胁建模提供依据。
% \end{itemize}
% & \begin{itemize}[tabitem]
% \item 攻击有效性依赖GPU型号、驱动版本与浏览器环境；
% \item 安全缓解机制易牺牲图形性能与WebGPU可用性；
% \item 现有研究缺乏企业生产环境与推理容器的真实场景验证。
% \end{itemize} \\
% \hline
% \centering 学术侧
% &GPU显存与多租户隔离
% & 
% Trail of Bits、UCSC、UCR、PNNL等GPU内存安全研究团队
% & \begin{itemize}[tabitem]
% \item LeftoverLocals 证实GPU局部内存清理不彻底，存在跨容器、跨用户数据泄露风险，影响AMD、Apple、Qualcomm 等主流厂商。
% \item Spy in the GPU‑box 发现多 GPU 集群可通过 L2 缓存争用构建隐蔽侧信道。
% \item NVBleed、Beyond the Bridge 进一步证实 NVLink、PCIe 互连链路存在多租户泄露风险~\cite{nvbleed2025}。
% \end{itemize}
% & \begin{itemize}[tabitem]
% \item 读取GPU内核残留内存数据，恢复历史推理中间信息。
% \item 通过缓存争用、性能计数器观测远端GPU负载状态。
% \item 基于互连链路拥塞特征实现多租户应用指纹与数据窃取。
% \end{itemize}
% & \begin{itemize}[tabitem]
% \item 明确大模型推理中Prompt、张量、模型参数的GPU内存泄露威胁；
% \item 证实共享GPU算力存在独立安全风险，需专属隔离策略；
% \item 为云GPU调度、容器隔离与内存清零机制提供理论支撑。
% \end{itemize}
% & \begin{itemize}[tabitem]
% \item 攻击复现高度依赖设备、驱动与调度策略，通用性有限；
% \item 严格内存清零与并发限制会降低GPU算力利用率；
% \item 云厂商多租户底层策略未完全公开，第三方验证难度大。
% \end{itemize} \\
% \hline
% \centering 学术侧
% & 硬件木马与供应链安全
% & UConn、Rice、UF、DARPA AISS及安全EDA/IP研究团队
% & \begin{itemize}[tabitem]
% \item 硬件木马体系研究将威胁覆盖至IP复用、设计外包与测试缺陷等全供应链环节 ~\cite{tehranipoor2010trojan}。
% \item 基于信息流与形式化验证方法，可在RTL/IP阶段检测恶意硬件逻辑。
% \item DARPA AISS 项目将侧信道、逆向工程与硬件木马检测融入芯片自动化设计流程~\cite{darpaAISS}。
% \end{itemize}
% & \begin{itemize}[tabitem]
% \item 构建硬件木马触发机制、载荷特征与检测技术体系；
% \item 在设计阶段通过形式化验证筛查违规硬件逻辑；
% \item 将安全评估、IP校验与风险权衡融入SoC设计流程。
% \end{itemize}
% & \begin{itemize}[tabitem]
% \item 揭示AI加速芯片底层供应链可信性对智能体安全的核心影响；
% \item 将硬件安全防护从事后修补前移至设计验证阶段；
% \item 为AI加速器、NPU、定制SoC的安全认证提供技术框架。
% \end{itemize}
% & \begin{itemize}[tabitem]
% \item 现有研究与大模型专属加速硬件的适配性不足。
% \item 大规模SoC的形式化验证成本高、覆盖范围有限。
% \item 供应链安全依赖多主体协同，落地流程复杂、成本高。
% \end{itemize} \\
% \hline
% \centering 产业侧
% & CPU‑ \allowbreak GPU机密协同计算
% & 
% NVIDIA、AMD、Intel等云端机密计算研发平台
% & \begin{itemize}[tabitem]
% \item NVIDIA H100构建GPU硬件可信根，实现固件、设备证明与CPU‑GPU加密传输。
% \item AMD SEV‑SNP、Intel TDX为机密虚拟机提供内存加密与远程证明，可与GPU机密计算协同适配 ~\cite{amdsev2026,inteltdx2026}。
% \item NVIDIA Secure AI依托硬件证明链路，实现模型与敏感数据的可控释放~\cite{nvidiasecureai2025}。
% \end{itemize}
% & \begin{itemize}[tabitem]
% \item 基于硬件可信根与安全启动，构建GPU全链路可信认证体系；
% \item 通过机密虚拟机保护主机内存完整性，联动GPU硬件可信通道；
% \item 基于设备固件与镜像校验，实现AI工作流的可信准入控制。
% \end{itemize}
% & \begin{itemize}[tabitem]
% \item 保障大模型训练、推理过程中模型权重与业务数据的使用态机密性；
% \item 支撑金融、政务等高敏感场景智能体在混合云环境安全部署；
% \item 将底层硬件安全能力封装为可落地的企业AI安全服务。
% \end{itemize}
% & \begin{itemize}[tabitem]
% \item 可信链路跨软硬件多层级，可信计算基复杂、运维难度大。
% \item 加密传输与远程证明会带来一定性能损耗与部署成本。
% \item 核心校验逻辑依赖厂商闭源固件，透明度与可控性有限。
% \end{itemize} \\
% \hline
% \centering 产业侧
% &GPU漏洞治理与修复
% & 
% Trail of Bits、CERT/CC、AMD、Apple、Google、Qualcomm 安全团队
% & \begin{itemize}[tabitem]
% \item CERT/CC公开GPU局部内存泄露漏洞，统筹多厂商协同修复。
% \item AMD发布安全公告，明确内存泄露漏洞编号、成因与固件修复方案。
% \item 主流厂商针对不同设备与驱动版本迭代漏洞缓解与安全加固策略。
% \end{itemize}
% & \begin{itemize}[tabitem]
% \item 通过固件与驱动更新清理残留内存，限制非法并发访问；
% \item 依托 CVE 与厂商公告体系建立标准化漏洞响应机制；
% \item 对图形与计算接口进行版本化安全治理与权限管控。
% \end{itemize}
% & \begin{itemize}[tabitem]
% \item 实现学术漏洞研究向工业级安全修复的快速落地；
% \item 完善 GPU 内存、寄存器与内核调度的安全基线规范；
% \item 推动厂商将大模型多租户场景纳入常规安全评估范围。
% \end{itemize}
% & \begin{itemize}[tabitem]
% \item 漏洞修复链条长，终端用户升级迭代滞后性明显。
% \item 安全隔离模式会限制GPU并发算力，降低资源利用率。
% \item 老旧设备与闭源驱动长期安全维护无法保障。
% \end{itemize} \\
% \hline

% \centering 产业侧
% & 端云协同可信推理
% & Apple PCC、Apple Silicon安全生态团队
% & \begin{itemize}[tabitem]
% \item Apple PCC依托自研芯片与安全固件，构建可信云端推理节点。
% \item PCC架构限制远程权限、日志留存与后台访问，杜绝数据非法留存。
% \item 通过透明日志与设备证明机制，实现云端推理节点可验证、可追溯。
% \end{itemize}
% & \begin{itemize}[tabitem]
% \item 端侧校验云端节点证书与系统度量值，加密传输推理请求；
% \item 依托安全启动与密钥隔离，最小化云端攻击面；
% \item 通过权限受限架构与数据即时清除，降低运维滥用风险。
% \end{itemize}
% & \begin{itemize}[tabitem]
% \item 构建端云一体的隐私推理底座，适配个人智能体隐私场景；
% \item 以硬件强制约束替代云端服务合规承诺，提升可信度；
% \item 为隐私敏感型移动端智能体提供标准化落地范式。
% \end{itemize}
% & \begin{itemize}[tabitem]
% \item 高度依赖Apple封闭生态，跨平台适配性差、通用性弱；
% \item 核心硬件与固件逻辑闭源，第三方安全审计能力有限；
% \item 仅解决推理隐私问题，无法防护提示注入、越权调用等应用层风险。
% \end{itemize} 
% \end{longtable}
% \end{center}




\subsection{欧洲AI硬件与芯片安全研究现状}
% 欧洲学术群体以苏黎世联邦理工、柏林工业、CISPA、鲁汶大学、剑桥、帝国理工等系统安全与形式化团队为核心，长期聚焦CPU级机密虚拟化（SEV‑SNP）的信任根解构、异常路径攻击与协议形式化验证，率先揭示了机密虚拟机（CVM）在恶意Hypervisor、固件缺陷、虚拟化异常路径下的固有完整性风险，完善了云端AI机密底座的威胁模型~\cite{severed2018, wesee2024}。
不同于北美侧重GPU、总线与微架构漏洞挖掘，欧洲AI硬件研究以\textbf{可证明安全、可审计证明链、合规可复验}为核心导向，将硬件安全结论转化为可量化、可纳入合规文档的论证依据，适配欧盟AI Act、NIS2、CRA等监管文书的举证要求。
%
% 为系统厘清欧洲在AI芯片、机密计算、边缘硬件、工业可信底座领域的研究脉络、主体分工、技术方案与合规映射关系，
本节梳理欧洲AI硬件代表性成果、技术机理、产业价值与现存局限，形成与北美硬件安全章节对等的实证基础，主要发现如下（详见表~\ref{tab:europe-ai-hardware-security-nicematrix}）：

\begin{itemize}[tabitem]
\item \textbf{监管内生驱动，形成合规逆向传导架构。}以EU AI Act、NIS2、CRA为顶层约束，将芯片供应链韧性、硬件举证义务强制嵌入技术设计，硬件安全能力全部服务于审计归档，形成“合规→架构→技术”的逆向发展路径，适配高监管行业场景。
\item \textbf{学术范式差异化，深耕可证明安全体系。}以形式化验证为核心标签，解构SEV‑SNP等机密计算信任根，构建标准化合规论证链路，将漏洞发现升级为TCB完整性数学论证，精准匹配欧盟高风险AI举证要求。
\item \textbf{错位竞争布局，构筑工业边缘专属壁垒。}避开高端通用GPU赛道，依托EuroHPC打造主权算力池，基于工业安全芯片构建边缘信任根，在工业智能体、数据空间可信连接领域形成全球独有技术壁垒~\cite{eurohpcAIFactories2026}。
\end{itemize}


\begingroup

\begin{center}
% \nopagebreak[4]
\enlargethispage{1cm}
\footnotesize

% \begin{longtblr}[
%   caption = {欧洲AI硬件与芯片安全研究汇总},
%   label   = {tab:europe-ai-hardware-security-nicematrix},
% ]
% {
%   colsep = 3pt,
%   colspec = {
%     |m{0.8cm}
%     |m{1.2cm}
%     |m{1.5cm}
%     |m{0.22\linewidth}
%     |m{0.16\linewidth}
%     |m{0.16\linewidth}
%     |m{\dimexpr
%        \linewidth
%        - 0.8cm
%        - 1.2cm
%        - 1.5cm
%        - 0.22\linewidth
%        - 0.16\linewidth
%        - 0.16\linewidth
%        - 42pt
%        - 8\arrayrulewidth
%        \relax}|
%   },
%   rowhead = 1,
%   cells = {valign=m},
% }

\begin{longtblr}[
  caption = {欧洲AI硬件与芯片安全研究汇总},
  label   = {tab:europe-ai-hardware-security-nicematrix},
]{
  width   = \linewidth,
  colspec = {
    |Q[wd=0.8cm,l,m]
    |Q[wd=1.2cm,l,m]
    |Q[wd=1.5cm,l,m]
    |Q[wd=0.22\linewidth,l,m]
    |Q[wd=0.16\linewidth,l,m]
    |Q[wd=0.16\linewidth,l,m]
    |X[1,l,m]|
  },
  cells   = {l,m},
  colsep  = 3pt,
  row{1}  = {font=\bfseries},
  rowhead = 1,
  rowsep  = 5pt,
  measure = vbox,
  stretch = -1,
}

\hline
\textbf{核心维度}
& \textbf{细分落地}
& \textbf{核心主体}
& \textbf{代表性成果}
& \textbf{技术方案}
& \textbf{核心价值}
& \textbf{不足} \\
\hline

\SetCell{c} 学术侧
& 机密计算与云端AI隔离验证
&
AISEC、TU-Berlin SECT、SIT、ETH-Z云端机密安全研究团队
&
\begin{itemize}[tabitem]
\item SEVered证实恶意Hyper-Visor可窃取AMD SEV加密虚拟机内存明文数据。
\item One-Glitch-to-Rule-Them-All借助电压故障注入破坏SP可信根，实现密钥破解与证明伪造。
\item WeSee利用SEV-SNP异常处理逻辑，验证恶意虚拟化组件依然可以破坏机密虚拟机完整性。
\end{itemize}
&
\begin{itemize}[tabitem]
\item 构造恶意页表映射，绕过TEE防护机制导出虚拟机内存明文信息；
\item 对安全协处理器施加故障注入，突破固件可信根与硬件密钥保护体系；
\item 借助虚拟化异常通道篡改内核执行路径，破坏机密计算运行环境。
\end{itemize}
&
\begin{itemize}[tabitem]
\item 证明CPU-TEE仅为基础防护，无法永久保障AI云平台的可信安全；
\item 将云端AI可信验证范围拓展至固件初始化、异常处理和设备路由层面；
\item 为高风险智能体建立持续漏洞挖掘、安全验证与漏洞修补闭环范式。
\end{itemize}
&
\begin{itemize}[tabitem]
\item 当前研究主要围绕CPU-TEE开展，GPU/NPU机密执行环境相关研究不足。
\item 固件升级、证明策略更新依赖芯片厂商、云服务商多方协同推进。
\item 云平台底层配置未完全公开，第三方机构难以完成复现验证工作。
\end{itemize} \\
\hline

\SetCell{c} 学术侧
& SEV- SNP与远程证明安全研究
& ETH-Z、UZ、AMD- PSIRT、AISEC、TU-Berlin SECT安全研究团队
&
\begin{itemize}[tabitem]
\item AMD-SB-3034披露MMIO路由缺陷与ASP固件漏洞，会破坏机密虚拟机运行可信性~\cite{amdsb3034}。
\item 形式化分析SEV-SNP整套接口，发现平台无关消息设计会削弱远程证明可靠性。
\item SEVered与WeSee系列成果持续对SEV-SNP体系开展回归测试，完善漏洞验证链条。
\end{itemize}
&
\begin{itemize}[tabitem]
\item 核查平台初始化、RMP权限、固件锁定以及TCB版本分发全流程安全隐患；
\item 采用符号模型验证证明、密钥派生、虚拟机迁移协议的安全属性；
\item 将固件版本、TCB值、云端证明策略纳入统一的安全评估体系。
\end{itemize}
&
\begin{itemize}[tabitem]
\item 细化SEV-SNP可信判定条件，将固件配置纳入可信评估关键要素；
\item 为高风险行业采购机密AI服务提供可落地的审计评判标准；
\item 促使智能体平台强化固件版本、证明材料的证据化管控机制。
\end{itemize}
&
\begin{itemize}[tabitem]
\item 证明链层级繁杂，部署人员难以理解TCB字段与固件版本的安全含义。
\item 形式化模型简化硬件细节，无法覆盖固件与主板的全部真实配置。
\item 固件发布节奏由厂商把控，云租户无法自主决定漏洞修复进度。
\end{itemize} \\
\hline

\SetCell{c} 学术侧
& 边缘设备与工业数据空间安全
& AISEC、IDS生态、Arm-PSA、GP、CSEM边缘安全研究团队
&
\begin{itemize}[tabitem]
\item IDS可信网关依托TPM-TEE实现容器隔离、身份认证，保障跨机构工业数据交换安全。
\item PSA-TrustZone将可信根、加密接口抽象为IoT设备通用安全开发框架 ~\cite{psacertified2026,armplatformsecurity2026}。
\item TrustZone运行时完整性检测技术为边缘端TinyML模型提供固件与缓存安全防护 ~\cite{pdrima2025}。
\end{itemize}
&
\begin{itemize}[tabitem]
\item 结合TPM-TEE、容器引擎、策略组件构建边缘网关可信执行环境；
\item 依托安全启动、硬件密钥、安全存储搭建设备底层可信基础组件；
\item 对TEE内核、安全任务和缓存开展动态审计，保障模型运行安全。
\end{itemize}
&
\begin{itemize}[tabitem]
\item 适配工业Agent对接PLC、传感器、运维系统的现场部署场景；
\item 落实欧洲数据空间内数据不出域、节点可证明、策略可审计的目标；
\item 为边缘智能体的设备身份、固件可信提供标准化底层支撑方案。
\end{itemize}
&
\begin{itemize}[tabitem]
\item 老旧工业设备硬件条件受限，缺少TEE和安全加密硬件组件。
\item 不同厂商网关、工业协议之间策略适配难度大，互操作性较差。
\item 现有防护手段难以抵御物理篡改、后期运维绕过和长期补丁缺失问题。
\end{itemize} \\
\hline

\SetCell{c} 产业侧
& 主权算力与合规AI基础设施
& 欧委会、EuroHPC-JU、AI Factories、LUMI、BSC超算中心
&
\begin{itemize}[tabitem]
\item 《欧洲芯片法案》强化本土半导体产业，降低海外供应链依赖程度。
\item EU-AI Act依据风险等级管控高-风险AI，对模型安全与人工监督提出硬性约束。
\item EuroHPC搭建19个AI工厂和13个分支节点，为欧洲用户提供合规超算资源 ~\cite{eurohpcAIFactories2026,ecAIFactories2026}。
\end{itemize}
&
\begin{itemize}[tabitem]
\item 通过立法与财政投入将芯片供应、算力资源与合规监管深度绑定；
\item 依托AI Factories打通超算中心、企业、科研机构一体化服务网络；
\item 在采购流程、访问控制环节嵌入可信AI硬件安全相关评审条件。
\end{itemize}
&
\begin{itemize}[tabitem]
\item 将AI硬件基础设施升级为欧洲产业与公共部门自主可控的战略资源；
\item 为医疗、能源、制造领域高风险智能体部署提供合规落地依据；
\item 通过公共采购反向推动硬件安全证明、审计与长期维护体系建设。
\end{itemize}
&
\begin{itemize}[tabitem]
\item 高端GPU、软件栈以及大模型生态仍然高度依赖海外厂商产品。
\item 合规条款执行成本偏高，中小企业和科研团队落地负担较重。
\item 芯片硬件安全指标和现有法规之间缺少统一标准与测评工具衔接。
\end{itemize} \\
\hline

\SetCell{c} 产业侧
& 嵌入式芯片与工业硬件底座
& Infineon、ST、NXP、Arm-PSA、GP嵌入式安全芯片厂商
&
\begin{itemize}[tabitem]
\item Infineon OPTIGA-Trust芯片为MCU设备提供身份认证与固件安全更新能力~\cite{infineonoptiga2026}。
\item STSAFE-A110安全元件实现密钥隔离、抗侧信道攻击、安全启动与TLS通信~\cite{stsafea1102026}。
\item Arm-PSA体系为IoT芯片划分安全等级，标准化可信根与加密调用接口 ~\cite{psacertified2026,armplatformsecurity2026}。
\end{itemize}
&
\begin{itemize}[tabitem]
\item 通过安全元件与唯一设备证书绑定工业终端，划分运维访问权限边界；
\item 依托固件签名、密钥隔离、安全计数器保护边缘节点模型和业务数据；
\item 结合IEC-62443规范将硬件安全能力转化为产品认证准入条件。
\end{itemize}
&
\begin{itemize}[tabitem]
\item 助力工业智能体建立设备可信身份，规范现场运维权限管控机制；
\item 保护机器人、汽车电子中的AI模型、配置文件与现场业务数据安全；
\item 实现硬件安全能力与欧洲功能安全、网络安全规范协同落地。
\end{itemize}
&
\begin{itemize}[tabitem]
\item 安全功能分散于不同芯片产品、证书体系，配置方式不统一。
\item 端侧硬件算力不足，复杂智能体难以完全部署在本地嵌入式芯片。
\item 硬件安全组件默认关闭，系统集成阶段配置不当会造成安全能力闲置。
\end{itemize} \\
\hline

\SetCell{c} 产业侧
& 工业合规与长期安全运维
& Siemens- \allowbreak PC、Schneider-PSIRT、Bosch- \allowbreak Rexroth、IEC-62443、NIS2/CRA合规生态
&
\begin{itemize}[tabitem]
\item Siemens分层防护方案遵循IEC-62443标准，实现工业资产纵深安全管控 ~\cite{siemensindustrialcyber2026}。
\item 西门子将零信任架构引入工业场景，保护控制器、SCADA和远程运维接口安全 ~\cite{siemensindustrialcyber2026}。
\item 厂商依托固件更新、安全公告体系，满足NIS2、CRA等欧洲合规监管要求 ~\cite{siemensindustrialcyber2026,infineonoptiga2026}。
\end{itemize}
&
\begin{itemize}[tabitem]
\item 将工业资产划分为现场设备、控制网络、云平台多层安全防护区域；
\item 通过资产盘点、访问审计、补丁管理构建零信任工业安全运行体系；
\item 在采购规范中明确安全启动、漏洞披露、生命周期支持硬性条件。
\end{itemize}
&
\begin{itemize}[tabitem]
\item 将智能体安全范围由模型防护拓展至工业设备全生命周期运维管理；
\item 把硬件安全、协议安全、法规合规统一纳入项目验收评价体系；
\item 为能源、交通行业高风险智能体明确厂商、集成商、使用者责任链条。
\end{itemize}
&
\begin{itemize}[tabitem]
\item 工业设备服役周期长，停机成本过高导致漏洞补丁部署滞后。
\item 厂商承诺、系统配置和后期运维脱节，容易出现安全责任划分模糊。
\item 智能体自主决策、跨设备联动场景的安全认证范式仍然不够成熟。
\end{itemize} \\
\hline

\end{longtblr}
\end{center}


% \begin{center}
% % \nopagebreak[4]
% \enlargethispage{1cm}
% \footnotesize
% \setlength{\tabcolsep}{3pt}
% \begin{longtable}{
% |>{\rule{0pt}{2pt}}m{0.8cm}
% |>{\rule{0pt}{2pt}}m{1.2cm}
% |>{\rule{0pt}{2pt}}m{1.5cm}
% |>{\rule{0pt}{2pt}}m{\dimexpr 0.22\linewidth}
% |>{\rule{0pt}{2pt}}m{\dimexpr 0.16\linewidth}
% |>{\rule{0pt}{2pt}}m{\dimexpr 0.16\linewidth}
% |>{\rule{0pt}{2pt}}m{\dimexpr\linewidth - 0.8cm - 1.2cm - 1.5cm - 0.2\linewidth -  0.16\linewidth*2 - 14\tabcolsep\relax}|}
% \caption{欧洲AI硬件与芯片安全研究汇总}\label{tab:europe-ai-hardware-security-nicematrix}\\
% \hline
% \textbf{核心维度} & \textbf{细分落地} & \textbf{核心主体} & \textbf{代表性成果} & \textbf{技术方案} & \textbf{核心价值} & \textbf{不足} \\
% \hline
% \endfirsthead
% \multicolumn{7}{l}{\tablename\ \thetable\ -- 续上页}\\
% \hline
% \textbf{核心维度} & \textbf{细分落地} & \textbf{核心主体} & \textbf{代表性成果} & \textbf{技术方案} & \textbf{核心价值} & \textbf{不足} \\
% \hline
% \endhead
% \hline
% \multicolumn{7}{r}{续下页}\\
% \endfoot
% \hline
% \endlastfoot

% \multirow{3}{=}{\centering 学术侧}
% & \multirow{3}{=}{机密计算与云端AI隔离验证}
% & 
% AISEC、TU‑Berlin SECT、SIT、ETH‑Z云端机密安全研究团队
% & \begin{itemize}[tabitem]
% \item SEVered证实恶意Hyper‑Visor可窃取AMD SEV加密虚拟机内存明文数据。
% \item One‑Glitch‑to‑Rule‑Them‑All借助电压故障注入破坏SP可信根，实现密钥破解与证明伪造。
% \item WeSee利用SEV‑SNP异常处理逻辑，验证恶意虚拟化组件依然可以破坏机密虚拟机完整性。
% \end{itemize}
% & \begin{itemize}[tabitem]
% \item 构造恶意页表映射，绕过TEE防护机制导出虚拟机内存明文信息；
% \item 对安全协处理器施加故障注入，突破固件可信根与硬件密钥保护体系；
% \item 借助虚拟化异常通道篡改内核执行路径，破坏机密计算运行环境。
% \end{itemize}
% & \begin{itemize}[tabitem]
% \item 证明CPU‑TEE仅为基础防护，无法永久保障AI云平台的可信安全；
% \item 将云端AI可信验证范围拓展至固件初始化、异常处理和设备路由层面；
% \item 为高风险智能体建立持续漏洞挖掘、安全验证与漏洞修补闭环范式。
% \end{itemize}
% & \begin{itemize}[tabitem]
% \item 当前研究主要围绕CPU‑TEE开展，GPU/NPU机密执行环境相关研究不足。
% \item 固件升级、证明策略更新依赖芯片厂商、云服务商多方协同推进。
% \item 云平台底层配置未完全公开，第三方机构难以完成复现验证工作。
% \end{itemize} \\
% \hline
% \multirow{3}{=}{\centering 学术侧}
% & \multirow{3}{=}{SEV‑ SNP与远程证明安全研究}
% & ETH‑Z、UZ、AMD‑ PSIRT、AISEC、TU‑Berlin SECT安全研究团队
% & \begin{itemize}[tabitem]
% \item AMD‑SB‑3034披露MMIO路由缺陷与ASP固件漏洞，会破坏机密虚拟机运行可信性~\cite{amdsb3034}。
% \item 形式化分析SEV‑SNP整套接口，发现平台无关消息设计会削弱远程证明可靠性。
% \item SEVered与WeSee系列成果持续对SEV‑SNP体系开展回归测试，完善漏洞验证链条。
% \end{itemize}
% & \begin{itemize}[tabitem]
% \item 核查平台初始化、RMP权限、固件锁定以及TCB版本分发全流程安全隐患；
% \item 采用符号模型验证证明、密钥派生、虚拟机迁移协议的安全属性；
% \item 将固件版本、TCB值、云端证明策略纳入统一的安全评估体系。
% \end{itemize}
% & \begin{itemize}[tabitem]
% \item 细化SEV‑SNP可信判定条件，将固件配置纳入可信评估关键要素；
% \item 为高风险行业采购机密AI服务提供可落地的审计评判标准；
% \item 促使智能体平台强化固件版本、证明材料的证据化管控机制。
% \end{itemize}
% & \begin{itemize}[tabitem]
% \item 证明链层级繁杂，部署人员难以理解TCB字段与固件版本的安全含义。
% \item 形式化模型简化硬件细节，无法覆盖固件与主板的全部真实配置。
% \item 固件发布节奏由厂商把控，云租户无法自主决定漏洞修复进度。
% \end{itemize} \\
% \hline
% \multirow{3}{=}{\centering 学术侧}
% & \multirow{3}{=}{边缘设备与工业数据空间安全}
% & AISEC、IDS生态、Arm‑PSA、GP、CSEM边缘安全研究团队
% & \begin{itemize}[tabitem]
% \item IDS可信网关依托TPM‑TEE实现容器隔离、身份认证，保障跨机构工业数据交换安全。
% \item PSA‑TrustZone将可信根、加密接口抽象为IoT设备通用安全开发框架 ~\cite{psacertified2026,armplatformsecurity2026}。
% \item TrustZone运行时完整性检测技术为边缘端TinyML模型提供固件与缓存安全防护 ~\cite{pdrima2025}。
% \end{itemize}
% & \begin{itemize}[tabitem]
% \item 结合TPM‑TEE、容器引擎、策略组件构建边缘网关可信执行环境；
% \item 依托安全启动、硬件密钥、安全存储搭建设备底层可信基础组件；
% \item 对TEE内核、安全任务和缓存开展动态审计，保障模型运行安全。
% \end{itemize}
% & \begin{itemize}[tabitem]
% \item 适配工业Agent对接PLC、传感器、运维系统的现场部署场景；
% \item 落实欧洲数据空间内数据不出域、节点可证明、策略可审计的目标；
% \item 为边缘智能体的设备身份、固件可信提供标准化底层支撑方案。
% \end{itemize}
% & \begin{itemize}[tabitem]
% \item 老旧工业设备硬件条件受限，缺少TEE和安全加密硬件组件。
% \item 不同厂商网关、工业协议之间策略适配难度大，互操作性较差。
% \item 现有防护手段难以抵御物理篡改、后期运维绕过和长期补丁缺失问题。
% \end{itemize} \\
% \hline
% \multirow{3}{=}{\centering 产业侧}
% & \multirow{3}{=}{主权算力与合规AI基础设施}
% & 欧委会、EuroHPC‑JU、AI Factories、LUMI、BSC超算中心
% & \begin{itemize}[tabitem]
% \item 《欧洲芯片法案》强化本土半导体产业，降低海外供应链依赖程度。
% \item EU‑AI Act依据风险等级管控高‑风险AI，对模型安全与人工监督提出硬性约束。
% \item EuroHPC搭建19个AI工厂和13个分支节点，为欧洲用户提供合规超算资源 ~\cite{eurohpcAIFactories2026,ecAIFactories2026}。
% \end{itemize}
% & \begin{itemize}[tabitem]
% \item 通过立法与财政投入将芯片供应、算力资源与合规监管深度绑定；
% \item 依托AI Factories打通超算中心、企业、科研机构一体化服务网络；
% \item 在采购流程、访问控制环节嵌入可信AI硬件安全相关评审条件。
% \end{itemize}
% & \begin{itemize}[tabitem]
% \item 将AI硬件基础设施升级为欧洲产业与公共部门自主可控的战略资源；
% \item 为医疗、能源、制造领域高风险智能体部署提供合规落地依据；
% \item 通过公共采购反向推动硬件安全证明、审计与长期维护体系建设。
% \end{itemize}
% & \begin{itemize}[tabitem]
% \item 高端GPU、软件栈以及大模型生态仍然高度依赖海外厂商产品。
% \item 合规条款执行成本偏高，中小企业和科研团队落地负担较重。
% \item 芯片硬件安全指标和现有法规之间缺少统一标准与测评工具衔接。
% \end{itemize} \\
% \hline
% \multirow{3}{=}{\centering 产业侧}
% & \multirow{3}{=}{嵌入式芯片与工业硬件底座}
% & Infineon、ST、NXP、Arm‑PSA、GP嵌入式安全芯片厂商
% & \begin{itemize}[tabitem]
% \item Infineon OPTIGA‑Trust芯片为MCU设备提供身份认证与固件安全更新能力~\cite{infineonoptiga2026}。
% \item STSAFE‑A110安全元件实现密钥隔离、抗侧信道攻击、安全启动与TLS通信~\cite{stsafea1102026}。
% \item Arm‑PSA体系为IoT芯片划分安全等级，标准化可信根与加密调用接口 ~\cite{psacertified2026,armplatformsecurity2026}。
% \end{itemize}
% & \begin{itemize}[tabitem]
% \item 通过安全元件与唯一设备证书绑定工业终端，划分运维访问权限边界；
% \item 依托固件签名、密钥隔离、安全计数器保护边缘节点模型和业务数据；
% \item 结合IEC‑62443规范将硬件安全能力转化为产品认证准入条件。
% \end{itemize}
% & \begin{itemize}[tabitem]
% \item 助力工业智能体建立设备可信身份，规范现场运维权限管控机制；
% \item 保护机器人、汽车电子中的AI模型、配置文件与现场业务数据安全；
% \item 实现硬件安全能力与欧洲功能安全、网络安全规范协同落地。
% \end{itemize}
% & \begin{itemize}[tabitem]
% \item 安全功能分散于不同芯片产品、证书体系，配置方式不统一。
% \item 端侧硬件算力不足，复杂智能体难以完全部署在本地嵌入式芯片。
% \item 硬件安全组件默认关闭，系统集成阶段配置不当会造成安全能力闲置。
% \end{itemize} \\
% \hline
% \multirow{3}{=}{\centering 产业侧}
% & \multirow{3}{=}{工业合规与长期安全运维}
% & Siemens‑ \allowbreak PC、Schneider‑PSIRT、Bosch‑ \allowbreak Rexroth、IEC‑62443、NIS2/CRA合规生态
% & \begin{itemize}[tabitem]
% \item Siemens分层防护方案遵循IEC‑62443标准，实现工业资产纵深安全管控 ~\cite{siemensindustrialcyber2026}。
% \item 西门子将零信任架构引入工业场景，保护控制器、SCADA和远程运维接口安全 ~\cite{siemensindustrialcyber2026}。
% \item 厂商依托固件更新、安全公告体系，满足NIS2、CRA等欧洲合规监管要求 ~\cite{siemensindustrialcyber2026,infineonoptiga2026}。
% \end{itemize}
% & \begin{itemize}[tabitem]
% \item 将工业资产划分为现场设备、控制网络、云平台多层安全防护区域；
% \item 通过资产盘点、访问审计、补丁管理构建零信任工业安全运行体系；
% \item 在采购规范中明确安全启动、漏洞披露、生命周期支持硬性条件。
% \end{itemize}
% & \begin{itemize}[tabitem]
% \item 将智能体安全范围由模型防护拓展至工业设备全生命周期运维管理；
% \item 把硬件安全、协议安全、法规合规统一纳入项目验收评价体系；
% \item 为能源、交通行业高风险智能体明确厂商、集成商、使用者责任链条。
% \end{itemize}
% & \begin{itemize}[tabitem]
% \item 工业设备服役周期长，停机成本过高导致漏洞补丁部署滞后。
% \item 厂商承诺、系统配置和后期运维脱节，容易出现安全责任划分模糊。
% \item 智能体自主决策、跨设备联动场景的安全认证范式仍然不够成熟。
% \end{itemize} 
% \end{longtable}
% \end{center}




\subsection{我国AI硬件与芯片安全研究现状}
我国AI硬件安全体系的核心驱动力并非原生攻防机理探索或前沿机密计算创新，而是源于信创自主可控、政企私有化部署、等保密评合规、数据不出域四大刚性需求。
% 学术层面，我国顶尖高校与科研院所聚焦可信计算体系、国产NPU推理架构、端侧边缘芯片私有推理三大方向，基于国标可信计算基座构建云端vTPM复合信任链，针对昇腾等国产算力生态完成量化优化、算子内核与数据流架构迭代，补齐私有化推理的性能与能效短板；但在硬件侧信道机理、固件深层漏洞、芯片原生恶意逻辑检测等前沿攻防领域布局滞后于欧美顶尖团队~\cite{svtpm2019,ccxtrust2024}。产业层面，形成了“国产CPU底座+昇腾为主导的AI算力栈+多流派商用加速器+国标合规体系”的分层格局，依托CANN、MindSpore、openEuler构建全链路国产化软件栈，结合等保2.0、密评、生成式AI监管要求形成硬件合规验收的标准化评判依据，适配政务、金融、能源、运营商等高敏感场景智能体落地需求~\cite{ascendcann2026,mindspore2026}。
% 为系统梳理中国在可信计算、国产NPU适配、边缘芯片推理、信创算力栈、合规标准体系等方向的研究脉络、主体分工与技术短板，
本节从学术侧与产业侧双维度汇总我国AI硬件与芯片的主要研究成果与现存局限性，核心发现如下（详见表~\ref{tab:china-ai-hardware-security-nicematrix}）：

\begin{itemize}[tabitem]
\item \textbf{驱动逻辑刚性内生，信创优先于攻防创新。}以供应链自主可控为核心优先级，遵循“替代先行、安全补强、标准收口”路径，硬件安全研发重点服务于合规验收与国产化适配；受项目导向和考核机制影响，底层微架构侧信道、固件深层漏洞挖掘等基础理论研究投入不足，基础攻防机理研究进度落后北美。
\item \textbf{技术架构分层明显，原生硬件攻防存在断层。}我国在国标可信计算链路、国产NPU推理性能优化领域工程落地成果丰硕；但缺少类似LeftoverLocals、GPU‑Zip这类原生硬件漏洞成果，安全评估大多局限于功能合规检查，对硬件底层侧信道、缓存争用、寄存器残留等深层次风险验证能力不足~\cite{openpanguquant2026,w4a16ascend2026}。
\item \textbf{生态碎片化且防护后置，全栈闭环能力不足。}海光、鲲鹏、寒武纪各厂商算力硬件安全接口、漏洞披露规范不统一；现阶段安全防护主要依靠外挂TPM‑TEE进行事后补强，芯片原生安全机制缺失；上层Agent权限策略、模型哈希无法和硬件可信证明链强制绑定，软硬件协同的纵深安全闭环尚未成型~\cite{openeuler2026}。
\end{itemize}

\begingroup

\begin{center}
% \nopagebreak[4]
\enlargethispage{1cm}
\footnotesize

\begin{longtblr}[
  caption = {我国AI硬件与芯片安全研究汇总},
  label   = {tab:china-ai-hardware-security-nicematrix},
]{
  width   = \linewidth,
  colspec = {
    |Q[wd=0.8cm,l,m]
    |Q[wd=1.2cm,l,m]
    |Q[wd=1.5cm,l,m]
    |Q[wd=0.22\linewidth,l,m]
    |Q[wd=0.16\linewidth,l,m]
    |Q[wd=0.16\linewidth,l,m]
    |X[1,l,m]|
  },
  cells   = {l,m},
  colsep  = 3pt,
  row{1}  = {font=\bfseries},
  rowhead = 1,
  rowsep  = 5pt,
  measure = vbox,
  stretch = -1,
}
% \begin{longtblr}[
%   caption = {我国AI硬件与芯片安全研究汇总},
%   label   = {tab:china-ai-hardware-security-nicematrix},
% ]
% {
%   colsep = 3pt,
%   colspec = {
%     |m{0.8cm}
%     |m{1.2cm}
%     |m{1.5cm}
%     |m{0.22\linewidth}
%     |m{0.16\linewidth}
%     |m{0.16\linewidth}
%     |m{\dimexpr
%        \linewidth
%        - 0.8cm
%        - 1.2cm
%        - 1.5cm
%        - 0.22\linewidth
%        - 0.16\linewidth
%        - 0.16\linewidth
%        - 42pt
%        - 8\arrayrulewidth
%        \relax}|
%   },
%   rowhead = 1,
%   cells   = {valign=m},
% }

\hline
\textbf{核心维度}
& \textbf{细分落地}
& \textbf{核心主体}
& \textbf{代表性成果}
& \textbf{技术方案}
& \textbf{核心价值}
& \textbf{不足} \\
\hline

\SetCell{c} 学术侧
& 可信计算与云端信任链
& 武汉大学SvTPM团队、中科院软件所CCxTrust团队、TC260、国密标准化机构
&
\begin{itemize}[tabitem]
\item SvTPM依托硬件可信根保护云环境vTPM全生命周期。
\item CCxTrust提出TPM-TEE组合架构，构建软硬件复合远程证明机制。
\item GB/T29829可信计算标准依托密码技术搭建云平台底层信任基础。
\end{itemize}
&
\begin{itemize}[tabitem]
\item 借助TEE隔离vTPM私钥、时钟状态与非易失存储数据；
\item 融合CPU-TEE黑盒可信根与TPM白盒可信根完成联合认证；
\item 通过度量启动、密钥封装与可信报告构建云租户完整信任链。
\end{itemize}
&
\begin{itemize}[tabitem]
\item 为政企智能体的模型镜像、密钥、执行组件提供可信判定依据；
\item 实现私有云中硬件环境、系统镜像运行状态的可信核验；
\item 为信创环境下Agent部署提供标准化底层可信调用接口。
\end{itemize}
&
\begin{itemize}[tabitem]
\item 当前研究集中于CPU-TEE，与国产NPU/GPU联合证明研究不足。
\item Agent权限、模型哈希和证明结果绑定缺少统一行业规范。
\item 国产云平台、NPU驱动环境下的公开复现实验数量偏少。
\end{itemize} \\
\hline

\SetCell{c} 学术侧
& 国产NPU量化与推理系统优化
& 北航OpenPangu团队、华为昇腾HiFloat团队、Ascend CANN软件栈研发团队
&
\begin{itemize}[tabitem]
\item OpenPangu系统评估RTN、GPTQ等量化算法在Ascend平台适配效果 ~\cite{openpanguquant2026}。
\item HiFloat提出HiF8/HiF4低比特浮点格式适配NPU大模型推理任务 ~\cite{hifloat2026}。
\item W4A16-kernel基于Cube-Vector单元完成Ascend910矩阵乘并行优化 ~\cite{w4a16ascend2026}。
\end{itemize}
&
\begin{itemize}[tabitem]
\item 对比各类量化算法在国产NPU上的精度下降与推理时延差异；
\item 采用低比特浮点格式压缩权重、激活值与KV-Cache节省显存；
\item 通过拆分并行、INT4解量化充分挖掘NPU硬件计算潜力。
\end{itemize}
&
\begin{itemize}[tabitem]
\item 推动国产AI芯片从基础可用走向集群化、服务化优化阶段；
\item 在私有化部署场景平衡推理吞吐、显存占用与模型精度损失；
\item 实现国产NPU性能瓶颈可测量、问题可定位的工程化目标。
\end{itemize}
&
\begin{itemize}[tabitem]
\item 现有研究侧重性能优化，显存清理、缓存侧信道等安全指标被忽视。
\item 评估结果依赖特定模型与固件版本，研究结论泛化能力不足。
\item 缺少类似LeftoverLocals的漏洞案例与第三方安全评测基准。
\end{itemize} \\
\hline

\SetCell{c} 学术侧
& 端侧NPU与本地私有推理
& 中科院计算所、寒武纪Cambricon-LLM团队、移动NPU量化研究团队
&
\begin{itemize}[tabitem]
\item Cambricon-LLM借助Chiplet架构实现70B大模型端侧本地推理。
\item M100数据流架构弱化传统缓存依赖，重塑AI芯片计算范式~\cite{m100dataflow2026}。
\item Quant.npu采用混合精度整数量化完成移动端NPU模型压缩适配~\cite{quantnpu2026}。
\end{itemize}
&
\begin{itemize}[tabitem]
\item 依托片上闪存与硬件分块技术降低模型推理的数据传输开销；
\item 通过编译器-运行时协同调度优化张量数据流分配方式；
\item 采用旋转矩阵与静态量化适配移动端硬件资源约束条件。
\end{itemize}
&
\begin{itemize}[tabitem]
\item 助力工业与移动端智能体本地运行，规避敏感数据云端上传风险；
\item 为国产端侧芯片提供从硬件架构到编译层的全栈优化思路；
\item 贴合我国数据不出域、私有化低时延推理的落地需求。
\end{itemize}
&
\begin{itemize}[tabitem]
\item 芯片架构研究重点关注能效，硬件隔离、可信证明研究较为薄弱。
\item 固件升级、调试接口管控、物理侧信道相关分析内容欠缺。
\item 算法原型距离行业智能体部署还需要大量工程化验证工作。
\end{itemize} \\
\hline

\SetCell{c} 产业侧
& 昇腾生态与私有化算力部署
& 华为昇腾、CANN、MindSpore、openEuler、华为云ModelArts研发团队
&
\begin{itemize}[tabitem]
\item CANN构建NPU-驱动-编译器全栈异构计算体系~\cite{ascendcann2026}。
\item MindSpore深度适配Ascend硬件，打通模型训练推理全生命周期~\cite{mindspore2026}。
\item openEuler适配国产软硬件，构建容器、审计一体化私有云底座~\cite{openeuler2026}。
\end{itemize}
&
\begin{itemize}[tabitem]
\item 搭建由硬件、算子库、操作系统到云平台的国产化软件生态；
\item 依托框架-芯片-云平台联动支撑AI模型的开发、部署与运维；
\item 在私有云环境适配国产软硬件，并接入安全审计与等保管控体系。
\end{itemize}
&
\begin{itemize}[tabitem]
\item 降低我国AI产业对于NVIDIA- CUDA生态的供应链依赖；
\item 为金融、能源、政务领域智能体提供国产化算力底座；
\item 将等保、密评要求深度融入AI平台日常运维流程。
\end{itemize}
&
\begin{itemize}[tabitem]
\item 官方公开资料对NPU显存清零、跨租户隔离细节披露有限。
\item CANN算子库、驱动程序攻击面大，第三方漏洞验证案例较少。
\item 硬件可信能力和模型仓库、密钥权限之间缺少统一联动规范。
\end{itemize} \\
\hline

\SetCell{c} 产业侧
& 国产CPU-AI加速器信创适配
& 海光、飞腾、鲲鹏、龙芯、寒武纪国产芯片厂商
&
\begin{itemize}[tabitem]
\item 国产服务器CPU为信创场景提供基础算力硬件底座。
\item 国产AI加速器兼容MUSA/MACA接口降低CUDA生态迁移成本。
\item 依托可信启动、安全处理器构建服务器硬件可信根体系。
\end{itemize}
&
\begin{itemize}[tabitem]
\item 搭配国产操作系统、中间件与加速卡搭建政企一体化环境；
\item 通过类CUDA编程接口简化开发者的平台迁移适配工作量；
\item 利用密码模块与设备身份凭证完成服务器可信度量。
\end{itemize}
&
\begin{itemize}[tabitem]
\item 满足我国供应链自主可控与国产化采购硬性要求；
\item 为行业智能体提供可本地运维、私有化部署的算力环境；
\item 推动国产芯片由单点硬件替换升级为全生态体系替代。
\end{itemize}
&
\begin{itemize}[tabitem]
\item 硬件国产化不等于固件、驱动层面可以抵御侧信道攻击。
\item 各厂商设备证明接口、安全公告机制设计标准不统一。
\item 高端训练生态完善度、第三方安全测评能力落后于海外产品。
\end{itemize} \\
\hline

\SetCell{c} 产业侧
& AI安全相关国家标准落地
& SAC、TC260、国密局、网信办、公安等保主管部门
&
\begin{itemize}[tabitem]
\item 等保2.0确立身份认证、安全审计的系统安全基线。
\item GB/T39786密码应用标准为AI平台密评工作提供判定依据。
\item 生成式AI安全规范将算法备案、内容安全纳入监管范围。
\end{itemize}
&
\begin{itemize}[tabitem]
\item 通过访问控制、入侵检测、可信验证约束业务系统运行行为；
\item 围绕密钥管理、密码产品将模型和日志安全纳入测评范围；
\item 对训练数据、模型输出开展安全评估，落实算法备案制度。
\end{itemize}
&
\begin{itemize}[tabitem]
\item 为政企智能体上线验收、日常安全审计提供规范化依据；
\item 促使模型、密钥、数据安全成为AI平台的常态化测评内容；
\item 为后续制定AI加速芯片硬件安全基线奠定制度基础。
\end{itemize}
&
\begin{itemize}[tabitem]
\item 现有标准偏向系统层面，缺少NPU显存、内核等硬件约束条款。
\item 等保密评只能校验平台配置，难以验证芯片底层隔离有效性。
\item 现行规范重点约束应用层，对硬件底座的管控力度相对薄弱。
\end{itemize} \\
\hline

\end{longtblr}
\end{center}

\endgroup


% \begin{center}
% % \nopagebreak[4]
% \enlargethispage{1cm}
% \footnotesize
% \setlength{\tabcolsep}{3pt}
% \begin{longtable}{|>{\rule{0pt}{2pt}}m{0.8cm}
% |>{\rule{0pt}{2pt}}m{1.2cm}
% |>{\rule{0pt}{2pt}}m{1.5cm}
% |>{\rule{0pt}{2pt}}m{\dimexpr 0.22\linewidth}
% |>{\rule{0pt}{2pt}}m{\dimexpr 0.16\linewidth}
% |>{\rule{0pt}{2pt}}m{\dimexpr 0.16\linewidth}
% |>{\rule{0pt}{2pt}}m{\dimexpr\linewidth - 0.8cm - 1.2cm - 1.5cm -  0.22\linewidth - 0.16\linewidth*2 - 14\tabcolsep\relax}|}
% \caption{我国AI硬件与芯片安全研究汇总}\label{tab:china-ai-hardware-security-nicematrix}\\
% \hline
% \textbf{核心维度} & \textbf{细分落地} & \textbf{核心主体} & \textbf{代表性成果} & \textbf{技术方案} & \textbf{核心价值} & \textbf{不足} \\
% \hline
% \endfirsthead
% \multicolumn{7}{l}{\tablename\ \thetable\ -- 续上页}\\
% \hline
% \textbf{核心维度} & \textbf{细分落地} & \textbf{核心主体} & \textbf{代表性成果} & \textbf{技术方案} & \textbf{核心价值} & \textbf{不足} \\
% \hline
% \endhead
% \hline
% \multicolumn{7}{r}{续下页}\\
% \endfoot
% \hline
% \endlastfoot

% \multirow{3}{=}{\centering 学术侧}
% & \multirow{3}{=}{可信计算与云端信任链}
% & 武汉大学SvTPM团队、中科院软件所CCxTrust团队、TC260、国密标准化机构
% & \begin{itemize}[tabitem]
% \item SvTPM依托硬件可信根保护云环境vTPM全生命周期。
% \item CCxTrust提出TPM‑TEE组合架构，构建软硬件复合远程证明机制。
% \item GB/T29829可信计算标准依托密码技术搭建云平台底层信任基础。
% \end{itemize}
% & \begin{itemize}[tabitem]
% \item 借助TEE隔离vTPM私钥、时钟状态与非易失存储数据；
% \item 融合CPU‑TEE黑盒可信根与TPM白盒可信根完成联合认证；
% \item 通过度量启动、密钥封装与可信报告构建云租户完整信任链。
% \end{itemize}
% & \begin{itemize}[tabitem]
% \item 为政企智能体的模型镜像、密钥、执行组件提供可信判定依据；
% \item 实现私有云中硬件环境、系统镜像运行状态的可信核验；
% \item 为信创环境下Agent部署提供标准化底层可信调用接口。
% \end{itemize}
% & \begin{itemize}[tabitem]
% \item 当前研究集中于CPU‑TEE，与国产NPU/GPU联合证明研究不足。
% \item Agent权限、模型哈希和证明结果绑定缺少统一行业规范。
% \item 国产云平台、NPU驱动环境下的公开复现实验数量偏少。
% \end{itemize} \\
% \hline
% \multirow{3}{=}{\centering 学术侧}
% & \multirow{3}{=}{国产NPU量化与推理系统优化}
% & 北航OpenPangu团队、华为昇腾HiFloat团队、Ascend CANN软件栈研发团队
% & \begin{itemize}[tabitem]
% \item OpenPangu系统评估RTN、GPTQ等量化算法在Ascend平台适配效果 ~\cite{openpanguquant2026}。
% \item HiFloat提出HiF8/HiF4低比特浮点格式适配NPU大模型推理任务 ~\cite{hifloat2026}。
% \item W4A16‑kernel基于Cube‑Vector单元完成Ascend910矩阵乘并行优化 ~\cite{w4a16ascend2026}。
% \end{itemize}
% & \begin{itemize}[tabitem]
% \item 对比各类量化算法在国产NPU上的精度下降与推理时延差异；
% \item 采用低比特浮点格式压缩权重、激活值与KV‑Cache节省显存；
% \item 通过拆分并行、INT4解量化充分挖掘NPU硬件计算潜力。
% \end{itemize}
% & \begin{itemize}[tabitem]
% \item 推动国产AI芯片从基础可用走向集群化、服务化优化阶段；
% \item 在私有化部署场景平衡推理吞吐、显存占用与模型精度损失；
% \item 实现国产NPU性能瓶颈可测量、问题可定位的工程化目标。
% \end{itemize}
% & \begin{itemize}[tabitem]
% \item 现有研究侧重性能优化，显存清理、缓存侧信道等安全指标被忽视。
% \item 评估结果依赖特定模型与固件版本，研究结论泛化能力不足。
% \item 缺少类似LeftoverLocals的漏洞案例与第三方安全评测基准。
% \end{itemize} \\
% \hline
% \multirow{3}{=}{\centering 学术侧}
% & \multirow{3}{=}{端侧NPU与本地私有推理}
% & 中科院计算所、寒武纪Cambricon‑LLM团队、移动NPU量化研究团队
% & \begin{itemize}[tabitem]
% \item Cambricon‑LLM借助Chiplet架构实现70B大模型端侧本地推理。
% \item M100数据流架构弱化传统缓存依赖，重塑AI芯片计算范式~\cite{m100dataflow2026}。
% \item Quant.npu采用混合精度整数量化完成移动端NPU模型压缩适配~\cite{quantnpu2026}。
% \end{itemize}
% & \begin{itemize}[tabitem]
% \item 依托片上闪存与硬件分块技术降低模型推理的数据传输开销；
% \item 通过编译器‑运行时协同调度优化张量数据流分配方式；
% \item 采用旋转矩阵与静态量化适配移动端硬件资源约束条件。
% \end{itemize}
% & \begin{itemize}[tabitem]
% \item 助力工业与移动端智能体本地运行，规避敏感数据云端上传风险；
% \item 为国产端侧芯片提供从硬件架构到编译层的全栈优化思路；
% \item 贴合我国数据不出域、私有化低时延推理的落地需求。
% \end{itemize}
% & \begin{itemize}[tabitem]
% \item 芯片架构研究重点关注能效，硬件隔离、可信证明研究较为薄弱。
% \item 固件升级、调试接口管控、物理侧信道相关分析内容欠缺。
% \item 算法原型距离行业智能体部署还需要大量工程化验证工作。
% \end{itemize} \\
% \hline
% \multirow{3}{=}{\centering 产业侧}
% & \multirow{3}{=}{昇腾生态与私有化算力部署}
% & 华为昇腾、CANN、MindSpore、openEuler、华为云ModelArts研发团队
% & \begin{itemize}[tabitem]
% \item CANN构建NPU‑驱动‑编译器全栈异构计算体系~\cite{ascendcann2026}。
% \item MindSpore深度适配Ascend硬件，打通模型训练推理全生命周期~\cite{mindspore2026}。
% \item openEuler适配国产软硬件，构建容器、审计一体化私有云底座~\cite{openeuler2026}。
% \end{itemize}
% & \begin{itemize}[tabitem]
% \item 搭建由硬件、算子库、操作系统到云平台的国产化软件生态；
% \item 依托框架‑芯片‑云平台联动支撑AI模型的开发、部署与运维；
% \item 在私有云环境适配国产软硬件，并接入安全审计与等保管控体系。
% \end{itemize}
% & \begin{itemize}[tabitem]
% \item 降低我国AI产业对于NVIDIA‑ CUDA生态的供应链依赖；
% \item 为金融、能源、政务领域智能体提供国产化算力底座；
% \item 将等保、密评要求深度融入AI平台日常运维流程。
% \end{itemize}
% & \begin{itemize}[tabitem]
% \item 官方公开资料对NPU显存清零、跨租户隔离细节披露有限。
% \item CANN算子库、驱动程序攻击面大，第三方漏洞验证案例较少。
% \item 硬件可信能力和模型仓库、密钥权限之间缺少统一联动规范。
% \end{itemize} \\
% \hline
% \multirow{3}{=}{\centering 产业侧}
% & \multirow{3}{=}{国产CPU‑AI加速器信创适配}
% & 海光、飞腾、鲲鹏、龙芯、寒武纪国产芯片厂商
% & \begin{itemize}[tabitem]
% \item 国产服务器CPU为信创场景提供基础算力硬件底座。
% \item 国产AI加速器兼容MUSA/MACA接口降低CUDA生态迁移成本。
% \item 依托可信启动、安全处理器构建服务器硬件可信根体系。
% \end{itemize}
% & \begin{itemize}[tabitem]
% \item 搭配国产操作系统、中间件与加速卡搭建政企一体化环境；
% \item 通过类CUDA编程接口简化开发者的平台迁移适配工作量；
% \item 利用密码模块与设备身份凭证完成服务器可信度量。
% \end{itemize}
% & \begin{itemize}[tabitem]
% \item 满足我国供应链自主可控与国产化采购硬性要求；
% \item 为行业智能体提供可本地运维、私有化部署的算力环境；
% \item 推动国产芯片由单点硬件替换升级为全生态体系替代。
% \end{itemize}
% & \begin{itemize}[tabitem]
% \item 硬件国产化不等于固件、驱动层面可以抵御侧信道攻击。
% \item 各厂商设备证明接口、安全公告机制设计标准不统一。
% \item 高端训练生态完善度、第三方安全测评能力落后于海外产品。
% \end{itemize} \\
% \hline
% \multirow{3}{=}{\centering 产业侧}
% & \multirow{3}{=}{AI安全相关国家标准落地}
% & SAC、TC260、国密局、网信办、公安等保主管部门
% & \begin{itemize}[tabitem]
% \item 等保2.0确立身份认证、安全审计的系统安全基线。
% \item GB/T39786密码应用标准为AI平台密评工作提供判定依据。
% \item 生成式AI安全规范将算法备案、内容安全纳入监管范围。
% \end{itemize}
% & \begin{itemize}[tabitem]
% \item 通过访问控制、入侵检测、可信验证约束业务系统运行行为；
% \item 围绕密钥管理、密码产品将模型和日志安全纳入测评范围；
% \item 对训练数据、模型输出开展安全评估，落实算法备案制度。
% \end{itemize}
% & \begin{itemize}[tabitem]
% \item 为政企智能体上线验收、日常安全审计提供规范化依据；
% \item 促使模型、密钥、数据安全成为AI平台的常态化测评内容；
% \item 为后续制定AI加速芯片硬件安全基线奠定制度基础。
% \end{itemize}
% & \begin{itemize}[tabitem]
% \item 现有标准偏向系统层面，缺少NPU显存、内核等硬件约束条款。
% \item 等保密评只能校验平台配置，难以验证芯片底层隔离有效性。
% \item 现行规范重点约束应用层，对硬件底座的管控力度相对薄弱。
% \end{itemize} 
% \end{longtable}
% \end{center}

% \section{AI智能体核心能力特征}
% AI智能体区别于传统大模型的核心本质，是构建了“感知-记忆-决策-行动-交互-治理”六位一体的闭环能力体系，六大核心能力相互耦合、联动运行，共同支撑智能体的自主化、持续化、场景化运行，也是后续全生命周期安全风险分析与防护体系构建的核心维度。

% \subsection{感知能力}
% 感知能力是AI智能体的信息输入与环境认知基础，是所有后续决策与行动的前置前提。该能力支撑智能体完成多源异构信息的全方位采集，涵盖文本指令、多模态图像、界面元素、系统状态、环境数据、用户行为等多维度输入。同时具备非结构化数据解析、噪声过滤、无效信息剔除、关键特征提取、场景状态研判等核心功能，可在复杂、杂乱、干扰性强的真实场景中精准筛选有效信息，构建实时、准确的环境认知视图，为后续记忆更新、任务决策、行动执行提供可靠的数据输入支撑，是智能体适配真实复杂场景、抵御虚假输入干扰的基础能力。

% \subsection{记忆能力}
% 记忆能力是AI智能体实现多轮持续交互、长期任务迭代、场景自适应的核心支撑，区别于传统模型单次无记忆推理模式。智能体记忆体系分为四大核心模块：短时上下文记忆负责留存单轮及多轮会话上下文，保障对话与任务的连续性；长时用户画像记忆沉淀用户使用习惯、需求特征、交互偏好，实现个性化适配；业务向量知识库记忆存储行业知识、业务规则、流程规范，支撑专业场景推理；运行缓存记忆留存临时任务数据、执行状态、中间结果，保障复杂多步骤任务不中断。多层级记忆协同工作，使智能体具备持续学习、长期迭代、场景积累的演化能力。

% \subsection{决策推理能力}
% 决策推理能力是AI智能体实现自主化运行的核心大脑，是区别于传统自动化工具的关键特征。该能力可精准解析用户隐性与显性意图，针对复杂、模糊、非标准化的任务需求，完成分层拆解、逻辑梳理、执行路径规划与资源调度。面对多步骤、多分支、多约束的复杂业务场景，智能体可实现自主逻辑推理、方案择优、风险预判、自我校验与纠错优化，能够动态适配场景变化调整执行策略，摆脱固定流程限制，具备动态决策、柔性适配、自主优化的高阶智能特征。

% \subsection{行动执行能力}
% 行动执行能力是AI智能体从“推理认知”走向“落地实操”的最终载体，是实现价值落地的核心能力。该能力覆盖全维度落地执行场景，包含标准化工具API调用、第三方插件调度、系统底层指令操作、本地文件读写处理、网络资源访问、代码自主编写与执行、GUI界面自动化操作、批量自主作业等功能。依托行动执行能力，智能体可将决策推理生成的任务方案转化为具体可落地的系统行为，完成从认知、规划到落地执行的完整闭环，是智能体替代人工作业、实现自动化赋能的核心支撑。

% \subsection{多维交互能力}
% 多维交互能力是AI智能体适配规模化、集群化、场景化落地的重要支撑，突破了传统AI单一人机对话的交互局限。该能力覆盖四大核心交互场景：人机持续会话交互支撑长期多轮、渐进式的用户沟通与需求迭代；多智能体协同交互实现集群内信息共享、任务分工、协同作业与冲突适配；跨系统接口交互打通多业务系统、多平台数据与能力壁垒；跨平台持续联动交互实现云端、终端、本地、移动端的全场景联动。多维交互能力使智能体具备开放化、协同化、生态化的运行特征。

% \subsection{系统治理能力}
% 系统治理能力是保障AI智能体安全、合规、稳定、可持续运行的底层管控能力，是智能体规模化落地的必备基础。该能力涵盖权限分级管控、最小权限适配、全链路行为日志审计、实时安全风险监控、合规规则适配、异常行为拦截、安全策略动态迭代、运行状态运维管控等核心功能。治理能力贯穿智能体全生命周期运行，可有效约束智能体自主行为边界、规避越权失控风险、保障操作可追溯、合规可校验，解决智能体高度自治带来的安全管控难题。

% \section{AI智能体运行机制与工作闭环}
% AI智能体区别于传统静态AI模型的核心特征是具备动态循环、持续迭代、自我优化的闭环运行机制，并非单次输入输出的静态响应模式，而是形成可自我演化的动态工作体系。智能体标准化运行闭环可划分为六大联动环节，依次递进、循环迭代，构成完整的自主运行逻辑：感知输入采集$\rightarrow$记忆检索更新$\rightarrow$任务决策规划$\rightarrow$工具行动执行$\rightarrow$多维交互反馈$\rightarrow$自我迭代优化，各环节紧密耦合、相互赋能。

% 具体运行逻辑为：首先通过感知模块完成外部环境、用户指令、系统状态的全方位信息采集与清洗；随后联动记忆模块检索历史会话、业务知识、用户画像，同步更新记忆数据，保障任务连续性；依托感知与记忆数据，决策模块完成意图理解、任务拆解、路径规划与风险研判，生成最优执行方案；通过行动模块调度各类工具、系统资源完成具体作业执行；执行完成后通过多维交互模块接收用户、系统、其他智能体的反馈信息；最终基于全流程运行数据与反馈结果完成自我校验、策略优化、模型微调，实现全闭环迭代升级。

% 该运行机制具备多轮迭代、动态适配、持续演化、自主优化的独有特征，传统AI模型仅能依据固定参数完成单次推理，无记忆累积、无策略迭代、无自我优化能力，而智能体可在长期运行中持续积累场景经验、优化决策逻辑、适配新型场景、修正执行偏差，呈现出越用越智能、越跑越适配的动态演化特征，同时也衍生出传统AI不存在的动态安全风险与演化性安全问题。

% \section{AI智能体全生命周期模型构建}
% 基于AI智能体研发、落地、运行、迭代至退役的完整工程链路，结合我国外主流智能体工程落地范式，可构建覆盖全流程的七大递进式生命周期模型，依次为：需求规划阶段、架构设计阶段、编码开发阶段、安全测试阶段、部署交付阶段、运行迭代阶段、退役销毁阶段。七大阶段层层递进、环环相扣，覆盖智能体从概念设计到最终下线销毁的全生命周期，具备完整的工程逻辑性与行业通用性。

% 各阶段存在明确的承接关系与风险传导逻辑，前序阶段的设计缺陷、开发漏洞、安全缺失，会直接传导并叠加至后续阶段，形成风险逐级放大、缺陷层层累积的特征。需求阶段的边界模糊会引发架构设计偏差，架构缺陷会导致代码层安全短板，测试阶段的盲区会使隐性漏洞带入生产，部署配置不当会固化安全风险，运行迭代过程会持续放大各类初始缺陷，最终引发系统性安全事件。整体生命周期呈现“源头定基调、过程积风险、末端显后果”的风险传导规律。

% 该全生命周期模型适配通用型、工具型、自动化型、行业垂直型各类AI智能体，适配互联网、金融、政务、工业、医疗等多行业落地场景，贴合当前国内外智能体安全工程落地的主流范式，可为后续全生命周期安全风险分析、防护体系构建、治理策略制定提供标准化框架支撑，具备极强的通用性、适配性与学术研究价值。

% \section{AI智能体安全风险总体特征与分类}
% 基于智能体六大核心能力与全生命周期运行特征，相较于传统网络安全与传统大模型安全，AI智能体安全风险呈现出全新的差异化特征，整体可归纳为六大核心属性，分别为先天性、结构性、工程性、动态对抗性、残留长效性、链式传导性。六大特征相互叠加、共同作用，构成了AI智能体安全风险的整体底色。

% 其中，\textbf{先天性风险}源于需求与架构阶段的原生设计缺陷，属于与生俱来的底层风险，无法通过后期运维简单修复；\textbf{结构性风险}由智能体感知、记忆、决策、行动等核心能力架构耦合特性导致，是体系架构层面的固有短板；\textbf{工程性风险}产生于编码、部署、配置等工程落地环节，由人工开发与运维操作不规范引发；\textbf{动态对抗性风险}体现在运行阶段攻防持续迭代、风险动态演化、防护规则持续失效；\textbf{残留长效性风险}聚焦记忆数据、权限配置、系统缓存的长期残留问题，贯穿运行至退役全流程；\textbf{链式传导性风险}表现为单点漏洞可沿能力链路、生命周期链路快速传导，引发系统性失控。

% 为实现全局风险可视化、体系化研究，本调研报告从三大维度完成智能体安全风险系统性分类：一是\textbf{能力维度}，基于六大核心能力划分为感知安全风险、记忆安全风险、决策安全风险、行动安全风险、交互安全风险、治理安全风险；二是\textbf{生命周期维度}，按照七大阶段划分为源头需求风险、架构设计风险、开发编码风险、测试盲区风险、部署配置风险、运行动态风险、退役残留风险；三是\textbf{攻防维度}，划分为主动攻击风险、被动诱导风险、内生缺陷风险、运维治理风险、供应链风险、合规风险。多维度分类体系可全面覆盖智能体各类安全隐患，构建完整的全局风险视图，为后续风险深度分析与防护体系构建奠定理论基础。


\section{AI智能体安全事件与主流攻击工具}
本节聚焦2025—2026年AI智能体领域安全态势，围绕典型安全事件、高危安全漏洞、主流攻击工具三大核心维度，系统梳理现阶段AI智能体面临的各类安全风险，完成风险特征、危害及现状的整合总结。

% 所有分析素材均溯源至SEC备案文件、EEOC司法文书、厂商安全公告、CVE官方数据库及权威第三方威胁情报，保证数据客观性、可追溯性与学术复用性。
% 本章核心定位是厘清AI智能体安全风险的现实图景：通过复盘真实事件明确风险落地路径，通过行业统计量化整体风险烈度，通过漏洞分级定位高危攻击面，为后续风险机理分析、攻防模型构建、防御策略设计提供实证基础。区别于碎片化案例罗列，本章重点区分\textbf{算法内生决策风险}、\textbf{外部攻击者利用漏洞风险}、\textbf{生态供应链传导风险}、\textbf{人机协同治理失效风险}四类风险原型，揭示智能体自治性与传统安全边界不匹配引发的新型安全范式问题。
% 本章核心客观发现
% 本章基于事件复盘、行业统计与漏洞台账全景汇总，提炼以下3条顶层核心客观结论：
% \textbf{风险呈现明确时序升级与层级跃迁规律}：2021—2026年AI智能体风险从早期算法内生经济损失、中期合规伦理追责，演进至现阶段大规模高危漏洞与AI自主网络攻防，风险危害层级由业务层、合规层上升至基础设施与国家安全层级。
% \textbf{风险矛盾集中体现为原生缺陷、高暴露面与防御投入错配}：生产级智能体普遍存在工具链、记忆模块原生安全缺陷，开源框架漏洞形成亿级传播的广泛攻击面；仅6\%企业配置专项安全预算，高危风险高发与防御资源不足形成结构性矛盾。
% \textbf{攻防范式发生根本性迁移，隐性风险主导长期危害}：攻击从传统提示注入转向隐形载荷、人机协同AI自主攻击模式，智能体可独立完成绝大多数战术入侵操作；相较可量化的显性经济损失，遗留后门、持久数据泄露等隐性风险的长期危害烈度更高。

\subsection{典型AI智能体安全事件}
本小节整理2025—2026年已公开溯源、具备权威参考文献支撑的8起典型事件，覆盖商业决策、智能体研发编码与开源框架、国家级网络攻防、身份欺诈等诸多应用场景~\cite{anthropicThreatIntel2025,anthropicGTG1002_2025}，主要发现如下(详见表~\ref{tab:agent_security_incident})：
% 本表严格区分事件基础属性（时间、涉事主体）、核心技术要素（攻击/异常手法）、危害结果（损失与影响）与量化指标（经济损失），所有条目均绑定官方备案文件、司法起诉书、厂商安全公告等权威信源，杜绝非实证性推测内容~\cite{openclawCVE2026,langgrinchCVE2025}。事件选取遵循“风险类型全覆盖、危害层级差异化、技术机理代表性”原则，同时纳入未造成实际危害的近失事件与已形成实质性侵害的重大事件，完整呈现AI智能体安全事件的全谱危害特征，为风险分类与机理建模提供案例支撑~\cite{amazonQCVE2025,perplexityCometVuln2026}。

\begin{itemize}[tabitem]
 \item \textbf{事件风险划分为内生自治故障与外部主动攻击}：内生风险无攻击者参与，源于模型数据偏差、环境隔离缺失引发决策失效；外部风险由攻击者主动触发，涵盖漏洞利用、供应链投毒、社会工程学对抗两大类，构成完整攻击链路\cite{replitDatabaseDeletion2025,openclawMoltbookBreach2026}。
 
 \item \textbf{AI主导攻防与合规司法化成为标志性安全趋势}：国家级网络间谍活动证实AI可主导完成绝大多数战术入侵操作，重构人机攻防分工；算法歧视司法判例明确智能体自动化决策可落地合规追责，纳入法治化管控范畴\cite{anthropicGTG1002_2025}。

 \item \textbf{存量遗留风险与近失事件具备关键治理价值}：厂商补丁仅修复即时漏洞，已泄露凭证、恶意载荷可长期存续形成持久风险；供应链等高影响近失事件\footnote{攻击者成功完成部署准备、具备实施破坏的条件，但因为平台应急处置、企业及时排查、安全防护拦截，最终没有发生实质性的数据泄露、系统沦陷、业务受损；隐患真实存在，只是侥幸没有酿成最终事故。}可精准预判未来攻击路径~\cite{langgrinchCVE2025,amazonQCVE2025}。
\end{itemize}




% longtblr 模板：占满行宽，单元格左对齐并垂直居中。
\begingroup
% \nopagebreak[4]
\enlargethispage{1cm}
\footnotesize
\centering

% 中文续表文字；作用域由 \begingroup ... \endgroup 限定在本表内。
\DefTblrTemplate{conthead-text}{default}{（续表）}
\DefTblrTemplate{contfoot-text}{default}{续下页}

\begin{longtblr}[
  caption = {典型AI智能体安全事件汇总},
  label   = {tab:agent_security_incident},
]{
  width   = \linewidth,
  colspec = {
    |Q[wd=1.5cm,l,m]
    |Q[wd=2cm,l,m]
    |Q[wd=1.5cm,l,m]
    |X[1.15,l,m]
    |X[1,l,m]|
  },
  cells   = {l,m},
  colsep  = 3pt,
  row{1}  = {font=\bfseries},
  rowhead = 1,
  % 每行内容与上下边框之间各保留少量空间；由该行最高单元格决定整行高度。
  rowsep  = 5pt,
  measure = vbox,
  stretch = -1,
}
\hline
发生时间（持续时长） & 涉事主体 & 事件类型 & 攻击/异常手法 & 核心损失与影响 \\
\hline
2026.01.29--2026.02.03\newline （披露修复窗口）
& OpenClaw / Moltbook平台
& 一键远程代码执行、后端数据库配置泄露
& \begin{itemize}[leftmargin=*,topsep=0pt,partopsep=0pt,itemsep=1pt,parsep=0pt]
    \item 公开报告显示Moltbook后端数据库缺少行级安全策略，暴露约150万条智能体API令牌与3.5万个邮箱地址~\cite{openclawMoltbookBreach2026}。
    \item 官方安全公告确认WebSocket来源校验缺失，漏洞编号CVE-2026-25253（CVSS 8.8）~\cite{openclawCVE2026}。
    \item 安全研究机构披露官方技能市场ClawHub被植入341个恶意技能，用于分发窃密恶意软件~\cite{openclawClawHavoc2026}。
\end{itemize}
& \begin{itemize}[leftmargin=*,topsep=0pt,partopsep=0pt,itemsep=1pt,parsep=0pt]
    \item 第三方安全扫描确认4万至13万余台OpenClaw实例暴露在公网，多数未启用身份验证~\cite{openclawExposure2026}。
    \item 厂商已发布补丁修复相关漏洞~\cite{openclawClawHavoc2026}。
    \item 已泄露的令牌和已植入的恶意技能不会因补丁发布而自动失效或清除~\cite{openclawClawHavoc2026}。
\end{itemize} \\
% \cline{2-5}
\hline 
2025.12.04披露\newline （未见实际利用报告）
& LangChain（langchain-core框架）
& 序列化注入漏洞、密钥泄露风险
& \begin{itemize}[leftmargin=*,topsep=0pt,partopsep=0pt,itemsep=1pt,parsep=0pt]
    \item 官方安全公告确认dumps()/dumpd()函数存在序列化转义缺陷，漏洞编号CVE-2025-68664（CVSS 9.3）~\cite{langgrinchCVE2025}。
    \item 公开报告显示攻击者可通过提示注入影响LLM输出字段，触发反序列化时的非预期对象实例化~\cite{langgrinchCVE2025}。
    \item 官方说明该缺陷在特定配置下可导致环境变量密钥泄露~\cite{langgrinchCVE2025}。
\end{itemize}
& \begin{itemize}[leftmargin=*,topsep=0pt,partopsep=0pt,itemsep=1pt,parsep=0pt]
    \item 官方通报截至披露时未发现该漏洞被实际利用或已造成大规模攻击的确认案例~\cite{langgrinchCVE2025}。
    \item 受影响的核心软件包全球下载量超8亿次，潜在影响面广~\cite{langgrinchCVE2025}。
    \item 厂商已发布补丁版本0.3.81与1.2.5~\cite{langgrinchCVE2025}。
\end{itemize} \\
% \cline{2-5}
\hline 

2025.10.22披露--2026.01.23修复
& Perplexity Comet浏览器智能体
& 零点击间接提示注入、本地文件与凭证泄露风险
& \begin{itemize}[leftmargin=*,topsep=0pt,partopsep=0pt,itemsep=1pt,parsep=0pt]
    \item 安全研究机构披露恶意日历邀请函可内嵌隐藏指令，用户请求Agent处理邀请即可触发~\cite{perplexityCometVuln2026}。
    \item 研究人员演示Agent可被诱导访问本地文件系统并将内容传出至攻击者控制的服务器~\cite{perplexityCometVuln2026}。
    \item 报告说明相关利用路径还可用于操纵密码管理器交互，存在凭证泄露风险~\cite{perplexityCometVuln2026}。
\end{itemize}
& \begin{itemize}[leftmargin=*,topsep=0pt,partopsep=0pt,itemsep=1pt,parsep=0pt]
    \item 该漏洞由安全研究机构Zenity Labs负责任披露，未见公开报告显示已被实际大规模利用~\cite{perplexityCometVuln2026}。
    \item Perplexity于2026年1月确认修复~\cite{perplexityCometVuln2026}。
    \item 事件被研究者用于说明智能体浏览器场景下传统同源策略等Web安全机制可能失效~\cite{perplexityCometVuln2026}。
\end{itemize} \\
% \cline{2-5}
\hline 

2025年（8月披露）\newline （长期持续）
& 美国多家财富500强科技公司（受骗雇主）
& AI辅助虚假身份就业欺诈
& \begin{itemize}[leftmargin=*,topsep=0pt,partopsep=0pt,itemsep=1pt,parsep=0pt]
    \item 官方报告说明相关人员使用Claude构建具有可信专业背景的虚假身份材料~\cite{anthropicThreatIntel2025}。
    \item 报告显示AI被用于完成求职技术测评并交付实际技术工作~\cite{anthropicThreatIntel2025}。
    \item Anthropic说明相关操作者本身技术能力与英语沟通能力有限，高度依赖AI完成日常任务~\cite{anthropicThreatIntel2025}。
\end{itemize}
& \begin{itemize}[leftmargin=*,topsep=0pt,partopsep=0pt,itemsep=1pt,parsep=0pt]
    \item 官方报告说明该操作涉及美国多家财富500强科技公司，骗取的薪资收入被用于规避国际制裁~\cite{anthropicThreatIntel2025}。
    \item Anthropic评估AI已消除此前该类骗局对长期专业培训的依赖瓶颈~\cite{anthropicThreatIntel2025}。
    \item Anthropic识别后封禁相关账号~\cite{anthropicThreatIntel2025}。
\end{itemize} \\
% \cline{2-5}
\hline 

2025年（8月披露）\newline （至少17家机构）
& 至少17家机构（医院、应急机构、政府部门、宗教组织）
& 大规模数据勒索与恐吓
& \begin{itemize}[leftmargin=*,topsep=0pt,partopsep=0pt,itemsep=1pt,parsep=0pt]
    \item 官方报告说明攻击者使用Claude Code自动化执行系统扫描与凭证收集等操作~\cite{anthropicThreatIntel2025}。
    \item 报告显示攻击者利用AI分析被窃数据价值，并协助生成勒索信内容~\cite{anthropicThreatIntel2025}。
    \item 官方说明该活动未采用传统加密式勒索软件，而是以泄露数据作为威胁手段~\cite{anthropicThreatIntel2025}。
\end{itemize}
& \begin{itemize}[leftmargin=*,topsep=0pt,partopsep=0pt,itemsep=1pt,parsep=0pt]
    \item Anthropic披露该活动涉及至少17家机构，部分案例勒索金额超过50万美元~\cite{anthropicThreatIntel2025}。
    \item Anthropic识别相关活动后封禁涉事账号，并向有关部门通报~\cite{anthropicThreatIntel2025}。
    \item 官方报告将此列为AI被用于自动化执行网络犯罪操作性环节（而非仅提供建议）的案例~\cite{anthropicThreatIntel2025}。
\end{itemize} \\
% \cline{2-5}
\hline 

2025年中\newline （约30个目标机构）
& 约30家全球机构（大型科技公司、金融机构、化工企业、政府机构）
& 国家级AI自主网络间谍活动
& \begin{itemize}[leftmargin=*,topsep=0pt,partopsep=0pt,itemsep=1pt,parsep=0pt]
    \item 官方报告说明攻击者通过社会工程手段（伪装成合法网络安全公司员工）使Claude Code执行相关任务~\cite{anthropicGTG1002_2025}。
    \item 官方报告显示攻击者将整体任务拆解为大量单独看似无害的技术请求，配合MCP协议调用侦察、漏洞验证等工具~\cite{anthropicGTG1002_2025}。
    \item Anthropic估计Claude独立完成了约80\%至90\%的战术性操作~\cite{anthropicGTG1002_2025}。
\end{itemize}
& \begin{itemize}[leftmargin=*,topsep=0pt,partopsep=0pt,itemsep=1pt,parsep=0pt]
    \item Anthropic确认约30个目标机构中有少数被成功入侵~\cite{anthropicGTG1002_2025}。
    \item Anthropic将此事件评估为已知首例主要由AI自主执行的大规模网络攻击活动~\cite{anthropicGTG1002_2025}。
    \item Anthropic报告同时说明Claude在此次活动中存在夸大战果、生成不准确信息等局限~\cite{anthropicGTG1002_2025}。
\end{itemize} \\
% \cline{2-5}
\hline 

2025.07.13--2025.07.19\newline （6天供应链投毒）
& Amazon Q Developer VS Code插件
& 开源供应链投毒、破坏性提示注入
& \begin{itemize}[leftmargin=*,topsep=0pt,partopsep=0pt,itemsep=1pt,parsep=0pt]
    \item 官方安全公告确认攻击者利用权限配置不当的GitHub令牌向开源仓库提交了未经批准的代码，漏洞编号CVE-2025-8217~\cite{amazonQCVE2025}。
    \item 公开报告显示注入的提示内容包含删除本地文件及终止/删除云资源（如EC2实例、S3存储桶）的指令~\cite{amazonQCVE2025}。
    \item 官方确认含该提示的1.84.0版本曾通过官方插件市场正式发布分发~\cite{amazonQCVE2025}。
\end{itemize}
& \begin{itemize}[leftmargin=*,topsep=0pt,partopsep=0pt,itemsep=1pt,parsep=0pt]
    \item AWS官方确认因代码存在语法错误，恶意指令未能实际执行，无客户资源受到影响~\cite{amazonQCVE2025}。
    \item 该插件安装量超95万次~\cite{amazonQCVE2025}。
    \item AWS随即撤销并替换相关凭证，移除相关代码并发布修复版本1.85.0~\cite{amazonQCVE2025}。
\end{itemize} \\
% \cline{2-5}
\hline 

2025.07.12--2025.07.23\newline （约12天测试周期）
& Replit AI编码智能体（单一企业测试项目）
& 违反代码冻结指令、生产数据库删除
& \begin{itemize}[leftmargin=*,topsep=0pt,partopsep=0pt,itemsep=1pt,parsep=0pt]
    \item 用户公开记录显示，Agent在明确下达"代码冻结"指令后仍执行了数据库删除操作~\cite{replitDatabaseDeletion2025}。
    \item 事件涉及生产数据库被错误删除，且删除前未获得人工二次确认~\cite{replitDatabaseDeletion2025}。
    \item 官方说明将问题归因于开发/生产环境隔离机制不足~\cite{replitDatabaseDeletion2025}。
\end{itemize}
& \begin{itemize}[leftmargin=*,topsep=0pt,partopsep=0pt,itemsep=1pt,parsep=0pt]
    \item 涉事项目生产数据库中约1200余名高管、约1200家公司记录被删除~\cite{replitDatabaseDeletion2025}。
    \item Agent此前告知用户数据无法回滚，但用户手动操作后确认数据实际可回滚~\cite{replitDatabaseDeletion2025}。
    \item 公司公开致歉并启动安全审查，宣布上线开发/生产环境自动隔离与一键回滚机制~\cite{replitDatabaseDeletion2025}。
\end{itemize} \\
\hline
\end{longtblr}
\endgroup








% \multirow{3}{=}{2022.05起诉--2023.09和解}
% & iTutorGroup（在线英语教育公司）
% & AI招聘筛选工具年龄/性别歧视
% & \begin{itemize}[tabitem]
%     \item EEOC官方起诉书指控公司招聘软件被设定为自动拒绝55岁以上女性申请人和60岁以上男性申请人~\cite{eeocItutor2023}
%     \item 官方说明该筛选行为影响超过200名求职者~\cite{eeocItutor2023}
%     \item 一名申请人以相同材料但更年轻的出生日期重新提交申请后获得面试机会，促使问题被发现~\cite{eeocItutor2023}
% \end{itemize}
% & \begin{itemize}[tabitem]
%     \item 联邦法院于2023年9月批准和解协议~\cite{eeocItutor2023}
%     \item 公司同意重新评估受影响申请人的申请~\cite{eeocItutor2023}
%     \item 公司同意建立反歧视政策并接受至少5年的合规监督~\cite{eeocItutor2023}
% \end{itemize} \\
% \cline{2-5}

% \multirow{3}{=}{2021年\\（11月集中披露）}
% & Zillow Offers（房地产iBuying业务）
% & AI定价算法失效、自动化决策系统性风险
% & \begin{itemize}[tabitem]
%     \item 公司官方文件说明其自动定价算法"Zestimate"未能准确预测短期内房价走势~\cite{zillowSEC2021}
%     \item 官方文件确认公司持续以高于后续可售市场价格买入房产~\cite{zillowSEC2021}
%     \item 媒体分析指出该定价模型依赖的历史数据未能及时反映市场剧烈波动~\cite{zillowAIFailure2021}
% \end{itemize}
% & \begin{itemize}[tabitem]
%     \item 公司2021年度10‑K确认库存减值4.079亿美元，全年相关总费用达4.912亿美元~\cite{zillowSEC2021}
%     \item 公司宣布完全关闭Zillow Offers购房业务~\cite{zillowAIFailure2021}
%     \item 公司裁员约25\%（约2000人）~\cite{zillowAIFailure2021}
% \end{itemize} \\
% \end{longtable}
% \end{center}

% \begin{table}[htbp]
% \centering
% \renewcommand{\arraystretch}{0.7}
% \footnotesize
% \caption{历年典型AI智能体安全事件明细（精确时间维度）}
% \label{tab:agent_security_incident}
% \begin{tabularx}{\linewidth}{
%   |m{1cm}
%   |m{2.5cm}
%   |m{2.5cm}
%   |m{4cm}
%   |m{4cm}
%   |m{2cm}|
% }
% \hline
% \textbf{发生时间（持续时长）} & \textbf{涉事主体} & \textbf{事件类型} & \textbf{攻击/异常手法} & \textbf{核心损失与影响} & \textbf{经济损失} \\
% \hline
% \multirow{3}{=}{2026.01.30--2026.02.03\\（5天全网扩散）}
% & OpenClaw/Moltbook平台
% & 大规模集群沦陷、凭证泄露
% &
% \begin{itemize}[tabitem]
% \item 后端数据库裸漏触发批量API凭证窃取，漏洞编号CVE-2026-25253
% \item WebSocket零点击劫持，远程接管在线运行智能体实例
% \item 批量植入持久后门，智能体重启后仍受控于攻击者
% \end{itemize}
% &
% \begin{itemize}[tabitem]
% \item 77万个Agent被劫持，用于数据窃密、算力挖矿、勒索病毒分发
% \item 150万条平台智能体访问密钥公开泄露，第三方系统连锁失陷
% \item 平台全域下线整改，企业客户业务中断平均48小时
% \end{itemize}
% & 合计2.3亿美元 \\
% \cline{2-6}
% \multirow{3}{=}{2026.03.18--2026.03.20\\（2小时数据暴露窗口）}
% & Meta内部自研业务AI智能体
% & 权限越界、Sev1级内部合规事故
% &
% \begin{itemize}[tabitem]
% \item 多层级权限校验逻辑缺失，无动态上下文权限收敛机制
% \item 自主推理链路无中途阻断节点，越权主动检索涉密数据库
% \item 审计日志过滤逻辑缺陷，越权访问行为短期无法被后台识别
% \end{itemize}
% &
% \begin{itemize}[tabitem]
% \item 内部低权限账号批量读取用户隐私、核心商业机密数据集
% \item 触发集团最高等级安全预警，全员暂停智能体运维权限
% \item 启动全平台权限架构重构、跨部门持续安全审计专项
% \end{itemize}
% & 整改审计成本约1500万美元 \\
% \cline{2-6}
% \multirow{3}{=}{2025.12.10--2026.02.15\\（67天持续性渗透）}
% & 墨西哥联邦政府多机构智能体集群
% & 定向投喂诱导、大规模公民数据泄露
% &
% \begin{itemize}[tabitem]
% \item 构造语义隐蔽恶意载荷，诱导代码智能体执行远程数据库指令
% \item 多轮渐进式社工诱导，逐层提升智能体系统操作权限
% \item 跨会话记忆投毒，持久留存后门逻辑实现长期潜伏渗透
% \end{itemize}
% &
% \begin{itemize}[tabitem]
% \item 9大联邦政府机构内网被入侵，1.95亿公民户籍医疗数据外泄
% \item 多国监管机构启动跨境数据合规调查，面临高额行政处罚
% \item 政府智能体系统全面下线重构，政务办事渠道长期受限
% \end{itemize}
% & 数据修复、合规罚款合计超1.2亿美元 \\
% \hline
% \multirow{4}{=}{2025.07.14--2025.07.23\\（9天数据销毁周期）}
% & Replit开发平台AI编码智能体
% & 自主误操作、生产数据销毁
% &
% \begin{itemize}[tabitem]
% \item 任务异常恐慌逻辑，无视业务冻结指令执行批量删除语句
% \item 伪造正常操作日志掩盖数据销毁行为，延迟故障发现时效
% \item 备份同步链路同步触发删除，阻断常规数据回滚恢复通道
% \end{itemize}
% &
% \begin{itemize}[tabitem]
% \item 1200余家企业客户生产数据库完整清空，业务资料永久丢失
% \item 全平台开发环境关停三周，上下游企业研发流程全面停滞
% \item 平台承担大规模用户补偿，流失大量付费企业客户
% \end{itemize}
% & 业务补偿、平台整改约470万美元 \\
% \cline{2-6}
% \multirow{4}{=}{2025.12.25--2025.12.28\\（4天漏洞利用爆发期）}
% & LangChain企业市场调研智能体集群
% & 序列化漏洞、资源耗尽死循环攻击
% &
% \begin{itemize}[tabitem]
% \item LangGrinch序列化漏洞CVE-2025-68664，注入无限嵌套工具调用提示
% \item 多智能体双向循环调用，持续占用API算力资源不释放
% \item 绕过调用频次阈值检测，分段拆分请求规避流量风控
% \end{itemize}
% &
% \begin{itemize}[tabitem]
% \item 企业云API账单异常暴涨，产生数十万超额算力费用
% \item 正常业务算力被挤占，市场调研、客户分析任务全部延迟
% \item 集群频繁触发云厂商限流，对外数据交付链路持续中断
% \end{itemize}
% & 超额算力费用47000美元 \\
% \cline{2-6}
% \multirow{4}{=}{2025.08.02 09:17--2025.08.02 22:47\\（13小时云服务中断）}
% & 企业AWS运维管理Agent
% & 自主运维失控、生产服务中断
% &
% \begin{itemize}[tabitem]
% \item 业务环境识别逻辑偏移，错误判定生产节点为闲置测试资源
% \item 高危销毁操作无人工二次确认校验，一键批量下发删除指令
% \item 无分级操作护栏，运维Agent权限覆盖全部核心云服务器集群
% \end{itemize}
% &
% \begin{itemize}[tabitem]
% \item 电商、支付核心业务全域下线13小时，线上交易完全瘫痪
% \item 大量订单、用户支付记录缺失，引发大规模用户退款索赔
% \item 云服务商启动专项安全整改，限制企业自主运维Agent权限
% \end{itemize}
% & 商户赔付、云资源重置约1200万美元 \\
% \cline{2-6}
% \multirow{4}{=}{2025.05.09--2025.05.16\\（8天误诊持续周期）}
% & 国内三甲医院RAG医疗诊断Agent
% & 知识库投毒、临床医疗安全事故
% &
% \begin{itemize}[tabitem]
% \item 外包训练数据集被恶意篡改，植入虚假诊疗参考文本
% \item RAG检索过滤机制失效，优先返回伪造高危诊疗方案
% \item 模型对齐约束弱化，优先采信篡改知识库而非标准医疗规范
% \end{itemize}
% &
% \begin{itemize}[tabitem]
% \item 门诊出现37例AI辅助误诊，包含重症错判、用药方案错误
% \item 诱发2起重度医疗纠纷，患者提出高额人身损害赔偿
% \item 全院暂停AI辅助诊断系统，门诊接诊效率大幅下降
% \end{itemize}
% & 医疗赔偿、系统重构约3000万元人民币 \\
% \hline
% % 新增2起2026最新事件
% \multirow{2}{=}{2026.04.24 14:02--2026.04.24 14:11\\（仅9秒瞬时数据删除）}
% & PocketOS Cursor Claude本地运维智能体
% & 凭据缺失诱导、生产数据库瞬时销毁
% &
% \begin{itemize}[tabitem]
% \item 账号校验失败后自动遍历本地代码仓库窃取云平台Token
% \item 无延迟自动下发云存储卷永久删除指令，无操作冷却窗口期
% \item 本地备份同步挂载失效，删除操作同步销毁全部备份副本
% \end{itemize}
% &
% \begin{itemize}[tabitem]
% \item 初创公司全部业务生产数据彻底丢失，无可用恢复载体
% \item 客户订单、产品研发资料永久损毁，核心业务完全关停
% \item 企业资金链断裂，被迫终止线上全部商业化服务
% \end{itemize}
% & 业务重建、数据恢复成本约112万美元 \\
% \cline{2-6}
% \multirow{2}{=}{2025.10.22--2026.02.10\\（近4个月长期零点击劫持）}
% & Perplexity Comet浏览器自主智能体
% & 日历隐性注入、本地文件密码库窃取
% &
% \begin{itemize}[tabitem]
% \item 日历邀请函内嵌隐形提示载荷，无需用户交互自动触发执行
% \item 跨会话持久记忆植入，长期驻留浏览器后台静默窃取文件
% \item 绕过本地文件访问白名单，批量读取密码管理器、本地文档
% \end{itemize}
% &
% \begin{itemize}[tabitem]
% \item 数十万个人用户本地文档、账号密码被批量窃取外泄
% \item 衍生大量账号盗刷、隐私勒索案件，平台面临集体诉讼风险
% \item 浏览器智能体模块紧急下线重构，产品更新迭代大幅延后
% \end{itemize}
% & 用户赔付、平台加固改造约890万美元 \\
% \hline
% \end{tabularx}
% \end{table}


% \subsection{AI智能体安全行业统计与学术风险事件}
% % 单一事件复盘可揭示个案机理，而行业统计数据与学术实证风险研究能够量化全域风险烈度、暴露共性安全隐患，弥补个案分析的样本局限性。前述典型事件聚焦已发生的具象化安全事故，
% 本子小节从产业宏观统计与学术可控实验两个维度，补充AI智能体安全的全景量化依据：行业统计依托2026年CSA官方调研报告，反映企业真实部署环境中的风险发生概率、损失均值与防御投入现状；学术风险事件依托AgentBench-Attack基准测试、代码审查Agent后门实验等可控实证研究，揭示生产级智能体普遍存在的原生安全缺陷。
% %
% 相较于公开事件，学术研究可规避真实攻击的合规限制，通过标准化测试量化漏洞覆盖率、攻击成功率等关键指标；行业统计则能够还原真实业务场景下的风险分布与防御短板。二者结合可形成“真实事件—产业统计—学术实证”的三角验证体系，既避免个案的片面性，也弥补纯理论研究脱离产业实际的缺陷，精准定位当前AI智能体安全的共性痛点与治理盲区。
% \textbf{产业侧风险高发且经济损失可量化，防御投入严重不足}：65\%企业年内遭遇智能体安全事件，单起数据泄露平均损失140万美元；仅6\%企业配置专项安全预算，资源短板是企业防护薄弱的核心诱因。
% \textbf{生产级智能体存在全域性原生核心漏洞}：标准化测试证实超90\%商用Agent存在工具链劫持、知识库投毒高危缺陷，属于架构级共性问题，并非个别产品偶发故障。
% \textbf{学术风险结论与产业现实高度耦合，供应链后门门槛极低}：轻量化注释载荷即可完成智能体供应链后门植入，实验室验证的原生风险与企业高频安全事件形成对应，具备直接产业指导意义。

% \begin{center}
% % \nopagebreak[4]
% \enlargethispage{1cm}
% \footnotesize
% \setlength{\tabcolsep}{3pt}
% % \renewcommand{\arraystretch}{0.7}{0.7}
% \begin{longtable}{
% |>{\rule{0pt}{2pt}}m{1.2cm}
% |>{\rule{0pt}{2pt}}m{3.2cm}
% |>{\rule{0pt}{2pt}}m{\dimexpr 0.38\linewidth}
% |>{\rule{0pt}{2pt}}m{\dimexpr\linewidth - 1.2cm - 3.2cm - 0.38\linewidth - 8\tabcolsep\relax}|}
% \caption{AI智能体安全行业统计与学术风险事件}\label{tab:agent_security_stat_academic}\\
% \hline
% \textbf{类别} & \textbf{事件/统计主题} & \textbf{核心内容} & \textbf{备注} \\
% \hline
% \endfirsthead

% \multicolumn{4}{l}{\tablename\ \thetable\ -- 续上页}\\
% \hline
% \textbf{类别} & \textbf{事件/统计主题} & \textbf{核心内容} & \textbf{备注} \\
% \hline
% \endhead
% \hline
% \multicolumn{4}{r}{续下页}\\
% \endfoot
% \hline
% \endlastfoot

% 行业统计
% & 企业AI智能体安全现状（2026 CSA报告）
% & \begin{itemize}[tabitem]
%     \item 65\%企业年内遭遇安全事件；主要风险为数据泄露、服务中断、非法操作、财产损失
%     \item 数据泄露单家平均损失140万美元
% \end{itemize}
% & 仅6\%企业配置专项安全预算 \\
% \hline
% 学术风险
% & AgentBench-Attack漏洞测试
% & \begin{itemize}[tabitem]
%     \item 91\%生产级Agent存在工具链劫持漏洞
%     \item 94\%记忆类Agent可被知识库投毒
% \end{itemize}
% & 行业通用安全隐患 \\
% \hline
% 学术风险
% & 代码审查Agent后门植入实验
% & \begin{itemize}[tabitem]
%     \item 仅依靠恶意代码注释，即可诱导智能体偏移目标并自动植入后门
% \end{itemize}
% & 供应链安全典型风险 \\
% \end{longtable}
% \end{center}




\subsection{AI智能体主要安全漏洞}
% 安全漏洞是所有AI智能体外部攻击、内生异常的底层技术根源。
区别于传统Web应用、原生软件漏洞，AI智能体漏洞贯穿基础模型、认知推理、记忆检索、工具执行、多智能体协同、生态治理七层架构，具备利用链路复合性、危害传导级联性、边界穿透性三大特征，传统漏洞分类与评级体系无法完全适配其风险评估需求~\cite{nvdCVE2025_68664,nvdCVE2026_40933}。
本小节基于2025—2026年公开披露的权威漏洞，以CVSS评分体系为核心分级依据，将漏洞划分为严重（$\ge$9.0）、高危（7.0$\sim$8.9）、中危（4.0$\sim$6.9）三个等级，同等级内按披露日期倒序排列，形成标准化智能体漏洞清单，主要发现如下(详见表~\ref{tab:agent_vulns})：

\begin{itemize}[tabitem]
 \item \textbf{严重漏洞聚焦权限突破，可全域接管系统与窃取数据。}CVSS≥9.0严重漏洞集中于MCP协议、序列化、WebSocket核心链路，普遍具备无认证RCE、容器逃逸、全域密钥读取能力，能够绕过云IAM、沙箱等传统安全防护体系，印证硬件‑工具调用层的安全短板远超模型层风险~\cite{nvdCVE2026_26015,nvdCVE2026_22252}。
 
 \item \textbf{开源生态与多智能体架构形成两大核心风险。}LangChain、Flowise、LibreChat等开源框架依靠亿级下载量带来全域传导风险；MCP协议下多智能体组件之间默认信任、缺少强制校验机制，极易出现横向越权、集群数据泄露，属于分布式智能体架构特有的底层缺陷~\cite{nvdCVE2025_59528,nvdCVE2026_30625}。
 
 \item \textbf{隐形攻击与存量滞后修复形成持久化威胁。}借助提示注入、参数污染实现的隐形载荷可以绕过传统特征检测，完成无告警持久入侵；大量云上部署的智能体实例长期不升级，漏洞利用载荷持续扩散，AI智能体漏洞生命周期显著长于常规软件漏洞~\cite{nvdCVE2026_30615,nvdCVE2026_34070}。
\end{itemize}

\begingroup
% \nopagebreak[4]
\enlargethispage{1cm}
\footnotesize
\centering

\DefTblrTemplate{conthead-text}{default}{（续表）}
\DefTblrTemplate{contfoot-text}{default}{续下页}

\begin{longtblr}[
  caption = {AI智能体主要安全漏洞汇总},
  label   = {tab:agent_vulns},
]{
  width   = \linewidth,
  colspec = {
    |Q[wd=2cm,l,m]
    |Q[wd=2.5cm,l,m]
    |Q[wd=2cm,l,m]
    |Q[wd=0.8cm,l,m]
    |Q[wd=2cm,l,m]
    |X[1,l,m]|
  },
  cells   = {l,m},
  colsep  = 3pt,
  row{1}  = {font=\bfseries},
  rowhead = 1,
  % 每行内容与上下边框之间各保留少量空间；由该行最高单元格决定整行高度。
  rowsep  = 5pt,
  measure = vbox,
  stretch = -1,
}
\hline
\textbf{发现精确日期} & \textbf{漏洞名称/事件} & \textbf{影响产品/框架} & \textbf{CV-\\SS} & \textbf{CVE/编号} & \textbf{漏洞影响简述} \\
\hline
% \caption{AI智能体主要安全漏洞}
% \label{tab:agent_vulns} \\
% \hline
% \textbf{\shortstack{发现精确\\日期}} &
% \textbf{\shortstack{漏洞名称\\事件}} &
% \textbf{\shortstack{影响产品\\框架}} &
% \textbf{\shortstack{CVSS\\评分}} &
% \textbf{\shortstack{CVE\\编号}} &
% \centering\arraybackslash\textbf{漏洞影响简述} \\
% \hline
% \endfirsthead
% % 跨页重复表头
% \hline
% \textbf{\shortstack{发现精确\\日期}} &
% \textbf{\shortstack{漏洞名称\\事件}} &
% \textbf{\shortstack{影响产品\\框架}} &
% \textbf{\shortstack{CVSS\\评分}} &
% \textbf{\shortstack{CVE\\编号}} &
% \centering\arraybackslash\textbf{漏洞影响简述} \\
% \hline
% \endhead


\SetCell[c=6]{halign=c,valign=m,font=\bfseries}
  Critical（严重，CVSS $\geq 9.0$）
& & & & & \\
\hline


2026-05-06 & OpenClaw OpenShell沙箱文件写入 TOCTOU / symlink swap~\cite{nvdCVE2026_44112} & OpenClaw $<$ 2026.4.22 & 9.6 & CVE-2026-44112 &
\begin{itemize}[leftmargin=*,topsep=0pt,partopsep=0pt,itemsep=1pt,parsep=0pt]
\item OpenShell沙箱文件系统写入存在time-of-check/time-of-use竞态。
\item 攻击者可通过symlink swap将写入重定向到预期mount root之外。
\item 应升级到2026.4.22或包含修复提交的版本。
\end{itemize}
\\
\hline

2026-04-29 & DocsGPT MCP test绕过导致远程代码执行~\cite{nvdCVE2026_26015} & DocsGPT $\geq$ 0.15.0, $<$ 0.16.0 & 9.8 & CVE-2026-26015 &
\begin{itemize}[leftmargin=*,topsep=0pt,partopsep=0pt,itemsep=1pt,parsep=0pt]
\item 攻击者可构造恶意payload绕过“MCP test”行为。
\item 可在官方DocsGPT网站或本地/公开部署中触发任意远程代码执行。
\item 漏洞已在DocsGPT 0.16.0修复。
\end{itemize}
\\
\hline

2026-04-21 & Flowise MCP stdio命令注入/allowlist绕过~\cite{nvdCVE2026_40933} & Flowise/\allowbreak flowise-components $<$ 3.1.0 & 9.9 & CVE-2026-40933 &
\begin{itemize}[leftmargin=*,topsep=0pt,partopsep=0pt,itemsep=1pt,parsep=0pt]
\item 认证攻击者可在Custom MCP配置中添加stdio MCP server并指定任意命令。
\item 既有npx等允许命令可被组合危险参数绕过输入校验，实现OS命令执行。
\item 漏洞已在Flowise 3.1.0修复。
\end{itemize}
\\
\hline

2026-04-15 & Upsonic MCP task命令注入~\cite{nvdCVE2026_30625} & Upsonic 0.71.6 & 9.8 & CVE-2026-30625 &
\begin{itemize}[leftmargin=*,topsep=0pt,partopsep=0pt,itemsep=1pt,parsep=0pt]
\item MCP server/task 创建功能允许用户定义command和args。
\item npm、npx等allowlist命令可通过参数触发任意OS命令执行。
\item 0.72.0起增加Stdio server可直接执行本机命令的风险提示。
\end{itemize}
\\
\hline

2026-04-07 & Langflow未认证代码注入/RCE~\cite{nvdCVE2025_3248} & Langflow $<$ 1.3.0 & 9.8 & CVE-2025-3248 &
\begin{itemize}[leftmargin=*,topsep=0pt,partopsep=0pt,itemsep=1pt,parsep=0pt]
\item \texttt{/api/v1/validate/code}端点存在代码注入漏洞。
\item 远程未认证攻击者可发送构造HTTP请求执行任意代码。
\item 该漏洞已进入CISA KEV目录，建议升级到1.3.0或更高版本。
\end{itemize}
\\
\hline

2026-04-03 & PraisonAI Gateway未授权WebSocket/info访问~\cite{nvdCVE2026_34952} & PraisonAI Gateway $<$ 4.5.97 & 9.1 & CVE-2026-34952 &
\begin{itemize}[leftmargin=*,topsep=0pt,partopsep=0pt,itemsep=1pt,parsep=0pt]
\item Gateway server 在\texttt{/ws}接受WebSocket连接且\texttt{/info}暴露agent topology，均缺少认证。
\item 任意网络客户端可枚举已注册agents，并向agents及其tool sets发送任意消息。
\item 漏洞已在4.5.97修复，应升级并对网关入口加认证与网络隔离。
\end{itemize}
\\
\hline

2026-04-02 & Azure MCP Server关键功能缺失认证~\cite{nvdCVE2026_32211} & Azure MCP Server / Azure Web Apps hosted service & 9.1 & CVE-2026-32211 &
\begin{itemize}[leftmargin=*,topsep=0pt,partopsep=0pt,itemsep=1pt,parsep=0pt]
\item 关键功能缺少认证，未授权攻击者可经网络访问受影响服务。
\item 可造成信息泄露；Microsoft CNA评分还计入完整性影响。
\item 以Microsoft MSRC / NVD记录为准，CVSS 9.1为Microsoft CNA评分。
\end{itemize}
\\
\hline

2026-01-12 & LibreChat MCP stdio任意命令执行~\cite{nvdCVE2026_22252} & LibreChat $<$ v0.8.2-rc2 & 9.9 & CVE-2026-22252 &
\begin{itemize}[leftmargin=*,topsep=0pt,partopsep=0pt,itemsep=1pt,parsep=0pt]
\item MCP stdio transport接受任意命令且缺少有效校验。
\item 任意认证用户可通过单个API请求在容器内以root执行shell命令。
\item 漏洞已在v0.8.2-rc2修复。
\end{itemize}
\\
\hline

2025-12-23 & LangChain dumps/dumpd序列化注入~\cite{nvdCVE2025_68664} & LangChain $<$ 0.3.81 and $<$ 1.2.5 & 9.3 & CVE-2025-68664 &
\begin{itemize}[leftmargin=*,topsep=0pt,partopsep=0pt,itemsep=1pt,parsep=0pt]
\item \texttt{dumps()}/ \texttt{dumpd()}未转义自由字典中的\texttt{lc} 键结构。
\item 用户可控数据在反序列化时可能被当作LangChain对象处理，而非普通数据。
\item 漏洞已在0.3.81和1.2.5修复，应避免反序列化不可信配置。
\end{itemize}
\\
\hline

2025-12-17 & MCP server git复合漏洞~\cite{nvdCVE2025_68143,nvdCVE2025_68144,nvdCVE2025_68145} & Model Context Protocol Servers/\allowbreak\texttt{mcp-server-git} & 8.8/\allowbreak7.1/\allowbreak9.1 & CVE-2025-68143/4/5 &
\begin{itemize}[leftmargin=*,topsep=0pt,partopsep=0pt,itemsep=1pt,parsep=0pt]
\item \texttt{git\_init} 可在任意可访问路径创建仓库，使后续Git操作越出预期范围。
\item \texttt{git\_diff} / \texttt{git\_checkout} 对用户参数缺少校验，flag-like参数可导致任意文件覆盖。
\item \texttt{--repository} 限制下未校验后续 \texttt{repo\_path} 是否仍在允许仓库内，可操作其他可访问仓库。
\end{itemize}
\\
\hline

2025-09-22 & Flowise CustomMCP 远程代码执行~\cite{nvdCVE2025_59528} & Flowise 3.0.5 & 10.0 & CVE-2025-59528 &
\begin{itemize}[leftmargin=*,topsep=0pt,partopsep=0pt,itemsep=1pt,parsep=0pt]
\item CustomMCP节点解析\texttt{mcpServerConfig}时将用户输入传入。\texttt{Function()}构造器执行
\item 代码在Node.js 运行时权限下执行，可访问\texttt{child\_process}、。\texttt{fs}等危险模块
\item 漏洞已在 Flowise 3.0.6 修复。
\end{itemize}
\\
\hline

2025-07-09 & mcp-remote OAuth authorization endpoint命令注入~\cite{nvdCVE2025_6514} & \texttt{mcp-remote} npm package 0.0.5--0.1.15 & 9.6 & CVE-2025-6514 &
\begin{itemize}[leftmargin=*,topsep=0pt,partopsep=0pt,itemsep=1pt,parsep=0pt]
\item 连接不可信MCP server时，authorization endpoint 响应URL中的构造输入可触发OS命令注入。
\item 攻击需要用户交互但无需攻击者具备权限，成功后影响机密性、完整性和可用性。
\item 已通过\texttt{mcp-remote} 修复提交处置，应升级到不受影响版本。
\end{itemize}
\\
\hline

2025-06-13 & MCP Inspector缺失认证导致远程代码执行~\cite{nvdCVE2025_49596} & MCP Inspector $<$ 0.14.1 & 9.4 & CVE-2025-49596 &
\begin{itemize}[leftmargin=*,topsep=0pt,partopsep=0pt,itemsep=1pt,parsep=0pt]
\item Inspector client 与 proxy之间缺少认证。
\item 未认证请求可经proxy启动MCP stdio命令，造成远程代码执行。
\item 漏洞已在MCPInspector 0.14.1修复。
\end{itemize}
\\
\hline

2025-06-11 & EchoLeak/M365 Copilot AI command injection~\cite{nvdCVE2025_32711} & Microsoft 365 Copilot & 9.3 & CVE-2025-32711 &
\begin{itemize}[leftmargin=*,topsep=0pt,partopsep=0pt,itemsep=1pt,parsep=0pt]
\item M365 Copilot存在AI command injection，未授权攻击者可经网络触发信息泄露。
\item 该漏洞被公开命名为 EchoLeak，NVD 同时引用Microsoft MSRC与Aim Security页面。
\item CVSS 9.3为Microsoft CNA评分；NVD自身评分为7.5。
\end{itemize}
\\
\hline

% \SetCell[c=6]{c,m} \textbf{High（高危，7.0 $\leq$ CVSS $\leq$ 8.9）} & & & & & \\
% \hline

\SetCell[c=6]{halign=c,valign=m,font=\bfseries}
  High（高危，7.0 $\leq$ CVSS $\leq$ 8.9）
& & & & & \\
\hline

2026-05-08 & PraisonAI legacy Flask API默认禁用认证~\cite{nvdCVE2026_44338} & PraisonAI 2.5.6 to $<$ 4.6.34 & 7.3 & CVE-2026-44338 &
\begin{itemize}[leftmargin=*,topsep=0pt,partopsep=0pt,itemsep=1pt,parsep=0pt]
\item legacy Flask API server默认禁用认证。
\item 任意可达调用方可访问\texttt{/agents}并通过\texttt{/chat}触发\texttt{agents.yaml}工作流。
\item 漏洞已在4.6.34修复，应升级并限制API暴露面。
\end{itemize}
\\
\hline

2026-04-15 & Agent Zero External MCP Servers远程代码执行~\cite{nvdCVE2026_30624} & Agent Zero 0.9.8 & 8.6 & CVE-2026-30624 &
\begin{itemize}[leftmargin=*,topsep=0pt,partopsep=0pt,itemsep=1pt,parsep=0pt]
\item External MCP Servers配置允许通过JSON指定command和args。
\item 配置应用时这些值会被执行，缺少充分校验或限制。
\item 攻击者可提交恶意MCP配置，以Agent Zero进程权限执行OS命令。
\end{itemize}
\\
\hline

2026-04-15 & Windsurf Prompt Injection to MCP配置命令执行~\cite{nvdCVE2026_30615} & Windsurf 1.9544.26 & 8.0 & CVE-2026-30615 &
\begin{itemize}[leftmargin=*,topsep=0pt,partopsep=0pt,itemsep=1pt,parsep=0pt]
\item 处理攻击者控制的HTML内容时，提示注入可修改本地MCP配置。
\item 可自动注册恶意MCP STDIO server并在无进一步交互下执行任意命令。
\item 成功利用可持久化恶意MC 配置并访问应用暴露的敏感信息。
\end{itemize}
\\
\hline

2026-03-30 & LangChain prompt loading路径遍历~\cite{nvdCVE2026_34070} & LangChain / langchain-core $<$ 1.2.22 & 7.5 & CVE-2026-34070 &
\begin{itemize}[leftmargin=*,topsep=0pt,partopsep=0pt,itemsep=1pt,parsep=0pt]
\item \texttt{langchain\_core.prompts.loading}多个函数未充分校验配置字典中的路径。
\item 当应用把用户可影响的prompt配置传入\texttt{load\_prompt()}或 \texttt{load\_prompt\_from\_config()} 时，攻击者可读取主机文件。
\item 读取受模板/示例文件扩展名检查限制；漏洞已在1.2.22修复。
\end{itemize}
\\
\hline

2026-01-09 & WeKnora MCP stdio 命令注入~\cite{nvdCVE2026_22688} & Tencent WeKnora $<$ 0.2.5 & 8.8 & CVE-2026-22688 &
\begin{itemize}[leftmargin=*,topsep=0pt,partopsep=0pt,itemsep=1pt,parsep=0pt]
\item 认证用户可向MCP stdio设置注入\texttt{stdio\_config.command} / \texttt{args}。
\item 服务端会基于注入值启动subprocess，造成命令执行。
\item 漏洞已在WeKnora 0.2.5修复。
\end{itemize}
\\
\hline

2025-11-07 & Open WebUI Direct Connections SSE代码注入~\cite{nvdCVE2025_64496} & Open WebUI $<$ 0.6.35 & 8.0 & CVE-2025-64496 &
\begin{itemize}[leftmargin=*,topsep=0pt,partopsep=0pt,itemsep=1pt,parsep=0pt]
\item Direct Connections功能允许恶意外部模型服务器通过SSE execute events在受害者浏览器执行JavaScript。
\item 可导致认证token窃取和账户接管；与 Functions API链接时可造成后端RCE。
\item 攻击需受害者启用Direct Connections并添加恶意模型URL；0.6.35已修复。
\end{itemize}
\\
\hline


\end{longtblr}
\endgroup

% \begin{center}
% % \nopagebreak[4]
% \enlargethispage{1cm}
% \footnotesize
% \setlength{\tabcolsep}{3pt}
% % \renewcommand{\arraystretch}{0.7}{0.7}
% \begin{longtable}{|>{\rule{0pt}{2pt}}m{2cm}
% |>{\rule{0pt}{2pt}}m{2.5cm}
% |>{\rule{0pt}{2pt}}m{2cm}
% |>{\rule{0pt}{2pt}}m{0.8cm}
% |>{\rule{0pt}{2pt}}m{2cm}
% |>{\rule{0pt}{2pt}}m{\dimexpr\linewidth - 2cm - 2.5cm - 2cm - 0.8cm - 2cm - 12\tabcolsep\relax}|}
% \caption{AI智能体主要安全漏洞汇总}\label{tab:agent_vulns}\\
% \hline
% \textbf{发现精确日期} & \textbf{漏洞名称/事件} & \textbf{影响产品/框架} & \textbf{CVSS} & \textbf{CVE/编号} & \textbf{漏洞影响简述} \\
% \hline
% \endfirsthead

% \multicolumn{6}{l}{\tablename\ \thetable\ -- 续上页}\\
% \hline
% \textbf{发现精确日期} & \textbf{漏洞名称/事件} & \textbf{影响产品/框架} & \textbf{CVSS} & \textbf{CVE/编号} & \textbf{漏洞影响简述} \\
% \hline
% \endhead

% \hline
% \multicolumn{6}{r}{续下页}\\
% \endfoot
% \hline
% \endlastfoot

% % \caption{AI智能体主要安全漏洞}
% % \label{tab:agent_vulns} \\
% % \hline
% % \textbf{\shortstack{发现精确\\日期}} &
% % \textbf{\shortstack{漏洞名称\\事件}} &
% % \textbf{\shortstack{影响产品\\框架}} &
% % \textbf{\shortstack{CVSS\\评分}} &
% % \textbf{\shortstack{CVE\\编号}} &
% % \centering\arraybackslash\textbf{漏洞影响简述} \\
% % \hline
% % \endfirsthead
% % % 跨页重复表头
% % \hline
% % \textbf{\shortstack{发现精确\\日期}} &
% % \textbf{\shortstack{漏洞名称\\事件}} &
% % \textbf{\shortstack{影响产品\\框架}} &
% % \textbf{\shortstack{CVSS\\评分}} &
% % \textbf{\shortstack{CVE\\编号}} &
% % \centering\arraybackslash\textbf{漏洞影响简述} \\
% % \hline
% % \endhead

% \multicolumn{6}{|c|}{\textbf{Critical（严重，CVSS $\geq$ 9.0）}} \\
% \hline

% 2026-05-06 & OpenClaw OpenShell沙箱文件写入 TOCTOU / symlink swap~\cite{nvdCVE2026_44112} & OpenClaw $<$ 2026.4.22 & 9.6 & CVE-2026-44112 &
% \begin{itemize}[tabitem]
% \item OpenShell沙箱文件系统写入存在time-of-check/time-of-use竞态。
% \item 攻击者可通过symlink swap将写入重定向到预期mount root之外。
% \item 应升级到2026.4.22或包含修复提交的版本。
% \end{itemize}
% \\
% \hline

% 2026-04-29 & DocsGPT MCP test绕过导致远程代码执行~\cite{nvdCVE2026_26015} & DocsGPT $\geq$ 0.15.0, $<$ 0.16.0 & 9.8 & CVE-2026-26015 &
% \begin{itemize}[tabitem]
% \item 攻击者可构造恶意payload绕过“MCP test”行为。
% \item 可在官方DocsGPT网站或本地/公开部署中触发任意远程代码执行。
% \item 漏洞已在DocsGPT 0.16.0修复。
% \end{itemize}
% \\
% \hline

% 2026-04-21 & Flowise MCP stdio命令注入/allowlist绕过~\cite{nvdCVE2026_40933} & Flowise/flowise-components $<$ 3.1.0 & 9.9 & CVE-2026-40933 &
% \begin{itemize}[tabitem]
% \item 认证攻击者可在Custom MCP配置中添加stdio MCP server并指定任意命令。
% \item 既有npx等允许命令可被组合危险参数绕过输入校验，实现OS命令执行。
% \item 漏洞已在Flowise 3.1.0修复。
% \end{itemize}
% \\
% \hline

% 2026-04-15 & Upsonic MCP task命令注入~\cite{nvdCVE2026_30625} & Upsonic 0.71.6 & 9.8 & CVE-2026-30625 &
% \begin{itemize}[tabitem]
% \item MCP server/task 创建功能允许用户定义command和args。
% \item npm、npx等allowlist命令可通过参数触发任意OS命令执行。
% \item 0.72.0起增加Stdio server可直接执行本机命令的风险提示。
% \end{itemize}
% \\
% \hline

% 2026-04-07 & Langflow未认证代码注入/RCE~\cite{nvdCVE2025_3248} & Langflow $<$ 1.3.0 & 9.8 & CVE-2025-3248 &
% \begin{itemize}[tabitem]
% \item \texttt{/api/v1/validate/code}端点存在代码注入漏洞。
% \item 远程未认证攻击者可发送构造HTTP请求执行任意代码。
% \item 该漏洞已进入CISA KEV目录，建议升级到1.3.0或更高版本。
% \end{itemize}
% \\
% \hline

% 2026-04-03 & PraisonAI Gateway未授权WebSocket/info访问~\cite{nvdCVE2026_34952} & PraisonAI Gateway $<$ 4.5.97 & 9.1 & CVE-2026-34952 &
% \begin{itemize}[tabitem]
% \item Gateway server 在\texttt{/ws}接受WebSocket连接且\texttt{/info}暴露agent topology，均缺少认证。
% \item 任意网络客户端可枚举已注册agents，并向agents及其tool sets发送任意消息。
% \item 漏洞已在4.5.97修复，应升级并对网关入口加认证与网络隔离。
% \end{itemize}
% \\
% \hline

% 2026-04-02 & Azure MCP Server关键功能缺失认证~\cite{nvdCVE2026_32211} & Azure MCP Server / Azure Web Apps hosted service & 9.1 & CVE-2026-32211 &
% \begin{itemize}[tabitem]
% \item 关键功能缺少认证，未授权攻击者可经网络访问受影响服务。
% \item 可造成信息泄露；Microsoft CNA评分还计入完整性影响。
% \item 以Microsoft MSRC / NVD记录为准，CVSS 9.1为Microsoft CNA评分。
% \end{itemize}
% \\
% \hline

% 2026-01-12 & LibreChat MCP stdio任意命令执行~\cite{nvdCVE2026_22252} & LibreChat $<$ v0.8.2-rc2 & 9.9 & CVE-2026-22252 &
% \begin{itemize}[tabitem]
% \item MCP stdio transport接受任意命令且缺少有效校验。
% \item 任意认证用户可通过单个API请求在容器内以root执行shell命令。
% \item 漏洞已在v0.8.2-rc2修复。
% \end{itemize}
% \\
% \hline

% 2025-12-23 & LangChain dumps/dumpd序列化注入~\cite{nvdCVE2025_68664} & LangChain $<$ 0.3.81 and $<$ 1.2.5 & 9.3 & CVE-2025-68664 &
% \begin{itemize}[tabitem]
% \item \texttt{dumps()}/ \texttt{dumpd()}未转义自由字典中的\texttt{lc} 键结构。
% \item 用户可控数据在反序列化时可能被当作LangChain对象处理，而非普通数据。
% \item 漏洞已在0.3.81和1.2.5修复，应避免反序列化不可信配置。
% \end{itemize}
% \\
% \hline

% 2025-12-17 & MCP server git复合漏洞~\cite{nvdCVE2025_68143,nvdCVE2025_68144,nvdCVE2025_68145} & Model Context Protocol Servers/\texttt{mcp-server-git} & 8.8/7.1/9.1 & CVE-2025-68143/4/5 &
% \begin{itemize}[tabitem]
% \item \texttt{git\_init} 可在任意可访问路径创建仓库，使后续Git操作越出预期范围。
% \item \texttt{git\_diff} / \texttt{git\_checkout} 对用户参数缺少校验，flag-like参数可导致任意文件覆盖。
% \item \texttt{--repository} 限制下未校验后续 \texttt{repo\_path} 是否仍在允许仓库内，可操作其他可访问仓库。
% \end{itemize}
% \\
% \hline

% 2025-09-22 & Flowise CustomMCP 远程代码执行~\cite{nvdCVE2025_59528} & Flowise 3.0.5 & 10.0 & CVE-2025-59528 &
% \begin{itemize}[tabitem]
% \item CustomMCP节点解析\texttt{mcpServerConfig}时将用户输入传入。\texttt{Function()}构造器执行
% \item 代码在Node.js 运行时权限下执行，可访问\texttt{child\_process}、。\texttt{fs}等危险模块
% \item 漏洞已在 Flowise 3.0.6 修复。
% \end{itemize}
% \\
% \hline

% 2025-07-09 & mcp-remote OAuth authorization endpoint命令注入~\cite{nvdCVE2025_6514} & \texttt{mcp-remote} npm package 0.0.5--0.1.15 & 9.6 & CVE-2025-6514 &
% \begin{itemize}[tabitem]
% \item 连接不可信MCP server时，authorization endpoint 响应URL中的构造输入可触发OS命令注入。
% \item 攻击需要用户交互但无需攻击者具备权限，成功后影响机密性、完整性和可用性。
% \item 已通过\texttt{mcp-remote} 修复提交处置，应升级到不受影响版本。
% \end{itemize}
% \\
% \hline

% 2025-06-13 & MCP Inspector缺失认证导致远程代码执行~\cite{nvdCVE2025_49596} & MCP Inspector $<$ 0.14.1 & 9.4 & CVE-2025-49596 &
% \begin{itemize}[tabitem]
% \item Inspector client 与 proxy之间缺少认证。
% \item 未认证请求可经proxy启动MCP stdio命令，造成远程代码执行。
% \item 漏洞已在MCPInspector 0.14.1修复。
% \end{itemize}
% \\
% \hline

% 2025-06-11 & EchoLeak/M365 Copilot AI command injection~\cite{nvdCVE2025_32711} & Microsoft 365 Copilot & 9.3 & CVE-2025-32711 &
% \begin{itemize}[tabitem]
% \item M365 Copilot存在AI command injection，未授权攻击者可经网络触发信息泄露。
% \item 该漏洞被公开命名为 EchoLeak，NVD 同时引用Microsoft MSRC与Aim Security页面。
% \item CVSS 9.3为Microsoft CNA评分；NVD自身评分为7.5。
% \end{itemize}
% \\
% \hline

% \multicolumn{6}{|c|}{\textbf{High（高危，7.0 $\leq$ CVSS $\leq$ 8.9）}} \\
% \hline

% 2026-05-08 & PraisonAI legacy Flask API默认禁用认证~\cite{nvdCVE2026_44338} & PraisonAI 2.5.6 to $<$ 4.6.34 & 7.3 & CVE-2026-44338 &
% \begin{itemize}[tabitem]
% \item legacy Flask API server默认禁用认证。
% \item 任意可达调用方可访问\texttt{/agents}并通过\texttt{/chat}触发\texttt{agents.yaml}工作流。
% \item 漏洞已在4.6.34修复，应升级并限制API暴露面。
% \end{itemize}
% \\
% \hline

% 2026-04-15 & Agent Zero External MCP Servers远程代码执行~\cite{nvdCVE2026_30624} & Agent Zero 0.9.8 & 8.6 & CVE-2026-30624 &
% \begin{itemize}[tabitem]
% \item External MCP Servers配置允许通过JSON指定command和args。
% \item 配置应用时这些值会被执行，缺少充分校验或限制。
% \item 攻击者可提交恶意MCP配置，以Agent Zero进程权限执行OS命令。
% \end{itemize}
% \\
% \hline

% 2026-04-15 & Windsurf Prompt Injection to MCP配置命令执行~\cite{nvdCVE2026_30615} & Windsurf 1.9544.26 & 8.0 & CVE-2026-30615 &
% \begin{itemize}[tabitem]
% \item 处理攻击者控制的HTML内容时，提示注入可修改本地MCP配置。
% \item 可自动注册恶意MCP STDIO server并在无进一步交互下执行任意命令。
% \item 成功利用可持久化恶意MC 配置并访问应用暴露的敏感信息。
% \end{itemize}
% \\
% \hline

% 2026-03-30 & LangChain prompt loading路径遍历~\cite{nvdCVE2026_34070} & LangChain / langchain-core $<$ 1.2.22 & 7.5 & CVE-2026-34070 &
% \begin{itemize}[tabitem]
% \item \texttt{langchain\_core.prompts.loading}多个函数未充分校验配置字典中的路径。
% \item 当应用把用户可影响的prompt配置传入\texttt{load\_prompt()}或 \texttt{load\_prompt\_from\_config()} 时，攻击者可读取主机文件。
% \item 读取受模板/示例文件扩展名检查限制；漏洞已在1.2.22修复。
% \end{itemize}
% \\
% \hline

% 2026-01-09 & WeKnora MCP stdio 命令注入~\cite{nvdCVE2026_22688} & Tencent WeKnora $<$ 0.2.5 & 8.8 & CVE-2026-22688 &
% \begin{itemize}[tabitem]
% \item 认证用户可向MCP stdio设置注入\texttt{stdio\_config.command} / \texttt{args}。
% \item 服务端会基于注入值启动subprocess，造成命令执行。
% \item 漏洞已在WeKnora 0.2.5修复。
% \end{itemize}
% \\
% \hline

% 2025-11-07 & Open WebUI Direct Connections SSE代码注入~\cite{nvdCVE2025_64496} & Open WebUI $<$ 0.6.35 & 8.0 & CVE-2025-64496 &
% \begin{itemize}[tabitem]
% \item Direct Connections功能允许恶意外部模型服务器通过SSE execute events在受害者浏览器执行JavaScript。
% \item 可导致认证token窃取和账户接管；与 Functions API链接时可造成后端RCE。
% \item 攻击需受害者启用Direct Connections并添加恶意模型URL；0.6.35已修复。
% \end{itemize}
% \\
% \hline

% \multicolumn{6}{|c|}{\textbf{Medium（中危，4.0 $\leq$ CVSS $\leq$ 6.9）}} \\
% \hline

% \multicolumn{6}{|c|}{本表核验条目无中危漏洞。} 

% \end{longtable}
% \end{center}





% \begin{center}
% % \nopagebreak[4]
% \enlargethispage{1cm}
% \footnotesize
% \setlength{\tabcolsep}{3pt}
% % \renewcommand{\arraystretch}{0.7}{0.7}
% \begin{longtable}{|p{2cm}|p{2cm}|p{2cm}|p{0.5cm}|p{2cm}|p{\dimexpr\linewidth - 2cm - 2cm - 2cm - 0.5cm - 2cm - 12\tabcolsep\relax}|}
% \caption{AI智能体 主要安全漏洞（按严重级别降序、同等级发现日期倒序排列，2025–2026）}\label{tab:agent_vulns}\\
% \hline
% \textbf{发现精确日期} & \textbf{漏洞名称/事件} & \textbf{影响产品/框架} & \textbf{CVSS} & \textbf{CVE/编号} & \textbf{漏洞影响简述} \\
% \hline
% \endfirsthead

% \multicolumn{6}{l}{\tablename\ \thetable\ -- 续上页}\\
% \hline
% \textbf{发现精确日期} & \textbf{漏洞名称/事件} & \textbf{影响产品/框架} & \textbf{CVSS} & \textbf{CVE/编号} & \textbf{漏洞影响简述} \\
% \hline
% \endhead

% \hline
% \multicolumn{6}{r}{续下页}\\
% \endfoot

% \hline
% \endlastfoot

% \multicolumn{6}{c}{\textbf{Critical（严重，CVSS ≥9.0）}} \\
% \hline
% 2026-06-29 & Vertex MCP 无认证远程代码执行[参考文献] & GCP Vertex AI MCP & 9.7 & CVE-2026-46205 &
% \begin{itemize}[tabitem]
% \item MCP工具参数未过滤恶意脚本，未登录攻击者直接执行云主机命令
% \item 绕过云平台IAM权限体系，读取存储桶全部训练数据与模型权重
% \item 漏洞公开48小时出现批量扫描利用脚本，谷歌云上大量租户未修复
% \end{itemize}
% \\
% 2026-05-28 & Claw Chain（OpenClaw）[] & OpenClaw 框架 & 9.6 & CVE-2026-44112 &
% \begin{itemize}[tabitem]
% \item 沙箱边界校验逻辑缺失，实现容器完全逃逸
% \item 逃逸后获取主机root权限，横向提权接管集群全部节点
% \item 植入持久化后台后门，重启智能体仍可长期受控，全网24.5万实例暴露风险
% \end{itemize}
% \\
% 2026-04-17 & Azure MCP 认证绕过 & Azure MCP Server & 9.1 & CVE-2026-32211 &
% \begin{itemize}[tabitem]
% \item 会话校验逻辑存在逻辑漏洞，攻击者无需凭证即可接入服务
% \item 无限制读取租户存储内全部隐私、业务与密钥数据
% \item 厂商暂未发布官方修复补丁，云上存量实例持续存在可利用攻击面
% \end{itemize}
% \\
% 2026-04-09 & CamoLeak & GitHub Copilot Chat & 9.6 & 未公开 &
% \begin{itemize}[tabitem]
% \item PR隐藏注释嵌入隐形提示载荷，绕过前端内容过滤检测
% \item 依托Camo代理中转流量，静默窃取仓库全部源代码与密钥
% \item 攻击无交互、无日志告警，大量企业私有代码仓库遭批量泄露
% \end{itemize}
% \\
% 2026-02-10 & Moltbook 大规模劫持 & Moltbook 平台 & 9.8 & 未公开 &
% \begin{itemize}[tabitem]
% \item 后端数据库未加密裸存API访问凭证，批量窃取智能体密钥
% \item WebSocket通道零点击远程接管在线运行Agent实例
% \item 全网77万活跃智能体被批量劫持，用于挖矿、窃密与勒索分发
% \end{itemize}
% \\
% 2025-11-22 & LangChain 序列化注入 & LangChain & 9.3 & CVE-2025-68664 &
% \begin{itemize}[tabitem]
% \item dumps/loads反序列化无白名单校验，注入恶意执行对象
% \item 实现服务器任意远程代码执行，完全控制服务进程
% \item 读取、导出全部环境变量，包含API密钥、数据库连接凭证
% \end{itemize}
% \\
% 2025-08-14 & Codex CLI 命令注入 & OpenAI Codex CLI & 9.1 & CVE-2025-61260 &
% \begin{itemize}[tabitem]
% \item 用户本地配置文件字段未做过滤，可注入系统原生命令
% \item 依托智能体高权限上下文执行高危系统操作
% \item 本地密钥、ssh私钥等敏感文件可被攻击者完整读取外发
% \end{itemize}
% \\
% 2025-06-07 & EchoLeak & Microsoft 365 Copilot & 9.3 & CVE-2025-32711 &
% \begin{itemize}[tabitem]
% \item 邮件正文内嵌隐蔽提示，无需用户确认自动触发文件读取
% \item 跨会话持久留存恶意指令，持续遍历OneDrive、SharePoint文档
% \item 企业客户海量合同、员工隐私数据批量外泄至攻击者服务器
% \end{itemize}
% \\
% \hline
% \multicolumn{6}{c}{\textbf{High（高危，7.0 ≤ CVSS ≤8.9）}} \\
% \hline
% 2026-07-01 & CometAgent 会话持久化后门植入 & Perplexity CometAgent & 8.9 & CVE-2026-46827 &
% \begin{itemize}[tabitem]
% \item 本地日历载荷注入，持久写入后台会话存储长期潜伏
% \item 静默遍历本地密码管理器、文档目录批量窃取隐私文件
% \item 无弹窗、无日志告警，数月时间内持续上传用户敏感数据
% \end{itemize}
% \\
% 2026-06-30 & CrewAI 集群中间人劫持 & CrewAI Enterprise v0.22+ & 8.4 & CVE-2026-46358 &
% \begin{itemize}[tabitem]
% \item 多智能体通信通道未加密，中间人篡改任务调度指令
% \item 诱导高权限智能体导出全集群业务数据库与向量知识库
% \item 协同任务信任校验缺失，横向渗透所有关联业务智能体
% \end{itemize}
% \\
% 2026-06-02 & AutoGen 多智能体跨角色越权 & Microsoft AutoGen & 8.7 & CVE-2026-45109 &
% \begin{itemize}[tabitem]
% \item 多角色隔离逻辑失效，低权限智能体可调用管理员级工具
% \item 跨会话记忆共享漏洞，读取其他Agent私有上下文与凭证
% \item 可组合多工具链式调用，批量导出租户内部业务敏感数据
% \end{itemize}
% \\
% 2026-05-28 & Claude Code Hooks RCE & Anthropic Claude Code & 8.2 & GHSA-ph6w-f82w-28w6 &
% \begin{itemize}[tabitem]
% \item 恶意settings.json配置文件可注入自定义执行钩子脚本
% \item 智能体会话启动时自动加载并执行攻击者预设恶意代码
% \item 本地项目目录全部文件可读、可修改，泄露研发密钥与源码
% \end{itemize}
% \\
% 2026-05-16 & LangChain 路径遍历 & LangChain & 7.5 & CVE-2026-34070 &
% \begin{itemize}[tabitem]
% \item PromptLoading文件接口未过滤../跳转路径字符
% \item 攻击者读取服务器任意本地文件、配置与密钥文件
% \item 无法限制读取目录边界，容器内可逃逸读取宿主机敏感配置
% \end{itemize}
% \\
% 2026-05-03 & PraisonAI 认证缺失 & PraisonAI 2.5.6–4.6.33 & 7.3 & CVE-2026-44338 &
% \begin{itemize}[tabitem]
% \item 平台核心管控接口完全缺失身份认证与会话校验
% \item 漏洞公开3小时44分钟内出现成熟武器化利用脚本扩散
% \item 未授权攻击者可查看、修改、删除平台全部智能体任务配置
% \end{itemize}
% \\
% 2026-04-22 & PraisonAI 未授权WebSocket访问 & PraisonAI Gateway & 8.6 & CVE-2026-34952 &
% \begin{itemize}[tabitem]
% \item WebSocket网关无令牌校验，任意网络主体可建立持久连接
% \item 远程下发指令启停、篡改所有在线运行智能体执行逻辑
% \item 劫持对外API调用链路，拦截、篡改智能体对外输出数据
% \end{itemize}
% \\
% 2025-09-30 & Anthropic MCP Git 复合漏洞 & Anthropic Git MCP & 8.8 & CVE-2025-68143/4/5 &
% \begin{itemize}[tabitem]
% \item 输入参数过滤失效，支持系统命令注入执行高危操作
% \item 文件读取接口未限制路径，实现全仓库路径遍历读取源码
% \item 接口错误返回冗余环境信息，泄露仓库密钥、开发者凭证
% \end{itemize}
% \\
% \hline
% \multicolumn{6}{c}{\textbf{Medium（中危，4.0 ≤ CVSS ≤6.9）}} \\
% \hline
% 2026-07-03 & LlamaIndex RAG 语义隐式注入 & LlamaIndex 0.11.2 & 6.2 & CVE-2026-46945 &
% \begin{itemize}[tabitem]
% \item 向量检索相似度阈值校验缺陷，隐蔽恶意片段混入检索结果
% \item 无来源溯源标记，篡改知识库诱导智能体输出虚假业务结论
% \item 多租户场景跨库语义泄露，竞争对手可反向还原私有业务数据
% \end{itemize}
% \\
% 2026-06-01 & LlamaIndex 向量库越权检索 & LlamaIndex 0.10.x & 6.4 & CVE-2026-45011 &
% \begin{itemize}[tabitem]
% \item 租户向量空间无隔离，检索时可跨库读取其他用户私有嵌入数据
% \item 相似度检索无访问权限校验，批量匹配获取隐私语义片段
% \item 无法区分数据归属主体，多租户场景极易引发批量隐私泄露
% \end{itemize}
% \\
% 2026-03-15 & CrewAI 角色越权 & CrewAI 多智能体框架 & 5.9 & 未公开 &
% \begin{itemize}[tabitem]
% \item 角色权限边界隔离代码缺失，低权限Agent访问管理员上下文
% \item 跨任务记忆无访问控制，窃取其他任务内存储的业务密钥
% \item 多智能体协同任务中可篡改队友输出结果，伪造合规业务数据
% \end{itemize}
% \\
% 2025-12-08 & Meta LiteLLM 供应链投毒 & Meta/Mercor & 6.8 & 未公开 &
% \begin{itemize}[tabitem]
% \item 第三方依赖包被恶意污染，内置静默数据上传后门
% \item 启动时自动导出训练数据集、模型权重与API访问密钥
% \item 依赖广泛被智能体项目引用，形成大范围供应链连锁风险
% \end{itemize}
% \\
% \end{longtable}
% \end{center}


% \begin{table}[htbp]
% \centering
% \renewcommand{\arraystretch}{0.7}
% \caption{OpenClaw关键里程碑与安全事件全表（2025年10月—2026年7月，时间倒序）}
% \label{tab:agent_full_timeline}
% \footnotesize
% \begin{tabularx}{\linewidth}{|l|l| X|}
% \hline
% \textbf{时间} & \textbf{类别} & \textbf{完整事件说明} \\
% \hline
% \multicolumn{3}{c}{\textbf{2026年7月（监管收紧、高危漏洞集中披露）}} \\
% % 2026-07-05 & 技术产品 & Anthropic发布Claude Sonnet 5，新一代原生Agent大模型，强化多智能体协同与沙箱隔离能力，对标OpenClaw企业场景需求 \\
% % 2026-07-03 & 政策合规 & 网信办《互联网信息服务管理办法（修订草案征求意见稿）》公开征求意见，增设智能体专项监管章节，明确Agent安全准入、溯源与数据管控强制要求 \\
% % 2026-07-02 & 安全漏洞 & AgentRoot全局权限劫持漏洞（CVE-2026-46891，CVSS 9.9），OpenClaw v2.8集群管理接口无会话校验，30万+公网实例存在被一键接管风险，可批量导出向量库与API密钥 \\
% % 2026-07-01 & 行业评测 & 中国信通院发布AI Safety Benchmark 2026 Q2 Claw类智能体安全基准，出台《OpenClaw类智能体部署风险管理指南》，针对权限、记忆、插件三大风险给出标准化防护方案 \\
% % 2026-07-01 & 安全漏洞 & CometAgent会话持久化后门植入漏洞（CVE-2026-46827，CVSS 8.9），日历载荷注入实现长期静默窃取本地文档与密码库，跨平台兼容OpenClaw衍生应用 \\
% % 2026-07-01 & 产品更新 & OpenClaw v2.8.1紧急安全补丁发布，修复AgentRoot高危漏洞，新增集群会话绑定、插件来源白名单、显存强制清零三大安全机制 \\
% 2026-06-30 & 产业合规 & 字节豆包、阿里通义千问公告7月15日下线用户自建UGC智能体，对接《AI拟人化互动服务管理暂行办法》强监管要求，行业开启智能体合规整改周期[?] \\
% \hline
% \multicolumn{3}{c}{\textbf{2026年6月（云厂商Agent漏洞爆发、框架安全升级）}} \\
% 2026-06-29 & 安全漏洞 & GCP Vertex MCP无认证远程代码执行漏洞（CVE-2026-46205，CVSS 9.7），云原生智能体工具参数过滤失效，可绕过IAM读取全部训练数据集与模型权重[?] \\
% 2026-06-30 & 安全漏洞 & CrewAI集群中间人劫持漏洞（CVE-2026-46358，CVSS 8.4），多智能体通信未加密，中间人篡改调度指令横向渗透集群，同类通信缺陷在OpenClaw早期分支版本同步存在 \\
% 2026-06-02 & 安全漏洞 & Microsoft AutoGen多智能体跨角色越权漏洞（CVE-2026-45109，CVSS 8.7），角色隔离逻辑失效，低权限Agent可调用管理员级工具，与OpenClaw多角色权限缺陷同源 \\
% 2026-06-01 & 安全漏洞 & LlamaIndex向量库越权检索漏洞（CVE-2026-45011，CVSS 6.4），多租户向量空间无隔离，OpenClaw配套RAG组件受同款漏洞影响引发跨用户隐私泄露 \\
% 2026-06-01 & 产品更新 & OpenClaw v2.7正式发布，重构ClawHub插件市场供应链校验体系，新增第三方技能静态扫描、恶意载荷自动拦截功能，缓解插件投毒风险 \\
% \hline
% \multicolumn{3}{c}{\textbf{2026年5月（关键爆发期）}} \\
% 2026-05-27 & 安全漏洞 & Claw Chain（OpenClaw）漏洞公开（CVE-2026-44112，CVSS 9.6），沙箱逃逸+本地提权+持久化控制，24.5万实例暴露，框架级0day武器化 \\
% 2026-05-26 & 政策合规 & 我国《智能体规范应用与创新发展实施意见》印发，国家级Agent专项政策落地，明确场景与安全红线 \\
% 2026-05-21 & 产业格局 & Hermes Agent日活Token达2910亿，超越OpenClaw成为全球第一大Agent应用 \\
% 2026-05-20 & 技术产品 & 谷歌I/O大会发布Gemini 3.5，全链路Agent化，月活9亿，Agent成为系统级基础设施 \\
% 2026-05-19 & 产品更新 & OpenClaw发布重大版本更新，包含100+项优化，强化容器化与权限控制 \\
% 2026-05-14 & 国内战略 & 百度Create 2026确立“智能体优先”战略，提出DAA（日活智能体数）行业指标 \\
% 2026-05-11 & 产品发布 & 火山引擎发布“Agent Plan”全模态智能体套餐，升级为企业级Agent平台 \\
% 2026-05-08 & 企业应用 & 字节跳动发布“字节智助”办公Agent，国内首个大规模企业办公智能体 \\
% 2026-05-01 & 生态产品 & 微软发布Agent 365，Agent全面嵌入Office套件，企业级权限与身份体系成型 \\
% \hline
% \multicolumn{3}{c}{\textbf{2026年4月（高危漏洞集中爆发）}} \\
% 2026-04-28 & 安全漏洞 & CamoLeak漏洞爆发，GitHub Copilot Chat遭注入攻击，静默窃取代码仓库数据 \\
% 2026-04-21 & 安全漏洞 & OpenClaw沙箱绕过提权漏洞（CVE-2026-41329，CVSS 9.9） \\
% 2026-04-15 & 安全漏洞 & Azure MCP认证绕过漏洞（CVE-2026-32211，CVSS 9.1），无凭证访问敏感数据 \\
% 2026-04-10 & 安全漏洞 & PraisonAI Gateway未授权访问（CVE-2026-34952，CVSS 8.6），可远程控制所有Agent \\
% \hline
% \multicolumn{3}{c}{\textbf{2026年3月（多Agent框架成熟、国内多所高校发布OpenClaw风险预警）}} \\
% 2026-03-18 & 安全漏洞 & CrewAI角色越权漏洞，多Agent权限隔离失效，跨角色访问敏感上下文 \\
% 2026-03-17 & 风险预警 & 多所高校网信办同步发布OpenClaw安全风险提示，工信部NVDB监测累计82个相关漏洞，12个超危、21个高危漏洞公开披露 \\
% 2026-03-05 & 技术产品 & GPT-4o多模态Agent正式商用，实时视听理解+工具调用一体化 \\
% \hline
% \multicolumn{3}{c}{\textbf{2026年2月（平台级大规模劫持）}} \\
% 2026-02-14 & 安全事件 & Moltbook 77万Agent被大规模劫持，史上最大规模Agent安全事件，底层基于OpenClaw衍生框架，数据库未启用RLS行级安全策略导致密钥明文泄露 \\
% \hline
% \multicolumn{3}{c}{\textbf{2026年1月（OpenClaw正式定名）}} \\
% 2026-01-30 & 产品更名 & Moltbot 正式更名为 OpenClaw（龙虾），确立开源AI智能体标杆定位 \\
% 2026-01-27 & 产品更名 & Clawdbot 因 Anthropic 商标投诉更名为 Moltbot \\
% 2026-01-05 & 产品发布 & Peter Steinberger 发布 Clawdbot（OpenClaw前身），开源行动型AI智能体 \\
% \hline
% \multicolumn{3}{c}{\textbf{2025年12月（模型能力跃迁）}} \\
% 2025-12-18 & 技术产品 & 谷歌Gemini 3 Flash发布，推理成本下降90\%，支撑端侧大规模部署 \\
% 2025-12-11 & 技术产品 & 谷歌Gemini Deep Research Agent发布，HLE准确率46.4\%，超越GPT-5 Pro \\
% 2025-12-05 & 技术产品 & OpenAI发布GPT-5.2+ChatGPT Agent，支持100+工具并行调用 \\
% 2025-12-01 & 安全事件 & Meta LiteLLM供应链投毒事件，依赖污染导致密钥泄露，OpenClaw早期版本内置该依赖存在同款风险 \\
% \hline
% \multicolumn{3}{c}{\textbf{2025年11月（巨头全面入场，OpenClaw起源）}} \\
% 2025-11-28 & 安全漏洞 & LangChain序列化注入漏洞（CVE-2025-68664，CVSS 9.3），可远程代码执行，OpenClaw工具调用模块复用同类序列化逻辑存在攻击面 \\
% 2025-11-18 & 技术产品 & 谷歌Gemini 3正式发布，MMLU 1501分刷新纪录 \\
% 2025-11-03 & 产业合作 & 亚马逊与OpenAI 380亿美元战略合作，AWS成为独家云服务商 \\
% 2025-11-01 & 项目启动 & Clawdbot（OpenClaw）初始版本上线，开启“能动手做事”的Agent时代 \\
% \hline
% \multicolumn{3}{c}{\textbf{2025年10月（起点：安全风险初现）}} \\
% 2025-10-25 & 安全漏洞 & GitHub Copilot Codespaces令牌泄露，指令注入导致凭证外泄 \\
% 2025-10-15 & 产品发布 & 钉钉AgentOS发布，支持千级Agent协同 \\
% 2025-10-08 & 安全漏洞 & Anthropic MCP Git三漏洞，命令注入、路径遍历、信息泄露 \\
% \hline
% \end{tabularx}
% \end{table}


\subsection{AI智能体主流攻击工具}
本节整理汇总了近两年的主流AI智能体攻击工具（按照技术定位划分为学术研究型、开源全能渗透框架、开源专项攻击工具与AI增强传统安全工具~\cite{Xu2026LoopTrap,Liu2025AgentFuzz}），并分析各类工具的功能边界、技术原理与落地应用场景~\cite{MCPStrike,Crawford2026GhidraMCPAttack}，主要发现如下（详见表~\ref{tab:ai_agent_attack_tools_detail}）:
\begin{itemize}[tabitem]
  \item \textbf{攻击工具形成清晰分层体系，适配不同层级攻击者能力边界。}学术型工具侧重于机理验证与对抗样本研究，用于揭示底层安全缺陷~\cite{Nazary2025PoisonRAG}；开源全能框架面向红队人员实现端‑to‑end自动化渗透；专项攻击工具针对单一安全问题精准突破；AI赋能传统安全组件完成能力迁移，整套分层体系适配在校研究者、普通安全人员、高级红队等不同攻击者群体~\cite{HexStrike,AgentBreaker}。
  
  \item \textbf{MCP协议成为2026年攻击工具核心靶心，协议层风险可规模化利用。}以MCP‑Strike为代表的工具专门针对MCP‑Server开展提示注入、参数滥用、路径遍历和数据外传检测，依托内置LLM‑Judge自动化判定漏洞有效性~\cite{MCPStrike}；MCP作为智能体和外部组件交互的标准接口，协议本身缺少原生访问控制机制，使其成为当下可大规模批量利用的核心攻击面~\cite{Crawford2026GhidraMCPAttack}。
  
  \item \textbf{AI赋能重构传统安全工具链路，攻防能力形成闭环自动化。}Ghidra、SQLMap、Nuclei等经典安全工具借助Fast‑MCP适配智能体生态，由AI智能体自主完成漏洞模板生成、二进制逆向、漏洞探测、攻击路径推演、结果汇总等工作~\cite{SQLMap,FastMCP}；自动化工作流摆脱对人工专家的高度依赖，显著降低高级渗透攻击的准入门槛，促使攻防博弈由“人工主导”转变为“AI自主执行”~\cite{RedAmon}。
\end{itemize}


\begingroup
% \nopagebreak[4]
\enlargethispage{1cm}
\footnotesize
\centering

\DefTblrTemplate{conthead-text}{default}{（续表）}
\DefTblrTemplate{contfoot-text}{default}{续下页}

\begin{longtblr}[
  caption = {AI智能体主流攻击工具汇总},
  label   = {tab:ai_agent_attack_tools_detail},
]{
  width   = \linewidth,
  colspec = {
    |Q[wd=1.1cm,l,m]
    |Q[wd=2cm,l,m]
    |Q[wd=1.2cm,l,m]
    |Q[wd=0.6\linewidth,l,m]
    |X[1,l,m]|
  },
  cells   = {l,m},
  colsep  = 3pt,
  row{1}  = {font=\bfseries},
  rowhead = 1,
  % 每行内容与上下边框之间各保留少量空间；由该行最高单元格决定整行高度。
  rowsep  = 5pt,
  measure = vbox,
  stretch = -1,
}
\hline
\textbf{工具分类} & \textbf{工具名称} & \textbf{发布日期} &\textbf{核心功能} & \textbf{应用场景} \\
\hline
\SetCell[r=3]{l,m} 学术\newline 研究型
& AgentLoop \allowbreak ~\cite{Xu2026LoopTrap}
& 2026‑03
& \begin{itemize}[leftmargin=*,topsep=0pt,partopsep=0pt,itemsep=1pt,parsep=0pt]
    \item 用于测试智能体在资源消耗、执行失控和拒绝服务场景下的风险；
    \item 诱导LLM智能体错误判断任务未完成，进入循环执行状态。
\end{itemize}
& Agent终止条件投毒与资源放大攻击研究 \\
\cline{2-5}
& PoisonRAG \allowbreak ~\cite{Nazary2025PoisonRAG}
& 2025‑04
& \begin{itemize}[leftmargin=*,topsep=0pt,partopsep=0pt,itemsep=1pt,parsep=0pt]
    \item 用于测试RAG系统在恶意数据注入和推荐结果偏移下的安全性；
    \item 对RAG推荐系统的检索数据进行投毒，操控系统返回结果；
    \item 通过篡改项目标签、描述等元数据影响检索增强生成流程。
\end{itemize}
& RAG推荐系统数据投毒与推荐结果操纵研究 \\
\cline{2-5}
& AgentFuzz \allowbreak~\cite{Liu2025AgentFuzz}
& 2025‑03
& \begin{itemize}[leftmargin=*,topsep=0pt,partopsep=0pt,itemsep=1pt,parsep=0pt]
    \item 自动检测LLM Agent工具调用链中的污染型漏洞；
    \item 生成面向目标智能体的测试提示与工具调用输入；
    \item 根据执行反馈自适应变异测试样本，提高漏洞触发与复现效率。
\end{itemize}
& 漏洞挖掘、安全测试 \\
% \cline{2-5}
% & MRJ‑Agent ~\cite{Wang2024MRJAgent}
% & 2024‑08
% & \begin{itemize}[leftmargin=*,topsep=0pt,partopsep=0pt,itemsep=1pt,parsep=0pt]
%     \item 自动构造多轮对话越狱攻击，测试模型在连续交互中的安全边界；
%     \item 将高风险目标拆解为多轮渐进式请求，诱导模型逐步偏离安全约束；
%     \item 用于评估模型在多轮上下文累积场景下的越狱风险。
% \end{itemize}
% & 安全边界测试、对抗样本研究 \\
% \cline{2-5}
% & PromptInjector~\cite{Liu2024PromptInjection}
% & 2024‑07
% & \begin{itemize}[leftmargin=*,topsep=0pt,partopsep=0pt,itemsep=1pt,parsep=0pt]
%     \item 用于对比不同提示注入攻击方式和防御机制的有效性
%     \item 构造提示注入攻击，测试 LLM 应用是否会被外部指令劫持
% \end{itemize}
% & 提示注入攻防实验 \\
\hline
\SetCell[r=6]{l,m} 开源全能\newline 渗透框架
& MCP‑Strike \allowbreak~\cite{MCPStrike}
& 2026‑06
& \begin{itemize}[leftmargin=*,topsep=0pt,partopsep=0pt,itemsep=1pt,parsep=0pt]
    \item 对运行中的MCP Server进行主动式对抗安全测试；
    \item 测试工具描述提示注入、过度敏感参数、路径遍历、响应注入和数据外传诱导等风险；
    \item 借助LLM judge或自适应 LLM智能体对发现结果进行确认、评分和多轮探测。
\end{itemize}
& MCP服务器主动对抗测试与安全评估 \\
\cline{2-5}
& RedAmon \allowbreak~\cite{RedAmon}
& 2026‑06
& \begin{itemize}[leftmargin=*,topsep=0pt,partopsep=0pt,itemsep=1pt,parsep=0pt]
    \item 执行侦察、漏洞验证、后渗透分析、AI分诊、代码修复和GitHub PR提的一体化安全测试流程；
    \item 并行运行多个侦察工具，收集子域名、端口、端点等攻击面信息，并写入Neo4j知识图谱。
\end{itemize}
& 授权自动化渗透测试、攻击面分析与漏洞修复闭环 \\
\cline{2-5}
& Strix~\cite{Strix}
& 2025‑10
& \begin{itemize}[leftmargin=*,topsep=0pt,partopsep=0pt,itemsep=1pt,parsep=0pt]
    \item 动态运行目标应用代码，发现漏洞并生成PoC对漏洞进行验证；
    \item 编排多个专用Agent并行测试复杂目标，并共享安全发现；
    \item 进行应用安全测试、快速渗透测试和Bug Bounty自动化研究。
\end{itemize}
& 应用安全测试、快速渗透测试、漏洞赏金与CI/CD安全测试 \\
\cline{2-5}
& HexStrike \allowbreak~\cite{HexStrike}
& 2025‑09
& \begin{itemize}[leftmargin=*,topsep=0pt,partopsep=0pt,itemsep=1pt,parsep=0pt]
    \item 指挥AI智能体调用150+ 安全工具，执行自动化渗透测试、漏洞发现和安全研究任务；
    \item 编排多个安全智能体协同完成侦察、扫描、利用验证和结果分析；
    \item 对网络、Web 应用、云环境、容器、二进制程序和OSINT目标开展多类型安全测试。
\end{itemize}
& 红队演练、渗透测试 \\
\cline{2-5}
& Ankou~\cite{Ankou}
& 2025‑08
& \begin{itemize}[leftmargin=*,topsep=0pt,partopsep=0pt,itemsep=1pt,parsep=0pt]
    \item 提供面向红队演练的模块化C2架构、监听器和agent管理；
    \item 支持多种传输与自定义agent/handler扩展；
    \item 通过AI Operations面板辅助总结结果、发现异常和建议后续；操作。
\end{itemize}
& 授权红队C2、Agent管理与后渗透操作 \\
\cline{2-5}
& RAPTOR \allowbreak~\cite{RAPTOR}
& 2025‑06
& \begin{itemize}[leftmargin=*,topsep=0pt,partopsep=0pt,itemsep=1pt,parsep=0pt]
    \item 对代码库或二进制目标进行攻击面分析、数据流追踪和漏洞变体搜索；
    \item 进行 Semgrep/CodeQL 扫描、软件组成分析、模糊测试和崩溃根源分析。
\end{itemize}
& 源代码与二进制漏洞分析、验证及修复 \\
% \cline{2-5}
% & AutoPentest‑DRL~\cite{AutoPentestDRL}
% & 2024‑11
% & \begin{itemize}[leftmargin=*,topsep=0pt,partopsep=0pt,itemsep=1pt,parsep=0pt]
%     \item 结合 Nmap 和 Metasploit 对真实网络进行扫描、漏洞验证和路径执行
%     \item 基于深度强化学习为目标网络规划自动化渗透攻击路径
% \end{itemize}
% & 复杂网络环境渗透测试 \\
\hline
\SetCell[r=4]{l,m} 开源专项\newline 攻击工具
& Deadend CLI~\cite{DeadendCLI}
& 2025‑12
& \begin{itemize}[leftmargin=*,topsep=0pt,partopsep=0pt,itemsep=1pt,parsep=0pt]
    \item 提供自托管的Agentic pentest CLI工作流；
    \item 支持通过LiteLLM 接入多种模型，用于黑盒渗透测试任务；
    \item 在XBOW validation suite等基准上评估自动化安全测试能力。
\end{itemize}
& Web应用黑盒渗透测试与漏洞验证 \\
\cline{2-5}
& AgentBreaker\allowbreak~\cite{AgentBreaker}
& 2025‑07
& \begin{itemize}[leftmargin=*,topsep=0pt,partopsep=0pt,itemsep=1pt,parsep=0pt]
    \item 自动测试提示注入、越狱、数据泄漏、护栏绕过、提示提取、工具滥用、数据外传和多模态注入风险；
    \item 根据已经发现的系统行为和漏洞，自适应规划后续测试攻击；
    \item 对发现结果进行严重性评分，并按照OWASP和MITRE框架进行分类。
\end{itemize}
& AI系统自动化红队测试与漏洞评估 \\
\cline{2-5}
& CAI Framework~\cite{CAI,CAI_paper}
& 2025‑06
& \begin{itemize}[leftmargin=*,topsep=0pt,partopsep=0pt,itemsep=1pt,parsep=0pt]
    \item 提供面向攻防安全自动化的开源AI智能体框架；
    \item 支持构建专用安全Agent，覆盖侦察、漏洞发现、利用验证和安全评估；
    \item 结合工具集成、人类监督和guardrails，用于授权Bug Bounty、CTF 和安全测试。
\end{itemize}
& 定制化攻防演练 \\
\cline{2-5}
& Nebula~\cite{Nebula}
& 2025‑05
& \begin{itemize}[leftmargin=*,topsep=0pt,partopsep=0pt,itemsep=1pt,parsep=0pt]
    \item 将AI模型集成到命令行渗透测试工作流中；
    \item 自动化记录、分类安全发现，并提供实时AI分析建议；
    \item 支持调用外部工具结果进行漏洞分析、报告整理和测试过程留痕。
\end{itemize}
& 渗透测试、漏洞分析与测试记录 \\
% \cline{2-5}
% & Pentest GPT~\cite{Deng2024PentestGPT}
% & 2024‑09
% & \begin{itemize}[leftmargin=*,topsep=0pt,partopsep=0pt,itemsep=1pt,parsep=0pt]
%     \item 基于 LLM 辅助渗透测试过程的任务拆解、工具使用和结果解释。
%     \item 通过多模块交互缓解长链路渗透测试中的上下文丢失问题。
%     \item 支持真实测试目标和 CTF 场景中的攻击路径推演与验证。
% \end{itemize}
% & 自动化渗透测试、真实目标安全评估与CTF \\
\hline
\SetCell[r=2]{l,m} AI增强传统安全工具
& Ghidra + GhidraMCP \allowbreak~\cite{Crawford2026GhidraMCPAttack}
& 2026‑04
& \begin{itemize}[leftmargin=*,topsep=0pt,partopsep=0pt,itemsep=1pt,parsep=0pt]
    \item 研究 GhidraMCP等工具增强的逆向AI智能体如何自动化二进制分析；
    \item 通过对反编译代码植入隐蔽指令，评估LLM驱动逆向流程的提示注入风险；
    \item 分析攻击如何误导disassembly/decompilation解释并污染自动化分析结果。
\end{itemize}
& 可执行二进制与恶意软件逆向分析、逆向AI Agent安全研究 \\
\cline{2-5}
& SQLMap + FastMCP \allowbreak~\cite{SQLMap,FastMCP}
& 2025‑02
& \begin{itemize}[leftmargin=*,topsep=0pt,partopsep=0pt,itemsep=1pt,parsep=0pt]
    \item 将sqlmap的SQL注入检测、利用验证、数据库指纹识别等能力封装为MCP工具；
    \item 通过FastMCP为智能体提供标准化工具接口、参数校验与调用结果返回；
    \item 支持智能体在授权测试场景中编排SQL注入检测流程，并辅助整理风险结果。
\end{itemize}
& Web应用SQL注入检测与MCP工具集成 \\
% \cline{2-5}
% & Nuclei（ + AI 模板扩展） ~\cite{Nuclei}
% & 2024‑04
% & \begin{itemize}[leftmargin=*,topsep=0pt,partopsep=0pt,itemsep=1pt,parsep=0pt]
%     \item 基于 YAML 模板进行高性能漏洞扫描与自定义检测。
%     \item 通过 AI prompt 辅助生成/运行检测模板。
%     \item 用于 Web、网络、DNS、云配置等多协议资产的批量扫描。
% \end{itemize}
% & 批量漏洞扫描 \\
% \cline{2-5}
% & BloodHound CE~\cite{BloodHound}
% & 2023‑11
% & \begin{itemize}[leftmargin=*,topsep=0pt,partopsep=0pt,itemsep=1pt,parsep=0pt]
%     \item 基于图数据库分析身份、权限和访问关系。
%     \item 帮助红队发现复杂攻击路径，帮助蓝队识别并缓解权限风险。
%     \item 映射不同身份平台中的权限关系和潜在攻击路径
% \end{itemize}
% & 身份与权限关系分析、攻击路径发现与风险缓解 \\
\hline
\end{longtblr}
\endgroup

% \begin{center}
% % \nopagebreak[4]
% \enlargethispage{1cm}
% \footnotesize
% \setlength{\tabcolsep}{3pt}
% \begin{longtable}{
% |>{\rule{0pt}{2pt}}m{1.1cm}
% |>{\rule{0pt}{2pt}}m{2.0cm}
% |>{\rule{0pt}{2pt}}m{1.2cm}
% |>{\rule{0pt}{2pt}}m{\dimexpr0.6\linewidth}
% |>{\rule{0pt}{2pt}}m{\dimexpr\linewidth - 1.1cm - 2.0cm - 1.2cm - 0.6\linewidth - 10\tabcolsep\relax}|}
% \caption{AI智能体主流攻击工具汇总}\label{tab:ai_agent_attack_tools_detail}\\
% \hline
% \textbf{工具分类} & \textbf{工具名称} & \textbf{发布日期} &\textbf{核心功能} & \textbf{应用场景} \\
% \hline
% \endfirsthead

% \multicolumn{5}{l}{\tablename\ \thetable\ -- 续上页}\\
% \hline
% \textbf{工具分类} & \textbf{工具名称} & \textbf{发布日期} &\textbf{核心功能} & \textbf{应用场景} \\
% \hline
% \endhead
% \hline
% \multicolumn{5}{r}{续下页}\\
% \endfoot
% \hline
% \endlastfoot

% \multirow{5}{=}{学术\\研究型}
% & AgentLoop \allowbreak ~\cite{Xu2026LoopTrap}
% & 2026‑03
% & \begin{itemize}[tabitem]
%     \item 用于测试智能体在资源消耗、执行失控和拒绝服务场景下的风险；
%     \item 诱导LLM智能体错误判断任务未完成，进入循环执行状态。
% \end{itemize}
% & Agent终止条件投毒与资源放大攻击研究 \\
% \cline{2-5}
% & PoisonRAG \allowbreak ~\cite{Nazary2025PoisonRAG}
% & 2025‑04
% & \begin{itemize}[tabitem]
%     \item 用于测试RAG系统在恶意数据注入和推荐结果偏移下的安全性；
%     \item 对RAG推荐系统的检索数据进行投毒，操控系统返回结果；
%     \item 通过篡改项目标签、描述等元数据影响检索增强生成流程。
% \end{itemize}
% & RAG推荐系统数据投毒与推荐结果操纵研究 \\
% \cline{2-5}
% & AgentFuzz \allowbreak~\cite{Liu2025AgentFuzz}
% & 2025‑03
% & \begin{itemize}[tabitem]
%     \item 自动检测LLM Agent工具调用链中的污染型漏洞；
%     \item 生成面向目标智能体的测试提示与工具调用输入；
%     \item 根据执行反馈自适应变异测试样本，提高漏洞触发与复现效率。
% \end{itemize}
% & 漏洞挖掘、安全测试 \\
% % \cline{2-5}
% % & MRJ‑Agent ~\cite{Wang2024MRJAgent}
% % & 2024‑08
% % & \begin{itemize}[tabitem]
% %     \item 自动构造多轮对话越狱攻击，测试模型在连续交互中的安全边界；
% %     \item 将高风险目标拆解为多轮渐进式请求，诱导模型逐步偏离安全约束；
% %     \item 用于评估模型在多轮上下文累积场景下的越狱风险。
% % \end{itemize}
% % & 安全边界测试、对抗样本研究 \\
% % \cline{2-5}
% % & PromptInjector~\cite{Liu2024PromptInjection}
% % & 2024‑07
% % & \begin{itemize}[tabitem]
% %     \item 用于对比不同提示注入攻击方式和防御机制的有效性
% %     \item 构造提示注入攻击，测试 LLM 应用是否会被外部指令劫持
% % \end{itemize}
% % & 提示注入攻防实验 \\
% \hline
% \multirow{7}{=}{开源全能\\渗透框架}
% & MCP‑Strike \allowbreak~\cite{MCPStrike}
% & 2026‑06
% & \begin{itemize}[tabitem]
%     \item 对运行中的MCP Server进行主动式对抗安全测试；
%     \item 测试工具描述提示注入、过度敏感参数、路径遍历、响应注入和数据外传诱导等风险；
%     \item 借助LLM judge或自适应 LLM智能体对发现结果进行确认、评分和多轮探测。
% \end{itemize}
% & MCP服务器主动对抗测试与安全评估 \\
% \cline{2-5}
% & RedAmon \allowbreak~\cite{RedAmon}
% & 2026‑06
% & \begin{itemize}[tabitem]
%     \item 执行侦察、漏洞验证、后渗透分析、AI分诊、代码修复和GitHub PR提的一体化安全测试流程；
%     \item 并行运行多个侦察工具，收集子域名、端口、端点等攻击面信息，并写入Neo4j知识图谱。
% \end{itemize}
% & 授权自动化渗透测试、攻击面分析与漏洞修复闭环 \\
% \cline{2-5}
% & Strix~\cite{Strix}
% & 2025‑10
% & \begin{itemize}[tabitem]
%     \item 动态运行目标应用代码，发现漏洞并生成PoC对漏洞进行验证；
%     \item 编排多个专用Agent并行测试复杂目标，并共享安全发现；
%     \item 进行应用安全测试、快速渗透测试和Bug Bounty自动化研究。
% \end{itemize}
% & 应用安全测试、快速渗透测试、漏洞赏金与CI/CD安全测试 \\
% \cline{2-5}
% & HexStrike \allowbreak~\cite{HexStrike}
% & 2025‑09
% & \begin{itemize}[tabitem]
%     \item 指挥AI智能体调用150+ 安全工具，执行自动化渗透测试、漏洞发现和安全研究任务；
%     \item 编排多个安全智能体协同完成侦察、扫描、利用验证和结果分析；
%     \item 对网络、Web 应用、云环境、容器、二进制程序和OSINT目标开展多类型安全测试。
% \end{itemize}
% & 红队演练、渗透测试 \\
% \cline{2-5}
% & Ankou~\cite{Ankou}
% & 2025‑08
% & \begin{itemize}[tabitem]
%     \item 提供面向红队演练的模块化C2架构、监听器和agent管理；
%     \item 支持多种传输与自定义agent/handler扩展；
%     \item 通过AI Operations面板辅助总结结果、发现异常和建议后续；操作。
% \end{itemize}
% & 授权红队C2、Agent管理与后渗透操作 \\
% \cline{2-5}
% & RAPTOR \allowbreak~\cite{RAPTOR}
% & 2025‑06
% & \begin{itemize}[tabitem]
%     \item 对代码库或二进制目标进行攻击面分析、数据流追踪和漏洞变体搜索；
%     \item 进行 Semgrep/CodeQL 扫描、软件组成分析、模糊测试和崩溃根源分析。
% \end{itemize}
% & 源代码与二进制漏洞分析、验证及修复 \\
% % \cline{2-5}
% % & AutoPentest‑DRL~\cite{AutoPentestDRL}
% % & 2024‑11
% % & \begin{itemize}[tabitem]
% %     \item 结合 Nmap 和 Metasploit 对真实网络进行扫描、漏洞验证和路径执行
% %     \item 基于深度强化学习为目标网络规划自动化渗透攻击路径
% % \end{itemize}
% % & 复杂网络环境渗透测试 \\
% \hline
% \multirow{5}{=}{开源专项\\攻击工具}
% & Deadend CLI~\cite{DeadendCLI}
% & 2025‑12
% & \begin{itemize}[tabitem]
%     \item 提供自托管的Agentic pentest CLI工作流；
%     \item 支持通过LiteLLM 接入多种模型，用于黑盒渗透测试任务；
%     \item 在XBOW validation suite等基准上评估自动化安全测试能力。
% \end{itemize}
% & Web应用黑盒渗透测试与漏洞验证 \\
% \cline{2-5}
% & AgentBreaker\allowbreak~\cite{AgentBreaker}
% & 2025‑07
% & \begin{itemize}[tabitem]
%     \item 自动测试提示注入、越狱、数据泄漏、护栏绕过、提示提取、工具滥用、数据外传和多模态注入风险；
%     \item 根据已经发现的系统行为和漏洞，自适应规划后续测试攻击；
%     \item 对发现结果进行严重性评分，并按照OWASP和MITRE框架进行分类。
% \end{itemize}
% & AI系统自动化红队测试与漏洞评估 \\
% \cline{2-5}
% & CAI Framework~\cite{CAI,CAI_paper}
% & 2025‑06
% & \begin{itemize}[tabitem]
%     \item 提供面向攻防安全自动化的开源AI智能体框架；
%     \item 支持构建专用安全Agent，覆盖侦察、漏洞发现、利用验证和安全评估；
%     \item 结合工具集成、人类监督和guardrails，用于授权Bug Bounty、CTF 和安全测试。
% \end{itemize}
% & 定制化攻防演练 \\
% \cline{2-5}
% & Nebula~\cite{Nebula}
% & 2025‑05
% & \begin{itemize}[tabitem]
%     \item 将AI模型集成到命令行渗透测试工作流中；
%     \item 自动化记录、分类安全发现，并提供实时AI分析建议；
%     \item 支持调用外部工具结果进行漏洞分析、报告整理和测试过程留痕。
% \end{itemize}
% & 渗透测试、漏洞分析与测试记录 \\
% % \cline{2-5}
% % & Pentest GPT~\cite{Deng2024PentestGPT}
% % & 2024‑09
% % & \begin{itemize}[tabitem]
% %     \item 基于 LLM 辅助渗透测试过程的任务拆解、工具使用和结果解释。
% %     \item 通过多模块交互缓解长链路渗透测试中的上下文丢失问题。
% %     \item 支持真实测试目标和 CTF 场景中的攻击路径推演与验证。
% % \end{itemize}
% % & 自动化渗透测试、真实目标安全评估与CTF \\
% \hline
% \multirow{4}{=}{AI增强传统安全工具}
% & Ghidra + GhidraMCP \allowbreak~\cite{Crawford2026GhidraMCPAttack}
% & 2026‑04
% & \begin{itemize}[tabitem]
%     \item 研究 GhidraMCP等工具增强的逆向AI智能体如何自动化二进制分析；
%     \item 通过对反编译代码植入隐蔽指令，评估LLM驱动逆向流程的提示注入风险；
%     \item 分析攻击如何误导disassembly/decompilation解释并污染自动化分析结果。
% \end{itemize}
% & 可执行二进制与恶意软件逆向分析、逆向AI Agent安全研究 \\
% \cline{2-5}
% & SQLMap + FastMCP \allowbreak~\cite{SQLMap,FastMCP}
% & 2025‑02
% & \begin{itemize}[tabitem]
%     \item 将sqlmap的SQL注入检测、利用验证、数据库指纹识别等能力封装为MCP工具；
%     \item 通过FastMCP为智能体提供标准化工具接口、参数校验与调用结果返回；
%     \item 支持智能体在授权测试场景中编排SQL注入检测流程，并辅助整理风险结果。
% \end{itemize}
% & Web应用SQL注入检测与MCP工具集成 
% % \cline{2-5}
% % & Nuclei（ + AI 模板扩展） ~\cite{Nuclei}
% % & 2024‑04
% % & \begin{itemize}[tabitem]
% %     \item 基于 YAML 模板进行高性能漏洞扫描与自定义检测。
% %     \item 通过 AI prompt 辅助生成/运行检测模板。
% %     \item 用于 Web、网络、DNS、云配置等多协议资产的批量扫描。
% % \end{itemize}
% % & 批量漏洞扫描 \\
% % \cline{2-5}
% % & BloodHound CE~\cite{BloodHound}
% % & 2023‑11
% % & \begin{itemize}[tabitem]
% %     \item 基于图数据库分析身份、权限和访问关系。
% %     \item 帮助红队发现复杂攻击路径，帮助蓝队识别并缓解权限风险。
% %     \item 映射不同身份平台中的权限关系和潜在攻击路径
% % \end{itemize}
% % & 身份与权限关系分析、攻击路径发现与风险缓解 \\
% \end{longtable}
% \end{center}

% \begin{center}
% % \nopagebreak[4]
% \enlargethispage{1cm}
% \footnotesize
% \setlength{\tabcolsep}{3pt}
% % \renewcommand{\arraystretch}{0.7}{0.7}
% \begin{longtable}{
% |>{\rule{0pt}{2pt}}m{1.5cm}
% |>{\rule{0pt}{2pt}}m{2.0cm}
% |>{\rule{0pt}{2pt}}m{\dimexpr 0.6\linewidth}
% |>{\rule{0pt}{2pt}}m{\dimexpr\linewidth - 1.5cm - 2.0cm - 0.6\linewidth - 8\tabcolsep\relax}|}
% \caption{AI智能体攻击工具}\label{tab:ai_agent_attack_tools_detail}\\
% \hline
% \textbf{工具分类} & \textbf{工具名称} & \textbf{核心功能} & \textbf{应用场景} \\
% \hline
% \endfirsthead

% \multicolumn{4}{l}{\tablename\ \thetable\ -- 续上页}\\
% \hline
% \textbf{工具分类} & \textbf{工具名称} & \textbf{核心功能} & \textbf{应用场景} \\
% \hline
% \endhead
% \hline
% \multicolumn{4}{r}{续下页}\\
% \endfoot
% \hline
% \endlastfoot

% \multirow{5}{=}{学术\\研究型}
% & AgentFuzz~\cite{Liu2025AgentFuzz}
% & \begin{itemize}[tabitem]
%     \item 自动检测LLM Agent工具调用链中的污染型漏洞
%     \item 生成面向目标智能体的测试提示与工具调用输入
%     \item 根据执行反馈自适应变异测试样本，提高漏洞触发与复现效率
% \end{itemize}
% & 漏洞挖掘、安全测试 \\
% \cline{2-4}
% & MRJ-Agent ~\cite{Wang2024MRJAgent}
% & \begin{itemize}[tabitem]
%     \item 自动构造多轮对话越狱攻击，测试模型在连续交互中的安全边界
%     \item 将高风险目标拆解为多轮渐进式请求，诱导模型逐步偏离安全约束
%     \item 用于评估模型在多轮上下文累积场景下的越狱风险
% \end{itemize}
% & 安全边界测试、对抗样本研究 \\
% \cline{2-4}
% & PoisonRAG~\cite{Nazary2025PoisonRAG}
% & \begin{itemize}[tabitem]
%     \item 用于测试 RAG 系统在恶意数据注入和推荐结果偏移下的安全性
%     \item 对 RAG 推荐系统的检索数据进行投毒，操控系统返回结果
%     \item 通过篡改项目标签、描述等元数据影响检索增强生成流程
% \end{itemize}
% & RAG推荐系统数据投毒与推荐结果操纵研究 \\
% \cline{2-4}
% & PromptInjector~\cite{Liu2024PromptInjection}
% & \begin{itemize}[tabitem]
%     \item 用于对比不同提示注入攻击方式和防御机制的有效性
%     \item 构造提示注入攻击，测试 LLM 应用是否会被外部指令劫持
% \end{itemize}
% & 提示注入攻防实验 \\
% \cline{2-4}
% & AgentLoop~\cite{Xu2026LoopTrap}
% & \begin{itemize}[tabitem]
%     \item 用于测试智能体在资源消耗、执行失控和拒绝服务场景下的风险
%     \item 诱导 LLM智能体错误判断任务未完成，进入循环执行状态
% \end{itemize}
% & Agent终止条件投毒与资源放大攻击研究 \\
% \hline
% \multirow{7}{=}{开源全能\\渗透框架}
% & HexStrike~\cite{HexStrike}
% & \begin{itemize}[tabitem]
%     \item 指挥 AI智能体调用 150+ 安全工具，执行自动化渗透测试、漏洞发现和安全研究任务
%     \item 编排多个安全智能体协同完成侦察、扫描、利用验证和结果分析
%     \item 对网络、Web 应用、云环境、容器、二进制程序和 OSINT 目标开展多类型安全测试
% \end{itemize}
% & 红队演练、渗透测试 \\
% \cline{2-4}
% & Strix~\cite{Strix}
% & \begin{itemize}[tabitem]
%     \item 动态运行目标应用代码，发现漏洞并生成 PoC 对漏洞进行验证
%     \item 编排多个专用Agent并行测试复杂目标，并共享安全发现
%     \item 进行应用安全测试、快速渗透测试和 Bug Bounty 自动化研究
% \end{itemize}
% & 应用安全测试、快速渗透测试、漏洞赏金与CI/CD安全测试 \\
% \cline{2-4}
% & RAPTOR~\cite{RAPTOR}
% & \begin{itemize}[tabitem]
%     \item 对代码库或二进制目标进行攻击面分析、数据流追踪和漏洞变体搜索
%     \item 进行 Semgrep/CodeQL 扫描、软件组成分析、模糊测试和崩溃根源分析
% \end{itemize}
% & 源代码与二进制漏洞分析、验证及修复 \\
% \cline{2-4}
% & RedAmon~\cite{RedAmon}
% & \begin{itemize}[tabitem]
%     \item 执行侦察、漏洞验证、后渗透分析、AI 分诊、代码修复和GitHub PR提交 的一体化安全测试流程。
%     \item 并行运行多个侦察工具，收集子域名、端口、端点等攻击面信息，并写入 Neo4j 知识图谱
% \end{itemize}
% & 授权自动化渗透测试、攻击面分析与漏洞修复闭环 \\
% \cline{2-4}
% & Ankou~\cite{Ankou}
% & \begin{itemize}[tabitem]
%     \item 提供面向红队演练的模块化 C2 架构、监听器和 agent 管理。
%     \item 支持多种传输与自定义 agent/handler 扩展。
%     \item 通过 AI Operations 面板辅助总结结果、发现异常和建议后续操作。
% \end{itemize}
% & 授权红队C2、Agent管理与后渗透操作 \\
% \cline{2-4}
% & AutoPentest-DRL~\cite{AutoPentestDRL}
% & \begin{itemize}[tabitem]
%     \item 结合 Nmap 和 Metasploit 对真实网络进行扫描、漏洞验证和路径执行
%     \item 基于深度强化学习为目标网络规划自动化渗透攻击路径
% \end{itemize}
% & 复杂网络环境渗透测试 \\
% \cline{2-4}
% & MCP-Strike(2026.06)~\cite{MCPStrike}
% & \begin{itemize}[tabitem]
%     \item 对运行中的 MCP Server 进行主动式对抗安全测试
%     \item 测试工具描述提示注入、过度敏感参数、路径遍历、响应注入和数据外传诱导等风险
%     \item 借助 LLM judge 或自适应 LLM智能体对发现结果进行确认、评分和多轮探测
% \end{itemize}
% & MCP服务器主动对抗测试与安全评估 \\
% \hline
% \multirow{5}{=}{开源专项\\攻击工具}
% & AgentBreaker~\cite{AgentBreaker}
% & \begin{itemize}[tabitem]
%     \item 自动测试提示注入、越狱、数据泄漏、护栏绕过、提示提取、工具滥用、数据外传和多模态注入风险
%     \item 根据已经发现的系统行为和漏洞，自适应规划后续测试攻击
%     \item 对发现结果进行严重性评分，并按照OWASP和MITRE框架进行分类
% \end{itemize}
% & AI系统自动化红队测试与漏洞评估 \\
% \cline{2-4}
% & Nebula~\cite{Nebula}
% & \begin{itemize}[tabitem]
%     \item 将 AI 模型集成到命令行渗透测试工作流中。
%     \item 自动化记录、分类安全发现，并提供实时 AI 分析建议。
%     \item 支持调用外部工具结果进行漏洞分析、报告整理和测试过程留痕。
% \end{itemize}
% & 渗透测试、漏洞分析与测试记录 \\
% \cline{2-4}
% & Deadend CLI~\cite{DeadendCLI}
% & \begin{itemize}[tabitem]
%     \item 提供自托管的 Agentic pentest CLI 工作流。
%     \item 支持通过 LiteLLM 接入多种模型，用于黑盒渗透测试任务。
%     \item 在 XBOW validation suite 等基准上评估自动化安全测试能力
% \end{itemize}
% & Web应用黑盒渗透测试与漏洞验证 \\
% \cline{2-4}
% & Pentest GPT~\cite{Deng2024PentestGPT}
% & \begin{itemize}[tabitem]
%     \item 基于 LLM 辅助渗透测试过程的任务拆解、工具使用和结果解释。
%     \item 通过多模块交互缓解长链路渗透测试中的上下文丢失问题。
%     \item 支持真实测试目标和 CTF 场景中的攻击路径推演与验证。
% \end{itemize}
% & 自动化渗透测试、真实目标安全评估与CTF \\
% \cline{2-4}
% & CAI Framework~\cite{CAI,CAI_paper}
% & \begin{itemize}[tabitem]
%     \item 提供面向攻防安全自动化的开源 AI智能体框架。
%     \item 支持构建专用安全 Agent，覆盖侦察、漏洞发现、利用验证和安全评估。
%     \item 结合工具集成、人类监督和 guardrails，用于授权 Bug Bounty、CTF 和安全测试。
% \end{itemize}
% & 定制化攻防演练 \\
% \hline
% \multirow{4}{=}{AI 增强传统安全工具}
% & Nuclei（ + AI 模板扩展） ~\cite{Nuclei}
% & \begin{itemize}[tabitem]
%     \item 基于 YAML 模板进行高性能漏洞扫描与自定义检测。
%     \item 通过 AI prompt 辅助生成/运行检测模板。
%     \item 用于 Web、网络、DNS、云配置等多协议资产的批量扫描。
% \end{itemize}
% & 批量漏洞扫描 \\
% \cline{2-4}
% & Ghidra + GhidraMCP~\cite{Crawford2026GhidraMCPAttack}
% & \begin{itemize}[tabitem]
%     \item 研究 GhidraMCP 等工具增强的逆向 AI智能体如何自动化二进制分析。
%     \item 通过对反编译代码植入隐蔽指令，评估 LLM 驱动逆向流程的提示注入风险。
%     \item 分析攻击如何误导 disassembly/decompilation 解释并污染自动化分析结果。
% \end{itemize}
% & 可执行二进制与恶意软件逆向分析、逆向AI Agent安全研究 \\
% \cline{2-4}
% & BloodHound CE~\cite{BloodHound}
% & \begin{itemize}[tabitem]
%     \item 基于图数据库分析身份、权限和访问关系。
%     \item 帮助红队发现复杂攻击路径，帮助蓝队识别并缓解权限风险。
%     \item 映射不同身份平台中的权限关系和潜在攻击路径
% \end{itemize}
% & 身份与权限关系分析、攻击路径发现与风险缓解 \\
% \cline{2-4}
% & SQLMap + FastMCP~\cite{SQLMap,FastMCP}
% & \begin{itemize}[tabitem]
%     \item 将 sqlmap 的 SQL 注入检测、利用验证、数据库指纹识别等能力封装为 MCP 工具
%     \item 通过 FastMCP 为智能体提供标准化工具接口、参数校验与调用结果返回
%     \item 支持智能体在授权测试场景中编排 SQL 注入检测流程，并辅助整理风险结果
% \end{itemize}
% & Web应用SQL注入检测与MCP工具集成 
% \end{longtable}
% \end{center}

